\documentclass[11pt,a4paper]{article}
\usepackage[T1]{fontenc}
\usepackage[utf8]{inputenc}
\usepackage{lmodern,amsmath,amssymb,amsthm,mathtools}
\usepackage[margin=27mm]{geometry}
\usepackage{microtype,booktabs,array,longtable,enumitem,listings,xcolor,xurl}
\usepackage[colorlinks=true,linkcolor=blue!35!black,citecolor=blue!35!black,urlcolor=blue!40!black,breaklinks=true]{hyperref}
\setlength{\emergencystretch}{3em}\setlength{\parskip}{3pt}
\setlist{itemsep=2pt,topsep=4pt}\raggedbottom
\setcounter{tocdepth}{1}
\newtheorem{theorem}{Theorem}[section]
\newtheorem{lemma}[theorem]{Lemma}\newtheorem{proposition}[theorem]{Proposition}
\newtheorem{corollary}[theorem]{Corollary}\newtheorem{mainthm}[theorem]{Theorem}
\theoremstyle{definition}\newtheorem{definition}[theorem]{Definition}
\theoremstyle{remark}\newtheorem{remark}[theorem]{Remark}
\newcommand{\N}{\mathbb N}\newcommand{\Z}{\mathbb Z}\newcommand{\Q}{\mathbb Q}
\newcommand{\R}{\mathbb R}\newcommand{\C}{\mathbb C}\newcommand{\F}{\mathbb F}
\newcommand{\PP}{\mathbb P}\newcommand{\W}{\mathcal W}\newcommand{\SaitoL}{\mathcal S}
\newcommand{\floor}[1]{\left\lfloor#1\right\rfloor}
\newcommand{\ceil}[1]{\left\lceil#1\right\rceil}
\newcommand{\xir}{\xi_r}\newcommand{\Nzero}{N_0}\newcommand{\Nb}{N_{\mathrm B}}
\newcommand{\No}{N_{\mathrm O}}\newcommand{\Nu}{N_{\mathrm U}}
\newcommand{\decrows}[2]{\left[\begin{gathered}\mathtt{#1}\\[-2pt]\mathtt{#2}\end{gathered}\right]_{10}}
\DeclareMathOperator{\Tr}{Tr}\DeclareMathOperator{\tr}{tr}\DeclareMathOperator{\Nm}{N}
\DeclareMathOperator{\Norm}{Norm}\DeclareMathOperator{\ord}{ord}\DeclareMathOperator{\Gal}{Gal}
\DeclareMathOperator{\lcm}{lcm}\DeclareMathOperator{\disc}{disc}\DeclareMathOperator{\Res}{Res}
\DeclareMathOperator{\Fix}{Fix}\DeclareMathOperator{\GL}{GL}\newcommand{\leg}[2]{\left(\frac{#1}{#2}\right)}
\lstset{language=Python,basicstyle=\ttfamily\scriptsize,keywordstyle=\bfseries,
commentstyle=\color{black!65},showstringspaces=false,columns=fullflexible,
keepspaces=true,breaklines=true,breakatwhitespace=false,frame=none,tabsize=4}
\hypersetup{pdftitle={Shifted Mills Constants and Cubic Pisot Trace Dynamics},pdfauthor={KEIO},pdfsubject={A unified treatment of affine schedules, recurring composite floors, and arithmetic barriers}}
\title{Shifted Mills Constants and Cubic Pisot Trace Dynamics}
\author{KEIO}\date{27 October 2026}
\begin{document}\maketitle
\begin{abstract}
We develop a unified treatment of least prime-representing constants on affine
ternary exponent schedules and of recurring composite floors for cubic Pisot
numbers. The positive-shift argument reduces algebraic candidates to three
units and eliminates them by exact local certificates and an infinite prime
continuation. Negative shifts admit fixed-point and multiplier extensions,
but leave a genuinely different even-multiplier problem. On cubic coefficient
lines we obtain simultaneous composite-floor returns for density-one sets,
including an improved logarithmic exceptional exponent $2/3$ and prime-relative bounds with exponent $5/3$, while explicit thin families
require adaptive projective witnesses. Finite-field trace-zero geometry,
cyclotomic norm bounds, and a two-sign three-adic depth theorem for arbitrary irreducible totally real cubics provide exact tests
for hypothetical classical Mills candidates. We also establish precise
limitations of finite congruence covers, skipped-level intervals, translated
prime patterns, and phase-blind Fourier estimates. The transcendence of the
classical Mills minimum and of the residual constant with exponents
$2\cdot3^k-1$ is not proved. Conditional reductions, proved restricted
families, finite computations, and unresolved arithmetic inputs are kept
separate throughout.
\end{abstract}
\tableofcontents
\clearpage
\section{Introduction and the structure of the argument}\label{sec:overview}
For a specified positive exponent schedule $C=(C_k)$, write
\[
 \mathcal W(C)=\{A>1:\floor{A^{C_k}}\text{ is prime for every specified }k\},
 \qquad \xi(C)=\min\mathcal W(C),
\]
whenever the minimum exists. The starting index is part of the schedule.
In particular, removing a finite prefix is not a harmless notational change.
Our central distinction is between a minimum of this entire admissible set,
a least completion of one feasible prime prefix, and an arbitrary cubic
Pisot number satisfying some of the resulting arithmetic constraints.
These three objects are not interchangeable.

A cubic Pisot number has one real conjugate $\beta>1$ and two conjugates
inside the unit disk. Integer traces then approximate its large powers.
When the small-conjugate contribution has the appropriate sign, the trace
is exactly the floor. The exponent schedules $3^k$ and $2\cdot3^k-1$
therefore turn a prime-floor question into a highly sparse recurrence
question. For a single auxiliary prime the resulting trace orbit is
finite, so a zero on a recurrent phase yields infinitely many composite
floors. The difficulty is not this implication, but ensuring a suitable
prime and a suitable recurrent phase for every fixed candidate.

\subsection{The main conclusions}
The following roadmap separates the results that close restricted cases
from the statements still needed for the principal problems.
\begin{enumerate}
\item For positive affine shifts, the exceptional unit classes left by
an earlier fixed-point calculation are removed by explicit divisibility
certificates and a bound on the actual infinite minimum
(Sections~\ref{src:P01} and~\ref{src:R07}). Exact scaling then yields
further equivalent schedules without treating them as independent results.
\item For negative shifts, the integral fixed-point method, the multiplier
thirteen argument, and the additional odd multiplier certificates give
restricted-family results using the structural reductions cited below (Sections~\ref{src:N01}, \ref{src:R02}, and~\ref{src:R04}). They do
not eliminate every candidate for $\xi(2\cdot3^k-1)$.
\item On each fixed coefficient line, recurrent divisors exist for a set
of coefficients of density one. Prime-relative estimates and the explicit
logarithmic exceptional-set saving are proved in
Sections~\ref{src:R08}--\ref{src:R11}. Density one does not mean that the
exceptional set is empty, nor that a prescribed minimum is covered.
\item Thin prime-power families and prime-leading examples require
pointwise certificates. Sections~\ref{src:R12}--\ref{src:R16} give exact
families, projective return tests, and limitations of bounded witness
strategies. A finite collection of successful residues is never promoted
to a universal covering theorem.
\item Hypothetical classical Mills candidates satisfy stronger order and
three-adic constraints in Sections~\ref{src:R16}--\ref{src:R17}.
Section~\ref{src:R18} identifies what three analytic routes do supply and
what each still lacks. No contradiction for every candidate is obtained.
\end{enumerate}

\subsection{Strengthenings obtained by combining the arguments}
The combined treatment yields stronger conclusions than a sequential
reading of the source papers alone.
\begin{itemize}
\item Scalar returns in cubic finite algebras enlarge the usable
Frobenius residues from one to four classes modulo $9$. This improves
the fixed-line logarithmic exceptional exponents from $1/6$ to $2/3$
for integers and from $7/6$ to $5/3$ for prime coefficients
(Theorem~\ref{R11:syn:projective-density-rate}).
\item The depth argument extends to both residue signs and to every
irreducible totally real cubic satisfying the stated congruences,
without assuming it is Pisot. A single cyclotomic argument then
controls both output signs
(Theorem~\ref{R17:thm:new-two-sign-depth} and
Corollary~\ref{R17:cor:new-two-sign-order}). Multiplicative independence
also makes the conditional shadow bounds strict.
\item The prime-$31$ family extends from one negative clock to every
even multiplier not divisible by $5$
(Corollary~\ref{R06:audit06:even-multipliers}). A different small-prime
witness synchronizes both clocks for the prime-seed outlier and an
entire coefficient progression
(Theorem~\ref{R13:syn:outlier-synchronization}).
\end{itemize}
These are strengthened restricted-family and necessary-condition results.
They do not change the unresolved status of the two main constants.

\subsection{Proof dependencies and interpretation}
The analytic-to-Pisot reductions use the cited work of Saito and the
explicitly identified structural lemmas in the source corpus. The
original source references are retained so that these dependencies are
visible rather than hidden inside a new title. Exact certificates establish
finite algebraic assertions; their periodic-return proofs explain when
those assertions imply infinitely many composite floors. A finite search
alone does not establish an infinite prime chain. Likewise, a theorem for
an auxiliary sequence of allowed integers is not a counterexample for the
actual primes.

This article combines the mathematical content of twenty source packages.
The sections are arranged by their role in the argument rather than by
delivery order. Repeated abstracts and title matter have been removed;
proofs, relevant finite tables, and verification procedures are retained. The density proof is replaced by its stronger projective version, and the positive-only depth proof by a two-sign totally real version; neither earlier proof is duplicated as a separate result.
The common convention is that a local symbol has the meaning assigned in its section. Some local results are intentionally stated under their own notation to
avoid silently changing hypotheses. A source map at the end records their
provenance. Earlier ``remaining'' cases are interpreted in light of the
later sections and the common final status, not as contradictory claims
about the final theorem list. No assertion of external peer review or
literature priority is made.

\clearpage\part{Affine schedules and their least prime representing constants}

\section{Scaling consequences local minima and normalized tails}\label{src:INIT}
% Integration-ready fragment. Requires amsmath, amssymb, amsthm and theorem,
% lemma, proposition, corollary environments. No private macros are used.
% Citation keys used: Saito, BHP, K031, K064, RamachandraSrinivasan1983.
% The first ZIP's intermediate 28/403 theorem and 46.2-trillion cutoff are
% superseded by positive01 and negative01 and are intentionally not restated.

\subsection{Scaling, prime-output consequences, and prefix minima}
\label{sec:init-consequences}

The results in this section consolidate the first archive's consequence
and proof-transfer notes. They distinguish an exact change of variables,
a reuse of the local-minimum proof, and a continuity statement that does
not preserve transcendence in its limit. All prescribed starting indices
are part of the definitions.
For a positive clock $C=(C_k)_{k\ge1}$, write
\[
 \mathcal W(C)=\{A>1:\lfloor A^{C_k}\rfloor\text{ is prime for all }k\ge1\}.
\]
Its least element, when it exists, is denoted by $\xi(C)$.

\begin{proposition}[Exact scaling]\label{init:scaling}
For any real $d>0$,
\[
 \mathcal W(dC)=\{A>1:A^d\in\mathcal W(C)\},\qquad
 \xi(dC)=\xi(C)^{1/d},
\]
whenever one, and hence both, of these minima exists.
For rational $d>0$, algebraicity of the two minima is equivalent.
In particular, if $\alpha=\xi((3^k-1)_{k\ge1})$, then
\[
 \xi\bigl((m(3^k-1)/2)_{k\ge1}\bigr)=\alpha^{2/m}
 \qquad(m\ge1).
\]
These are all positive rational scalings of $3^k-1$ that are integral
at every prescribed index.
\end{proposition}
\begin{proof}
The map $A\mapsto A^d$ is an increasing bijection on $(1,\infty)$,
and $(A^d)^{C_k}=A^{dC_k}$ term by term. A positive rational power
of a positive real number is algebraic if and only if that number is
algebraic. Finally, $\gcd_{k\ge1}(3^k-1)=2$: necessity follows from
the first exponent $2$, and every exponent is even.
\end{proof}

Thus parameter hypotheses should be stated, where possible, for the
primitive pair $(R,S)=(r/d,s/d)$, $d=\gcd(r,s)$. For example,
$3\nmid S$ is equivalent to $v_3(s)\le v_3(r)$; it does not require
$3\nmid s$. This is a strengthening of the literal unnormalized
statements, but not a new independent transcendence mechanism.

\subsubsection{A uniform consequence for all nonzero affine shifts}

The next theorem extends the first archive's output consequences from
$3^k-1$ to every transcendental affine-clock base. Primality is not needed
for its proof; it applies in particular to each transcendental minimum
proved in the subsequent sections.

\begin{theorem}[Laurent limits and escape from fixed curves]
\label{init:affine-outputs}
Let $b\ge2$ and $r\ge1$ be integers, let $\sigma\in\mathbb Z\setminus\{0\}$,
and suppose $C_k=rb^k+\sigma>0$ for $k\ge1$. Fix a transcendental
real number $A>1$ and put $p_k=\lfloor A^{C_k}\rfloor$.
For every fixed $h\ge1$,
\[
 \frac{p_{k+h}}{p_k^{b^h}}\longrightarrow
 L_h:=A^{\sigma(1-b^h)}.
\]
For any fixed integers $a_0,\ldots,a_m$, put
$U=\sum_{j=0}^m a_jb^j$ and $V=\sum_{j=0}^m a_j$. Then
\[
 \lim_{k\to\infty}\prod_{j=0}^m p_{k+j}^{a_j}
 =\begin{cases}
 +\infty,&U>0,\\
 0,&U<0,\\
 A^{\sigma V},&U=0.
 \end{cases}
\]
The finite nonzero limits, as all fixed integral Laurent monomials are
allowed, are exactly
\[
 \{A^{\sigma(b-1)n}:n\in\mathbb Z\}.
\]
Only $1$ in this set is algebraic. Moreover, for every nonzero
$F\in\overline{\mathbb Q}[X,Y]$ and every fixed $h\ge1$,
\[
 F(p_k,p_{k+h})\ne0\qquad\text{for all sufficiently large }k.
\]
Consequently no fixed polynomial recurrence with algebraic coefficients
can hold for these pairs on an infinite subsequence.
\end{theorem}
\begin{proof}
For every fixed finite window,
$p_{k+j}=A^{rb^{k+j}+\sigma}(1+o(1))$, and hence
\[
 \prod_{j=0}^m p_{k+j}^{a_j}
 =A^{rb^kU+\sigma V}(1+o(1)).
\]
This proves the three alternatives, including negative exponents.
If $U=0$, reduction modulo $b-1$ gives $V\equiv0\pmod{b-1}$.
Conversely $(a_0,a_1)=(bn,-n)$ gives $U=0$ and $V=(b-1)n$.
A nonzero integral power of a transcendental positive real number is
transcendental, proving the arithmetic assertion.
The ratio formula is the special case $(a_0,a_h)=(-b^h,1)$.

Write $F=\sum c_{ij}X^iY^j$, put $d=b^h$, and let
$D=\max\{i+dj:c_{ij}\ne0\}$. Since $p_k\to\infty$,
\[
 \frac{F(p_k,p_{k+h})}{p_k^D}
 \longrightarrow\sum_{i+dj=D}c_{ij}L_h^j.
\]
The right side is a nonzero polynomial evaluated at the transcendental
number $L_h$: for each $j$, the equation $i+dj=D$ determines at most
one $i$, so distinct terms do not cancel as formal powers of $L_h$.
The limit is therefore nonzero.
\end{proof}

For the ternary clock this yields exactly the limit group
$\{A^{2\sigma n}:n\in\mathbb Z\}$. For $\alpha$ and
$p_k=\lfloor\alpha^{3^k-1}\rfloor$, in particular,
\[
 \frac{p_{k+1}}{p_k^3}\to\alpha^2,\qquad
 \frac{p_{k+1}-p_k^3}{p_k^3}\to\alpha^2-1,
\]
and the exact telescoping product is
\[
 \alpha=p_1^{1/3}\prod_{k\ge1}
       \left(\frac{p_{k+1}}{p_k^3}\right)^{1/3^{k+1}}.
\]
Its finite products equal $p_n^{1/3^n}$.
The curve-escape conclusion means that every infinite subset of the
pair sequence is Zariski dense over $\overline{\mathbb Q}$; it does
not say that individual integer pairs are algebraically independent.
Nor does the two-coordinate proof automatically extend to arbitrary
polynomials in three or more output coordinates.

\subsubsection{Every feasible-prefix minimum for the clock $3^k-1$}

A \emph{local left minimum} of $\mathcal W(C)$ is a member $A$ for
which $(A-\eta,A)\cap\mathcal W(C)=\varnothing$ for some $\eta>0$.
The following result uses more than transcendence of the global minimum.
Its explicit inherited inputs are the local Type B reduction of
\cite[section ``Local Type B and factorial perturbations'']{K031},
and the negative-fixed-point norm exclusion of \cite{K064}.
The localization in the former is the Saito prime-chain replacement
argument with every fixed finite prefix retained. These are unpublished
corpus results, not additional external hypotheses about primes.

\begin{theorem}[Local and prefix minima]\label{init:prefix-minima}
Every local left minimum of
$\mathcal W((3^k-1)_{k\ge1})$ is transcendental.
These points are exactly the least completions of feasible finite prime
prefixes. They form a countable dense subset of the admissible set.
\end{theorem}
\begin{proof}
Put $C_k=3^k-1$. For each $m$, a later proper multiple of $C_m$
occurs among the $C_k$, by taking $k-m$ a positive multiple of
$\operatorname{ord}_{C_m}(3)$. Thus every admissible $A$ satisfies
$A^{C_m}\notin\mathbb N$: otherwise the prime at level $m$ would
become a composite proper prime power at that later level.

Suppose an algebraic admissible $A$ were a local left minimum.
The stated local reduction applies because $C_{k+1}/C_k>3$,
these ratios tend to $3$, and the recurring divisibility just verified
preserves the initial endpoints. It gives a cubic Pisot number
$\beta=A^g$, $g\mid2$, and, for all sufficiently large $k$,
\[
 n_k=(3^k-1)/g,\qquad
 T_k=\operatorname{Tr}(\beta^{n_k})=\lfloor A^{C_k}\rfloor\text{ prime},
 \qquad 0\le\beta^{n_k}-T_k\le T_k^{-17/40}.
\]
Here the degree bound is $3\le\deg\beta\le1+40/17<4$.
The norm-exclusion lemma of \cite{K064} gives
$3\nmid\operatorname{N}(\beta)$.

For every $k$ divisible by $12$, $208\mid3^k-1$, so $104\mid n_k$.
Each residue of a conjugate of $\beta$ in an algebraic closure of
$\mathbb F_3$ is nonzero and has degree at most three over
$\mathbb F_3$. Its multiplicative order therefore divides one of
$2,8,26$, hence divides $104$. Reducing the trace at these exponents
modulo a prime above $3$ gives $T_k\equiv1+1+1=0\pmod3$.
The traces tend to infinity, contradicting eventual primality.

For the remaining topological assertion, a feasible prefix defines a
nonempty cylinder $I=[a,b)$, the finite intersection of its prime cells.
Let $A=\inf(\mathcal W(C)\cap I)$. A sequence of completions tending
to $A$ from the right, and right continuity of each fixed floor, show
that $A\in\mathcal W(C)$. A chosen completion gives $A<b$.
Also $A\ne a$, since $a$ is a lower prime-cell endpoint and integral
powers were excluded. Thus $(a,A)\cap\mathcal W(C)=\varnothing$,
so this least completion is a local left minimum.
Conversely, a sufficiently long prefix cylinder about any local left
minimum has width smaller than its left gap; its least completion is
therefore that point. There are countably many finite prefixes, and
successive prefix cylinders at any admissible point have widths tending
to zero. Their minima converge to that point, proving density.
\end{proof}

This theorem and Proposition~\ref{init:scaling} also transfer local
left minima and feasible-prefix minima throughout the rational scaling
class of $3^k-1$. They make no claim that every admissible base is
transcendental.

\subsection{Continuity at the unshifted clock, without arithmetic transfer}
\label{sec:init-continuity}

Extend $\xi$ to positive real exponent clocks and put
\[
 F(t)=\xi((3^k-t)_{k\ge1})\quad(-1\le t\le1),\qquad
 M=F(0).
\]
All these minima exist by Saito's existence theorem
\cite[Theorem 1.3]{Saito}, since the consecutive ratios are at least
$5/2$. The proof below needs only a uniform sufficiently-large
prime-gap threshold, not its numerical value.

\begin{lemma}[Uniform large-prime continuation]\label{init:uniform-seed}
There is a prime threshold $P_0$ such that, for every positive clock
with consecutive ratios in $[5/2,4]$, every prime $p\ge P_0$ at any
one level admits an infinite prime continuation inside its closed cell,
whose subsequent closed cells lie strictly inside their parents.
\end{lemma}
\begin{proof}
The Baker--Harman--Pintz prime-interval theorem \cite{BHP} supplies
constants $K,X_0>0$ such that every real $x\ge X_0$ has a prime
$q$ with $x<q\le x+Kx^{21/40}$ (enlarging $K,X_0$ if necessary).
Let $c\in[5/2,4]$. The mean value theorem gives
\[
 (p+1)^c-p^c\ge cp^{c-1},\qquad
 \frac{cp^{c-1}}{p^{21c/40}}
 =cp^{19c/40-1}\ge\frac52p^{3/16}.
\]
Thus a single sufficiently large $P_0$ ensures, uniformly in $c$,
$p^c+Kp^{21c/40}+1<(p+1)^c$. The prime just supplied satisfies
$p^c<q$ and $q+1<(p+1)^c$. Iteration gives nested compact root
cells, and any point in their intersection has the required floors:
the strict upper containment at the next level removes every open
right-endpoint ambiguity. The same threshold remains valid because
the primes strictly increase.
\end{proof}

\begin{theorem}[Two-sided continuity and normalized tails]
\label{init:continuity}
One has $F(t)\to M$ as $t\to0$. Consequently
\[
 \xi((3^{k+h}-1)_{k\ge1})^{3^h}\longrightarrow M,
 \qquad
 \xi((3^{k+h}+1)_{k\ge1})^{3^h}\longrightarrow M
 \qquad(h\to\infty).
\]
The corresponding unnormalized minima tend to $1$.
For each fixed finite set of output indices, the two normalized
families eventually reproduce the prime outputs of $M$ at those indices,
when evaluated on their own perturbed clocks.
\end{theorem}
\begin{proof}
The ratios $(3^{k+1}-t)/(3^k-t)$ lie in $[5/2,4]$ for $|t|\le1$.
A fixed sufficiently large first prime and
Lemma~\ref{init:uniform-seed} give a uniform upper bound on $F(t)$;
also $F(t)\ge2^{1/4}>1$.

Take a sequence $t_n\to0$ for which $F(t_n)\to A$.
If $A^{3^j}=m\in\mathbb N$ for some $j$, then $m\ge2$ and
\[
 F(t_n)^{3^{j+2}-t_n}\longrightarrow m^9.
\]
Its floor, for all sufficiently large $n$, is either $m^9-1$ or $m^9$.
Both are composite; indeed
$m^9-1=(m^3-1)(m^6+m^3+1)$.
This contradicts admissibility. Therefore each $A^{3^j}$ is
nonintegral, and continuity at each fixed index gives
$\lfloor A^{3^j}\rfloor$ prime. Thus $A\ge M$, proving
$\liminf_{t\to0}F(t)\ge M$.

Write $q_j=\lfloor M^{3^j}\rfloor$. No $M^{3^j}$ is integral,
because otherwise the next required power would be composite.
Moreover the closed classical prime cells strictly nest:
\[
 q_{j+1}>q_j^3,\qquad q_{j+1}+1<(q_j+1)^3.
\]
The lower equality is composite; the only possible upper equality
would give $q_{j+1}=(q_j+1)^3-1$, also composite.
Given $\varepsilon>0$, choose $n$ with $q_n\ge P_0$ and
$(q_n+1)^{1/3^n}<M+\varepsilon$.
Its closed cell lies strictly inside every earlier prime cell.
By continuity of finitely many endpoints, this containment persists
with exponents $3^j-t$ whenever $|t|$ is sufficiently small.
Continue from the same prime $q_n$ by
Lemma~\ref{init:uniform-seed}. The resulting admissible base
preserves the earlier prefix, so
\[
 F(t)\le(q_n+1)^{1/(3^n-t)}.
\]
Taking the upper limit and then $\varepsilon\downarrow0$ proves
continuity.

Exact scaling gives
$\xi(3^{k+h}-1)^{3^h}=F(3^{-h})$ and
$\xi(3^{k+h}+1)^{3^h}=F(-3^{-h})$.
These converge to $M>1$, and taking $3^h$-th roots gives the raw
limits $1$. Finally, $F(t)^{3^j-t}\to M^{3^j}$ at every fixed $j$;
nonintegrality of the latter makes the finitely many floors stable.
\end{proof}

The ninth power in the lower-limit proof is essential to this simple
argument: $2^3-1=7$ is prime. The displayed limits do not give an
identity between a full minimum and a tail minimum. Although the
shifted approximants here are transcendental by the shifted-clock
results, convergence of transcendental numbers does not determine the
arithmetic nature of $M$.

\subsection{Classical exponential-value consequences}
\label{sec:init-exponentials}

The following bound is the theorem of Ramachandra and Srinivasan
\cite{RamachandraSrinivasan1983}, included with its short
six-exponentials proof. Its use here is an application of an existing
theorem, not a new transcendence theorem for individual exponential
values.

\begin{theorem}[Sparse algebraic exponential values]
\label{init:six-exponentials}
Let $A>1$ be transcendental and fix $z\in\mathbb C\setminus\{0\}$.
In any interval of $N\ge1$ consecutive positive integer indices,
at most
\[
 \frac{1+\sqrt{16N-15}}2
\]
of the numbers $\exp(zA^n)$ are algebraic.
For real $A$, the same exceptional-index bound holds simultaneously
for algebraicity of $\sin(A^n)$ and $\cos(A^n)$.
\end{theorem}
\begin{proof}
Let $E=\{n\ge1:\exp(zA^n)\in\overline{\mathbb Q}\}$.
For a fixed positive difference $d$, there are at most two indices
$n$ with $n,n+d\in E$. Otherwise three such distinct indices
$n_1,n_2,n_3$ give three $\mathbb Q$-linearly independent numbers
$zA^{n_i}$ and two $\mathbb Q$-linearly independent numbers $1,A^d$.
Their six exponentials would all be algebraic, contrary to the
six-exponentials theorem. The independence follows directly from
transcendence of $A$.
If an interval of length $N$ contains $m$ elements of $E$, its
$\binom m2$ pairwise positive differences have at most $2(N-1)$
occurrences in total. Solving this quadratic inequality gives the bound.
For $z=i$, put $u=e^{ix}$ with $x=A^n$ real.
Algebraicity of $u$, $\cos x$, and $\sin x$ are equivalent:
$u$ satisfies $u^2-2\cos(x)u+1=0$ and
$u^2-2i\sin(x)u-1=0$, while the reverse implications follow from
$\cos x=(u+u^{-1})/2$ and $\sin x=(u-u^{-1})/(2i)$.
\end{proof}

For every fixed real $b>0$, $b\ne1$, one may take $z=\log b$;
no algebraicity assumption on $b$ is needed. A separate, weaker-input
consequence is that for any irrational real $A$, at most two distinct
primes $p$ can have $p^A$ algebraic. Indeed, the logarithms of three
distinct primes are $\mathbb Q$-linearly independent, and applying
six exponentials with $1,A$ would contradict algebraicity of all six
values. More generally the multiplicative group
\[
 \{b\in\overline{\mathbb Q}\cap\mathbb R_{>0}:
                      b^A\in\overline{\mathbb Q}\}
\]
has $\mathbb Q$-rank at most two; it need not have only two elements.
None of these assertions identifies $e^\alpha$, $\log\alpha$, or
$\alpha^{\sqrt2}$ individually, or gives algebraic independence of
an infinite family of exponential values.

\section{Positive shifts and elimination of the exceptional unit candidates}\label{src:P01}
The next argument completes the positive unit elimination that was only partial in the initial notes. Its finite-ring certificates must be combined with a bound for the actual infinite minimum; divisibility tests alone do not finish the argument.
\medskip
\begingroup




\subsection{Statement and dependencies}

For a positive sequence $C=(C_k)_{k\geq1}$ define
\[
 \mathcal W(C)=\{A>1:\floor{A^{C_k}}\text{ is prime for every }k\geq1\}.
\]
When this set has a least element, denote it by $\xi(C)$.
Saito's existence theorem \cite[Theorem 1.3]{Saito} applies if
$C_{k+1}/C_k\geq2$. For $C_k=r3^k+s$ this condition follows from
$s\leq3r$.

\begin{theorem}\label{P01:thm:main}
Let $r,s\geq1$ be integers satisfying $3\nmid s$ and $s\leq3r$.
Then
\[
             \xi\bigl((r3^k+s)_{k\geq1}\bigr)
\]
is transcendental.
\end{theorem}

Theorem~\ref{P01:thm:main} removes the exceptional divisibility classes in the
positive-shift theorem of \cite{K064}. It concerns the actual least element
of $\mathcal W(C)$, with the schedule starting at $k=1$.
We use Saito's algebraic-minimum reduction
\cite[Proposition 3.1]{Saito}, based on the prime-interval exponent
$21/40$ of Baker, Harman and Pintz \cite{BHP}, and the theoretical
ternary trace results of \cite{K064}. The latter is an unpublished
research manuscript; the present argument is not an independent replacement
for all of its theoretical lemmas. We isolate exactly the assertions used
below. The finite candidate computation is also reproduced by a separate
unit-only enumeration, rather than relying on the original enumeration of
all cubic Pisot numbers below $9$.

The additional argument has two parts. Finite modular certificates make
any surviving normalized coefficient $R=r/s$ very large. A separate
infinite construction makes the \emph{global} minimum small for all such
$R$. Neither part alone removes every $R$.

\subsection{Reduction to three algebraic candidates}\label{P01:sec:reduction}

A Pisot number is a real algebraic integer $\beta>1$ whose other
conjugates have absolute value less than $1$.
All field traces below are over $\Q$.

\begin{proposition}[Inherited algebraic-minimum and trace inputs]
\label{P01:prop:inputs}
Under the assumptions of Theorem~\ref{P01:thm:main}, suppose that
$\xi=\xi(r3^k+s)$ is algebraic. There is an integer $g\mid2s$ such that
$\beta=\xi^g$ is a cubic Pisot number and
\[
 n_k=\frac{r3^k+s}{g}\in\N
\]
for every sufficiently large $k$. Write $x_2,x_3$ for its other
conjugates and $\lambda=s/g$. Then:
\begin{enumerate}[label=\textup{(\roman*)}]
\item $T_k=\Tr(\beta^{n_k})=\floor{\xi^{r3^k+s}}$ is prime for all
sufficiently large $k$.
\item For every fixed $\varepsilon>0$,
\[
 |x_2^{n_k}+x_3^{n_k}|
     =O\bigl(\beta^{-(17/40-\varepsilon)n_k}\bigr).
\]
\item $\beta^{\lambda}<3$, hence $\beta<9$.
\item If $\beta$ is a unit, then the possible pairs $(\lambda,\beta)$
and the necessary eventual exponent residues are precisely among the
five rows of Table~\ref{P01:tab:unit}.
\end{enumerate}
\end{proposition}

Items (i) and (ii) follow from Saito's Proposition 3.1, after excluding
integer powers and the quadratic alternative as in the Mills-shadow
corollary of \cite{K064}. For clarity, $3\nmid s$ makes every $C_m$
coprime to $3$, and $C_m\mid C_{m+h}$ whenever
$3^h\equiv1\pmod{C_m}$. An integer power $\xi^{C_m}$ would therefore
give composite outputs at infinitely many later indices. Likewise, a
putative sequence of prime integers $C_k/g$ would repeat one of its
prime divisors along an affine return, contrary to its strict growth.
The surviving degree is $3$, since
$1+(3(1-21/40)-1)^{-1}<4$.
Finally, $g$ divides every sufficiently late $C_k$ and therefore divides
$C_{k+1}-3C_k=-2s$.

Items (iii) and (iv) are the positive-fixed-point consequence and the
necessary finite table of \cite{K064}. Its affine trace-return theorem
forces each relevant twisted $3$-adic limit to equal $\pm1$.
Transporting that algebraic identity to $\C$ gives
$\beta^\lambda<1+1+1=3$. The exact unit-only check of (iv) is
described in Section~\ref{P01:sec:enumeration}.
All Saito numbering here refers to arXiv version 3; in particular,
the reduction is Proposition \emph{3.1}, not 3.2.

\begin{lemma}[Unit reduction]\label{P01:lem:unit}
The number $\beta$ in Proposition~\ref{P01:prop:inputs} has
$|\Nm(\beta)|=1$.
\end{lemma}
\begin{proof}
Take $\varepsilon=1/40$ in Proposition~\ref{P01:prop:inputs}(ii), so that
the approximation exponent is $\kappa=2/5$.
The approximation-magnitude lemma of \cite{K064} gives
\[
 |\Nm(\beta)|\leq\beta^{1-2\kappa}
       =\beta^{1/5}<9^{1/5}<2.
\]
The norm is a nonzero integer, proving the claim.
This use of the magnitude lemma is valid for both totally real and
nonreal cubic Pisot numbers. In the nonreal case its proof uses the
polynomial lower bound for a nonzero linear form in logarithms
\cite{BW}; it does not assume a uniform lower bound on an oscillating
cosine.
\end{proof}

Let $\delta$ and $\tau$ denote the unique real roots greater than $1$ of
\[
 f_\delta(X)=X^3-X^2-1,
 \qquad f_\tau(X)=X^3-X^2-X-1,
\]
respectively. Their relevant powers have minimal polynomials
\[
 f_{\delta^2}(X)=X^3-X^2-2X-1,
 \qquad f_{\tau^2}(X)=X^3-3X^2-X-1.
\]

\begin{table}[htbp]
\centering
\small
\begin{tabular}{cclcl}
\toprule
$\lambda$ & $\beta$ & Minimal polynomial of $\beta$
 & $\xi^s$ & Necessary exponent residue \\
\midrule
$1$ & $\delta$ & $X^3-X^2-1$ & $\delta$ & $n_k\equiv1\pmod8$\\
$1$ & $\tau$ & $X^3-X^2-X-1$ & $\tau$ & $n_k\equiv1\pmod{13}$\\
$1$ & $\delta^2$ & $X^3-X^2-2X-1$ & $\delta^2$ & $n_k\equiv1\pmod4$\\
$1/2$ & $\tau^2$ & $X^3-3X^2-X-1$ & $\tau$ & $n_k\equiv7\pmod{13}$\\
$2$ & $\delta$ & $X^3-X^2-1$ & $\delta^2$ & $n_k\equiv2\pmod8$\\
\bottomrule
\end{tabular}
\caption{The reduced unit table. Each listed residue is a necessary
condition, not an assertion that the corresponding trace is prime.}
\label{P01:tab:unit}
\end{table}

\begin{proposition}[Normalized candidate reduction]\label{P01:prop:normalized}
If $\xi(r3^k+s)$ is algebraic under the hypotheses of
Theorem~\ref{P01:thm:main}, then $s\mid r$. Put
\[
 R=r/s,\qquad \eta=\xi(r3^k+s)^s.
\]
Then
\[
 \eta=\xi(R3^k+1)\in\{\delta,\delta^2,\tau\},
\]
and $\Tr_{\Q(\eta)/\Q}(\eta^{R3^k+1})$ is prime for every sufficiently
large $k$. Moreover, $8\mid R$ for the $\delta$ branch,
$4\mid R$ for the $\delta^2$ branch, and $13\mid R$ for the $\tau$ branch.
\end{proposition}
\begin{proof}
For $\lambda=1$ we have $g=s$, so the integrality of $n_k$ and
$3\nmid s$ give $s\mid r$. Table~\ref{P01:tab:unit} then gives
$8\mid R$, $13\mid R$, or $4\mid R$, respectively.
For $\lambda=1/2$ we have $g=2s$, so again $s\mid r$, and
$n_k=(R3^k+1)/2\equiv7\pmod{13}$ implies $13\mid R$.
For $\lambda=2$ we have $g=s/2$; the last residue condition says
\[
       \frac{2r3^k}{s}\equiv0\pmod8.
\]
Consequently $4s\mid r3^k$, and since $3\nmid4s$, we get
$4s\mid r$. These arguments establish the stated integrality and
divisibility conditions in every row.

The possible values of $\eta=\beta^\lambda$ are displayed in the table.
A Pisot number and any positive integer power generate the same field:
an embedding fixing that power must send the Pisot number to a conjugate
of the same absolute value, and only the dominant conjugate has that
absolute value. Thus the trace field is unchanged in the $\tau^2$ and
$\delta^2$ rows, and
\[
 \Tr_{\Q(\beta)/\Q}(\beta^{n_k})
       =\Tr_{\Q(\eta)/\Q}(\eta^{R3^k+1}).
\]
Lastly, for any positive integer $d$,
\[
 \mathcal W(dC)=\{A>1:A^d\in\mathcal W(C)\},
 \qquad \xi(dC)^d=\xi(C).
\]
Applying this exact minimum-preserving identity with $d=s$ proves
$\eta=\xi(R3^k+1)$, not just $\eta\in\mathcal W(R3^k+1)$.
\end{proof}

\subsection{Exact large divisors from finite rings}\label{P01:sec:finite}

For a cubic unit with minimal polynomial
$f(X)=X^3-aX^2+bX-1$, put $T_n=\Tr(\eta^n)$.
Newton's identities give
\begin{equation}\label{P01:eq:newton}
 T_0=3,\quad T_1=a,\quad T_2=a^2-2b,\qquad
 T_n=aT_{n-1}-bT_{n-2}+T_{n-3}\quad(n\geq3).
\end{equation}

\begin{lemma}[Finite-ring obstruction]\label{P01:lem:finite}
Let $\eta$ be a cubic Pisot unit, and let $p,t$ be primes with
$t\neq3$. Suppose
\begin{enumerate}[label=\textup{(\alph*)}]
\item $\ord_t(3)=t-1$;
\item $X^t=1$ in $\F_p[X]/(f)$;
\item for some $n\not\equiv1\pmod t$, one has $T_n\equiv0\pmod p$.
\end{enumerate}
If $T_{R3^k+1}$ is prime for every sufficiently large $k$, where
$R\in\N$, then $t\mid R$.
\end{lemma}
\begin{proof}
If $t\nmid R$, then $R3^k$ runs periodically through every nonzero
residue modulo $t$, so that $R3^k+1\equiv n\pmod t$ infinitely often.
Condition (b) makes the power traces periodic modulo $p$ with period
dividing $t$. Hence $p\mid T_{R3^k+1}$ for infinitely many $k$.
Since $\eta$ is Pisot,
$T_m=\eta^m+o(1)\longrightarrow+\infty$ as $m\longrightarrow\infty$.
These divisible terms eventually exceed $p$, contradicting their
primality.
Neither irreducibility of $f$ modulo $p$ nor exact order $t$ of $X$
in the quotient ring is required.
\end{proof}

\begin{table}[htbp]
\centering
\begin{tabular}{rrrr}
\toprule
$p$ & $t$ & $n$ for $\delta$ & $n$ for $\delta^2$\\
\midrule
2 & 7 & 3 & 3\\
5 & 31 & 4 & 2\\
41 & 1723 & 42 & 21\\
101 & 10303 & 58 & 29\\
1217 & 1482307 & 1218 & 609\\
1373 & 1886503 & 86 & 43\\
2729 & 7450171 & 2730 & 1365\\
\bottomrule
\end{tabular}
\caption{Certificates for $\delta$ and $\delta^2$.}
\label{P01:tab:delta}
\end{table}

\begin{table}[htbp]
\centering
\begin{tabular}{rrr}
\toprule
$p$ & $t$ & $n$ for $\tau$\\
\midrule
5 & 31 & 11\\
1193 & 1424443 & 1620\\
1373 & 1886503 & 78\\
1973 & 3894703 & 331\\
\bottomrule
\end{tabular}
\caption{Certificates for $\tau$.}
\label{P01:tab:tau}
\end{table}

Every row of Tables~\ref{P01:tab:delta} and~\ref{P01:tab:tau} satisfies
Lemma~\ref{P01:lem:finite}; in each row $t=p^2+p+1$.
Section~\ref{P01:app:certificate} gives a complete exact verifier.
It checks primality by trial division, the order of $3$ using every
prime divisor of $t-1$, quotient-ring powering, and the zero trace
both by powering and by the independent recurrence~\eqref{P01:eq:newton}.
A separate companion-matrix implementation reproduces these checks.

\begin{proposition}\label{P01:prop:divisors}
If $\Tr(\eta^{R3^k+1})$ is eventually prime for
$\eta\in\{\delta,\delta^2,\tau\}$ and $R\in\N$, then
\begin{align*}
 80254725161999161502908681243&\mid R
       &&\text{if }\eta\in\{\delta,\delta^2\},\\
 324443153856491628397&\mid R
       &&\text{if }\eta=\tau.
\end{align*}
In particular, $R>10^{20}$ in all three cases.
\end{proposition}
\begin{proof}
Apply Lemma~\ref{P01:lem:finite} to every row for the given candidate.
Within each table the moduli $t$ are distinct primes, so their product
divides $R$. Their exact products are the displayed integers.
\end{proof}

These certificates assert necessary divisibility, not compositeness for
every $R$. In particular, multiples of all listed moduli survive this
finite sieve. The next section supplies the separate global bound that
eliminates this escape.

\subsection{An explicit prime interval}\label{P01:sec:gap}

We give a deliberately conservative consequence of the explicit
estimates in Cully-Hugill \cite{CH}. The large threshold keeps the
verification elementary and leaves ample room below the divisors in
Proposition~\ref{P01:prop:divisors}.

\begin{lemma}\label{P01:lem:gap}
For every real $u\geq X_0:=\exp(10^{18})$, there is a prime $p$ with
\[
                  u<p\leq u+u^{2/3}.
\]
\end{lemma}
\begin{proof}
First let $x$ be a half-odd-integer at least $X_0$, and put
\[
 y=\log x,\qquad \alpha=\frac{17}{50},\qquad
 T=x^\alpha,\qquad h=\floor{x^{2/3}/2},\qquad
 w=\log\frac{x}{T^{8/3}(\log T)^2}.
\]
Thus $x+h$ is also a half-odd-integer, $50<T<x$, and
$x^{2/3}/3\leq h\leq x^{2/3}/2$.

Write $\theta(x)=\sum_{p\leq x}\log p$.
Theorem 2, Table 2, and equations (2), (9), (15), and (19)--(21) of
\cite{CH}, with $\sigma_1=3/5$, imply
\begin{equation}\label{P01:eq:theta}
 \theta(x+h)-\theta(x)
 >h\left(1-F(x)-\frac{K G(x,h)}{Th}+\frac{E(x)}h\right),
\end{equation}
where $K=3.4747$,
\begin{align*}
 G(x,h)&=(x+h)\log(x+h)\log\log(x+h)+x\log x\log\log x,\\
 E(x)&=0.999\sqrt{x}+x^{1/3}
                   -a_1\sqrt{x+h}-a_2(x+h)^{1/3},
\end{align*}
with $a_1,a_2<2$. Taking
$C_1=17.418$, $C_2=5.272$, and
\[
 v=\frac{1}{57.54(\alpha y)^{2/3}(\log(\alpha y))^{1/3}},
\]
one may use the upper estimate
\begin{equation}\label{P01:eq:F}
 F(x)\leq\frac{\alpha}{\pi}y e^{-(1-\alpha-3/5)y}
   +\frac{2C_1\alpha^3y^4}{w}e^{-vw}
   +2C_2\alpha^2y^2e^{-vy}.
\end{equation}
Here the negative subtraction involving
$\exp((\sigma_1-1)w)$ has been discarded, and the source's zero-free
width has been replaced by the smaller width $v$.

For $y\geq10^{18}$, elementary inequalities give
\[
 w=\frac7{75}y-2\log(\alpha y)\geq\frac y{12},
 \quad 0<w<y,\quad \log y\leq y^{1/10}.
\]
Consequently
\[
 v\geq\frac1{60y^{7/10}},\qquad
 vw\geq\frac{y^{3/10}}{720}=:D(y),\qquad vy\geq D(y).
\]
Since $24C_1\alpha^3<20$ and $2C_2\alpha^2<2$,
\begin{equation}\label{P01:eq:Fsmall}
 F(x)<y e^{-3y/50}+(20y^3+2y^2)e^{-D(y)}<\frac1{100}.
\end{equation}
For completeness, at $y=10^{18}$ one has
$D(y)>300$ and $\log(22y^3)<130$.
These follow from $\log10<7/3$, $\log22<4$, and
$10^{2/5}>5/2$; the last inequality follows on taking fifth powers.
For $1\leq j\leq3$, the logarithmic derivative of
$y^j e^{-D(y)}$ with respect to $\log y$ is
$j-(3/10)D(y)<0$. The first term of~\eqref{P01:eq:Fsmall} also decreases.
This proves the bound throughout the asserted range, not just at the
left endpoint. The auxiliary bound $\log y\leq y^{1/10}$ and the
estimate for $w$ follow by the same elementary endpoint and monotonicity
checks.

Because $x+h<2x$, we have $G(x,h)<9xy\log y$. Hence
\[
 \frac{KG(x,h)}{Th}
       <108y\log y\,e^{-y/150}<\frac1{100}.
\]
This expression decreases on $y\geq10^{18}$, and its value at the
left endpoint is below $1/100$ without any delicate numerical estimate.
Also $E(x)>-6\sqrt x$, so
\[
                  \frac{E(x)}h>-18e^{-y/6}>-\frac1{100}.
\]
The bracket in~\eqref{P01:eq:theta} is therefore greater than $97/100$,
proving the existence of a prime in $(x,x+h]$.

For an arbitrary real $u\geq X_0$, choose the least half-odd-integer
$x\geq u$. Then $0\leq x-u<1$, and the prime just obtained satisfies
\[
 u<p\leq x+\tfrac12x^{2/3}
       \leq u+1+\tfrac12(u+1)^{2/3}<u+u^{2/3}.
\]
The last inequality holds already for $u\geq8$, using
$(u+1)^{2/3}\leq u^{2/3}+1$.
\end{proof}

\subsection{A bound for the actual infinite minimum}\label{P01:sec:upper}

\begin{proposition}\label{P01:prop:upper}
For every integer $R\geq10^{19}$,
\[
                     \xi(R3^k+1)<\frac43.
\]
\end{proposition}
\begin{proof}
Put $C_k=R3^k+1$, $B=4/3$, and
\[
 c_{k+1}=\frac{C_{k+1}}{C_k}=3-\frac2{C_k}\geq\frac52.
\]
Bertrand's postulate gives a prime $p_1$ with $X_0<p_1<2X_0$;
one can apply the integer version to $\floor{X_0}$.
The upper root endpoint satisfies
\[
 \frac{\log(p_1+1)}{C_1}
  <\frac{10^{18}+\log3}{3R+1}<\frac1{20}
  <\log\frac43.
\]
For the last inequality, $\log(1+t)>t/(1+t)$ at $t=1/3$
gives $\log(4/3)>1/4$.

Suppose a prime $p_k$ has been chosen with
$(p_k+1)^{1/C_k}<B$, and set $c=c_{k+1}$ and $x=p_k^c$.
Then $x\geq X_0$ and
\begin{equation}\label{P01:eq:width}
 \frac{c p_k^{c-1}}{x^{2/3}}
   =c p_k^{-2/(3C_k)}
   >\frac52 B^{-2/3}>1.
\end{equation}
Taylor's theorem gives
\[
 (p_k+1)^c\geq p_k^c+c p_k^{c-1}
                +\frac{c(c-1)}2p_k^{c-2}
           >p_k^c+c p_k^{c-1}+1,
\]
since $c\geq5/2$ and $p_k\geq2$.
Applying Lemma~\ref{P01:lem:gap} at $x$ therefore supplies a prime
$p_{k+1}$ with
\begin{equation}\label{P01:eq:nest}
 p_k^c<p_{k+1}\leq x+x^{2/3},\qquad
             p_{k+1}+1<(p_k+1)^c.
\end{equation}
The second inequality, including its strictness and the extra $1$,
is what is needed for the floor endpoints.

Thus the compact intervals
\[
 I_k=\bigl[p_k^{1/C_k},(p_k+1)^{1/C_k}\bigr]
\]
are nested; each child has its upper endpoint strictly below its
parent's upper endpoint. Their intersection is nonempty.
If $A$ lies in this intersection, then for every $k$ the containment
$A\in I_{k+1}$ and~\eqref{P01:eq:nest} imply
\[
                  p_k\leq A^{C_k}<p_k+1.
\]
Hence $A\in\mathcal W(C)$, and
\[
 \xi(C)\leq A\leq(p_1+1)^{1/C_1}<B.
\]
This constructs an entire infinite prime chain. It makes no assumption
about extending an arbitrary prime prefix below $X_0$.
\end{proof}

\subsection{Completion of the transcendence proof}

\begin{proof}[Proof of Theorem~\ref{P01:thm:main}]
Suppose $\xi(r3^k+s)$ were algebraic. By
Proposition~\ref{P01:prop:normalized}, $R=r/s\in\N$ and
\[
 \eta=\xi(r3^k+s)^s=\xi(R3^k+1)
                 \in\{\delta,\delta^2,\tau\},
\]
with eventually prime traces along $R3^k+1$.
Proposition~\ref{P01:prop:divisors} gives $R>10^{20}$, so
Proposition~\ref{P01:prop:upper} yields $\eta<4/3$.

On the other hand, $f_\delta(4/3)=-11/27<0$, and
$f_\delta$ is strictly increasing on $(1,\infty)$, so $\delta>4/3$.
Also $\delta^2>\delta$, while
$f_\tau(\delta)=-\delta<0$ and $f_\tau$ is strictly increasing on
$(1,\infty)$, whence $\tau>\delta$.
Every possible $\eta$ is therefore greater than $4/3$, a contradiction.
\end{proof}

\subsection{The finite unit enumeration}\label{P01:sec:enumeration}

Here are the exact checks underlying Table~\ref{P01:tab:unit}.
They reproduce the necessary unit rows of \cite{K064} and can be
carried out without floating-point root tests.
Write the minimal polynomial of a cubic Pisot unit below $9$ as
\[
                     X^3-aX^2+bX-c.
\]
The root bounds give the safe exhaustive ranges
\[
                0\leq a\leq10,\quad |b|\leq20,\quad c\in\{-1,1\}.
\]
Reject polynomials having a root at $1$ or $-1$; for these monic
cubics with unit constant term, this is the full rational-root test.
For negative discriminant, a root in $(1,9)$ and unit norm force the
nonreal pair to have modulus $\beta^{-1/2}<1$.
For positive discriminant, use exact Sturm counts to require one root
in $(1,9)$ and two in $(-1,1)$.
Exactly $144$ cubic Pisot units remain.

Since $g\mid2s$, the fixed point is a positive half-integer.
For $\lambda=\ell/2$, $1\leq\ell\leq7$, the positive-fixed-point bound
is $\beta^\ell<9$. Rational bisection tests this inequality exactly.
The numbers of candidates are
\[
                           144,\ 12,\ 4,\ 2,\ 2,\ 1,\ 1.
\]
There are no candidates for $\lambda\geq4$: Siegel's smallest-Pisot
theorem \cite{Siegel} gives $\beta\geq\rho$, where
$\rho^3-\rho-1=0$ and $\rho>3^{1/4}$.
The impossible values $\lambda=3/2,3$ may harmlessly be included in
the finite check.

To test the twisted-limit condition, fix $J\in\{12,22,30\}$ and
$K=J+3$. Let $n_0$ be the least nonnegative residue of
$\ell/2$ modulo $3^K$, and consider
\[
                 n_t=n_0+t3^K,\qquad 0\leq t<104.
\]
These integers cover every residue modulo $104$.
Compute $\Tr(\beta^{n_t})$ in the exact ring
$(\Z/3^J\Z)[X]/(f)$ by binary powering, and retain the residue
$n_t\bmod104$ only when the trace equals $1$ or $-1$ modulo $3^J$.
The principal-unit estimate from \cite{K064} shows that these
computed traces equal the corresponding twisted limits modulo $3^J$.
Indeed, each root lies in a local field of ramification index at most
$3$; for its principal-unit factor $u$, cubing gives
$v_3(u^{3^K}-1)\geq K$.
The Teichm\"uller orders divide
$\operatorname{lcm}(2,8,26)=104$.

For all three precisions, the same five rows remain, with exactly the
residue sets displayed in Table~\ref{P01:tab:unit}. A genuine limit
$\pm1$ must pass this test. Conversely, passing a finite-precision
test is \emph{not} used to prove a limit identity or eventual primality.
The program \texttt{unit\_table\_independent.py}, supplied with this
draft's certificates, performs the enumeration using exact Sturm counts
and rational arithmetic and records the complete residue sets.
The full logical direction needed here is exclusion: every unlisted
candidate fails an exact necessary congruence.

\subsection{Scope and remaining questions}

The hypothesis $3\nmid s$ is used in the algebraic-minimum reduction,
in affine returns, and in the $3$-adic fixed-point analysis.
The hypothesis $s\leq3r$ guarantees the consecutive-ratio condition
from the first prescribed exponent onward. We do not remove either
hypothesis here, although exact common-factor scaling can transport an
individual result to an equivalent schedule.

The global minimum is essential in
Proposition~\ref{P01:prop:upper}. Producing a single prime-representing
$A<4/3$ bounds the least element but need not bound any other local
minimum. Thus this proof does not establish transcendence of all local
minima or all prime-representing constants for the same schedule.
Likewise, the trace argument starts from the eventual floor--trace
identity of an algebraic minimum. It does not identify arbitrary Pisot
floors with traces; the small-conjugate sum may have either sign.

For the classical Mills schedule $C_k=3^k$, the fixed point of the
affine exponent dynamics is $0$, not positive. The principal-unit
factors then disappear from the transferred identity, so the
positive-fixed-point candidate bound used here is unavailable.
No unconditional transcendence statement about the classical Mills
constant follows from Theorem~\ref{P01:thm:main}. The negative-shift
families left open in \cite{K064} also require separate arguments.
Earlier finite-certificate work in \cite{K031} motivates the modular
step, but its negative-shift cutoffs are not used in the present proof.


\subsection{An exact verifier for the large divisors}\label{P01:app:certificate}

The following Python 3 program uses only the standard library.
All arithmetic affecting acceptance is integral. Its final assertions
check both product values and the inequality used in the proof.
The trace verification is performed in two distinct ways.

\begin{lstlisting}
from math import isqrt, prod

def prime(n):
    return n >= 2 and all(n % d for d in range(2, isqrt(n) + 1))

def prime_factors(n):
    out, d = [], 2
    while d * d <= n:
        if n % d == 0:
            out.append(d)
            while n % d == 0:
                n //= d
        d += 1
    if n > 1:
        out.append(n)
    return out

# f = X^3 - a X^2 + b X - 1; tuples are coefficients of 1,X,X^2.
def mul(u, v, a, b, p):
    z = [0] * 5
    for i in range(3):
        for j in range(3):
            z[i+j] += u[i] * v[j]
    for i in (4, 3):
        z[i-1] += a * z[i]
        z[i-2] -= b * z[i]
        z[i-3] += z[i]
    return tuple(z[i] % p for i in range(3))

def xp(n, a, b, p):
    v, x = (1, 0, 0), (0, 1, 0)
    while n:
        if n & 1:
            v = mul(v, x, a, b, p)
        x = mul(x, x, a, b, p)
        n //= 2
    return v

def trace_power(n, a, b, p):
    x, y, z = xp(n, a, b, p)
    return (3*x + a*y + (a*a - 2*b)*z) % p

def trace_recurrence(n, a, b, p):
    v = [3 % p, a % p, (a*a - 2*b) % p]
    for j in range(3, n + 1):
        v.append((a*v[-1] - b*v[-2] + v[-3]) % p)
    return v[n]

DATA = {
    'delta': (1, 0, [
        (2,7,3), (5,31,4), (41,1723,42), (101,10303,58),
        (1217,1482307,1218), (1373,1886503,86),
        (2729,7450171,2730)]),
    'delta_square': (1, -2, [
        (2,7,3), (5,31,2), (41,1723,21), (101,10303,29),
        (1217,1482307,609), (1373,1886503,43),
        (2729,7450171,1365)]),
    'tau': (1, -1, [
        (5,31,11), (1193,1424443,1620),
        (1373,1886503,78), (1973,3894703,331)])
}

products = {}
for name, (a, b, rows) in DATA.items():
    moduli = []
    for p, t, n in rows:
        assert prime(p) and prime(t)
        assert t == p*p + p + 1 and t != 3
        assert pow(3, t-1, t) == 1
        assert all(pow(3, (t-1)//q, t) != 1
                   for q in prime_factors(t-1))
        assert xp(t, a, b, p) == (1, 0, 0)
        assert 0 < n < t and n != 1
        assert trace_power(n, a, b, p) == 0
        assert trace_recurrence(n, a, b, p) == 0
        moduli.append(t)
    assert len(moduli) == len(set(moduli))
    products[name] = prod(moduli)
    print(name, products[name])

assert products['delta'] == 80254725161999161502908681243
assert products['delta_square'] == products['delta']
assert products['tau'] == 324443153856491628397
assert all(value > 10**20 for value in products.values())
print('ALL EXACT CERTIFICATES PASSED')
\end{lstlisting}



\endgroup

\section{Full prefix positive extensions and cube rescaling}\label{src:R07}
The preceding elimination concerns its stated starting index. We now extend the full-prefix range and isolate the exact cube rescaling that yields the clock $3^k+3$. This keeps tail minima separate from minima with every initial constraint.
\medskip
\begingroup




\subsection{Statement, indexing, and dependencies}

For a positive clock $C=(C_i)$ define
\[
 \mathcal W(C)=\{A>1:\floor{A^{C_i}}\text{ is prime at every prescribed index}\}.
\]
If this set has a least element, write $\xi(C)=\min\mathcal W(C)$.
The starting index is part of the definition.

\begin{theorem}[Full-clock extension]\label{R07:thm:full}
Let $r,s$ be integers satisfying
\[
                 r\ge s\ge1,\qquad 3\nmid s,\qquad 3\nmid(r+s).
\]
Then the minimum
\[
                  B=\xi\bigl((r3^j+s)_{j\ge0}\bigr)
\]
exists and is transcendental. Consequently
\[
                  \xi\bigl((r3^k+3s)_{k\ge1}\bigr)=B^{1/3}
\]
exists and is transcendental, as does the minimum for every positive
integer scaling of either clock.
\end{theorem}

\begin{corollary}\label{R07:cor:small}
For every integer $q\ge1$,
\[
                  \xi\bigl((q(3^k+3))_{k\ge1}\bigr)
\]
exists and is transcendental.
\end{corollary}

The full clock $(3^j+1)_{j\ge0}$ is $2,4,10,28,\ldots$;
it is not the shorter clock $4,10,28,\ldots$.
\emph{Transcendence of a tail's minimum does not, by itself, imply
transcendence of the full clock's minimum.} No such inference is made here.
Only after applying Saito's theorem to the full minimum do we use
an eventual statement about its trace orbit.

The inputs are Saito's Theorem~1.3 and Proposition~3.1 in
arXiv version~3 \cite{Saito}, the prime-interval exponent $21/40$
of Baker--Harman--Pintz \cite{BHP}, and the positive trace theory of
corpus manuscript [064] \cite{K064}. The unit table, large-divisor
certificates, and explicit prime-interval lemma are inherited from
\cite{P01}; the finite checks are reproduced with this note.
The two preceding manuscripts are unpublished research manuscripts.
This note proves the new full-prefix bridge and infinite-chain adaptation;
it does not independently reprove all their trace or analytic lemmas.

For later use, the increasing bijection $A\mapsto A^d$ gives, for
any positive integer $d$,
\begin{equation}\label{R07:eq:scaling}
 \mathcal W(dC)=\{A>1:A^d\in\mathcal W(C)\},\qquad
                  \xi(dC)=\xi(C)^{1/d}.
\end{equation}
This is an exact identity of global minima whenever they exist.

\subsection{The full-clock algebraic reduction}\label{R07:sec:bridge}

A Pisot number is a real algebraic integer greater than $1$ whose other
conjugates have absolute value less than $1$.
All traces and norms below are over $\Q$.

\begin{proposition}\label{R07:prop:bridge}
Under the hypotheses of Theorem~\ref{R07:thm:full}, the full minimum $B$
exists. If $B$ is algebraic, there is $g\mid2s$ such that
$\beta=B^g$ is a cubic Pisot number. For every sufficiently large $j$,
\[
 n_j=\frac{r3^j+s}{g}\in\N,\qquad
 \Tr(\beta^{n_j})=\floor{B^{r3^j+s}}\text{ is prime}.
\]
If $x_2,x_3$ are the other conjugates of $\beta$, then, for every fixed
$0<\kappa<17/40$,
\[
                  |x_2^{n_j}+x_3^{n_j}|
                         =O(\beta^{-\kappa n_j}).
\]
\end{proposition}
\begin{proof}
Put $D_j=r3^j+s$ for $j\ge0$. The consecutive ratios are
\[
 c_j:=\frac{D_{j+1}}{D_j}=3-\frac{2s}{D_j}\ge2,
 \qquad c_j\longrightarrow3.
\]
Thus Saito's Theorem~1.3 gives existence of the minimum for this
\emph{full} clock, with $C_{j+1}=D_j$ in his indexing.

Every $D_m$ is coprime to $3$: for $m\ge1$ this follows from
$3\nmid s$, and for $m=0$ it is exactly the additional condition
$3\nmid(r+s)$. Let $h=\ord_{D_m}(3)$. Then
\[
                 D_{m+h}\equiv D_m\equiv0\pmod{D_m},
                 \qquad D_{m+h}>D_m.
\]
If $B^{D_m}$ were an integer, its prime floor would make it a prime
$p$, and the later output $B^{D_{m+h}}=p^{D_{m+h}/D_m}$ would be
composite. Hence
\begin{equation}\label{R07:eq:nonintegral}
                     B^{D_m}\notin\N\quad\text{for every }m\ge0.
\end{equation}
This verifies the initial boundary as well as every later one.

Take $\theta=21/40$ and $\varepsilon_0=1/10$ in Saito's
Proposition~3.1. Since $D_j\ge4s$ for $j\ge1$, we have $c_j\ge5/2$
there. The selected set
\[
          I=\{j\ge0:c_j\ge40/19+1/10\}
\]
therefore contains every $j\ge1$, possibly also $j=0$.
The proposition gives $\beta=B^g$ Pisot, with
\[
        2\le\deg\beta\le1+(3\cdot19/40-1)^{-1}
                            =1+40/17<4,
\]
and $g\mid D_j$ at all sufficiently large indices.
Taking consecutive indices shows
$g\mid D_{j+1}-3D_j=-2s$, so $g\mid2s$ and $3\nmid g$.

If the degree were $2$, Saito's quadratic alternative would require
$n_j=D_j/g$ to be an odd prime for every sufficiently large $j$.
Choose one such prime $p=n_j>g$ and put $h=\ord_p(3)$; this order
exists since $3\nmid D_j$. The congruence
$D_{j+h}\equiv D_j\pmod p$ and $p\nmid g$ imply
$p\mid n_{j+h}$, while $n_{j+h}>p$. This contradicts primality.
Thus the degree is $3$, and Saito's trace identity gives the asserted
prime traces.

Finally Saito's fractional-part estimate is
\[
 \{B^{D_j}\}\le\frac{2}{c_j}
       \floor{B^{D_j}}^{-(1-\theta)c_j+1}
              \quad\text{for all sufficiently large }j\in I.
\]
Since $(1-\theta)c_j-1\to17/40$, this implies the stated estimate
for every fixed $\kappa<17/40$, using the trace identity and
$\floor{B^{D_j}}\sim\beta^{n_j}$. We do not assert the limiting
exponent $17/40$ itself.
\end{proof}

Saito's proposition is used on the original minimum, not on a tail
minimum. Its finite-prefix preservation is part of the primary proof
\cite[Lemma~5.2, equations (5.5)--(5.14)]{Saito}; the selected set need
not include the first coordinate. The numbering here is that of
version~3: the reduction is Proposition~\emph{3.1}, not 3.2.

\subsection{A short proof for the clock \texorpdfstring{$3^k+3$}{3\string^k+3}}

We isolate the part of [064] needed for this family.

\begin{lemma}[Inherited positive trace rigidity]\label{R07:lem:rigidity}
Let $\beta$ be a cubic Pisot number and let $r,s,g\in\N$ with
$3\nmid gs$ and $n_j=(r3^j+s)/g\in\N$ for all sufficiently large $j$.
If $\Tr(\beta^{n_j})$ is eventually prime, then
$\beta^{s/g}<3$. Moreover, if $4s\nmid r$ and $13s\nmid r$,
no such eventually prime trace orbit exists.
\end{lemma}

This is the positive-fixed-point part (iv) of the ternary trace-rigidity
theorem in \cite{K064}. It is a theorem about eventual \emph{traces},
not about transferring minima between clocks. Its proof uses affine
finite-field return, twisted $3$-adic limits, Galois transfer, and a
finite candidate calculation.

\begin{proof}[Proof of Corollary~\ref{R07:cor:small}, independently of the large-height bound]
Let $E_j=3^j+1$ for $j\ge0$, and let $\eta=\xi((E_j)_{j\ge0})$.
Proposition~\ref{R07:prop:bridge}, with $r=s=1$, gives existence.
If $\eta$ were algebraic, the same proposition would give a cubic
Pisot number $\beta=\eta^g$, $g\mid2$, with eventually prime traces
at $(3^j+1)/g$. Lemma~\ref{R07:lem:rigidity} rules this out, since neither
$4$ nor $13$ divides $1$.

In this case the initial boundary check is especially simple:
$\eta^2=p\in\N$ would force the next output $\eta^4=p^2$.
The selected set omits just this first coordinate, because
$c_0=2$ and $c_1=5/2$; it nevertheless belongs to the minimum problem
throughout the proof.
Since $q(3^k+3)=3qE_{k-1}$, equation~\eqref{R07:eq:scaling} gives
\[
              \xi\bigl((q(3^k+3))_{k\ge1}\bigr)=\eta^{1/(3q)}.
\]
An algebraic value on the left would make $\eta$ algebraic, proving
the corollary. No explicit large prime-gap threshold is used here.
\end{proof}

\subsection{Inherited unit candidates and their normalization}

For the broader theorem we state the further inherited inputs precisely.
The approximation-magnitude lemma of \cite{K064} says that a cubic
Pisot number satisfying
$|x_2^n+x_3^n|\le A_0\beta^{-\kappa n}$ for infinitely many $n$,
with $0<\kappa<1/2$, satisfies
\begin{equation}\label{R07:eq:magnitude}
                       |\Nm(\beta)|\le\beta^{1-2\kappa}.
\end{equation}
For nonreal conjugates, its proof uses a polynomial lower bound from
linear forms in logarithms \cite{BW}, not a uniform lower bound on
an oscillating cosine.

Apply this to Proposition~\ref{R07:prop:bridge} with $\kappa=2/5$.
Lemma~\ref{R07:lem:rigidity} and $g\le2s$ give $\beta<9$, whence
\[
                     |\Nm(\beta)|\le\beta^{1/5}<2.
\]
The nonzero integral norm is therefore $\pm1$.

Write $\lambda=s/g$. The exact unit-only table of \cite{P01},
using the necessary twisted-limit identities of \cite{K064}, leaves
only Table~\ref{R07:tab:units}. Define $\delta,\tau>1$ by
\[
        \delta^3-\delta^2-1=0,\qquad
        \tau^3-\tau^2-\tau-1=0.
\]

\begin{table}[htbp]
\centering\small
\begin{tabular}{cclcl}
\toprule
$\lambda$ & $\beta$ & Minimal polynomial & $B^s$ & Necessary eventual residue\\
\midrule
$1$ & $\delta$ & $X^3-X^2-1$ & $\delta$ & $n_j\equiv1\pmod8$\\
$1$ & $\tau$ & $X^3-X^2-X-1$ & $\tau$ & $n_j\equiv1\pmod{13}$\\
$1$ & $\delta^2$ & $X^3-X^2-2X-1$ & $\delta^2$ & $n_j\equiv1\pmod4$\\
$1/2$ & $\tau^2$ & $X^3-3X^2-X-1$ & $\tau$ & $n_j\equiv7\pmod{13}$\\
$2$ & $\delta$ & $X^3-X^2-1$ & $\delta^2$ & $n_j\equiv2\pmod8$\\
\bottomrule
\end{tabular}
\caption{Necessary unit candidates; no row asserts eventual primality.}
\label{R07:tab:units}
\end{table}

The hypotheses of this inherited table are exactly those now available:
a cubic Pisot unit $\beta<9$, $g\mid2s$, $3\nmid gs$, the bound
$\beta^{s/g}<3$, and eventual prime traces along $(r3^j+s)/g$.
Its finite verification is described in Section~\ref{R07:app:unit}.

\begin{proposition}\label{R07:prop:normalize}
If the full minimum $B$ is algebraic, then $s\mid r$ and, with
$R=r/s$ and $\eta=B^s$,
\begin{equation}\label{R07:eq:fullnormalize}
 \boxed{\ \eta=\xi\bigl((R3^j+1)_{j\ge0}\bigr)
                       \in\{\delta,\delta^2,\tau\}.\ }
\end{equation}
Furthermore $\Tr(\eta^{R3^j+1})$ is prime for all sufficiently large $j$.
\end{proposition}
\begin{proof}
In the rows $\lambda=1$ and $1/2$, integrality of $n_j$ and
$3\nmid s$ imply $s\mid r$.
In the row $\lambda=2$, the residue condition gives
\[
                    \frac{2r3^j}{s}=n_j-2\equiv0\pmod8.
\]
Thus $4s\mid r3^j$, and $3\nmid4s$ implies $4s\mid r$.
Every row therefore has integral $R=r/s$, and the possible values
of $\eta=\beta^\lambda$ are those displayed.

A Pisot number and any positive integer power generate the same field:
an embedding fixing that power must send the Pisot number to a conjugate
of the same absolute value, and no other conjugate has that value.
This handles both the $\tau^2$ and $\delta^2$ rows and gives
\[
       \Tr_{\Q(\beta)/\Q}(\beta^{n_j})
           =\Tr_{\Q(\eta)/\Q}(\eta^{R3^j+1}).
\]
Finally, apply the exact scaling identity~\eqref{R07:eq:scaling} to
$D_j=s(R3^j+1)$ at \emph{every} $j\ge0$. This proves the equality
of minima in~\eqref{R07:eq:fullnormalize}, with initial exponent $R+1$;
no shorter-clock minimum appears.
\end{proof}

\subsection{Exact divisibility obstructions}

For $\eta\in\{\delta,\delta^2,\tau\}$ write its minimal polynomial
as $f(X)=X^3-aX^2+bX-1$, and put $T_n=\Tr(\eta^n)$.
Newton's identities give
\begin{equation}\label{R07:eq:recurrence}
 T_0=3,\quad T_1=a,\quad T_2=a^2-2b,\qquad
 T_n=aT_{n-1}-bT_{n-2}+T_{n-3}\quad(n\ge3).
\end{equation}

\begin{lemma}\label{R07:lem:ring}
Suppose $p,t$ are primes, $t\ne3$, $\ord_t(3)=t-1$,
$X^t=1$ in $\F_p[X]/(f)$, and $T_n\equiv0\pmod p$ for some
$n\not\equiv1\pmod t$. If $T_{R3^j+1}$ is eventually prime,
then $t\mid R$.
\end{lemma}
\begin{proof}
If $t\nmid R$, the residues $R3^j$ run periodically through all
nonzero classes modulo $t$, so $R3^j+1\equiv n\pmod t$ infinitely
often. The quotient-ring identity makes the traces periodic modulo $p$
with period dividing $t$. Thus these trace values are divisible by $p$.
They tend to infinity because $T_m=\eta^m+o(1)$, so they eventually
exceed $p$, contradicting primality. Neither irreducibility of $f$
modulo $p$ nor exact order $t$ of $X$ in the quotient is required.
\end{proof}

\begin{table}[htbp]
\centering
\begin{tabular}{rrrr}
\toprule
$p$ & $t$ & $n$ for $\delta$ & $n$ for $\delta^2$\\
\midrule
2 & 7 & 3 & 3\\
5 & 31 & 4 & 2\\
41 & 1723 & 42 & 21\\
101 & 10303 & 58 & 29\\
1217 & 1482307 & 1218 & 609\\
1373 & 1886503 & 86 & 43\\
2729 & 7450171 & 2730 & 1365\\
\bottomrule
\end{tabular}
\qquad
\begin{tabular}{rrr}
\toprule
$p$ & $t$ & $n$ for $\tau$\\
\midrule
5 & 31 & 11\\
1193 & 1424443 & 1620\\
1373 & 1886503 & 78\\
1973 & 3894703 & 331\\
\bottomrule
\end{tabular}
\caption{Exact inherited finite-ring certificates. Each row has $t=p^2+p+1$.}
\label{R07:tab:divisors}
\end{table}

Every row of Table~\ref{R07:tab:divisors} satisfies Lemma~\ref{R07:lem:ring}.
The supplied programs check primality, the primitive-root condition,
quotient-ring powering, and the zero trace by two methods; a separate
companion-matrix implementation reproduces the identities.
The distinct prime $t$'s in each branch have products
\begin{align}
 80254725161999161502908681243&\mid R
                   &&(\eta=\delta\text{ or }\delta^2),\label{R07:eq:divdelta}\\
 324443153856491628397&\mid R
                   &&(\eta=\tau).\label{R07:eq:divtau}
\end{align}
In particular, $R>10^{20}$ in every branch.
These certificates use only eventual traces, so adding the initial
coordinate does not alter their conclusions. Multiples of these
products survive the finite sieve; the global bound below eliminates
them separately.

\subsection{A full-prefix infinite prime chain}\label{R07:sec:global}

\begin{lemma}[Inherited explicit prime-interval bound]\label{R07:lem:gap}
For every real $u\ge X_{\mathrm{gap}}:=\exp(10^{18})$, there is a prime
\[
                         u<p\le u+u^{2/3}.
\]
\end{lemma}

This is the explicit lemma proved in \cite[Section~4]{P01}
from Cully-Hugill \cite{CH}, using Theorem~2, Table~2, and
equations (2), (9), (15), (19)--(21) of that paper.
The substituted parameters are
\[
 T=x^{17/50},\quad \sigma_1=3/5,\quad
 K=3.4747,\quad C_1=17.418,\quad C_2=5.272,
\]
with the zero-free width
\[
 \frac{1}{57.54(\log T)^{2/3}(\log\log T)^{1/3}}.
\]
The derivation treats half-odd-integer endpoints first and then arbitrary
real $u$. This explicit lemma is inherited as stated; no numerical
constant is inferred from an asymptotic $O$-bound.

\begin{proposition}[The initial exponent is $R+1$]\label{R07:prop:upper}
For every integer $R\ge10^{19}$,
\[
                   \xi\bigl((R3^j+1)_{j\ge0}\bigr)<\frac43.
\]
\end{proposition}
\begin{proof}
Put $C_j=R3^j+1$ for $j\ge0$ and $B_0=4/3$.
Bertrand's postulate, applied to $\floor{X_{\mathrm{gap}}}$,
gives a prime $p_0$ with
$X_{\mathrm{gap}}<p_0<2X_{\mathrm{gap}}$.
As $p_0+1<3X_{\mathrm{gap}}$, we have
\begin{equation}\label{R07:eq:start}
 \frac{\log(p_0+1)}{C_0}
 <\frac{10^{18}+\log3}{R+1}
 <\frac15<\frac14<\log\frac43.
\end{equation}
Here $C_0=R+1$, not $3R+1$. The last inequality follows from
$\log(1+t)>t/(1+t)$ with $t=1/3$.

Suppose $p_j$ has been chosen with $(p_j+1)^{1/C_j}<B_0$.
Set
\[
                c=\frac{C_{j+1}}{C_j}=3-\frac2{C_j}\ge\frac52,
                \qquad u=p_j^c\ge X_{\mathrm{gap}}.
\]
Since $p_j^{1/C_j}<B_0$, the first-order width satisfies
\begin{equation}\label{R07:eq:width}
 \frac{cp_j^{c-1}}{u^{2/3}}
       =cp_j^{-2/(3C_j)}
       >\frac52 B_0^{-2/3}>1.
\end{equation}
Taylor's theorem, with increasing second derivative for $c>2$, gives
\[
 (p_j+1)^c\ge p_j^c+cp_j^{c-1}
                   +\frac{c(c-1)}2p_j^{c-2}
             >u+cp_j^{c-1}+1>u+u^{2/3}+1.
\]
The remainder exceeds $1$ since $c\ge5/2$ and $p_j\ge2$.
By Lemma~\ref{R07:lem:gap}, choose a prime $p_{j+1}$ with
\begin{equation}\label{R07:eq:nest}
 p_j^c<p_{j+1}\le u+u^{2/3},\qquad
                        p_{j+1}+1<(p_j+1)^c.
\end{equation}
Its upper root endpoint remains below $B_0$, so this constructs
an entire infinite sequence.

The compact intervals
\[
                  I_j=[p_j^{1/C_j},(p_j+1)^{1/C_j}]
\]
are nested, with each child upper endpoint strictly below its parent's.
Their intersection is nonempty. For any point $A$ in the intersection,
use $A\in I_{j+1}$ and~\eqref{R07:eq:nest} to obtain
\[
                       p_j\le A^{C_j}<p_j+1
                             \quad\text{for every }j\ge0.
\]
Thus every floor is prime, including the first floor at exponent $R+1$,
and $A\le(p_0+1)^{1/C_0}<B_0$ by~\eqref{R07:eq:start}.
Saito's existence theorem supplies the minimum, so
$\xi((C_j)_{j\ge0})\le A<B_0$ as asserted.
This constructs a full infinite chain; it does not presume that an
arbitrarily chosen finite prime prefix extends.
\end{proof}

\subsection{Completion and scope}

\begin{proof}[Proof of Theorem~\ref{R07:thm:full}]
Existence and the algebraic reduction are established in
Proposition~\ref{R07:prop:bridge}. If $B$ were algebraic,
Proposition~\ref{R07:prop:normalize} would give
\[
       \eta=B^s=\xi\bigl((R3^j+1)_{j\ge0}\bigr)
                  \in\{\delta,\delta^2,\tau\},\qquad R=r/s\in\N.
\]
Equations~\eqref{R07:eq:divdelta}--\eqref{R07:eq:divtau} force $R>10^{20}$,
so Proposition~\ref{R07:prop:upper} gives $\eta<4/3$.
But $f_\delta(4/3)=-11/27<0$ and $f_\delta$ is strictly increasing
on $(1,\infty)$, so $\delta>4/3$. Also $\delta^2>\delta$ and
$f_\tau(\delta)=-\delta<0$, whence $\tau>\delta$.
Every candidate exceeds $4/3$, a contradiction.

Finally, $r3^k+3s=3D_{k-1}$ for $k\ge1$.
Equation~\eqref{R07:eq:scaling} proves the translated statement and all
positive integer scalings, preserving global minima exactly.
\end{proof}

The condition $3\nmid(r+s)$ is retained in Theorem~\ref{R07:thm:full}:
it excludes an integral power at the newly imposed coordinate.
Without that condition the following necessary-only statement is valid.

\begin{proposition}[The only possible algebraic initial boundary]
\label{R07:prop:boundary}
Let $r\ge s\ge1$ and $3\nmid s$, without assuming $3\nmid(r+s)$.
The full minimum $B=\xi((r3^j+s)_{j\ge0})$ exists.
If $B$ is algebraic, then for some prime $p$,
\[
 B^{r+s}=p,\qquad B=p^{1/(r+s)},\qquad
 \xi\bigl((r3^k+3s)_{k\ge1}\bigr)=p^{1/[3(r+s)]}.
\]
This is a necessary condition only: it asserts neither the existence of
such an algebraic minimum nor the admissibility or minimality of any
specified prime radical.
\end{proposition}
\begin{proof}
The ratio bound still gives existence. For every $j\ge1$, the exponent
$r3^j+s$ is coprime to $3$, so the return argument in
Proposition~\ref{R07:prop:bridge} excludes an integral $B^{r3^j+s}$.
If an algebraic $B$ also had $B^{r+s}\notin\N$, all the hypotheses of
Saito's reduction would hold. The remainder of the full-clock candidate,
divisibility, and global-bound argument above would then give the same
contradiction: none of those later steps uses $3\nmid(r+s)$.
Consequently $B^{r+s}$ must be an integer, hence a prime because its
floor is prime. The last identity follows from~\eqref{R07:eq:scaling}.
\end{proof}

The condition $r\ge s$ supplies the ratio bound from the first coordinate.
No claim is made about every admissible base, every local minimum,
the classical Mills constant, or the unresolved negative-shift families.

\paragraph{Relation to the earlier corpus.}
The earlier additive-clock theorem
\cite[Section~6, Theorem~6.1]{R07-K008} proves transcendence of the full
minima for $(ba^k\pm k)_{k\ge1}$, with $a\ge3$ and $b\ge1$.
Its perturbation is the varying term $\pm k$, not a constant shift.
A stronger antecedent relevant here is
\cite[Section~27, Theorem~27.2]{K031}: for
$C_n=P(a^n)+Q(n)$, with $P$ of positive leading coefficient and
positive degree, $P(0)+Q$ not the zero polynomial, and
$L=a^{\deg P}>40/19$, every sufficiently late tail minimum is
transcendental. Constant shifts are included. Taking
$a=3$, $P(X)=rX$, and $Q(X)=s$ gives $L=3$ and applies to sufficiently
late tails of the present clocks; taking $P(X)=X$, $Q(X)=3$ applies
to sufficiently late tails of $3^k+3$.

The distinction is the prescribed starting coordinate. At $L=3$,
that proof uses the fact that the tail minima tend to $1$ as the
starting index tends to infinity. It does not give the corresponding
claim for an arbitrary fixed full prefix. The same theorem's stronger
conclusion for every tail left-gap point (a local left minimum), and
its full-clock extension under recurring divisibility, are stated for
$L\ge4$. Moreover, the finite-prefix dichotomy in
\cite[Section~27, Theorem~27.1]{K031} requires transcendence of
\emph{every} left-gap point of the underlying tail set, not merely
of its global minimum. Thus the already established $L=3$ tail-minimum
result alone does not justify a finite-prefix transfer here.
The present argument instead applies Saito directly to the specified
full minimum and keeps its initial coordinate throughout. It is a
full-prefix extension of the existing framework; no historical
priority claim is made.


\subsection{Reproducing the exact finite checks}\label{R07:app:reproduce}

The accompanying \texttt{certificates/} directory contains four portable
Python~3 programs. Both \texttt{unit\_table\_independent.py} and
\texttt{positive\_prefix\_certificate.py}, which imports the unit-table
module, require SymPy. The two large-divisor programs use only the
standard library. All acceptance tests use exact integer or rational
arithmetic.
Run from the corresponding extracted source archive's directory:
\begin{lstlisting}[language=bash]
python3 certificates/large_divisor_certificate.py
python3 certificates/independent_matrix_audit.py
python3 certificates/unit_table_independent.py
python3 certificates/positive_prefix_certificate.py
\end{lstlisting}
The programs assert their claimed conclusions. Expected terminal lines are
\begin{lstlisting}[language={}]
ALL EXACT CERTIFICATES PASSED
ALL INDEPENDENT MATRIX CHECKS PASSED
ALL EXACT UNIT TABLE CHECKS PASSED
PASS: all full-prefix algebraic candidates excluded by exact finite residues
\end{lstlisting}
Rerun logs are supplied alongside them. The companion-matrix verifier
reads only the literal certificate records from the first program and
implements its powering checks separately.

For the large divisors, primality of all $p,t$ is checked by deterministic
trial division. The test $\ord_t(3)=t-1$ uses every prime divisor of
$t-1$. Quotient-ring multiplication uses $X^3=aX^2-bX+1$;
the zero trace is checked both by powering and by~\eqref{R07:eq:recurrence}.
These finite identities prove the implications of Lemma~\ref{R07:lem:ring};
they do not certify eventual primality for any surviving class.

\subsubsection{The necessary unit table}\label{R07:app:unit}
Write a cubic Pisot unit below $9$ as a root of
$X^3-aX^2+bX-c$. Safe exhaustive bounds are
\[
                     0\le a\le10,\qquad |b|\le20,
                     \qquad c\in\{-1,1\}.
\]
The rational-root test rejects roots $\pm1$. Exact Sturm counts handle
the totally real case; in the nonreal case the unit norm and the unique
root $\beta\in(1,9)$ give the other roots modulus $\beta^{-1/2}<1$.
Exactly $144$ cubic Pisot units remain.

As $g\mid2s$, write $\lambda=\ell/2$ with $\ell\ge1$.
The bound $\beta^\lambda<3$ becomes $\beta^\ell<9$ and is tested by
rational bisection. For $\ell=1,\ldots,7$ the respective candidate
counts are
\[
                         144,\ 12,\ 4,\ 2,\ 2,\ 1,\ 1.
\]
There are no candidates with $\lambda\ge4$, by Siegel's smallest-Pisot
theorem \cite{Siegel}: its smallest value $\rho$, satisfying
$\rho^3-\rho-1=0$, exceeds $3^{1/4}$.
The impossible values $\lambda=3/2,3$ can harmlessly be tested too.

For each $J\in\{12,22,30\}$ set $K=J+3$. Let
$n_0\equiv\ell/2\pmod{3^K}$ be the least nonnegative representative
and test $n_t=n_0+t3^K$, $0\le t<104$.
These cover every residue modulo $104$.
Binary powering in $(\Z/3^J\Z)[X]/(f)$ computes the trace exactly;
retain a residue only if that trace equals $1$ or $-1$ modulo $3^J$.
The principal-unit estimate and twisted-limit theorem of \cite{K064}
make this a necessary test: the Teichm\"uller orders divide
$104=\operatorname{lcm}(2,8,26)$, and principal-unit powers at $3^K$
are sufficiently close to $1$ modulo $3^J$.
Every genuine twisted limit $\pm1$ passes, whereas passing is not
claimed sufficient. All three precisions give exactly the five necessary
rows of Table~\ref{R07:tab:units}.

For the small family, the additional program intersects these necessary
residues with the full orbits
\[
                  \left\{\frac{3^j+1}{g}\bmod104:j\ge0\right\},
                  \qquad g=1,2.
\]
Their periods are $6$ and $12$, respectively. At precision $J=22$,
all $12$ bounded unit candidates for $g=1$ and all $144$ for $g=2$
are excluded. Equivalently, the first three relevant rows would require
$3^j$ divisible by $8,4$, or $13$; the $g=2$ row would require
$(3^j+1)/2\equiv7\pmod{13}$, hence $3^j\equiv0\pmod{13}$.
This separately checks the finite part of the small-family argument.
It does not independently reprove the inherited trace or gap lemmas.

\subsubsection{The portable large-divisor verifier}
For convenience, the complete first verifier is included below; the
separate unit and matrix programs are supplied as companion files.
\begin{lstlisting}
#!/usr/bin/env python3
"""Exact finite certificates for trace(eta^(R*3^k+1)) eventual primality.
No nonstandard dependencies; all primality/order checks are deterministic.
Polynomials are x^3-a*x^2+b*x-1.
"""
from math import isqrt, prod

def prime(n):
    if n<2:return False
    return all(n%d for d in range(2,isqrt(n)+1))

def prime_factors(n):
    ans=[];d=2
    while d*d<=n:
        if n%d==0:
            ans.append(d)
            while n%d==0:n//=d
        d+=1
    if n>1:ans.append(n)
    return ans

def mul(u,v,a,b,p):
    z=[0]*5
    for i in range(3):
        for j in range(3):z[i+j]+=u[i]*v[j]
    for i in (4,3):
        z[i-1]+=a*z[i]
        z[i-2]-=b*z[i]
        z[i-3]+=z[i]
    return tuple(z[i]%p for i in range(3))

def xp(n,a,b,p):
    v=(1,0,0);x=(0,1,0)
    while n:
        if n&1:v=mul(v,x,a,b,p)
        x=mul(x,x,a,b,p);n//=2
    return v

def trace_pow(n,a,b,p):
    x,y,z=xp(n,a,b,p)
    return (3*x+a*y+(a*a-2*b)*z)%p

def trace_recurrence(n,a,b,p):
    v=[3%p,a%p,(a*a-2*b)%p]
    for j in range(3,n+1):v.append((a*v[-1]-b*v[-2]+v[-3])%p)
    return v[n]

DATA={
'delta':(1,0,[(2,7,3),(5,31,4),(41,1723,42),(101,10303,58),(1217,1482307,1218),(1373,1886503,86),(2729,7450171,2730)]),
'delta_square':(1,-2,[(2,7,3),(5,31,2),(41,1723,21),(101,10303,29),(1217,1482307,609),(1373,1886503,43),(2729,7450171,1365)]),
 'tau':(1,-1,[(5,31,11),(1193,1424443,1620),(1373,1886503,78),(1973,3894703,331)])}

def verify():
    for name,(a,b,rows) in DATA.items():
        ts=[]
        for p,t,n in rows:
            assert prime(p) and prime(t)
            assert t==p*p+p+1
            assert pow(3,t-1,t)==1
            factors=prime_factors(t-1)
            assert all(pow(3,(t-1)//ell,t)!=1 for ell in factors)
            assert xp(t,a,b,p)==(1,0,0)
            assert 0<n<t and n!=1
            assert trace_pow(n,a,b,p)==0
            assert trace_recurrence(n,a,b,p)==0
            ts.append(t)
            print(name,'p',p,'t',t,'zero exponent',n,'factors(t-1)',factors)
        assert len(ts)==len(set(ts))
        L=prod(ts)
        print('NECESSARY DIVISOR OF R:',name,L)
    assert prod(row[1] for row in DATA['delta'][2])==80254725161999161502908681243
    assert prod(row[1] for row in DATA['tau'][2])==324443153856491628397
    print('ALL EXACT CERTIFICATES PASSED')

if __name__=='__main__':verify()

\end{lstlisting}




\endgroup

\section{Negative shifts and integral fixed point rigidity}\label{src:N01}
The negative shift changes the fixed-point argument. The results below establish a structural reduction and explicit families; the multiplier-thirteen and other odd-multiplier cases are subsequently strengthened in Sections~\ref{src:R02} and~\ref{src:R04}. Statements of what remains at this stage are superseded only by those explicitly later theorems.
\medskip
\begingroup



\subsection{Main results and scope}

For an integer sequence $C=(C_k)_{k\ge1}$ with $C_{k+1}\ge2C_k>0$, write
\[
 \xi(C)=\min\{A>1:\floor{A^{C_k}}\text{ is prime for every }k\ge1\}.
\]
Existence is covered by Saito's general existence theorem
\cite[Theorem~1.3]{Saito}. Throughout this note,
\begin{equation}\label{N01:eq:scope}
 r,s\in\N,\qquad 3\nmid s,\qquad s<3r,\qquad C_k=r3^k-s.
\end{equation}
Thus $C_k>0$, $C_{k+1}=3C_k+2s$, and the consecutive ratios tend to $3$
from above. We abbreviate $\xi(C)$ to $\xi(r3^k-s)$.

\begin{mainthm}[Extension to even shifts]\label{N01:thm:extension}
Assume \eqref{N01:eq:scope} and put $d=\gcd(r,s)$. If $r/d$ is odd and
$13\nmid r/d$, then $\xi(r3^k-s)$ is transcendental.
\end{mainthm}

The odd--odd case is Theorem~B of \cite{K064}. The added cases have
$r/d$ odd and $s/d$ even; for example, $\xi(3^k-2)$ is covered.
The hypothesis $3\nmid s$ has \emph{not} been removed.

\begin{mainthm}[A uniform cutoff for the shift minus one]\label{N01:thm:cutoff}
For every integer $r\ge12\,230\,000\,000\,000$, the number
$\xi(r3^k-1)$ is transcendental.
\end{mainthm}

Theorem~\ref{N01:thm:cutoff} includes even $r$ and odd multiples of $13$.
Its proof uses a finite computation only after obtaining a global bound on
\emph{every} hypothetical algebraic minimum in the stated range. It does not
infer a theorem for arbitrary $r$ from a finite prime prefix or a bounded
search for Pisot numbers.

Section~\ref{N01:app:general} gives a further sufficient condition for
general $-s$: with $R=r/d$, $S=s/d$ and the parity gcd $D\in\{1,2\}$,
transcendence holds if $(9R-S)/D\ge2\,924\,000\,000\,000\,000$.
In particular, the still-unresolved \emph{normalized} parameter pairs
form a finite set; this finiteness is not a proof that the set is empty.

The new structural result is Lemma~\ref{N01:lem:inverse}. It gives
$\Tr(\beta^{-s})=\pm1$ without an upper bound on $\beta$ and without
assuming that $\beta$ is a unit or is non-totally-real. A positive-shift
argument cannot simply be reused: at a negative fixed point the terms in
the limiting identity are inverse powers of the conjugates. The proof
below treats this sign directly.

\subsection{Inputs and normalization}\label{N01:sec:inputs}

We use the following prior results, with their provenance made explicit.
The initial reduction is Saito's \cite[Proposition~3.1]{Saito}; its
specialization and the subsequent $3$-adic arguments are from \cite{K064}.
The numbering of Saito here refers to arXiv version~3. In particular, the
general reduction is Proposition~3.1 in that version, rather than the
Proposition~3.2 citation used in \cite{K064}.

\begin{enumerate}[label=\textup{(I\arabic*)}]
\item\label{N01:input:shadow}
\emph{Cubic Pisot shadow.}
If $\xi=\xi(C)$ is algebraic, the Mills-shadow corollary in
\cite[Section~2]{K064} provides $g\mid D$, where
$D=\lim_m\gcd(C_m,C_{m+1},\ldots)$, and a cubic Pisot number
$\beta=\xi^g$. For all sufficiently large $k$, putting $n_k=C_k/g$ gives
\begin{equation}\label{N01:eq:shadow}
 T_k=\Tr(\beta^{n_k})=\floor{\xi^{C_k}}\text{ prime},\qquad
 0\le\beta^{n_k}-T_k\le T_k^{-17/40}.
\end{equation}
Here $D\mid2s$, and $3\nmid D$.

\item\label{N01:input:norm}
\emph{Norm bounds.}
The approximation-magnitude lemma in \cite[Section~3.7]{K064} gives
\begin{equation}\label{N01:eq:norm}
 |\Nm(\beta)|\le\beta^{3/20}.
\end{equation}
The negative-fixed-point norm-exclusion lemma in
\cite[Section~4, lemma titled $3\nmid\Nm(\beta)$ at negative fixed
points]{K064} gives $3\nmid\Nm(\beta)$ when \eqref{N01:eq:shadow} holds
and $g\le2s$. The magnitude lemma uses the Pisot power-field property
and, for a nonreal conjugate pair, a lower bound for a nonzero linear
form in logarithms.

\item\label{N01:input:limits}
\emph{Actual twisted limit identities.}
Write the $3$-adic unit conjugates as $\beta_i=\omega_i u_i$, with
$\omega_i\in\mu_{104}$ and $u_i$ a principal unit. The affine-return
lemma and the limit-identity proposition in
\cite[Sections~3.4--3.5]{K064} imply that, in each phase class of $k$,
\begin{equation}\label{N01:eq:general-limit}
 \sum_i\omega_i^{n_k}u_i^{-s/g}=\varepsilon_k,
 \qquad\varepsilon_k\in\{1,-1\}.
\end{equation}
When $3\nmid\Nm(\beta)$, all three conjugates occur. The phases are
those of the actual orbit; the coefficients are not freely chosen
roots of unity.

\item\label{N01:input:odd}
\emph{Odd primitive coefficients.}
Theorem~B of \cite{K064} settles coprime odd $r,s$ with $13\nmid r$.
For coprime odd $r,s$ with $13\mid r$, the theorem in
\cite[Section~7.2]{K064}, together with its immediately following
general-$s$ remark, implies that an algebraic minimum $\xi$ itself is
a cubic Pisot number, with nonreal conjugates, $3$ inert in $\Q(\xi)$,
and $\Tr(\xi^{-s})=\pm1$.
\end{enumerate}

We use these as prior results; their full proofs are not replaced by the
new argument in this note. In contrast, all the additional structural
steps and the implications of the finite certificates are proved below.

\begin{lemma}[Normalization]\label{N01:lem:normalization}
For $d\in\N$, one has $\xi(dC)=\xi(C)^{1/d}$. If
$R=r/d$, $S=s/d$ with $d=\gcd(r,s)$, then the asymptotic gcd of
$R3^k-S$ equals $1$ when $R,S$ have opposite parity, and equals $2$
when both are odd.
\end{lemma}
\begin{proof}
The first assertion follows from
$\floor{A^{dC_k}}=\floor{(A^d)^{C_k}}$ and monotonicity of $A\mapsto A^d$.
For the second, the asymptotic gcd divides $2S$ and every $R3^k-S$.
No odd prime can divide it: such a prime would divide $S$ and $R3^k$,
and cannot be $3$ because $3\nmid S$. This contradicts $\gcd(R,S)=1$.
Opposite parity therefore gives gcd $1$. If both are odd, all exponents
are even, but they are not all divisible by $4$: multiplication of the
odd residue $R3^k$ by $3$ interchanges its two residues modulo $4$.
Thus the gcd is exactly $2$.
\end{proof}

\subsection{Inverse trace rigidity at integral negative fixed points}

\begin{lemma}[Uniform inverse trace]\label{N01:lem:inverse}
Let $\beta$ be a cubic Pisot number with $3\nmid\Nm(\beta)$, and let
$r,s\in\N$ with $3\nmid s$. Suppose that
$\Tr(\beta^{r3^k-s})$ is prime for every sufficiently large $k$.
Then
\[
 \Tr(\beta^{-s})=\pm1.
\]
Moreover, in every phase class of the orbit, the three actual twists
$\omega_i^{r3^k}$ are all one common sign $\pm1$.
\end{lemma}

\begin{proof}
Work inside a fixed algebraic closure of $\Q_3$ and identify the
algebraic numbers appearing there with an algebraic closure of $\Q$.
All three $\beta_i$ are $3$-adic units. Put
\[
 a_i=\beta_i^{-s},\qquad q=a_1+a_2+a_3\in\Q,
 \qquad F=\Q(\mu_{104}).
\]
Although $q$ need not initially be an integer, it is rational.
The $a_i$ are pairwise distinct. Indeed, for a Pisot number
$\Q(\beta^s)=\Q(\beta)$: otherwise two embeddings over
$\Q(\beta^s)$ would send $\beta$ to conjugates of the same modulus
$\beta>1$, contrary to the Pisot property. The same degree statement
holds for $\beta^{-s}$.

By input~\ref{N01:input:limits}, the identity at the integral fixed point
$-s$ takes the form
\begin{equation}\label{N01:eq:integral}
 \zeta_1a_1+\zeta_2a_2+\zeta_3a_3=\varepsilon,
 \qquad \zeta_i=\omega_i^{r3^k},\qquad \varepsilon\in\{1,-1\}.
\end{equation}
This follows directly from
$\omega_i^{n_k}u_i^{-s}=\omega_i^{n_k+s}\beta_i^{-s}$.
We prove the conclusion first in the generic Galois case, then treat
the exceptional cubic field.

\paragraph{A cyclic normal-basis argument over $F$.}
Let $M$ be the splitting field of $\beta$ and put $E=MF$.
Unless $M$ is a cyclic cubic subfield of $F$, the group
$\Gal(E/F)$ contains an element $\tau$ cyclically permuting the
$\beta_i$. If $\Gal(M/\Q)=S_3$, the intersection $M\cap F$ is
abelian and Galois over $\Q$, hence is at most quadratic; its fixing
subgroup contains $A_3$. If $\Gal(M/\Q)=C_3$, the intersection is
$\Q$ unless $M\subseteq F$. These assertions also show that the
required three-cycle extends while fixing $F$.

Apply $1,\tau,\tau^2$ to \eqref{N01:eq:integral} and add. Since $\tau$
fixes the $\zeta_i$, this gives
\begin{equation}\label{N01:eq:qnonzero}
 (\zeta_1+\zeta_2+\zeta_3)q=3\varepsilon,
\end{equation}
so $q\ne0$. Consider the relation space
\[
 \mathcal R=\{(v_1,v_2,v_3)\in F^3:\textstyle\sum_i v_i a_i=0\}.
\]
It is stable under cyclic permutation. The field $F$ contains no
primitive cube root of unity; for example, $3$ is unramified in
$\Q(\mu_{104})$, whereas it ramifies in $\Q(\mu_3)$.
Consequently $X^2+X+1$ is irreducible over $F$. The permutation
representation decomposes as
\[
 F^3=U\oplus V,\qquad
 U=F(1,1,1),\qquad V=\{(v_i):\textstyle\sum_i v_i=0\},
\]
where $U$ and $V$ are nonisomorphic simple modules for the three-cycle.
Equivalently, the projections onto these summands are polynomials in
that cycle and its restriction to $V$ has irreducible minimal
polynomial $X^2+X+1$. Every invariant subspace is therefore one of
$0,U,V,F^3$. Equation~\eqref{N01:eq:qnonzero} excludes $U\subseteq\mathcal R$,
and $a_1\ne a_2$ excludes $V\subseteq\mathcal R$. Hence
$\mathcal R=0$.

Subtracting $(\varepsilon/q)(a_1+a_2+a_3)=\varepsilon$ from
\eqref{N01:eq:integral} now gives $\zeta_i=\varepsilon/q$ for every $i$.
This common value is a rational root of unity, hence is $1$ or $-1$.
Thus $q=\pm1$, with the asserted common signed twists.

\paragraph{The exceptional cyclic field of conductor $13$.}
The group
\[
 \Gal(F/\Q)\simeq (\Z/104\Z)^\times\simeq C_2\times C_2\times C_{12}
\]
has a unique index-$3$ subgroup. Its cubic fixed field $K$ is the
real cubic subfield of $\Q(\mu_{13})$, fixed by
$H=\{1,5,8,12\}\subseteq(\Z/13\Z)^\times$.
The class of $3$ has order $3$ and is not in $H$, so $3$ is
unramified and inert in $K$.

We must allow the possibility that the integral generator $\beta$
has nonmaximal index at $3$. Its residue in $\F_{27}$ either belongs
to $\F_3$ or generates $\F_{27}$. In the former case all three
conjugates have the same nonzero residue $e$, and
$\Tr(\beta^n)\equiv3e^n=0\pmod3$ for every positive $n$.
Since these traces tend to infinity, this contradicts eventual
primality. Thus the residue generates $\F_{27}$, and, after labeling,
\[
 \omega_i=\omega^{3^{i-1}},\qquad \ord(\omega)\in\{13,26\}.
\]
If $13\mid r$, each $\omega_i^{r3^k}$ is the same sign, and
\eqref{N01:eq:integral} immediately proves the conclusion.

Suppose instead that $13\nmid r$. Put $\eta=\omega^{r3^k}$.
The three successive phase identities have cyclically permuted
coefficient rows, because $\eta^{27}=\eta$. Adding them yields
\begin{equation}\label{N01:eq:period}
 (\eta+\eta^3+\eta^9)q=\varepsilon_0+\varepsilon_1+\varepsilon_2.
\end{equation}
The right side is an odd, nonzero integer. The Gaussian period on
the left is irrational. For $\eta$ of order $13$, a rational value
would give a polynomial $X^9+X^3+X-c$ over $\Q$ annihilating
$\eta$, impossible because its degree is $12$. For order $26$,
replace $\eta$ by $-\eta$, of order $13$; all three exponents are
odd, so the same argument applies with a common minus sign.
Equation~\eqref{N01:eq:period} first implies $q\ne0$ and then contradicts
$q\in\Q$. This exhausts the exceptional case and proves the lemma.
\end{proof}

\subsection{The extended negative shift theorem}

\begin{proof}[Proof of Theorem~\ref{N01:thm:extension}]
By Lemma~\ref{N01:lem:normalization}, it is enough to take
$\gcd(r,s)=1$ with $r$ odd. For odd $s$, apply input~\ref{N01:input:odd}.
Suppose $s$ is even and that the minimum is algebraic.
Lemma~\ref{N01:lem:normalization} gives $D=1$, so input~\ref{N01:input:shadow}
provides $\beta=\xi$ itself as a cubic Pisot number.
Input~\ref{N01:input:norm} gives $3\nmid\Nm(\beta)$, and
Lemma~\ref{N01:lem:inverse} shows that the three twists
$\omega_i^{r3^k}$ are one common sign. Since $r3^k$ is odd, inspect
the factorization of the minimal polynomial modulo $3$.

\begin{enumerate}[label=\textup{(\roman*)}]
\item If it has three linear factors, allowing repetitions, every
$\omega_i$ is a sign. Equality of their odd powers makes all three
signs equal. Every trace is then $0$ modulo $3$, impossible.
\item If it has a linear factor and an irreducible quadratic factor,
the latter's Teichm\"uller roots are $\omega,\omega^3$, of order
$4$ or $8$. Equality of their $r3^k$-th powers requires
$\omega^{2r3^k}=1$, impossible for odd $r3^k$.
\item If it is irreducible, its roots have Teichm\"uller lifts
$\omega,\omega^3,\omega^9$, of order $13$ or $26$.
Equality requires $13\mid r$, contrary to the hypothesis.
\end{enumerate}
These are all cubic factorization types with nonzero constant term
over $\F_3$. The contradiction proves the theorem.
\end{proof}

\begin{corollary}[A uniform necessary algebraic form]\label{N01:cor:form}
Assume \eqref{N01:eq:scope} and write $d=\gcd(r,s)$, $R=r/d$, $S=s/d$.
If $\xi(r3^k-s)$ is algebraic, then
$\alpha=\xi(r3^k-s)^d=\xi(R3^k-S)$ itself is a cubic Pisot number and
\begin{equation}\label{N01:eq:necessary}
 3\nmid\Nm(\alpha),\qquad
 \Tr(\alpha^{-S})=\pm1,\qquad
 |\Nm(\alpha)|\le\alpha^{3/20}.
\end{equation}
If $R$ is odd, then $13\mid R$, the prime $3$ is inert in $\Q(\alpha)$,
and the two small conjugates are nonreal.
\end{corollary}
\begin{proof}
For opposite parity, the asymptotic gcd is $1$, so the assertions
follow from inputs~\ref{N01:input:shadow}--\ref{N01:input:limits} and
Lemma~\ref{N01:lem:inverse}. If $R$ is odd and $S$ even, the preceding
modulo-$3$ classification forces an irreducible reduction and
$13\mid R$. Irreducibility modulo $3$ gives inertness. Total reality
would make each small conjugate's inverse even power greater than
$1$, so $\Tr(\alpha^{-S})>2$, a contradiction.

If $R,S$ are both odd, input~\ref{N01:input:odd} gives the conclusion
except possibly the displayed norm bound. Its degree descent uses
$\alpha^S$ and $\alpha^2$ in the same cubic field and $\gcd(S,2)=1$;
thus it is a statement about $\alpha$, not merely $\alpha^2$.
Apply \eqref{N01:eq:norm} to $\beta=\alpha^2$ and take square roots to
obtain the bound for $\alpha$.
\end{proof}

In particular, for $S=1$, a hypothetical algebraic minimum has
minimal polynomial
\begin{equation}\label{N01:eq:polynomial}
 f(X)=X^3-aX^2+bX-c,\qquad b/c=\Tr(\alpha^{-1})=\pm1.
\end{equation}
The condition $b=\pm c$ is uniform, not an extrapolation from the
$\beta<9$ table in \cite{K064}.

\subsection{Finite ring certificates and the global cutoff}

\subsubsection{The certificate implication}

For integers $A,B$ put
\[
 F_{A,B}(X)=X^3-AX^2+BX-1,\qquad
 M_{A,B}=\begin{pmatrix}0&0&1\\1&0&-B\\0&1&A\end{pmatrix}.
\]
This is the matrix of multiplication by $X$ in the basis $1,X,X^2$.
Its power traces are the Newton traces of $F_{A,B}$.

\begin{lemma}[A modular certificate]\label{N01:lem:certificate}
Suppose $p,q$ are primes, $d$ is a nonnegative integral power of $3$,
$t=dq$, and $n\ge0$ satisfy
\begin{align}
 &\ord_q(3)=q-1,\qquad M_{A,B}^{t}=I\pmod p,\label{N01:eq:cert1}\\
 &\Tr(M_{A,B}^{n})=0\pmod p,\qquad d\mid n+1,
 \qquad q\nmid n+1.\label{N01:eq:cert2}
\end{align}
If the absolute values of the traces at exponents $r3^k-1$ tend to
infinity and are eventually prime, then $q\mid r$.
\end{lemma}
\begin{proof}
Assume $q\nmid r$. Since $3$ generates $(\Z/q\Z)^\times$, there are
infinitely many $k$, in one residue class modulo $q-1$, with
$r3^k\equiv n+1\pmod q$. For sufficiently large such $k$,
$r3^k\equiv0\equiv n+1\pmod d$ as well. As $q\ne3$,
the Chinese remainder theorem gives $r3^k-1\equiv n\pmod t$.
Equations~\eqref{N01:eq:cert1}--\eqref{N01:eq:cert2} make the corresponding
trace divisible by $p$. Its absolute value eventually exceeds $p$,
contradicting primality.
\end{proof}

\begin{proposition}[Exact finite certificate]\label{N01:prop:finite}
Set $R_0=400\,300\,000\,000\,000$. For every integer pair
\[
 -102\le A\le102,\qquad B\in\{-1,1\},
\]
other than $(-3,1)$, $(-1,-1)$ and $(1,1)$, the supplementary
certificate file contains tuples satisfying
\eqref{N01:eq:cert1}--\eqref{N01:eq:cert2} whose $q$-values have least common
multiple greater than $R_0$. Thus no irreducible polynomial in this
range can have eventually prime absolute traces along $r3^k-1$
with $1\le r\le R_0$ and traces tending in absolute value to infinity.
\end{proposition}
\begin{proof}[Exact verification and deduction]
There are $410$ pairs and $2\,202$ recorded tuples. The independent
verifier described in Section~\ref{N01:app:verification} checks every
condition by integer arithmetic and modular matrix multiplication.
Exactly $407$ rows have the required least common multiple; the
smallest of these least common multiples is
$404\,048\,787\,244\,321>R_0$.
The three exceptions factor as
\begin{align*}
 F_{-3,1}&=(X+1)(X^2+2X-1),\\
 F_{-1,-1}&=(X-1)(X+1)^2,\\
 F_{1,1}&=(X-1)(X^2+1).
\end{align*}
They cannot be minimal polynomials of cubic numbers.
For every other row, Lemma~\ref{N01:lem:certificate} forces each certified
$q$ to divide $r$, hence forces their least common multiple to divide
$r$. This is impossible for $1\le r\le R_0$.
\end{proof}

\subsubsection{Why the finite certificate covers all large parameters}

Put
\begin{equation}\label{N01:eq:constants}
 E=e^{32.76},\qquad x_0=e^{3E},\qquad M=2e^E+1,
 \qquad r_*=12\,230\,000\,000\,000.
\end{equation}
Saito's Type~C theorem gives the following bound for a hypothetical
algebraic minimum:
\begin{equation}\label{N01:eq:global}
 \xi(r3^k-1)\le M^{1/(3r-1)}.
\end{equation}
Indeed, \cite[Theorem~2.6 and its application in Section~2.3]{Saito}
gives $\xi^g\le(2x_0^{1/c_2}+1)^{g/c_1}$, with
$c_1=3r-1$ and $c_2=(9r-1)/(3r-1)>3$; taking the $g$-th root
gives \eqref{N01:eq:global}. This is a global prime-chain bound, not a
finite-prefix estimate.

\begin{proof}[Proof of Theorem~\ref{N01:thm:cutoff}]
For $r\ge R_0$, Saito's established large-$r$ result already applies
\cite[Theorem~1.9(C)]{Saito}. Let $r_*\le r<R_0$, and suppose that
$\alpha=\xi(r3^k-1)$ is algebraic. Corollary~\ref{N01:cor:form} makes
$\alpha$ itself cubic Pisot with \eqref{N01:eq:necessary}.
The rational estimates in Section~\ref{N01:app:rounding} prove
\[
 (3r_*-1)\log100>\log M.
\]
Therefore \eqref{N01:eq:global} gives $\alpha<100$.
Its norm is a nonzero integer and
\[
 |\Nm(\alpha)|\le\alpha^{3/20}<100^{3/20}<2,
 \qquad\text{since }100^3<2^{20}.
\]
Thus $c=\pm1$ in \eqref{N01:eq:polynomial}, and $b=\pm1$.
If $x_2,x_3$ are the small conjugates, then
$-1<\alpha+x_2+x_3<102$, so the integral coefficient satisfies
$0\le a\le101$.

For $c=1$, the polynomial already has the form $F_{a,b}$.
For $c=-1$, put $\gamma=-\alpha$; its polynomial is $F_{-a,b}$,
and
\[
 \Tr(\gamma^n)=(-1)^n\Tr(\alpha^n).
\]
Consequently all zero-trace certificates transfer, regardless of
exponent parity. Both signs are covered by Proposition~\ref{N01:prop:finite}.
The traces of $\alpha^{r3^k-1}$ are the eventual prime floors from
input~\ref{N01:input:shadow}; for the odd--odd descent the field is
unchanged, so these traces are unchanged as well. They tend to
infinity. The finite certificate now contradicts $r<R_0$.
\end{proof}

The unrounded sufficient threshold in this argument is
\[
 \frac{\log M}{3\log100}+\frac13.
\]
The stated integer cutoff is a conservative rationally certified
rounding, not an optimized endpoint.

\subsection{What remains open}

Neither theorem settles all parameters in \eqref{N01:eq:scope}.
For $s=1$, the cases not settled by the present conclusions can
still include even $r<r_*$ and odd multiples of $13$ below $r_*$.
This section alone does not settle $\xi(2\cdot3^k-1)$ or $\xi(13\cdot3^k-1)$; the latter is settled in Section~\ref{src:R02}, while the former remains open.
The next proposition explains one limitation of simply extending
a fixed modular sieve.

Let $\alpha_a$ be the dominant root of
$F_a=X^3-aX^2-X-1$ for $a\ge1$. These are cubic Pisot units with
nonreal conjugates and $\Tr(\alpha_a^{-1})=-1$; see
\cite[Section~7.1]{K064}. Put $T_n(a)=\Tr(\alpha_a^n)$.
The residual families in \cite[Section~7]{K064} use
$a\equiv2\pmod3$ for even $r$, and $a\equiv1\pmod3$ for odd
$r$ divisible by $13$. They satisfy the cited $3$-adic necessary
conditions; this does not assert prime traces or the full shadow condition.

\begin{proposition}[A finite fixed-prime sieve cannot cover the residual family]
Fix an even $r$, or an odd $r$ divisible by $13$. For every finite
set of primes $\mathcal P$, infinitely many parameters $a$ in the
corresponding residual family have
\[
 p\nmid T_{r3^k-1}(a)\qquad
 \text{for all }p\in\mathcal P\text{ and all }k\ge1.
\]
\end{proposition}
\begin{proof}
The exact specialization
$F_{-1}=(X-1)(X+1)^2$ gives, by Newton traces,
\begin{equation}\label{N01:eq:barrier}
 T_n(a)\equiv1+2(-1)^n\pmod{a+1}.
\end{equation}
For even $r$, choose infinitely many positive
$a\equiv-1\pmod{3\prod_{p\in\mathcal P}p}$. Then
$a\equiv2\pmod3$ and $r3^k-1$ is odd, so every tested trace is
$-1$ modulo every tested prime.

For odd $13\mid r$, use the Chinese remainder theorem to choose
$a\equiv1\pmod3$ and $a\equiv-1\pmod p$ for each
$p\in\mathcal P\setminus\{3\}$. The exponents are even, so the
traces are $3$ modulo those primes. Modulo $3$ the polynomial is
$X^3-X^2-X-1$, irreducible, and $X^{13}=1$ in its residue field.
Hence $13\mid r$ gives
$T_{r3^k-1}(a)\equiv T_{-1}(a)=-1\pmod3$ as well.
All required divisibilities fail.
\end{proof}

This is a limitation of finite \emph{fixed-prime divisibility covers},
not a counterexample to transcendence. Parameter-dependent primes,
finer use of the real shadow, or additional global bounds remain
possible. The finite-ring certificate above is compatible with this
barrier because it is applied only after bounding the candidate
coefficients.


\subsection{Reproducible finite verification}\label{N01:app:verification}

The proof-critical supplementary files are
\begin{itemize}
\item \texttt{finite\_modular\_scan\_signed102.json}, containing all
$410$ rows and $2\,202$ tuples;
\item \texttt{bounds\_verify\_certificates.py}, an independent verifier;
\item \texttt{bounds\_roundoff.py}, containing the rational numerical checks.
\item \texttt{cutoff\_audit\_checks.py}, checking completeness of the
signed grid and the small Pisot gap used in Section~\ref{N01:app:general}.
\end{itemize}
For each row, the verifier performs the following exhaustive operations.
It tests $p$ and $q$ by trial division; checks $t=dq$ and that repeated
division of $d$ by $3$ ends at $1$; factors $q-1$ and checks
\[
 3^{q-1}\equiv1\pmod q,\qquad
 3^{(q-1)/\ell}\not\equiv1\pmod q
 \quad(\ell\text{ prime},\ \ell\mid q-1);
\]
and uses binary matrix powering modulo $p$ to check
$M_{A,B}^{t}=I$ and $\Tr(M_{A,B}^{n})=0$. It then checks the two
conditions on $n+1$, recomputes the least common multiple of the
$q$-values, and verifies that the only failed rows are the three
explicitly reducible rows of Proposition~\ref{N01:prop:finite}.
The additional check verifies that the row keys are exactly the full
$410$-element signed grid, with no duplicates or omissions.
No floating-point root approximation or probable-prime test is used.
The matrix implementation is separate from the generating program's
quotient-ring representation.

The SHA-256 digest of the proof-critical JSON data is
\begin{center}\small\ttfamily
c97b6693c8057e73c90641f5f30d6ad0\\
2167f3e83b82b4c428c342866f1840a1.
\end{center}
The line break and final punctuation are not part of the digest.
The verifier's terminal certificate is
\begin{quote}\small\ttfamily
INDEPENDENT MATRIX CHECK PASSED 410 rows 2202 certificates;\newline
only reducible failures [(-3, 1), (-1, -1), (1, 1)].
\end{quote}
These are exact computational checks, not a claim of external peer review.

As a nonessential cross-check, a separate exact Sturm enumeration of
a safe superset of cubic Pisot units below $100$ produced $10\,277$
candidates. Necessary twisted-limit tests at precisions $3^{12}$
and $3^{22}$ left the same $394$ odd-residue candidates, all with
$|b|=1$, and all mapped to certified signed rows. The proof above
does not depend on inferring any exact limit identity from such
finite-precision agreement.

\subsection{Rational certification of the cutoff}\label{N01:app:rounding}

Set $x=819/25$ and $N=250$. The Taylor sum with geometric tail bound
\[
 e^x\le \sum_{j=0}^{N}\frac{x^j}{j!}
 +\frac{x^{N+1}}{(N+1)!}\frac{1}{1-x/(N+2)}
 <168\,845\,000\,000\,000
\]
is verified over $\Q$. For $0<z<1$, the positive expansion
\[
 \log\frac{1+z}{1-z}
 =2\sum_{j=0}^{\infty}\frac{z^{2j+1}}{2j+1}
\]
gives lower bounds by truncation. Taking $20$ terms at $z=1/3$ and
$z=1/9$, and using $\log100=6\log2+2\log(5/4)$, proves
$\log100>460517/100000=4.60517$.
Since $E=e^{32.76}>2$, one has $2+e^{-E}<e$, and therefore
$\log M=E+\log(2+e^{-E})<E+1$.
Finally the exact positive difference is
\[
 (3r_*-1)\frac{460517}{100000}
 -168\,845\,000\,000\,001
 =\frac{11\,868\,729\,999\,439\,483}{100000}>0.
\]
Combining these strict inequalities gives
$(3r_*-1)\log100>\log M$, as required.

\subsection{A sufficient bound for general negative shifts}\label{N01:app:general}

\begin{theorem}\label{N01:thm:general}
Let $r,s\in\N$, $3\nmid s$, and $s<3r$. Put
$d=\gcd(r,s)$, $R=r/d$, $S=s/d$, and
\[
 D=\begin{cases}2,&R,S\text{ both odd},\\1,&R,S\text{ of opposite parity}.
 \end{cases}
\]
If
\begin{equation}\label{N01:eq:general-cutoff}
 \frac{9R-S}{D}\ge Q_0:=2\,924\,000\,000\,000\,000,
\end{equation}
then $\xi(r3^k-s)$ is transcendental. In particular, all normalized
$R\ge975\,000\,000\,000\,000$ are covered; for opposite parity,
$R\ge487\,500\,000\,000\,000$ suffices.
\end{theorem}

\begin{proof}
By normalization, take the coprime schedule $R3^k-S$ and write
$A=\xi(R3^k-S)$. Set
\[
 H=3R-S>0,\quad J=9R-S,\quad N=H/D\in\N,
 \quad c_2=J/H,\quad L=\log x_0=3e^{32.76}.
\]
Suppose $A$ is algebraic. The reduction gives a cubic Pisot number
$A^g$ with $g\mid D$, so $A^D$ and $A^H$ are cubic Pisot numbers
by the power-field property. This observation does not require
the degree descent in Corollary~\ref{N01:cor:form}.

The exact rational checks give
\[
 Q_0\frac{693147}{1000000}
 >12\cdot168\,845\,000\,000\,000,
 \qquad \log2>\frac{693147}{1000000}.
\]
Together with Section~\ref{N01:app:rounding}, condition
\eqref{N01:eq:general-cutoff} therefore implies
\[
 \frac JD>\frac{4L}{\log2},\qquad
 x_0^{1/c_2}=\exp(LH/J)<2^{N/4}.
\]
Let $\kappa$ be the smallest Pisot number, the real root of
$X^3-X-1$ \cite{Siegel}; one has $\kappa>k_0:=331/250$.

For $N=4,5,6,7,8,9,10$, respectively choose the prime
\[
 p=2,3,3,5,5,5,7.
\]
Exact rational inequalities give $p^4\ge2^N$ and
$p+1<k_0^N<\kappa^N$. For $N\ge11$, Bertrand's postulate gives
$2^{N/4}<p<2\cdot2^{N/4}$. The rational inequalities
\[
 3^4 2^{11}<k_0^{44},\qquad 2<k_0^4
\]
imply $p+1<3\cdot2^{N/4}<\kappa^N$ for every such $N$.
In either case $p>x_0^{1/c_2}$. Saito's nested prime-chain
construction \cite[Lemma~9.1]{Saito} produces a prime-representing
number with initial prime $p$, so minimality gives
\[
 A^D\le(p+1)^{D/H}=(p+1)^{1/N}<\kappa,
\]
contradicting the Pisot property.

It remains to handle $1\le N\le3$, where a crude initial-prime
bound would not suffice. Here $p=2>x_0^{1/c_2}$ and $c_2>100$.
Choose the first prime child $q$ using the explicit prime interval:
\[
 2^{c_2}\le q\le2^{c_2}+3\cdot2^{2c_2/3}.
\]
Continue this same child by the nested construction, rather than
replacing the original problem by an unrelated tail. It gives
\[
 2\le A^H\le(q+1)^{1/c_2}<\frac{21}{10}.
\]
For the last strict inequality, for every $c\ge100$,
\[
 1+3\cdot2^{-c/3}+2^{-c}<2<(21/20)^c.
\]
But there is no cubic Pisot number in $[2,21/10]$. To verify this
small gap exhaustively, the root bounds leave only
\[
 f=X^3-aX^2+bX-c,\qquad
 0\le a\le4,\quad |b|\le5,\quad c\in\{-2,-1,1,2\}.
\]
The supplementary rational Sturm verifier checks all $220$ triples.
It first rejects rational roots (possible integral roots divide
$c$), then counts roots in $(2,21/10)$ and $(-1,1)$.
For negative discriminant, a root $\beta\in(2,21/10)$ can be Pisot
only for $c>0$, and then its conjugate-pair modulus is
$\sqrt{c/\beta}<1$; for positive discriminant it must have two
roots in $(-1,1)$. No irreducible polynomial passes either test.
The endpoints cannot be roots of an irreducible monic cubic.
An independent Fraction-only implementation and a separate symbolic
Sturm implementation both verify the empty list. This contradicts
the fact that $A^H$ is cubic Pisot.

Finally $J=9R-S>6R$ and $D\le2$, giving the stated simpler bounds
on $R$ by direct comparison with $Q_0$.
\end{proof}

Since $1\le S<3R$, a bound on the remaining normalized $R$ leaves
only finitely many coprime schedule pairs $(R,S)$ within the stated scope. This is not a finite enumeration of all cubic Pisot candidates. Arbitrarily large common
scalings $d$ are still possible, but scaling preserves algebraicity.
The theorem does not enumerate or resolve that finite residual set.



\endgroup

\section{The multiplier thirteen and a residual branch exclusion}\label{src:R02}
This section closes the multiplier-thirteen case left by the basic negative theorem. Its final branch exclusion is only one part of the even-multiplier residual problem, not a proof for all candidates.
\medskip
\begingroup



\subsection{Main results and their scope}

For a sequence $C=(C_k)_{k\ge1}$ of positive integers, put
\[
 \mathcal W(C)=\{A>1:\ \floor{A^{C_k}}\text{ is prime for every }k\ge1\}.
\]
If $C_{k+1}\ge2C_k$, this set is nonempty and has a least element
\cite[Theorem~1.3]{Saito}; denote it by $\xi(C)$.
Our principal schedule is
\begin{equation}\label{R02:eq:schedule}
 n_k=13\cdot3^k-1\qquad(k\ge1),
 \qquad n_1=38,\quad n_2=116.
\end{equation}
Its consecutive ratios tend to $3$ from above.

\begin{mainthm}\label{R02:thm:main}
The number $\xi((13\cdot3^k-1)_{k\ge1})$ is transcendental.
\end{mainthm}

The proof uses the structural theorem for odd multiples of $13$ in
\cite[Section~7.2, Theorem~\textup{\texttt{D13}}]{K064}, whose precise
consequences are recorded in Section~\ref{R02:sec:reduction}. That theorem
is a result of the author's earlier, unpublished manuscript; it is an
explicit dependency, not an assertion proved anew in this paper.
The underlying algebraic-minimum reduction is Saito's
\cite[Proposition~3.1]{Saito}. Our citation numbers for Saito refer to
arXiv version~3; the corresponding reference in [064] is numbered
Proposition~3.2 there.

The added steps are an analytic bound from an initial prime floor and a
complete modular exclusion. The analytic bound applies to every
hypothetical algebraic minimum, not merely to a numerical search
range. The modular exclusion proves infinitely recurring divisibility,
not merely compositeness in a finite initial segment. No prime-gap
conjecture is assumed. The additional finite computation is independent
of numerical root approximations and probabilistic primality tests.

We also obtain the following result, which does not require the
structural theorem for multiples of $13$.

\begin{mainthm}\label{R02:thm:fifth}
Let $\beta>1$ be a cubic Pisot number with minimal polynomial
\[
 f(X)=X^3-aX^2+cX-c,\qquad a,c\in\Z,\quad c>0.
\]
If $c\le\beta^{3/20}$, then $\floor{\beta^5}$ is composite.
\end{mainthm}

The exponent $5$ is the first exponent of $2\cdot3^k-1$.
Thus Theorem~\ref{R02:thm:fifth} excludes the entire positive-norm branch
$b=c$ in the algebraic-minimum problem for this schedule, including
nonunits. Section~\ref{R02:sec:r2} specifies the remaining branches. We
make no claim here that $\xi(2\cdot3^k-1)$, the classical Mills
constant, or all negative-shift minima have been resolved.

\subsection{The structural reduction and its residue-field consequence}
\label{R02:sec:reduction}

A Pisot number is a real algebraic integer $\beta>1$ all of whose other
conjugates have modulus strictly less than one. Write its cubic minimal
polynomial and its power traces as
\begin{equation}\label{R02:eq:poly}
 f(X)=X^3-aX^2+bX-c,\qquad
 T_n=\Tr_{\Q(\beta)/\Q}(\beta^n).
\end{equation}
Newton's identities give
\begin{equation}\label{R02:eq:recurrence}
 \begin{gathered}
 T_0=3,\qquad T_1=a,\qquad T_2=a^2-2b,\\
 T_n=aT_{n-1}-bT_{n-2}+cT_{n-3}\qquad(n\ge3).
 \end{gathered}
\end{equation}

\begin{proposition}[Structural input from {\cite{K064}}]
\label{R02:prop:input}
If $\xi=\xi((n_k)_{k\ge1})$ is algebraic, then $\beta=\xi$ itself
is a cubic Pisot number with nonreal conjugates $u,v=\overline u$.
For its minimal polynomial \eqref{R02:eq:poly}, one has
\begin{equation}\label{R02:eq:structural}
 b=\varepsilon c,\qquad \varepsilon\in\{-1,1\},\qquad
 0<c\le\beta^{3/20},\qquad 3\nmid c.
\end{equation}
The polynomial $f$ is irreducible modulo $3$, and
\begin{equation}\label{R02:eq:eventual}
 T_{n_k}=\floor{\beta^{n_k}}\text{ is prime for every sufficiently large }k.
\end{equation}
Independently of the eventual identity, the definition of $\xi$ gives
\begin{equation}\label{R02:eq:allfloors}
 \floor{\beta^{n_k}}\text{ is prime for every }k\ge1.
\end{equation}
\end{proposition}

\begin{proof}[Explanation of the cited reduction and the mod $3$ bridge]
The Mills-shadow corollary in \cite[Section~2]{K064}, based on Saito's
\cite[Proposition~3.1]{Saito}, first produces the cubic Pisot number
$B=\xi^2$, with eventual prime traces at exponents $n_k/2$.
The theorem for odd multiples of $13$ in
\cite[Section~7.2]{K064} proves that $\xi$ itself is cubic Pisot,
with nonreal conjugates, and that
\[
 \Q(\xi)=\Q(B),\qquad \Tr(\xi^{-1})=\pm1,
 \qquad \xi\ge|\Nm(\xi)|^{20/3}.
\]
Since $\Tr(\xi^{-1})=b/c$ and
$c=\beta u\overline u>0$, these statements give the coefficient and
norm assertions in \eqref{R02:eq:structural}. Equality of the fields gives
\[
 \Tr_{\Q(B)/\Q}(B^{n_k/2})
 =\Tr_{\Q(\xi)/\Q}(\xi^{n_k}),
\]
so the eventual prime-trace identity transfers to \eqref{R02:eq:eventual}
with $\beta=\xi$, not merely with $\beta=\xi^2$.

There is one residue-field detail that is essential for the finite
coefficient cover. In the proof of the cited theorem, the minimal
polynomial of $B=\xi^2$ is shown to be irreducible modulo $3$.
Consequently $3$ is inert in $K=\Q(B)=\Q(\xi)$, its residue field is
$\F_{27}$, and the residue $\overline B$ generates $\F_{27}$ over
$\F_3$. The algebraic integer $\xi$ has a residue $\overline\xi$ in
this same field and $\overline\xi^{\,2}=\overline B$. Hence
\[
 \F_3(\overline B)\subseteq\F_3(\overline\xi)\subseteq\F_{27},
\]
and both inclusions are equalities. The reduction of the monic cubic
minimal polynomial of $\xi$ annihilates an element of residue degree
$3$, so it is irreducible. In particular its constant coefficient is
nonzero modulo $3$, giving $3\nmid c$.
This uses the stronger square-generator statement from the cited proof:
inertness alone would not justify irreducibility modulo $3$ for an
arbitrary integral generator.
\end{proof}

Theorem~\ref{R02:thm:main} uses only this input from [064] and Saito,
together with the arguments below. In particular it does not depend on
the more general inverse-trace theorem in \cite{N01}, which is used
only to describe the $r=2$ application in Section~\ref{R02:sec:r2}.

\subsection{Parity and the second prime floor}

We first isolate the parity obstruction. It will apply to both
schedules used below.

\begin{lemma}\label{R02:lem:parity}
Let $\beta$ be a cubic Pisot number as in \eqref{R02:eq:poly}, with
$b\equiv c\pmod2$. Suppose that, for $r\in\{2,13,39\}$, the traces
$T_{r3^k-1}$ are prime for all sufficiently large $k$. Then $a$ is odd,
and $T_n$ is odd for every $n\ge0$.
\end{lemma}
\begin{proof}
The Pisot property gives $T_n=\beta^n+o(1)$, so these traces tend to
infinity. If $a$ and $c$ are even, then $b$ is even, and
\eqref{R02:eq:recurrence} makes every $T_n$, $n\ge1$, even. This contradicts
eventual primality.

If $a$ is even and $c$ is odd, then $b$ is odd, and the polynomial
modulo $2$ is $X^3+X+1$. It is irreducible over $\F_2$, because it
has no root there. Its nonzero roots belong to $\F_8^\times$, a group
of order $7$. Thus its power traces have period dividing $7$; the
residues for $n=0,\ldots,6$ are
\[
 (T_n\bmod2)_{n=0}^{6}=(1,0,0,1,0,1,1).
\]
For $r=13$, $r3^k-1\equiv4\pmod7$ whenever $k\equiv2\pmod6$.
For $r=2$, the same residue occurs whenever $k\equiv3\pmod6$,
and for $r=39$ it occurs whenever $k\equiv1\pmod6$.
In every case infinitely many arbitrarily large required traces are
even, again impossible.

Therefore $a$ is odd. The initial values $T_0,T_1,T_2$ are all odd.
Since $a-b+c\equiv1\pmod2$, recurrence \eqref{R02:eq:recurrence} now
proves by induction that every trace is odd.
\end{proof}

Suppose henceforth that the hypothetical algebraic minimum of
Proposition~\ref{R02:prop:input} exists. Put
\begin{equation}\label{R02:eq:qs}
 q=uv=\frac c\beta,\qquad s=u+v=a-\beta.
\end{equation}
Vieta's identity $b=\beta s+q$, together with $b=\varepsilon c$, gives
\begin{equation}\label{R02:eq:sformula}
 s=q\left(\varepsilon-\frac1\beta\right),
 \qquad 0<q\le\beta^{-17/20}<1.
\end{equation}
Choose $\theta\in(0,\pi)$ so that
$u=\sqrt q\,e^{i\theta}$ and $v=\sqrt q\,e^{-i\theta}$. Then
\begin{equation}\label{R02:eq:shadow}
 \cos\theta=\frac{\sqrt q}{2}
       \left(\varepsilon-\frac1\beta\right),\qquad
 S_n:=u^n+v^n=2q^{n/2}\cos(n\theta),\qquad
 \beta^n=T_n-S_n.
\end{equation}

\begin{proposition}\label{R02:prop:bound}
The algebraic minimum in Proposition~\ref{R02:prop:input} would satisfy
\[
 \beta<8192.
\]
\end{proposition}
\begin{proof}
Suppose instead that $\beta\ge8192=2^{13}$. The norm bound gives
\[
 0<q\le\beta^{-17/20}\le2^{-221/20}<\frac1{2048}.
\]
As $45^2=2025<2048$, this implies $\sqrt q<1/45$ and hence
\[
 t:=|\cos\theta|
 \le\frac{\sqrt q}{2}\left(1+\frac1\beta\right)
 <\frac{1+1/8192}{90}<\frac1{89}.
\]
Set $\delta=|\theta-\pi/2|=\arcsin t$. By integrating the increasing
function $(1-x^2)^{-1/2}$ on $[0,t]$, we obtain
\[
 \delta\le\frac{t}{\sqrt{1-t^2}}<\frac1{88},
\]
where $\sqrt{1-1/89^2}>88/89$ was used. Consequently
\[
 116\delta<\frac{116}{88}=\frac{29}{22}<\frac32<\frac\pi2.
\]
The exponent $116$ is divisible by $4$, so
$\cos(116\theta)=\cos(116\delta)>0$. Formula \eqref{R02:eq:shadow}
therefore gives the exact inequalities
\[
 0<S_{116}=2q^{58}\cos(116\delta)<2\left(\frac1{2048}\right)^{58}<1.
\]
Since $T_{116}$ is an integer,
\begin{equation}\label{R02:eq:floor116}
 \floor{\beta^{116}}=T_{116}-1.
\end{equation}
Lemma~\ref{R02:lem:parity}, applied to the eventual prime traces from
\eqref{R02:eq:eventual}, makes this floor even. It exceeds $2$, since
$\beta\ge8192$. But $116=n_2$ is an actual required prime exponent,
contradicting \eqref{R02:eq:allfloors}.
\end{proof}

\begin{remark}
The proof does not assume the eventual identity
$\floor{\beta^{n_k}}=T_{n_k}$ at $k=2$. Instead, it proves the
opposite rounding identity \eqref{R02:eq:floor116} directly under the
hypothesis $\beta\ge8192$. Eventual prime traces justify the parity
lemma; the definition of the minimum justifies primality of this
particular early floor. These are different uses of the reduction.
\end{remark}

\begin{corollary}\label{R02:cor:grid}
The minimal polynomial of the hypothetical algebraic minimum belongs
to the finite set
\begin{equation}\label{R02:eq:grid}
 \mathcal G=\left\{X^3-aX^2+bX-c:
 \begin{array}{l}
 1\le a\le8191,\quad a\text{ odd},\\
 c\in\{1,2\},\quad b\in\{-c,c\},\\
 f\bmod3\text{ irreducible}
 \end{array}\right\}.
\end{equation}
This set contains exactly $8192$ polynomials.
\end{corollary}
\begin{proof}
From Proposition~\ref{R02:prop:bound},
\[
 1\le c\le\beta^{3/20}<8192^{3/20}=2^{39/20}<4.
\]
Since $3\nmid c$, necessarily $c=1$ or $c=2$.
If $b=-c$, then \eqref{R02:eq:sformula} gives $s<0$ and $a<\beta<8192$.
If $b=c$, it gives $0<s=q(1-1/\beta)<1$, so
$a<\beta+1<8193$. In both cases oddness forces $a\le8191$.
Moreover $|s|<2$, because both small conjugates have modulus less
than one. Hence $a=\beta+s>\beta-2>-1$. As an odd integer, $a\ge1$.
Irreducibility modulo $3$ was established in
Proposition~\ref{R02:prop:input}.

For a cubic over a field, reducibility is equivalent to having a root.
Testing $0,1,2$ over $\F_3$ gives the residue classes and counts in
Table~\ref{R02:tab:grid}. Their sum is $8192$.
\end{proof}

\begin{table}[ht]
\centering
\begin{tabular}{rrrr}
\toprule
$c$ & $b$ & Allowed $a\pmod3$ & Odd $a$ in $[1,8191]$\\
\midrule
$1$ & $-1$ & $0,1$ & $2731$\\
$1$ & $1$ & $2$ & $1365$\\
$2$ & $-2$ & $1$ & $1366$\\
$2$ & $2$ & $0,2$ & $2730$\\
\midrule
\multicolumn{3}{r}{Total} & $8192$\\
\bottomrule
\end{tabular}
\caption{The complete coefficient set after irreducibility modulo $3$.}
\label{R02:tab:grid}
\end{table}

No further root test, norm cutoff, or test of the Pisot property is used
to shrink this grid. Eliminating this deliberately larger set suffices.

\subsection{Finite certificates for infinitely recurring divisors}
\label{R02:sec:certificates}

For any monic cubic $f=X^3-aX^2+bX-c$ with integer coefficients,
let $T_n(f)$ denote the integer sequence defined by
\eqref{R02:eq:recurrence}. It agrees with the sum of the $n$th powers of
the three roots, counted with multiplicity. Over $\F_p$, put
\[
 R_p=\F_p[X]/(f),\qquad x=X\bmod f.
\]
The matrix for multiplication by $x$ in the ordered basis
$1,x,x^2$ is
\begin{equation}\label{R02:eq:matrix}
 M_f=\begin{pmatrix}0&0&c\\1&0&-b\\0&1&a\end{pmatrix}\pmod p,
 \qquad \tr(M_f^n)\equiv T_n(f)\pmod p.
\end{equation}
The trace identity follows, for example, from triangularizing the
matrix over an algebraic closure, with algebraic multiplicities. It
does not require $f\bmod p$ to be squarefree.

\begin{lemma}[Return certificate]\label{R02:lem:return}
Let $p$ be a prime. Suppose there are positive integers $E,U,k,h$
and an integer $v\ge0$ such that
\begin{equation}\label{R02:eq:certconditions}
 \begin{gathered}
 E=3^vU,\quad \gcd(U,3)=1,\quad k\ge v,\quad k\ge1,\\
 M_f^E=I\pmod p,\qquad 3^h\equiv1\pmod U,\\
 n\equiv13\cdot3^k-1\pmod E,\qquad
 \tr(M_f^n)\equiv0\pmod p,
 \quad 0\le n<E.
 \end{gathered}
\end{equation}
Then
\begin{equation}\label{R02:eq:return}
 p\mid T_{13\cdot3^{k+jh}-1}(f)\qquad\text{for every }j\ge0.
\end{equation}
\end{lemma}
\begin{proof}
For every $j\ge0$,
\[
 (13\cdot3^{k+jh}-1)-(13\cdot3^k-1)
 =13\cdot3^k(3^{jh}-1).
\]
Its right-hand side is divisible by $3^v$, since $k\ge v$, and by
$U$, since $3^h\equiv1\pmod U$. The two factors $3^v$ and $U$ are
coprime, so it is divisible by $E$. The identity $M_f^E=I$ therefore
identifies the corresponding matrix powers with $M_f^n$. Taking the
trace proves \eqref{R02:eq:return}.
\end{proof}

\begin{remark}
The hypothesis $M_f^E=I$ is checked in the full matrix algebra. It
allows repeated factors modulo $p$ and needs no exclusion of primes
dividing the discriminant. Nor is it necessary that $p\nmid13$.
The witness prime $13$ is as valid as the other primes in the cover.
\end{remark}

\begin{theorem}[Exact finite exclusion]\label{R02:thm:finite}
For every polynomial $f\in\mathcal G$, there are a prime
\[
 p\in\{5,7,11,13,17,19,23,29,31,37,41,43,47\}
\]
and positive integers $k,h$ satisfying \eqref{R02:eq:return}.
Consequently no polynomial in $\mathcal G$ is the minimal polynomial
of a Pisot number whose traces on \eqref{R02:eq:schedule} are eventually
prime.
\end{theorem}
\begin{proof}
The accompanying exact data file \texttt{r13\_scan.json} contains
$230$ certificates. Each stores
\[
 (p,a_0,b,c,k,E,U,n,h),\qquad a_0\equiv a\pmod p,
\]
with $h$ stored under the field name \texttt{period}.
For every certificate the verifier checks primality of $p$ by trial
division, the coefficient conditions, and all identities in
\eqref{R02:eq:certconditions}. Exponentiation is repeated squaring using
integer arithmetic modulo $p$ or modulo $U$.

The file also assigns each polynomial of $\mathcal G$ to one
certificate with exactly the same reduced coefficients modulo its
witness prime. The verifier independently constructs $\mathcal G$
from \eqref{R02:eq:grid}, checks the coefficient match for every
assignment, rejects duplicate triples, and checks equality of the
expected and assigned sets. Thus it verifies a complete cover,
not just the validity of a listed collection of local witnesses.
Section~\ref{R02:app:verification} gives a complete matrix verifier.
Its output is
\begin{quote}\ttfamily
PASS: 230 independent matrix certificates cover all 8192 candidates.
\end{quote}
Lemma~\ref{R02:lem:return} proves the first assertion.
If $f$ were the minimal polynomial of a Pisot number $\beta>1$,
then $T_n(f)=\beta^n+o(1)\to\infty$. Along the infinite progression
in \eqref{R02:eq:return}, all sufficiently late traces are larger than
$p$ and divisible by $p$. They are therefore composite, contradicting
eventual primality.
\end{proof}

\begin{table}[ht]
\centering
\begin{tabular}{r|rrrrrrr}
\toprule
Witness prime & $5$ & $7$ & $11$ & $13$ & $17$ & $19$ & $23$\\
Remaining & $3549$ & $1911$ & $1113$ & $664$ & $391$ & $230$ & $112$\\
\midrule
Witness prime & $29$ & $31$ & $37$ & $41$ & $43$ & $47$ & \\
Remaining & $77$ & $33$ & $18$ & $8$ & $1$ & $0$ & \\
\bottomrule
\end{tabular}
\caption{Successive coverage counts in the exact modular exclusion.}
\label{R02:tab:sieve}
\end{table}

\subsubsection{Discovery and verification}

For clarity, the search procedure and the mathematical verification
have separate roles. The search uses quotient-ring coefficient triples
and iterates
\[
 z_k=x^{13\cdot3^k},\qquad z_{k+1}=z_k^3.
\]
If $z_k=z_0+z_1x+z_2x^2$, then $p\nmid c$ gives
\[
 \Tr_{R_p/\F_p}(x^{-1}z_k)
 =\frac bc z_0+3z_1+az_2.
\]
Here $b/c=\pm1$. A zero of this expression suggests a certificate.
The search chooses
\[
 E=p^2\lcm(p-1,p^2-1,p^3-1),\qquad
 U=E/3^{v_3(E)},\qquad h=\varphi(U).
\]
Neither formula is trusted as an unverified premise: the verification
checks $M_f^E=I$ and $3^h\equiv1\pmod U$ directly for every witness.
Similarly, stopping the search at a finite number of indices causes
no logical problem, because every candidate receives a verified
infinite-return certificate. The largest seed index used is $81$.

Three separately implemented checks of the same finite cover are
included in the supplement:
\begin{enumerate}[label=\textup{(\roman*)}]
\item \texttt{r13\_verify.py} uses $3\times3$ matrix multiplication.
\item \texttt{audit\_r13\_certificates.py} uses polynomial multiplication
followed by explicit long division. It also checks the first two
returned exponent residues and raises explicit exceptions on failure.
\item \texttt{r13\_verify\_direct\_cycles.py} enumerates the triple-state
cycles of recurrence \eqref{R02:eq:recurrence} modulo each witness prime,
without matrix or quotient-ring exponentiation. It checks the zero
trace and the affine return against those cycle periods, and builds
the coefficient grid from Table~\ref{R02:tab:grid}.
\end{enumerate}
All three checks pass. The recurrence verification enumerates
$662665$ trace states in total, with largest period $50652$.
These are exact reproducibility checks, not formal verification in a
proof assistant or external peer review.

\subsubsection{Two explicit certificates}

For $(a,b,c)=(1857,2,2)$, one certificate is
\[
 p=5,\quad a_0=2,\quad k=2,\quad
 E=18600=3\cdot6200,\quad U=6200,\quad n=116,\quad h=2400.
\]
The matrix \eqref{R02:eq:matrix} satisfies
\[
 M_f^{18600}=I\pmod5,\qquad \tr(M_f^{116})=0\pmod5,
 \qquad 3^{2400}\equiv1\pmod{6200}.
\]
Thus $5$ divides the trace at every index $k=2+2400j$.
For $(a,b,c)=(935,1,1)$, the same $p,E,U,h$ work with
$k=5$ and $n=3158$. These examples illustrate why a certificate
excludes eventual prime traces even if its seed is small.

\subsection{Proof of the transcendence theorem}

\begin{proof}[Proof of Theorem~\ref{R02:thm:main}]
Suppose that $\xi=\xi((13\cdot3^k-1)_{k\ge1})$ is algebraic.
By Proposition~\ref{R02:prop:input}, $\beta=\xi$ is a cubic Pisot number
with the structural properties \eqref{R02:eq:structural}, irreducible
minimal polynomial modulo $3$, and eventual prime traces
$T_{13\cdot3^k-1}$.

Lemma~\ref{R02:lem:parity} forces $a$ to be odd.
Proposition~\ref{R02:prop:bound} uses the actual prime floor at exponent
$116$ to force $\beta<8192$. Corollary~\ref{R02:cor:grid} then puts the
minimal polynomial in the exhaustive finite set $\mathcal G$.
Theorem~\ref{R02:thm:finite} supplies a fixed prime dividing infinitely
many of the required traces. Those traces tend to infinity and so
cannot all be prime from any point onward. This contradicts
Proposition~\ref{R02:prop:input} and proves transcendence.
\end{proof}

The proof has two distinct finite ingredients. The exponent $116$
is used analytically to bound all possible algebraic minima. The
finite cover is then used arithmetically to exclude that entire bounded
set by infinite progressions of divisors. Neither ingredient is a
substitute for the other.

\begin{corollary}\label{R02:cor:39}
The number $\xi((39\cdot3^k-1)_{k\ge1})$ is transcendental.
\end{corollary}
\begin{proof}
The structural theorem for odd multiples of $13$ in
\cite[Section~7.2]{K064} applies also to $39$. If this minimum were
algebraic, it would itself be a cubic Pisot number with nonreal conjugates with
$b=\pm c$, $0<c\le\beta^{3/20}$, and eventual prime traces on
$39\cdot3^k-1$. The same square-generator argument as in
Proposition~\ref{R02:prop:input} gives irreducibility modulo $3$.
Lemma~\ref{R02:lem:parity} makes $a$ odd.

Here $116=39\cdot3^1-1$ is the \emph{first} required exponent.
Thus the proof of Proposition~\ref{R02:prop:bound} applies unchanged:
if $\beta\ge8192$, the actual floor at exponent $116$ is an even
integer greater than $2$. Hence $\beta<8192$, and
Corollary~\ref{R02:cor:grid} places its polynomial in the same set
$\mathcal G$.

For that polynomial, Theorem~\ref{R02:thm:finite} gives a prime divisor
of $T_{13\cdot3^{k_0+jh}-1}$ for every $j\ge0$. For all sufficiently
large $j$, put $k=k_0+jh-1\ge1$. Then
\[
 13\cdot3^{k_0+jh}-1=39\cdot3^k-1.
\]
The resulting traces on the $39$ schedule tend to infinity and are
divisible by the same fixed prime, contradicting eventual primality.
\end{proof}

This is not a deduction from the transcendence of the minimum for
$r=13$: deleting an initial floor condition changes the minimum.
The proof instead reuses the global bound and the infinite-return
certificates. It does not settle every schedule
$13\cdot3^h\cdot3^k-1$, because for $h\ge2$ the exponent $116$ is
absent and the same initial-floor bound is not available.

\subsection{A branch exclusion for the exponents \texorpdfstring{$2\cdot3^k-1$}{2 times powers of three minus one}}
\label{R02:sec:r2}

We now prove Theorem~\ref{R02:thm:fifth}. Its hypotheses do not require
any prime-trace assumption or irreducibility modulo $3$.

\begin{proof}[Proof of Theorem~\ref{R02:thm:fifth}]
Let $u,v$ be the small conjugates and put $q=uv=c/\beta$ and
$s=u+v=a-\beta$. As in \eqref{R02:eq:sformula},
\[
 s=q(1-1/\beta)>0,\qquad 0<q\le\beta^{-17/20}<1.
\]
Set $z=s^2/q=q(1-1/\beta)^2$. Then $0<z<1$, so
$s^2-4q=q(z-4)<0$. In particular $u,v$ are necessarily nonreal;
there is no omitted totally real case.

We first show $\beta>3/2$. Otherwise the norm bound forces $c=1$,
and Vieta's identity gives
\[
 a=g(\beta),\qquad g(x)=x+\frac1x-\frac1{x^2}.
\]
For $x>1$, $g'(x)=1-x^{-2}+2x^{-3}>0$, and therefore
\[
 1<a\le g(3/2)=\frac{31}{18}<2,
\]
contradicting integrality. If $c=1$, this gives $q=1/\beta<2/3$.
If $c\ge2$, the norm bound gives $\beta\ge c^{20/3}$ and hence
\[
 q=\frac c\beta\le c^{-17/3}\le2^{-17/3}<\frac12.
\]
Thus $q<2/3$ in every case.

The two-root Newton identity gives
\[
 S_5=u^5+v^5=s^5-5qs^3+5q^2s
 =sq^2(z^2-5z+5)>0.
\]
Indeed $z^2-5z+5>1$ for $0<z<1$. Since the conjugate pair has
modulus $\sqrt q$,
\[
 |S_5|\le2q^{5/2}<2(2/3)^{5/2}<1;
\]
the last inequality follows on squaring from $128/243<1$.
Consequently
\begin{equation}\label{R02:eq:fifthfloor}
 \floor{\beta^5}=T_5-1.
\end{equation}
Newton's identities factor the integer on the right exactly:
\begin{equation}\label{R02:eq:factor}
 \begin{split}
 T_5-1&=(a-1)Q(a,c),\\
 Q(a,c)&=a^4+a^3+a^2+a+1-5a^2c+5c^2.
 \end{split}
\end{equation}
For verification, expanding the product yields
$a^5-5a^3c+5a^2c+5ac^2-5c^2-1$, the fifth power trace minus one
for $b=c$.

Since $a=\beta+s>\beta>3/2$, the integer $a$ is at least $2$.
If $a=2$, then $\beta<2$ and $c\le\beta^{3/20}<2$, so $c=1$ and
\eqref{R02:eq:factor} equals $16$.
If $a\ge3$, then
\[
 c\le\beta^{3/20}<a^{3/20}\le a/2,
\]
where the last inequality follows from
$a^{17}\ge3^{17}>2^{20}$. Therefore
\[
 Q(a,c)\ge a^4-\frac32a^3+a^2+a+1>1.
\]
Both factors $a-1$ and $Q(a,c)$ are integers greater than one.
Equations \eqref{R02:eq:fifthfloor}--\eqref{R02:eq:factor} prove compositeness.
\end{proof}

\begin{corollary}\label{R02:cor:r2}
If $\alpha=\xi((2\cdot3^k-1)_{k\ge1})$ is algebraic, then, using
the structural reduction of \cite{K064} and the inverse-trace lemma
of \cite[Section~3]{N01}, its cubic minimal polynomial
$X^3-aX^2+bX-c$ must belong to one of the two sign branches
\[
 \begin{array}{ll}
 c>0,\quad b=-c,&\text{with a nonreal small conjugate pair};\\
 c<0,\quad b=c,&\text{with two real small conjugates of opposite signs}.
 \end{array}
\]
In both branches $a$ is odd, $3\nmid c$, and
$|c|\le\alpha^{3/20}$.
\end{corollary}
\begin{proof}
For $2\cdot3^k-1$, the eventual gcd of the exponents is $1$, so the
Mills-shadow reduction gives $\alpha$ itself as cubic Pisot. The
negative-fixed-point norm exclusion and magnitude bound of
\cite{K064} give $3\nmid c$ and $|c|\le\alpha^{3/20}$.
The inverse-trace lemma of \cite{N01} gives
$\Tr(\alpha^{-1})=\pm1$, hence $b=\pm c$.
The same reduction gives eventual floor--trace equality and eventual
prime traces, so Lemma~\ref{R02:lem:parity} makes $a$ odd.
The case $c>0,b=c$ is excluded by Theorem~\ref{R02:thm:fifth}, because
$5=2\cdot3^1-1$ is an actual required prime exponent.

It remains to inspect the signs of the small conjugates $u,v$.
Write $q=c/\alpha$ and
$s=q(\varepsilon-1/\alpha)$, with $\varepsilon=b/c$.
If $c>0$ and $b=-c$, then
\[
 \frac{s^2}{q}=q(1+1/\alpha)^2<4,
\]
since $q\le\alpha^{-17/20}<1$ and $1+1/\alpha<2$.
Thus the small roots are nonreal.
If $c<0$, their product is negative, so they are real with opposite
signs. When $b=-c$, one has $s>0$, and the positive root has larger
absolute value. For every odd $n$, therefore, $u^n+v^n>0$.
For all sufficiently large odd $n$ this sum is less than one, so
$\floor{\alpha^n}=T_n-1$, contradicting eventual floor--trace equality
on the entirely odd schedule $2\cdot3^k-1$.
Hence $b=c$ when $c<0$. In that case $s<0$, so the negative small
root has larger absolute value, as claimed.
\end{proof}

This corollary is a branch reduction, not a transcendence theorem for
$r=2$. Theorem~\ref{R02:thm:fifth} uses primality at the first required
floor. It does not rule out eventual prime traces for every polynomial
in its branch by itself. Conversely, the $r=13$ modular cover rules
out eventual prime traces but is justified only after its global
coefficient bound has been proved.


\subsection{Certificate data and a complete verifier}
\label{R02:app:verification}

The supplementary file \texttt{r13\_scan.json} is part of the
computer-assisted proof of Theorem~\ref{R02:thm:finite}. Its SHA-256 digest,
written on two consecutive lines for readability, is
\begin{center}\ttfamily\small
b14b061fb261ebff87c6fe27412bfd5d\\
bc0d917f45bbc4b9f3de96133845d957.
\end{center}
The two strings are to be concatenated, without a line break or the
sentence-ending period. The data include $230$ local certificates,
$8192$ candidate-to-certificate assignments, and an empty remainder.
The digest fixes the exact evidence being referred to; the verifier's
arithmetic and set-equality checks establish its mathematical content.

The following Python~3 program uses only the standard library. Run it
in the same directory as \texttt{r13\_scan.json}. All failed checks
raise an exception; no assertion can be disabled by an optimized
interpreter. An independent long-division verifier and the sequential
recurrence verifier are also supplied.

\begin{lstlisting}[language=Python]
import json
from math import isqrt
from pathlib import Path

D = json.loads(Path("r13_scan.json").read_text())
def check(ok):
    if not ok:
        raise ValueError("Certificate verification failed")

def prime(p):
    return p >= 2 and all(p % d for d in range(2, isqrt(p)+1))

I = tuple(tuple(int(i == j) for j in range(3)) for i in range(3))
def mul(A, B, p):
    return tuple(tuple(sum(A[i][t]*B[t][j] for t in range(3)) % p
                       for j in range(3)) for i in range(3))

def power(A, n, p):
    B = I
    while n:
        if n & 1:
            B = mul(B, A, p)
        A = mul(A, A, p)
        n //= 2
    return B

check(D["r"] == 13 and D["a_max"] == 8191)
check(D["remaining"] == [])
certs = D["certificates"]
check(len(certs) == 230)
for w in certs:
    p,a,b,c,k,E,U,n,h = [w[t] for t in
        ("p","a_mod","b","c","k","E","U","n","period")]
    check(prime(p) and 0 <= a < p and c in (1,2))
    check(b in (-c,c) and c % p != 0)
    check(k >= 1 and E >= 1 and U >= 1 and h >= 1)
    check(0 <= n < E)
    U0, v = E, 0
    while U0 % 3 == 0:
        U0 //= 3
        v += 1
    check(U == U0 and k >= v and pow(3,h,U) == 1 % U)
    check(n == (13*pow(3,k,E)-1) % E)
    M = ((0,0,c % p), (1,0,(-b) % p), (0,1,a))
    check(power(M,E,p) == I)
    B = power(M,n,p)
    check(sum(B[i][i] for i in range(3)) % p == 0)

expected = {(a,b,c) for a in range(1,8192,2)
            for c in (1,2) for b in (-c,c)
            if all((t**3-a*t*t+b*t-c) % 3 for t in range(3))}
actual = set()
for row in D["excluded"]:
    a,b,c = [row[t] for t in ("a","b","c")]
    key = (a,b,c)
    check(key not in actual)
    actual.add(key)
    j = row["certificate"]
    check(isinstance(j,int) and 0 <= j < len(certs))
    w = certs[j]
    check(b == w["b"] and c == w["c"])
    check(a % w["p"] == w["a_mod"])
check(len(expected) == 8192 == D["initial_count"])
check(actual == expected)
print("PASS: 230 independent matrix certificates cover all 8192 candidates.")
\end{lstlisting}



\endgroup

\section{Odd multiplier extensions}\label{src:R04}
The same certificate mechanism extends to further odd multipliers. The finite coefficient grids are exhaustive only under the height and structural hypotheses proved here; they are not a finite enumeration of every cubic Pisot candidate.
\medskip
\begingroup



\subsection{The results and their dependencies}

For a sequence $C=(C_k)_{k\ge1}$ of positive integers, write
\[
 \mathcal W(C)=\{A>1:\ \floor{A^{C_k}}\text{ is prime for every }k\ge1\},
 \qquad \xi(C)=\min\mathcal W(C).
\]
Existence of this minimum when $C_{k+1}\ge2C_k$ is supplied by
\cite[Theorem~1.3]{Saito}. All schedules in this paper satisfy that
inequality. We abbreviate
\begin{equation}\label{R04:eq:schedule}
 n_k(r)=r3^k-1\quad(k\ge1),\qquad
 \xi_r=\xi((n_k(r))_{k\ge1}).
\end{equation}

\begin{mainthm}\label{R04:thm:main}
The numbers $\xi_{65}$, $\xi_{91}$ and $\xi_{195}$ are transcendental.
\end{mainthm}

Our explicit structural dependency is the theorem for odd multiples of
$13$ in \cite[Section~7.2, Theorem~\textup{\texttt{D13}}]{K064}.
That theorem is proved in the author's earlier unpublished manuscript;
we do not claim to reprove it here. Its consequences, including the
residue-field detail needed by the computation, are recorded in
Proposition~\ref{R04:prop:input}. Subject to that cited input, the new
analytic and finite-cover arguments below are exact and unconditional:
no conjecture about primes, floating-point root test, or extrapolation
from a finite prime prefix is used.

The analytic part is uniform, rather than a search cutoff chosen for
convenience. Define, for odd $r\ge1$,
\begin{equation}\label{R04:eq:N}
 j(r)=\begin{cases}1,&r\equiv3\pmod4,\\2,&r\equiv1\pmod4,\end{cases}
 \qquad N(r)=r3^{j(r)}-1.
\end{equation}
Then $4\mid N(r)$ and $N(r)\ge8$.

\begin{mainthm}\label{R04:thm:boundmain}
If $r$ is odd and $\xi_r$ is algebraic, then the structural reduction
from \cite{K064} either already gives a contradiction, when $13\nmid r$,
or gives $\beta=\xi_r$ with positive integral norm $c$ satisfying
\[
 \beta<\left(\frac{N(r)}2\right)^{40/17},\qquad
 c<\left(\frac{N(r)}2\right)^{6/17}.
\]
The latter bound alone does not exclude the remaining coefficients
for every odd multiple of $13$.
\end{mainthm}

The first two new cases are settled by finite covers of entire necessary
coefficient grids, not just of numerically identified Pisot roots.
The third case reuses a cover only after proving that the same grid
still contains every hypothetical algebraic minimum.

\subsection{The structural input and the trace sequence}
\label{R04:sec:input}

A Pisot number is a real algebraic integer $\beta>1$ whose other
conjugates have absolute value less than one. In the cubic case put
\begin{equation}\label{R04:eq:f}
 f(X)=X^3-aX^2+bX-c,\qquad
 T_n=\Tr_{\Q(\beta)/\Q}(\beta^n).
\end{equation}
Newton's identities give the integer recurrence
\begin{equation}\label{R04:eq:trace}
 \begin{gathered}
 T_0=3,\quad T_1=a,\quad T_2=a^2-2b,\\
 T_n=aT_{n-1}-bT_{n-2}+cT_{n-3}\qquad(n\ge3).
 \end{gathered}
\end{equation}

\begin{proposition}[Structural input from {\cite{K064}}]\label{R04:prop:input}
Let $r$ be odd with $13\mid r$. If $\xi_r$ is algebraic, then
$\beta=\xi_r$ itself is a cubic Pisot number, with nonreal conjugates
$u,v=\overline u$. Its minimal polynomial \eqref{R04:eq:f} satisfies
\begin{equation}\label{R04:eq:structural}
 b=\varepsilon c,\quad \varepsilon\in\{-1,1\},\qquad
 0<c\le\beta^{3/20},\qquad 3\nmid c,
 \qquad f\text{ irreducible over }\F_3.
\end{equation}
Furthermore,
\begin{align}
 T_{n_k(r)}=\floor{\beta^{n_k(r)}}\text{ is prime}
 &\quad\text{for all sufficiently large }k,\label{R04:eq:tail}\\
 \floor{\beta^{n_k(r)}}\text{ is prime}
 &\quad\text{for every }k\ge1.\label{R04:eq:floorall}
\end{align}
The equality in \eqref{R04:eq:tail} is not asserted at every initial index.
\end{proposition}

\begin{proof}[Explanation of the cited input]
The algebraic-minimum reduction in \cite[Section~2]{K064}, based on
\cite[Proposition~3.1]{Saito}, first gives the cubic Pisot number
$B=\xi_r^2$ with eventual prime traces at exponents $n_k(r)/2$.
Theorem~\texttt{D13} proves that $\xi_r$ itself is cubic Pisot with
nonreal conjugates and
\[
 \Q(\xi_r)=\Q(B),\qquad
 \Tr(\xi_r^{-1})=\pm1,\qquad
 \xi_r\ge\lvert\operatorname{N}_{\Q(\xi_r)/\Q}(\xi_r)\rvert^{20/3}.
\]
Since $c=\beta u\overline u>0$ and $\Tr(\beta^{-1})=b/c$, these give
the coefficient and norm assertions. Equality of the fields gives
\[
 \Tr_{\Q(B)/\Q}(B^{n_k(r)/2})
 =\Tr_{\Q(\beta)/\Q}(\beta^{n_k(r)}),
\]
which transfers the eventual prime-trace statement to $\beta=\xi_r$.
The initial floor condition \eqref{R04:eq:floorall} follows separately
from the definition of the minimum.

For completeness, irreducibility modulo $3$ requires more than the
statement that $3$ is inert. The proof of \texttt{D13} shows that the
minimal polynomial of $B=\beta^2$ is irreducible modulo $3$.
Thus $3$ is inert in $K=\Q(B)=\Q(\beta)$, with residue field $\F_{27}$,
and the residue $\overline B$ generates that field over $\F_3$.
The integral element $\beta$ has a residue in the same field, and
\[
 \F_3(\overline B)\subseteq\F_3(\overline\beta)\subseteq\F_{27},
 \qquad \overline\beta^{\,2}=\overline B.
\]
Therefore $\overline\beta$ has degree three over $\F_3$.
Its monic cubic polynomial $\overline f$ is irreducible; in particular
$3\nmid c$. This square-generator argument avoids an unjustified
inference from inertness for an arbitrary integral generator.
\end{proof}

For $13\nmid r$, the transcendence result already follows from
\cite[Theorem~B]{K064}. All further arguments concern the residual
case of Proposition~\ref{R04:prop:input}.

\subsection{Scheduled parity and the uniform height bound}
\label{R04:sec:bound}

\begin{lemma}\label{R04:lem:parity}
Let $r\ge1$ and let $\beta$ be a cubic Pisot number with
$b\equiv c\pmod2$. If $T_{n_k(r)}$ is eventually prime, then every
scheduled trace $T_{n_k(r)}$, $k\ge1$, is odd. If $7\nmid r$, then
$a$ is odd. If $7\mid r$, only the simultaneous conditions
$a\equiv c\equiv0\pmod2$ are excluded by this parity argument.
\end{lemma}
\begin{proof}
The Pisot property gives $T_n=\beta^n+o(1)$, so eventually prime
scheduled traces eventually exceed two. There are three cases.
If $c$ is even, the roots modulo two have multiset $\{0,0,a\}$,
and $T_n\equiv a^n\pmod2$ for $n\ge1$.
If $c$ and $a$ are both odd, then
$\overline f=(X+1)^3$ over $\F_2$, and every trace is odd.
If $c$ is odd and $a$ is even, then
$\overline f=X^3+X+1$, whose roots are nonzero elements of $\F_8$
with order seven. In this case
\begin{equation}\label{R04:eq:parities}
 (T_n\bmod2)_{n=0}^{6}=(1,0,0,1,0,1,1),
 \qquad T_{n+7}\equiv T_n\pmod2.
\end{equation}
The sequence $r3^k-1$ modulo seven is periodic from its first index,
with period dividing six. Consequently the parity of the scheduled
traces is purely periodic in all three cases. An eventually odd,
purely periodic sequence is odd from the first index.

If $7\nmid r$, the powers of three run through all nonzero residues
modulo seven, and multiplication by $r$ permutes them. Thus the
scheduled exponents include residues $1,2,4$ infinitely often;
\eqref{R04:eq:parities} excludes the third case. The first case excludes
even $a$ as well. If $7\mid r$, all scheduled exponents are $6$ modulo
seven, and the third case has odd scheduled traces. It must not be
excluded.
\end{proof}

In particular, $r=91$ does not force $a$ to be odd. The conclusion
needed for the initial floor is only oddness on the actual schedule.

\begin{proposition}\label{R04:prop:bound}
Suppose $r$ is odd, and $\beta$ satisfies
Proposition~\ref{R04:prop:input}. With $N=N(r)$ from \eqref{R04:eq:N},
\[
 \beta<(N/2)^{40/17},\qquad c<(N/2)^{6/17}.
\]
\end{proposition}
\begin{proof}
Put $q=uv=c/\beta$. Vieta's identities and $b=\varepsilon c$ give
\begin{equation}\label{R04:eq:angle}
 u+v=q\left(\varepsilon-\frac1\beta\right),\qquad
 0<q\le\beta^{-17/20}<1.
\end{equation}
Write $u=\sqrt q\,e^{i\theta}$ and $v=\sqrt q\,e^{-i\theta}$ with
$0<\theta<\pi$. Then
\begin{equation}\label{R04:eq:shadow}
 \cos\theta=\frac{\sqrt q}{2}\left(\varepsilon-\frac1\beta\right),
 \qquad S_n=u^n+v^n=2q^{n/2}\cos(n\theta),
 \qquad \beta^n=T_n-S_n.
\end{equation}
Suppose $\beta\ge(N/2)^{40/17}$. Since $N\ge8$, we have $\beta>4$.
For $t=|\cos\theta|$, \eqref{R04:eq:angle} yields
\[
 0<t<\frac58\beta^{-17/40}
 \le\frac5{4N}\le\frac5{32}<\frac12.
\]
The strict positivity follows from $q>0$, $\beta>1$ and
$\varepsilon\in\{-1,1\}$. With $\delta=|\theta-\pi/2|=\arcsin t$,
\[
 \delta\le\frac{t}{\sqrt{1-t^2}}
 <\frac65t
 <\frac34\beta^{-17/40}
 \le\frac3{2N}.
\]
Here $2/\sqrt3<6/5$ follows by squaring from $100<108$.
It follows that $N\delta<3/2<\pi/2$. Since $4\mid N$,
$\cos(N\theta)=\cos(N\delta)>0$.
Moreover $q\le(2/N)^2$, and therefore
\[
 0<S_N\le2q^{N/2}\le2(2/N)^N<1.
\]
The integer $T_N$ consequently satisfies
\[
 \floor{\beta^N}=T_N-1.
\]
By Lemma~\ref{R04:lem:parity}, $T_N$ is odd. This floor is even and,
since $\beta>4$, greater than two. But $N=n_{j(r)}(r)$ is an actual
required exponent, so \eqref{R04:eq:floorall} says that this same floor
is prime. This contradiction proves the bound on $\beta$.
Finally $c\le\beta^{3/20}<(N/2)^{6/17}$.
\end{proof}

Proposition~\ref{R04:prop:bound}, together with the cited result for
$13\nmid r$, proves Theorem~\ref{R04:thm:boundmain}. Notice that no
initial floor/trace equality was assumed: the sign and size of $S_N$
determined the initial floor directly.

\subsection{The two exhaustive coefficient grids}
\label{R04:sec:grids}

For $r=65$ and $r=91$, the parameters are
\[
 \begin{array}{c|c|c|c}
 r&j(r)&N(r)&B\\\hline
 65&2&584&2^{20}\\
 91&1&272&2^{17}
 \end{array}
\]
The following integer comparisons, checked by both verifiers, imply
$\beta<B$ and the required norm cutoffs:
\begin{equation}\label{R04:eq:integerbounds}
 \begin{aligned}
 292^{40}&<(2^{20})^{17},&292^6&<8^{17},\\
 136^{40}&<(2^{17})^{17},&136^6&<6^{17}.
 \end{aligned}
\end{equation}
For example, $292^2=85\,264<2^{17}$ proves the first inequality;
$136^5=46\,525\,874\,176<2^{36}$ proves the third after raising to
the eighth power. The norm cutoff also follows from
$c<2^3=8$ in the first case and
$c<2^{51/20}<6$, with $2^{51}<6^{20}$, in the second.
As $c$ is a positive integer not divisible by three, the norm sets are
\[
 C_{65}=\{1,2,4,5,7\},\qquad C_{91}=\{1,2,4,5\}.
\]

There is a safe coefficient bound without approximating any roots.
From \eqref{R04:eq:angle}, $a=\beta+u+v>\beta-2>-1$, hence $a\ge0$.
If $\varepsilon=-1$ then $a<\beta<B$; if $\varepsilon=1$ then
$0<u+v<1$, so $a<\beta+1<B+1$. Thus in all cases $a\le B$.
When $a$ is odd and $B$ even this improves to $a\le B-1$.

Define $\mathcal G_{65}$ and $\mathcal G_{91}$ to be the sets of
integer triples $(a,b,c)$ with $b\in\{-c,c\}$, with
$X^3-aX^2+bX-c$ irreducible over $\F_3$, and with the following
additional conditions:
\begin{equation}\label{R04:eq:grids}
 \begin{aligned}
 \mathcal G_{65}:\;&1\le a\le2^{20}-1,\quad a\text{ odd},
       \quad c\in C_{65};\\
 \mathcal G_{91}:\;&0\le a\le2^{17},\quad c\in C_{91},
       \quad\text{not both }a,c\text{ even}.
 \end{aligned}
\end{equation}
Every algebraic minimum for the corresponding schedule yields a
triple in its grid. In particular, the even-$a$, odd-$c$ branch and
both endpoints of the $r=91$ interval are retained whenever the
modulo-three condition permits them. We do not impose any numerical
Pisot, discriminant, or additional norm filter.

For a cubic over $\F_3$, irreducibility is exactly the absence of
roots in $\{0,1,2\}$. Direct enumeration, or counting the admissible
$a$ classes modulo six, gives Table~\ref{R04:tab:counts}.
\begin{table}[htbp]
\centering
\caption{Exact grid sizes separated by the coefficients $b,c$.}
\label{R04:tab:counts}
\begin{tabular}{rrr@{\qquad}rr}
\toprule
&&&\multicolumn{2}{c}{Number of admissible $a$}\\
$c$&$b$&&$\mathcal G_{65}$&$\mathcal G_{91}$\\
\midrule
1&$-1$&&349526&87382\\
1&1&&174762&43691\\
2&$-2$&&174763&21846\\
2&2&&349525&43690\\
4&$-4$&&349526&43691\\
4&4&&174762&21845\\
5&$-5$&&174763&43691\\
5&5&&349525&87382\\
7&$-7$&&349526&---\\
7&7&&174762&---\\
\midrule
\multicolumn{3}{l}{Total}&2621440&393218\\
\bottomrule
\end{tabular}
\end{table}

\subsection{Exact certificates for infinitely recurring divisibility}
\label{R04:sec:certificates}

For any integer coefficients $a,b,c$, define $T_n$ by
\eqref{R04:eq:trace}; equivalently it is the sum of the $n$th powers of
the roots of $f$, counted with multiplicity. The companion matrix
\begin{equation}\label{R04:eq:matrix}
 M(a,b,c)=\begin{pmatrix}0&0&c\\1&0&-b\\0&1&a\end{pmatrix}
 \quad\text{satisfies}\quad \tr(M^n)=T_n\qquad(n\ge0).
\end{equation}
This remains true after reduction modulo any prime, including when
$f$ becomes reducible or inseparable.

\begin{lemma}[Recurrent-divisor certificate]\label{R04:lem:cert}
Fix $r\ge1$ and integers $b,c$. Let
\[
 (p,a_0,b,c,k,E,U,n,h)
\]
be a tuple of integers satisfying the following tests:
\begin{enumerate}[label=\textup{(\roman*)}]
\item $p$ is prime, $p\nmid c$, and $0\le a_0<p$;
\item $E=3^vU>0$ with $v\ge0$, $3\nmid U$, $k\ge\max(v,1)$,
$h\ge1$, and $3^h\equiv1\pmod U$;
\item $0\le n<E$ and $n\equiv r3^k-1\pmod E$;
\item modulo $p$, $M(a_0,b,c)^E=I$ and
$\tr(M(a_0,b,c)^n)=0$.
\end{enumerate}
Then every integer $a\equiv a_0\pmod p$ satisfies
\[
 p\mid T_{r3^{k+th}-1}\qquad\text{for every integer }t\ge0.
\]
\end{lemma}
\begin{proof}
The factor $3^v$ divides $3^k$, while $U$ divides $3^{th}-1$.
Since $3^v$ and $U$ are coprime,
\[
 E\mid r3^k(3^{th}-1).
\]
Thus $r3^{k+th}-1\equiv n\pmod E$, and the matrix identities in
(iv) give the claimed trace divisibility. Reduction of $a$ to $a_0$
does not change the matrix modulo $p$.
\end{proof}

\begin{proposition}[Complete finite covers]\label{R04:prop:cover}
Every triple in $\mathcal G_{65}$ or $\mathcal G_{91}$ admits a
certificate from Lemma~\ref{R04:lem:cert}, for its corresponding $r$.
The supplied exact certificate files have the following sizes:
\[
\begin{array}{c|r|r|r|r}
 r&\text{grid size}&\text{certificates}&\max p&\text{uncovered}\\\hline
 65&2\,621\,440&1\,447&107&0\\
 91&393\,218&774&89&0
\end{array}
\]
\end{proposition}
\begin{proof}[Finite verification and its meaning]
The data consist of the tuples of Lemma~\ref{R04:lem:cert}, not one row
for each coefficient triple. For each fixed $(b,c)$ a tuple covers
exactly the grid elements satisfying $a\equiv a_0\pmod p$.
The matrix verifier independently reconstructs \eqref{R04:eq:grids},
tests every tuple, and removes these arithmetic progressions.
It verifies the initial counts in Table~\ref{R04:tab:counts}, all
per-prime survivor counts, and an empty final survivor list in each
case. Primality of every witness prime is checked by trial division;
all other calculations are exact integer or finite-field operations.

A separately authored verifier changes both central implementations:
it checks identities using generic polynomial long division in
$\F_p[X]/(f)$ instead of matrices, and reconstructs and subtracts
coefficient sets using Python sets instead of bytearrays. It also
checks the seed exponent and two explicit returns for every tuple.
It verifies the same two empty final sets. The certificate files,
verifier sources, commands, and their identifying hashes are specified
in Section~\ref{R04:app:reproduce}. Thus the computational assertion is
an exhaustive finite check with independently checkable witnesses,
not an assumption that a divisor observed once must recur.
Lemma~\ref{R04:lem:cert} proves the required infinite recurrence.
\end{proof}

The discovery program chooses
\begin{equation}\label{R04:eq:E}
 E=p^2\lcm(p-1,p^2-1,p^3-1),\qquad
 U=E/3^{v_3(E)},\qquad h=\varphi(U).
\end{equation}
The proof does not depend on a claim that this is a minimal period.
The verifiers test $M^E=I$ and $3^h\equiv1\pmod U$ directly.
Thus reducible or repeated factors modulo $p$ cause no gap; primes
dividing $r$ are also harmless. A prime dividing $c$ is not used for
a unit-return certificate.

\subsection{Transcendence and the limits of the transfer}
\label{R04:sec:conclusion}

\begin{proof}[Proof of Theorem~\ref{R04:thm:main}]
Suppose first that $\xi_r$ is algebraic for $r=65$ or $r=91$.
Proposition~\ref{R04:prop:input} gives $\beta=\xi_r$ and eventual prime
traces. Proposition~\ref{R04:prop:bound} and Section~\ref{R04:sec:grids}
place its coefficients in $\mathcal G_r$.
Proposition~\ref{R04:prop:cover} and Lemma~\ref{R04:lem:cert} supply a fixed
prime $p$ dividing infinitely many of its scheduled traces.
But $T_n=\beta^n+u^n+v^n\longrightarrow+\infty$.
At all sufficiently large certified indices, $T_n>p$ and is divisible
by $p$, hence composite. This contradicts eventual primality.

Now let $r=195=3\cdot65$ and suppose $\xi_{195}$ algebraic.
The structural input still applies. Since $7\nmid195$, $a$ is odd.
Its first exponent is
\[
 N(195)=3\cdot195-1=584=N(65).
\]
Therefore the same height and norm estimates put its coefficients
in the identical grid $\mathcal G_{65}$. A certificate for that grid
provides divisibility at exponents $65\cdot3^{k+th}-1$. These equal
\[
 195\cdot3^{k+th-1}-1,
\]
and $k+th-1\ge1$ for all sufficiently large $t$.
Thus infinitely many actual $r=195$ traces are divisible by the same
prime and eventually exceed it. This again contradicts eventual
primality.
\end{proof}

\begin{corollary}\label{R04:cor:100}
For every odd integer $r$ with $1\le r\le100$, $\xi_r$ is transcendental.
\end{corollary}
\begin{proof}
Theorem~B of \cite{K064} covers $13\nmid r$.
The only positive odd multiples of $13$ at most $100$ are
$13,39,65,91$. The first two are proved in \cite{R02}; the last two
are Theorem~\ref{R04:thm:main}.
\end{proof}

The $r=195$ proof is not a transfer of transcendence merely because a
schedule is a tail of another. Deleting an initial prime-floor
condition changes the minimization problem. We instead reestablished
the same finite coefficient bound and then transferred infinite
recurring-divisor witnesses. For example, the exponent $272$ used
for $r=91$ is absent from the $r=273$ schedule, so the present proof
does not automatically settle $r=273$.

The uniform bound leaves a finite problem for each residual odd
multiplier, but this paper proves complete covers only for the stated
grids. No result for all odd multiples of $13$ is asserted.
The constant $\xi_2=\xi((2\cdot3^k-1)_{k\ge1})$ remains unresolved
by this argument, as does the transcendence of the classical Mills
constant with exponents $3^k$.


\subsection{Reproducible finite verification}
\label{R04:app:reproduce}

All files in Table~\ref{R04:tab:files} are supplied together in the
accompanying \texttt{certificates} directory. Only the two JSON files and
either verifier are needed to check the finite-cover assertion.
The verifiers use the Python standard library; SymPy is needed only
to rerun discovery. The data encode exactly the tests in
Lemma~\ref{R04:lem:cert} and the grids in \eqref{R04:eq:grids}.

\begin{table}[htbp]
\centering\small
\caption{Reproduction files and their roles.}\label{R04:tab:files}
\begin{tabular}{@{}p{0.54\linewidth}p{0.40\linewidth}@{}}
\toprule
File&Role\\\midrule
\texttt{odd65\_91\_r65.json}&1,447 local certificates and grid metadata\\
\texttt{odd65\_91\_r91.json}&774 local certificates and grid metadata\\
\texttt{odd65\_91\_scan.py}&Exact quotient-ring discovery program\\
\texttt{odd65\_91\_verify.py}&Matrix and bytearray cover verifier\\
\texttt{audit\_odd65\_91\_verify.py}&Long-division and set cover verifier\\
\texttt{odd65\_91\_verification.json}&Matrix verifier terminal results\\
\texttt{audit\_odd65\_91\_verification.json}&Independent audit terminal results\\
\texttt{odd65\_91\_certificates.sha256}&Frozen certificate digests\\
\bottomrule
\end{tabular}
\end{table}

The following commands verify the existing frozen data without
rerunning discovery. Run them without Python's optimization flag,
so that the matrix verifier's assertion checks remain enabled.
\begin{lstlisting}
cd certificates
python3 odd65_91_verify.py
python3 audit_odd65_91_verify.py
sha256sum -c odd65_91_certificates.sha256
\end{lstlisting}
Both verifiers report $2\,621\,440$ candidates, $1\,447$ certificates,
and zero survivors for $r=65$; $393\,218$ candidates, $774$
certificates, and zero survivors for $r=91$.
The largest seed indices are respectively $290$ and $230$.

For clarity, the complete logical structure of the finite check is:
\begin{lstlisting}
for r in (65, 91):
    G = exact_grid(r)  # stated bounds, parity, no root in F_3
    assert all_pair_counts(G) == supplied_initial_counts
    for (p, a0, b, c, k, E, U, n, h) in certificates[r]:
        verify_all_recurrent_divisor_certificate_conditions()
        G[b,c] = {a for a in G[b,c] if a % p != a0}
    assert every coefficient set G[b,c] is empty
\end{lstlisting}
Here \texttt{exact\_grid} uses only integer inequalities and the
three evaluations $x^3-ax^2+bx-c\pmod3$ for $x=0,1,2$.
Matrix powers are computed by binary exponentiation modulo $p$.
The alternative verifier computes powers by polynomial multiplication
and long division by $f$; for a residue $z_0+z_1X+z_2X^2$, its trace is
$3z_0+az_1+(a^2-2b)z_2$ modulo $p$.
These tests completely specify the finite proposition independently
of the method used to discover the witnesses.

To rediscover the covers, work in a separate copy of the certificates
folder, because discovery rewrites the JSON and timing metadata:
\begin{lstlisting}
python3 odd65_91_scan.py 65 1000
python3 odd65_91_scan.py 91 1000
\end{lstlisting}
The upper search limit $1000$ is not a mathematical hypothesis:
the scans stop once coverage is complete, at witness primes $107$
and $89$. Their search updates $X^{r3^k}$ by cubing in
$\F_p[X]/(f)$ and uses the identity
\[
 \Tr\!\left(X^{-1}(z_0+z_1X+z_2X^2)\right)
 =(b/c)z_0+3z_1+az_2.
\]
The verification does not reuse these discovery formulas.

\paragraph{Frozen SHA256 digests.}
Line breaks below are typographical only; each digest is one
64-character hexadecimal string.
\begin{lstlisting}
odd65_91_r65.json
73a813b489d134e6e3c7d085c768ebc5715ea88dc98f6ebdc2e00221c13d96a6
odd65_91_r91.json
56aa4b60938822f0d00c4df7f5f9e8c4c45d14792599a9feefb3e71c9206b2ce
odd65_91_scan.py
7aa615056aa2119c53e0fbacf3749ee1b24677f0050c74371a427903aa605984
odd65_91_verify.py
ef5b9c5a084d468da95024a0ebc5e2e57409efd537d045b25c0b0ed151aa4f77
audit_odd65_91_verify.py
fbd28dece660e180d94b0108503eee0eb4630012b6cfd86cd10e56bf927e315a
\end{lstlisting}



\endgroup

\clearpage\part{Prime cells and finite to infinite continuation}

\section{Prime cells height bounds and Pell families}\label{src:R03}
We now separate finite prime prefixes from their infinite continuations. The prime-cell formalism supports the height and Pell-family arguments below and provides the precise interface needed for the subsequent prime-forest criteria.
\medskip
\begingroup



\subsection{Statements and dependencies}

For positive exponents $n_k\to\infty$, define
\[
 \W(n_k)=\{A>1:\floor{A^{n_k}}\in\PP\text{ for every }k\ge1\}.
\]
When a least element exists, write it as $\xi(n_k)$. The two schedules
in this report are
\begin{equation}\label{R03:eq:clocks}
 n_k=2\cdot3^k-1=5,17,53,\ldots\quad(k\ge1),
 \qquad M_r=3^r\quad(r\ge1).
\end{equation}
The prime-representing sets for both schedules are nonempty and have
least elements by \cite[Theorem~1.3]{Saito}. Their arithmetic problems
must nevertheless be kept separate: a statement about $3^r$ does not
automatically transfer to $2\cdot3^k-1$.

A \emph{cubic Pisot number} is a real algebraic integer $\beta>1$ of
degree three whose other conjugates have modulus less than one. Write
\[
 T_n=\Tr_{\Q(\beta)/\Q}(\beta^n),\qquad
 S_n=T_n-\beta^n.
\]
Thus the sign of $S_n$ determines which of $T_n$ and $T_n-1$ is the
floor once $|S_n|<1$. We reserve $V_r=T_{3^r}$ for the classical
schedule in Section~\ref{R03:sec:real}.

The shifted structural reduction used in earlier work
\cite{K064,N01,R02} leaves the following two families:
\begin{equation}\label{R03:eq:residual}
 \begin{gathered}
 \begin{aligned}
 f_+(X)&=X^3-aX^2-CX-C &&\text{with a nonreal conjugate pair},\\
 f_-(X)&=X^3-aX^2-CX+C &&\text{with three real roots},
 \end{aligned}\\
 a\text{ odd},\quad C\in\Z_{\ge1},\quad C\le\beta^{3/20}.
 \end{gathered}
\end{equation}
Here $\beta$ is required to be cubic Pisot. Our first theorem is
stated directly under these hypotheses; its proof does not reprove
the earlier reduction, and does not assume that every polynomial
with coefficients as displayed is a Pisot polynomial.

Put $\theta=21/40$. For a real number $X\ge8$, let $(\mathrm P_X)$ mean
\begin{equation}\label{R03:eq:PX}
 [x-x^\theta,x]\cap\PP\ne\varnothing
 \qquad\text{for every real }x\ge X.
 \tag{$\mathrm P_X$}
\end{equation}
The Baker--Harman--Pintz theorem supplies such a threshold
\cite[Theorem~1]{BHP}.

\begin{mainthm}\label{R03:thm:height}
Suppose $\beta$ satisfies \eqref{R03:eq:residual}, every
$\floor{\beta^{2\cdot3^k-1}}$ is prime, and $\beta$ is the global
minimum of $\W(2\cdot3^k-1)$. For every $X\ge8$ satisfying
\eqref{R03:eq:PX},
\[
 \beta<\max\{4,(X+2)^{1/17}\}.
\]
The prime competitor used in the proof has a complete infinite
continuation and retains the first prime floor.
\end{mainthm}

\begin{mainthm}\label{R03:thm:pell}
Define
\begin{equation}\label{R03:eq:pell-main}
 \binom{m_0}{y_0}=\binom37,\qquad
 \binom{m_{t+1}}{y_{t+1}}
 =\begin{pmatrix}577&204\\1632&577\end{pmatrix}
  \binom{m_t}{y_t}-\binom{468}{1326},\qquad a_t=2m_t-1.
\end{equation}
For every $t\ge0$, the real root $\beta_t$ of
$X^3-a_tX^2-X-1$ is a cubic Pisot unit and has infinitely many
composite values $\floor{\beta_t^{2\cdot3^k-1}}$. The five primes
$5,19,29,31,59$ suffice to certify this assertion for the entire ray.
No minimum or conjugate-phase assumption is needed.
The same conclusion holds for the totally real reflected family in
Corollary~\ref{R03:cor:real-reflection}.
\end{mainthm}

For the classical schedule, let $\SaitoL$ consist of totally real
cubic Pisot numbers with conjugates
\begin{equation}\label{R03:eq:S}
 \beta,-\rho,\gamma,\qquad
 0<\sigma=|\gamma|<\rho<1,\qquad
 \rho\le\min\{\sigma^{17/23},\beta^{-17/40}\}.
\end{equation}
Saito's reduction says that a hypothetical algebraic classical Mills
constant has a positive integral power in this locus
\cite[Theorem~1.7]{Saito}.

\begin{mainthm}\label{R03:thm:finite}
The locus $\SaitoL$ has the exact coefficient description in
Theorem~\ref{R03:thm:locus}. Among its polynomials
$X^3-aX^2-BX-c$ with $1\le a\le1000$, there are exactly $32,232$
coefficient triples. Every one has a prime $p\le59$ dividing
$V_r=\Tr(\beta^{3^r})$ on an infinite arithmetic progression of
indices. Consequently every such $\beta$ has infinitely many composite
classical floors $\floor{\beta^{3^r}}$.
\end{mainthm}

Theorem~\ref{R03:thm:finite} is a computer-assisted finite theorem;
its exhaustive coefficient bounds, exact comparisons, and certificate
semantics are given below. It does not bound the trace of an arbitrary
member of $\SaitoL$. The norm bound, redundancy of one inequality in
\eqref{R03:eq:S}, and fixed-norm density-one results are already in
manuscript~[031] \cite{K031}; they are background, not new results
of this report. The bounded theorem was obtained in the preceding
round and is presented here with an independent integer-only
verification and a uniform floor-to-trace onset.

\subsection{Exact prime cells and minimum criteria}\label{R03:sec:cells}

\subsubsection{Viability includes the infinite continuation}

Assume throughout this subsection that $n_k$ are positive and tend to
infinity. A prime prefix $p_1,\ldots,p_k$ determines the half-open cell
\begin{equation}\label{R03:eq:prefix}
 J_k=[L_k,U_k),\qquad
 L_k=\max_{j\le k}p_j^{1/n_j},\quad
 U_k=\min_{j\le k}(p_j+1)^{1/n_j}.
\end{equation}
A prime $q$ is a \emph{viable next label} if some member of
$\W(n_k)$ lies in
\[
 J_k\cap[q^{1/n_{k+1}},(q+1)^{1/n_{k+1}}).
\]
The first label is viable if its cell contains a member of $\W(n_k)$.
A nonempty finite intersection alone is not a viability certificate.
An element $A\in\W(n_k)$ is a \emph{local left minimum} if some
$\varepsilon>0$ satisfies $(A-\varepsilon,A)\cap\W(n_k)=\varnothing$.

\begin{proposition}\label{R03:prop:viable}
An all-prime base is the global minimum if and only if its first label
and its next label after every prefix are the least viable choices.
If additionally $A^{n_j}\notin\Z$ for every $j$, then $A$ is a local
left minimum if and only if its next label is the least viable choice
after every sufficiently long prefix.
\end{proposition}
\begin{proof}
Two distinct bases have a first different label: for $1<A<B$,
$B^{n_k}-A^{n_k}\to\infty$. At that first difference, the smaller
base has the smaller label. This proves both global implications.

The cell containing a fixed $A>1$ has width tending to zero, by the
mean value theorem (or directly from its endpoints). Infinitely many
smaller viable siblings therefore produce smaller all-prime bases
sharing successively longer prefixes and converging to $A$; a local
left gap is impossible. Conversely, if all choices after a prefix
of length $K$ are least viable, nonintegrality puts $A$ in the
interior of that entire finite prefix. Any smaller all-prime base
in a sufficiently small left neighborhood has a first different
label after $K$, contrary to leastness. Equal labels forever would
force the bases to be equal.
\end{proof}

\subsubsection{The affine boundary child cannot be omitted}

For the shifted schedule, put $n=n_k$, $m=n_{k+1}=3n+2$ and $d=m/n$.
Suppose for the moment that the current prefix cell is the full
$p$-cell. Since $\gcd(n,m)=1$ and $n\ge5$, the number $p^d$ is not
an integer. The integer
\begin{equation}\label{R03:eq:boundary}
 q_0=\floor{p^d}
\end{equation}
can be prime, and its cell can meet the parent in the truncated cell
\begin{equation}\label{R03:eq:truncated}
 [p^{1/n},(q_0+1)^{1/m}).
\end{equation}
Indeed, $q_0<p^d<q_0+1<(p+1)^d$. Labels below $q_0$ cannot meet
the parent. The child \eqref{R03:eq:truncated} has no analogue for the
integer cubic ratio: there the boundary value is the composite
integer $p^3$. Thus ``least prime strictly above $p^d$'' alone is
not an exact minimum criterion for the affine clock.

\begin{lemma}\label{R03:lem:nesting}
Let $\beta$ be a cubic Pisot number with all shifted floors prime.
Suppose on the scheduled tail that $p_k=T_{n_k}$ and
$e_{n_k}=\beta^{n_k}-T_{n_k}>0$. Then, for every sufficiently large
$k$, the full current prime cell equals the entire prefix cell $J_k$,
and its successor has closure strictly inside it.
\end{lemma}
\begin{proof}
Let $\rho<1$ be the largest modulus of the two small conjugates.
The algebraic integer $\beta^n-T_n$ has nonzero integral norm.
To see nonvanishing, a rational positive power of $\beta$ would
have the same value under every embedding, contradicting the small
conjugates. Its other two conjugates have modulus below $T_n+1$.
Consequently, for large scheduled $n$,
\begin{equation}\label{R03:eq:normlower}
 \tfrac14\beta^{-2n}<(T_n+1)^{-2}<e_n\le2\rho^n.
\end{equation}
For $\delta_n=\beta-T_n^{1/n}$, the mean value theorem yields
\[
 \delta_n\ge\frac{\beta}{4n}\beta^{-3n},\qquad
 \delta_m\le\frac{4\beta}{m}(\rho/\beta)^m.
\]
Since $m=3n+2$, their ratio is at most
$16(n/m)\beta^{-2}\rho^m$, which tends to zero. Hence the lower
endpoints strictly increase eventually. The upper distances satisfy
\[
 (T_n+1)^{1/n}-\beta\sim\frac{\beta}{n\beta^n},
\]
and their successive ratio tends to zero, so the upper endpoints
strictly decrease. Every fixed earlier prefix contains $\beta$ in
its interior, again by nonintegrality of Pisot powers. The shrinking
current cells are eventually contained in that whole earlier prefix.
This proves the assertion without a lower bound on a lacunary cosine.
\end{proof}

\begin{corollary}\label{R03:cor:affineiff}
Under Lemma~\ref{R03:lem:nesting}'s hypotheses, $\beta$ is a local left
minimum if and only if, for every sufficiently large $k$, both of
the following hold:
\begin{enumerate}[label=(\roman*)]
\item $(p_k^{n_{k+1}/n_k},p_{k+1})$ contains no prime;
\item the boundary label $q_0=\floor{p_k^{n_{k+1}/n_k}}$ is composite,
or its truncated cell \eqref{R03:eq:truncated} has no infinite
all-prime continuation.
\end{enumerate}
\end{corollary}
\begin{proof}
Every sufficiently large full prime child extends indefinitely by
Lemma~\ref{R03:lem:extension} below. A prime strictly between
$p_k^{n_{k+1}/n_k}$ and $p_{k+1}$ has a full cell inside the parent:
its upper endpoint is bounded by $p_{k+1}^{1/n_{k+1}}$.
Lemma~\ref{R03:lem:nesting} makes $q_0$ a strictly smaller child for
large $k$. These full children and the one boundary child are all
possible smaller labels. Apply Proposition~\ref{R03:prop:viable}.
\end{proof}

Global leastness requires the earlier least viable choices as well.
In particular, eventual prime-free windows by themselves prove
neither global leastness nor the viability failure of a boundary child.

\subsection{The first affine window and a height theorem}\label{R03:sec:height}

\subsubsection{Exact fifth-power cancellation}

Assume \eqref{R03:eq:residual}, all scheduled floors prime, and $\beta\ge4$.
Put $h=C/\beta$. Then
\begin{equation}\label{R03:eq:h}
 h\le\beta^{-17/20}<\tfrac13.
\end{equation}
The sums and products of the two small roots are
\[
 (s_+,q_+)=(-h(1+\beta^{-1}),h),\qquad
 (s_-,q_-)=(-h(1-\beta^{-1}),-h).
\]
The two-root Newton identity $S_5=s^5-5qs^3+5q^2s$ gives
\begin{equation}\label{R03:eq:e5}
 \begin{aligned}
 e_{5,+}&=h^3(1+\beta^{-1})(5-5z_++z_+^2),
 &z_+&=h(1+\beta^{-1})^2,\\
 e_{5,-}&=h^3(1-\beta^{-1})(5+5z_-+z_-^2),
 &z_-&=h(1-\beta^{-1})^2.
 \end{aligned}
\end{equation}
Here $e_n=-S_n$. Since $z_+<25/48$, the first bracket exceeds
$5-125/48>2$. In the second formula the product of the outside
factor $1-\beta^{-1}$ and the bracket exceeds $5(3/4)>2$.
Therefore
\begin{equation}\label{R03:eq:e5lower}
 e_5>2C^3\beta^{-3}\ge2\beta^{-3}.
\end{equation}
Also $e_5<1$. For the nonreal pair use
$|S_5|\le2h^{5/2}<1$. For the real pair put
$\delta=h(1-\beta^{-1})<h$; then
\[
 e_5=\delta^5+5h\delta^3+5h^2\delta
 <h^5+5h^4+5h^3<61/243<1.
\]

For exponent $17$, the nonreal modulus estimate again gives
$|S_{17}|<1$. In the real case the larger small modulus is below
$\sqrt h+h<11/12$, so $|S_{17}|<2(11/12)^{17}<1$.
All traces are odd: Newton's recurrence modulo $2$ has odd initial
traces, an odd leading coefficient, and equal lower coefficients.
If $S_{17}>0$, the actual floor would be $T_{17}-1$, an even integer
greater than two, contradicting primality. No $S_n$ can vanish,
since that would make $\beta^n$ integral. We have therefore proved
\begin{equation}\label{R03:eq:initial-primes}
 p=T_5=\beta^5-e_5,\qquad P=T_{17}=\beta^{17}-e_{17},
 \qquad 0<e_5,e_{17}<1.
\end{equation}
These are the actual initial floors, not just tail traces.

The mean value theorem now gives
\[
 \begin{aligned}
 \beta^{17}-(\beta^5-e_5)^{17/5}
 &\ge\frac{17}{5}(\beta^5-1)^{12/5}e_5\\
 &>\frac{34}{5}\beta^9(1-\beta^{-5})^{12/5}
 >6\beta^9.
 \end{aligned}
\]
The last inequality follows from Bernoulli's inequality already at
$\beta=4$. Subtracting $e_{17}<1$ proves
\begin{equation}\label{R03:eq:firstwindow}
 \boxed{T_{17}-T_5^{17/5}>5\beta^9.}
\end{equation}
The relevant exponent relative to $\beta^{17}$ is $9/17$, not
$19/34$. The cancellation in \eqref{R03:eq:e5} is essential.

\subsubsection{A full infinite continuation}

\begin{lemma}\label{R03:lem:extension}
Assume \eqref{R03:eq:PX}. If $r\ge2$ is prime, $d\ge3$, and $r^d\ge X$,
there is a prime $r'$ such that
\[
 r^d<r',\qquad r'+1<(r+1)^d.
\]
Consequently every sufficiently large full prime cell extends
indefinitely along the shifted schedule.
\end{lemma}
\begin{proof}
Put $y=r^d$ and $z=y+2y^\theta$. Apply \eqref{R03:eq:PX} at $z$.
Since $z\le3y$ and $3^\theta<2$, the supplied prime satisfies
\[
 y<z-z^\theta\le r'\le z=y+2y^\theta.
\]
For $r\ge2$ and $d\ge3$,
\[
 d r^{(1-\theta)d-1}>3>2+r^{-\theta d},
\]
because $(1-\theta)d-1\ge17/40>0$. Thus convexity gives
\[
 (r+1)^d-r^d\ge d r^{d-1}>2r^{\theta d}+1.
\]
This proves strict inclusion of the child closure inside the parent.
Iterate. The compact cell closures are nested, their diameters tend
to zero, and strict one-step upper inclusion places the limit inside
every original half-open prime cell. Hence the limit is an actual
all-prime base, not merely a limit satisfying weak endpoint inequalities.
\end{proof}

\begin{proof}[Proof of Theorem~\ref{R03:thm:height}]
Suppose $\beta\ge4$ and $\beta^{17}\ge X+2$. In
\eqref{R03:eq:initial-primes}, take $x=P-1>\beta^{17}-2\ge X$.
The prime $q$ supplied by \eqref{R03:eq:PX} satisfies $q\le P-1<P$.
Moreover \eqref{R03:eq:firstwindow} gives
\[
 q-p^{17/5}>5\beta^9-1-\beta^{357/40}>0,
\]
since $357/40<9$. Also
\[
 q>\beta^{17}-2-\beta^{357/40}>\tfrac12\beta^{17},
 \qquad q^3>X.
\]
For the middle inequality, divide by $\beta^{17}$ and use
$2\cdot4^{-17}+4^{-8}<1/2$.

The full $q$-cell at exponent $17$ is strictly inside the first
$p$-cell, because
\[
 p^{17/5}<q,\qquad q+1\le P<\beta^{17}<(p+1)^{17/5}.
\]
Lemma~\ref{R03:lem:extension} then constructs every subsequent child
using the same $X$, since all later ratios are greater than $3$ and
all later prime labels are at least $q$. Its limit $B$ has the
original first prime floor and all later floors prime, while
\[
 B^{17}<q+1\le P<\beta^{17}.
\]
Thus $B<\beta$, contradicting global minimality. This proves the
strict bound.
\end{proof}

Applying the prime theorem at $P$, rather than $P-1$, would not
suffice: it could return the already known prime $P$. Likewise a
finite prime prefix or a tail constant would not suffice without
the complete nesting argument above.

\subsubsection{What the height theorem does and does not provide}

Theorem~\ref{R03:thm:height} is sensitive to global minimality: it
changes the second prime floor, so it does not exclude every
large algebraic local minimum. BHP's original theorem states that
its threshold can in principle be made effective, but does not
print a numerical value at that theorem \cite[p.~532]{BHP}.
No numerical $X$ is established by the present argument. It therefore
does not supply an executable finite coefficient bound by itself.

An enormous explicit upper bound for the global least base already
follows from Saito's construction \cite[Section~2.3 and
Lemma~9.1]{Saito}. For example, his recorded sufficient prime-interval
threshold allows
\[
 X_0=\exp(3\exp(32.76)),\qquad
 \xi(2\cdot3^k-1)<(2X_0^{5/17}+1)^{1/5}.
\]
Thus finiteness in principle is not a new conclusion. The exponent
$2/3$ in that explicit interval estimate is too large to force a
prime into \eqref{R03:eq:firstwindow}; substituting that estimate into
the thin window would be invalid.

There is also a limitation on repeating the argument later. For
fixed $C$ and a fixed scheduled odd exponent $n$, the two-root
power-sum identity gives
\[
 e_n\sim nC^{(n+1)/2}\beta^{-(n+1)/2}\qquad(\beta\to\infty)
\]
within either family. For the transition $n\mapsto3n+2$ the
window consequently has order
\[
 (3n+2)C^{(n+1)/2}\beta^{3(n+1)/2},\qquad
 \eta_n=\frac{3n+3}{6n+4}=\frac12+\frac1{6n+4}
\]
relative to $x=\beta^{3n+2}$. Although $\eta_5=9/17>21/40$,
$\eta_{17}=27/53<21/40$. This is a fixed-index, large-height
comparison, not a fixed-$\beta$, large-index asymptotic.

\subsection{An infinite Pell family with composite shifted floors}
\label{R03:sec:pell}

\subsubsection{The family and its surviving primitive character}

The recurrence \eqref{R03:eq:pell-main} preserves
\begin{equation}\label{R03:eq:pell}
 8m_t^2-13m_t+16=y_t^2.
\end{equation}
Indeed, with $x=16m-13$, it is multiplication of $x+y\sqrt{32}$
by $577+102\sqrt{32}$, a unit of norm one, on
$x^2-32y^2=-343$. Both coordinates increase from $(3,7)$, and
$m_t\equiv3\pmod{24}$ and oddness of $y_t$ are invariant. Hence
\[
 a_t\equiv5\pmod{48},\qquad
 a_t=5,5381,6208949,7165121045,\ldots.
\]
For $a\ge5$, the only possible rational roots of
$X^3-aX^2-X-1$ are $\pm1$, and neither is a root. Its discriminant is
\begin{equation}\label{R03:eq:disc-pell}
 \Delta=-4m(8m^2-13m+16)=-4my^2<0.
\end{equation}
The real root lies in $(a,a+1)$; the other two roots have common
modulus $\beta^{-1/2}<1$. Thus every member is a cubic Pisot unit.

This ray is useful because a natural quadratic-character obstruction
does not exclude it. Modulo $a+1=2m$,
\[
 X^3-aX^2-X-1\equiv(X+1)(X^2-1),
 \qquad T_n\equiv-1\pmod{2m}\quad(n\text{ odd}).
\]
For scheduled $n\equiv1\pmod4$, Newton's recurrence also gives
$T_n\equiv a\equiv1\pmod4$. Since $m\equiv3\pmod8$, reciprocity
therefore gives, for positive such traces,
\begin{equation}\label{R03:eq:primitive}
 \left(\frac{-m}{T_n}\right)
 =\left(\frac{T_n}{m}\right)
 =\left(\frac{-1}{m}\right)=-1.
\end{equation}
The symbol in \eqref{R03:eq:primitive} is the primitive quadratic-field
character attached to the square class $-m$; the coprimality needed
here follows from $T_n\equiv-1\pmod m$.
The earlier projective prime-return test requires this negative sign
for eventual prime traces in the complex unit branch $a\equiv2\pmod3$
\cite{R02}. Thus the ray survives that particular necessary test.
This says nothing about primes dividing the discarded square factor
$y^2$: the full polynomial-discriminant symbol can be zero there.
It also says nothing by itself about primality or minimality.

\subsubsection{Seven exact return certificates}

For a prime $p$ and coefficient residue $a_0$, work in
\[
 R=\F_p[X]/(X^3-a_0X^2-X-1),\qquad x=X\bmod f.
\]
The element $x$ is a unit with $x^{-1}=x^2-a_0x-1$. Its algebra trace is
\begin{equation}\label{R03:eq:trace-ring}
 \tau(v_0+v_1x+v_2x^2)=3v_0+a_0v_1+(a_0^2+2)v_2.
\end{equation}
For a seed $k_0$, put $z=x^{2\cdot3^{k_0}}$. The two finite identities
\begin{equation}\label{R03:eq:certificate}
 \tau(x^{-1}z)=0,\qquad z^{3^H}=z
\end{equation}
imply
$p\mid T_{2\cdot3^{k_0+jH}-1}$ for every $j\ge0$.
This remains valid for non-squarefree reductions: multiplication
trace in the quotient algebra gives the usual power sum with
multiplicities.

\begin{table}[ht]
\centering
\caption{A complete Pell-index cover for Theorem~\ref{R03:thm:pell}}
\label{R03:tab:pell}
\begin{tabular}{rccrrr}
\toprule
$p$ & condition on $t$ & $a_0\pmod p$ & $k_0$ & $H$ & $z=(z_0,z_1,z_2)$\\
\midrule
5  & $0\pmod3$   & 0  & 1  & 2   & $(1,2,1)$\\
5  & $1\pmod3$   & 1  & 10 & 30  & $(4,2,3)$\\
19 & $0\pmod5$   & 5  & 4  & 126 & $(1,11,5)$\\
19 & $2\pmod5$   & 15 & 9  & 126 & $(18,1,1)$\\
29 & $3\pmod5$   & 14 & 7  & 66  & $(22,6,21)$\\
31 & $11\pmod{15}$ & 23 & 8 & 330 & $(23,13,2)$\\
59 & $4\pmod5$   & 23 & 6  & 28  & $(0,10,50)$\\
\bottomrule
\end{tabular}
\end{table}

All entries in Table~\ref{R03:tab:pell} satisfy \eqref{R03:eq:certificate}.
For completeness, the coefficient cycles from the initial Pell
state modulo the five primes are
\[
\begin{array}{rcl}
p=5&:&0,1,4;\\
p=19&:&5,4,15,18,11;\\
p=29&:&5,16,20,14,17;\\
p=31&:&5,18,21,29,20,11,19,22,4,30,13,23,17,27,10;\\
p=59&:&5,12,25,34,23.
\end{array}
\]
The full $(m,y)$ states return with periods $3,5,5,15,5$, respectively;
one need not infer a state period from an accidentally repeated
coefficient. These small cycles and the seven quotient-ring identities
are finite exact certificates. A separate matrix implementation checks
\begin{equation}\label{R03:eq:matrix-pell}
 M=\begin{pmatrix}0&0&1\\1&0&1\\0&1&a_0\end{pmatrix},\qquad
 \tr(M^{2\cdot3^{k_0}-1})=0,\quad
 M^{2\cdot3^{k_0}(3^H-1)}=I\pmod p.
\end{equation}
Binary exponentiation makes these checks short even when $H=330$;
no enormous integer trace is factored.

\begin{proof}[Proof of Theorem~\ref{R03:thm:pell}]
The first two rows cover $t\not\equiv2\pmod3$. Of the remaining
indices, the rows for $19,29,59$ cover the four residues
$0,2,3,4\pmod5$. The only uncovered condition is
$t\equiv2\pmod3$, $t\equiv1\pmod5$, namely $t\equiv11\pmod{15}$,
and the row for $31$ covers it. Thus a fixed prime from the table
divides infinitely many scheduled traces for every fixed $t$.
Those traces tend to infinity and are eventually larger than that
prime, so they are composite.

To pass from traces to floors without a phase assumption, note that
all traces are odd: $T_0=3$, $T_1=a_t$, $T_2=a_t^2+2$ are odd and
$T_n=a_tT_{n-1}+T_{n-2}+T_{n-3}$ preserves oddness. For $n\ge5$,
\[
 |S_n|\le2\beta_t^{-n/2}<1,\qquad
 \floor{\beta_t^n}=
 \begin{cases}T_n,&S_n\le0,\\T_n-1,&S_n>0.\end{cases}
\]
At each sufficiently large certified return, the first option exceeds
its fixed prime divisor and is composite. The second is an even
integer greater than two and is composite as well. Hence infinitely
many scheduled floors are composite, regardless of the sign of $S_n$.
\end{proof}

\subsubsection{The totally real reflected family}

\begin{corollary}\label{R03:cor:real-reflection}
Start at $(m_0,y_0)=(165,469)$ and iterate
\[
 \binom{m_{t+1}}{y_{t+1}}
 =\begin{pmatrix}577&204\\1632&577\end{pmatrix}\binom{m_t}{y_t}
 +\binom{468}{1326},\qquad a_t=2m_t+1.
\]
For every $t\ge0$, the Pisot root of $X^3-a_tX^2-X+1$ has infinitely
many composite shifted floors. The same five primes as in
Theorem~\ref{R03:thm:pell} suffice; the coefficients begin
$331,382699,\ldots$.
\end{corollary}
\begin{proof}
Extend the complex recurrence \eqref{R03:eq:pell-main} to all integer
indices; its affine matrix is integrally invertible. Its index $-1$
state is $(-165,469)$. Reversing the recurrence and changing
$(m,y)$ to $(-m,y)$ gives the displayed real recurrence, and hence
\[
 a_t^{\mathrm R}=-a_{-1-t}^{\mathrm C}.
\]
Writing $f_a^{\mathrm C}(X)=X^3-aX^2-X-1$, the identity
$X^3-aX^2-X+1=-f_{-a}^{\mathrm C}(-X)$ shows that every odd real
trace is the negative of the corresponding complex trace.
The modular state cycles in Table~\ref{R03:tab:pell} hold also at negative
indices, since their maps are invertible. Reflect each Pell-index
class $r\pmod P$ to $-1-r\pmod P$ and negate its coefficient residue.
The same seven return identities therefore give a complete real
cover modulo $15$.

For $a\ge331$, polynomial signs place the roots in
$(-1,0),(0,1),(a,a+1)$, proving irreducibility and the Pisot property.
Writing the small roots as $-\rho,\sigma$ gives
$0<\sigma<\rho<1$, since $\beta>a$.
At every positive odd $n$,
$0<\beta^n-T_n=\rho^n-\sigma^n<1$, so the floor equals the trace.
Every recurring divisor thus supplies infinitely many composite
floors. The reflected seven identities are also checked directly by
the independent polynomial and matrix verifiers, without relying on
the reflection argument.
\end{proof}

\subsubsection{Why a finite fixed-prime cover cannot solve every Pell ray}

The following obstruction is structural, not the outcome of a capped
search. For $s\in\{1,-1\}$, start at $(m_0,y_0)=(0,4)$ and set
\begin{equation}\label{R03:eq:boundary-rays}
 \binom{m_{t+1}}{y_{t+1}}=
 \begin{pmatrix}577&204\\1632&577\end{pmatrix}\binom{m_t}{y_t}
 -s\binom{468}{1326},\qquad a_t=2m_t-s.
\end{equation}
For $s=1$ this is another complex ray with positive-index coefficients
$695,801311,\ldots$; for $s=-1$ it is a real ray with coefficients
$2569,2965345,\ldots$. The index-zero polynomials are reducible and
are discarded. Every positive-index member is an irreducible cubic
Pisot unit of the indicated type. The Pell identities are respectively
$8m^2-13m+16=y^2$ and $8m^2+13m+16=y^2$.

\begin{proposition}\label{R03:prop:nogo}
For either ray \eqref{R03:eq:boundary-rays} and every finite set of primes
$S$, there are infinitely many positive indices $t$ such that
\[
 T_n\equiv-s\pmod p\qquad(p\in S, n\text{ any positive odd integer}).
\]
In particular, no finite set of fixed recurring prime divisors covers
all members of either ray.
\end{proposition}
\begin{proof}
The affine Pell map has linear determinant
$577^2-204\cdot1632=1$. Modulo each prime it is a permutation of the
finite state space, so its orbit from $(0,4)$ is purely periodic.
Let $L$ be a common multiple of these periods for $p\in S$.
At every positive multiple $t$ of $L$, one has $m_t\equiv0\pmod p$
and $a_t\equiv-s\pmod p$. Therefore
\[
 X^3-a_tX^2-X-s\equiv(X+s)(X^2-1)\pmod p.
\]
Its odd power sum, with multiplicities, is $-s+1-1=-s$.
This is nonzero modulo every prime, proving the assertion.
\end{proof}

For the real ray this obstruction occurs even inside the full
shrinking-conjugate region. If $-\rho$ is the larger small root,
\[
 \rho<\beta^{-1/2}+\beta^{-1}<\beta^{-17/40}
 \quad(\beta>2569),
\]
since $1+\beta^{-1/2}<17/16<3/2<\beta^{3/40}$.
It follows neither that any trace is prime nor that any member is
minimal. Adaptive primes, noncongruence arguments, and global
competitors remain possible. Theorem~\ref{R03:thm:pell} and
Corollary~\ref{R03:cor:real-reflection} solve their two specified rays,
not all Pell rays or all unit candidates.

\subsection{An exact bounded theorem for the classical schedule}
\label{R03:sec:real}

\subsubsection{The coefficient locus in both directions}

Use the convention
\begin{equation}\label{R03:eq:realpoly}
 f(X)=X^3-aX^2-BX-c,\qquad q=|c|\ge1,
\end{equation}
so $c$ is the signed norm. Put
\begin{equation}\label{R03:eq:realdisc}
 D=a^2B^2+4B^3-4a^3c-27c^2-18aBc.
\end{equation}

\begin{theorem}\label{R03:thm:locus}
The polynomial \eqref{R03:eq:realpoly} has roots
$\beta>1,-\rho,\gamma$, with $0<|\gamma|<\rho<1$, if and only if
$a,B,q$ are positive integers and one of the following holds:
\begin{enumerate}[label=(\roman*)]
\item $c=-q$ and $a\ge B+q$;
\item $c=q$, $D>0$, $B\le a+q$, and $B<2a+3$.
\end{enumerate}
It is then irreducible, and $a<\beta<a+2$.
Among these polynomials, membership in $\SaitoL$ is equivalent to
\begin{equation}\label{R03:eq:strict-shrink}
 \rho^{40}\beta^{17}<1.
\end{equation}
Writing $\delta=\beta-a$ and $z=\beta^{1/40}$, an equivalent test
using only the large root is
\begin{equation}\label{R03:eq:large-root-test}
 \begin{gathered}
 \delta=B/\beta+c/\beta^2,\qquad
 \rho=\frac{\delta+\sqrt{\delta^2-4c/\beta}}2,\\
 z^{17}(z^{40}-a)<2,\qquad
 c+z^6+az^{23}-z^{63}>0.
 \end{gathered}
\end{equation}
\end{theorem}
\begin{proof}
For negative norm, the stated inequality gives
$f(-1)<0<f(0)$ and $f(1)<0$, placing one root in each of
$(-1,0),(0,1),(1,\infty)$. Also $f(a)=-aB+q<0$, so
$\beta>a$ and the negative small root has larger modulus.
Conversely that configuration gives
$B=\beta(\rho-\gamma)+\rho\gamma>0$ and
$f(-1)<0$, hence $a\ge B+q$ by integrality.

For positive norm the two small roots must be negative. The negative
critical point $u=(a-\sqrt{a^2+3B})/3$ lies in $(-1,0)$ exactly when
$B<2a+3$. Positive discriminant gives a positive local maximum at
$u$. The additional inequality $B\le a+q$ is equivalent to
$f(-1)<0$; together with $f(0)<0$, these locate both negative roots
in $(-1,0)$. The positive root exceeds one and $a$. Conversely the
same signs and critical-point location follow from the root configuration.
The formula $B=\beta(\rho+\sigma)-\rho\sigma$ is positive.

There is no missing zero-trace case. In the positive-norm configuration,
$a=\beta-\rho-\sigma>-1$, so $a\ge0$. If $a=0$, then
$\beta<2$, whence $q=1$ and $B\le1$, giving
$D=4B^3-27<0$, a contradiction. In the negative-norm case $a>0$
directly. Reducibility of a monic cubic would give an integer root;
neither small nonzero root is integral. If $\beta$ were that integer
root, the remaining monic integral quadratic would have nonzero
constant term of absolute value below one, also impossible.
Finally $\beta-a=\rho-\gamma\in(0,2)$.

We recall the norm/redundancy argument from \cite{K031}. If
$\rho\le\beta^{-17/40}$, then
\[
 1\le q=\beta\rho\sigma<\beta\rho^2\le\beta^{3/20},
 \qquad
 \rho^{40}\le\beta^{-17}\le q^{17}\beta^{-17}
 =\rho^{17}\sigma^{17}.
\]
Thus $\rho\le\sigma^{17/23}$ follows. Equality
$\rho^{40}\beta^{17}=1$ is impossible: in the totally real splitting
field, apply an automorphism sending $\beta$ to its negative conjugate
in $\beta^{17}(-\rho)^{40}=1$. The transformed left side is negative.
This proves \eqref{R03:eq:strict-shrink}.

The two small roots give the displayed quadratic formula for $\rho$.
For $R=\beta^{-17/40}$, the larger quadratic root is below $R$ if
and only if $R>\delta/2$ and $R^2-\delta R+c/\beta>0$.
Substituting $\beta=z^{40}$ gives the two inequalities in
\eqref{R03:eq:large-root-test}. Both are needed to select the correct
side of the quadratic.
\end{proof}

\subsubsection{Exhaustiveness and exact arithmetic}

Every member of $\SaitoL$ satisfies the necessary bounds
\begin{equation}\label{R03:eq:coeff-bounds}
 q^{20}<(a+2)^3,\qquad B^{40}<2^{40}(a+2)^{23}.
\end{equation}
The first follows from $q<\beta^{3/20}$ and $\beta<a+2$.
For positive norm, $B<2\beta\rho<2\beta^{23/40}$; for negative
norm the sharper $B<\beta\rho<\beta^{23/40}$ holds. Hence
\eqref{R03:eq:coeff-bounds} bounds \emph{every} possible $B$ and $q$
for a specified trace $a$, not just parameters chosen for a scan.

The independent verifier enumerates all triples satisfying
\eqref{R03:eq:coeff-bounds} with $1\le a\le1000$, then applies the exact
integer root-locus conditions of Theorem~\ref{R03:thm:locus}. It isolates
$\beta\in(a,a+2)$ using integer polynomial signs. For $\rho$ it uses
\[
 g(t)=-f(-t)=t^3+at^2-Bt+c.
\]
When $c<0$, $g$ has one positive root in $(0,1)$. When $c>0$, exact
bisection of $g'(t)=3t^2+2at-B$ first identifies its positive minimum;
the larger zero is then isolated on the increasing branch. This
avoids confusing the two positive roots of $g$.

With denominator $d=2^{32}$, each bracket
$b_l/d<\beta<b_h/d$, $r_l/d<\rho<r_h/d$ is classified by one of
\begin{equation}\label{R03:eq:integer-classification}
 r_h^{40}b_h^{17}<d^{57},\qquad
 r_l^{40}b_l^{17}>d^{57}.
\end{equation}
The first certifies membership and the second nonmembership. Every
one of the $166,052$ root-locus candidates is decided at this precision;
there are no unclassified boundary cases. The implementation permits
adaptive refinement, although none is needed here. All arithmetic is
integral, including the root isolation; floating-point proposals are
not an input to this independent verification.

\begin{table}[ht]
\centering
\caption{Exact counts in the bounded classical locus}
\label{R03:tab:counts}
\begin{tabular}{lr}
\toprule
Stage or signed norm & Number of triples\\
\midrule
Coefficient bounds with $1\le a\le1000$ & 265,126\\
Exact negative-dominant real root locus & 166,052\\
Shrinking locus $\SaitoL$ & 32,232\\
\midrule
$c=-1$ in $\SaitoL$ &20,018\\
$c=1$ in $\SaitoL$ &4,759\\
$c=-2$ in $\SaitoL$ &6,936\\
$c=2$ in $\SaitoL$ &519\\
\bottomrule
\end{tabular}
\end{table}

\subsubsection{A recurrent divisor is a finite certificate of infinitely many floors}

For a prime $p$, work in the finite ring
\[
 R_p=\F_p[X]/(X^3-aX^2-BX-c),\qquad
 u_0=X,\quad u_{r+1}=u_r^3.
\]
The linear algebra trace is
\begin{equation}\label{R03:eq:tau-real}
 \tau(v_0+v_1X+v_2X^2)=3v_0+av_1+(a^2+2B)v_2\pmod p.
\end{equation}
Therefore $\tau(u_r)\equiv V_r\pmod p$. Iterate until an exact
polynomial state repeats. If a zero trace occurs on the eventual
cycle, it occurs on an infinite arithmetic progression of indices.
A transient zero alone is not enough. This argument works also when
$p\mid cD$; neither invertibility nor squarefreeness is assumed.

This verification is independent of the original symmetric Newton
map, whose state transformation is
\[
 (T,U,C)\longmapsto
 (T^3-3TU+3C,\ U^3-3TUC+3C^2,\ C^3).
\]
The independent implementation uses direct polynomial multiplication
and reduction, rather than importing that graph calculation.

\begin{table}[ht]
\centering
\caption{Recurring-divisor certificates for all $32,232$ triples}
\label{R03:tab:prime-counts}
\begin{tabular}{rrr@{\qquad}rrr}
\toprule
$p$ & Newly certified & Remaining & $p$ & Newly certified & Remaining\\
\midrule
2  &18,105&14,127 &29&34&35\\
3  &4,711&9,416 &31&15&20\\
5  &5,196&4,220 &37&11&9\\
7  &2,583&1,637 &41&4&5\\
11 &834&803 &43&3&2\\
13 &411&392 &47&1&1\\
17 &165&227 &53&0&1\\
19 &73&154 &59&1&0\\
23 &85&69 &&&\\
\bottomrule
\end{tabular}
\end{table}

The complete accepted coefficient set, not merely its cardinality,
agrees with the original enumeration. Table~\ref{R03:tab:prime-counts}
records the successive removals, read down the left half and then
the right half. The final two certificates are particularly compact:
\[
 \begin{array}{ccl}
 (a,B,c)=(917,63,1),&p=47:&47\mid V_{8+88j},\\
 (a,B,c)=(845,65,1),&p=59:&59\mid V_{3+28j},
 \end{array}\qquad j\ge0.
\]
Both cycles begin at index $1$. Another illustrative case is
$(a,B,c)=(211,7,-1)$, for which both $7$ and $31$ divide
$V_{2+4j}$ for all $j\ge0$. Thus nonreturn of a selected prime seed,
as in the earlier example in [031], does not prevent a different
prime from giving infinitely recurring divisibility.

\begin{lemma}\label{R03:lem:uniform-floor}
Every $\beta\in\SaitoL$ satisfies $\beta>2$, and
\[
 \floor{\beta^{3^r}}=V_r\qquad(r\ge1).
\]
For $r\ge2$ this integer is at least $512$.
\end{lemma}
\begin{proof}
For negative norm, $a\ge B+q\ge2$ and $\beta>a$.
For positive norm, if $\beta\le2$, then $a=1$ and
$\rho+\sigma=\beta-1$, implying
\[
 1\le q=\beta\rho\sigma
 \le\frac{\beta(\beta-1)^2}{4}\le\frac12,
\]
a contradiction. For $M=3^r$, oddness of $M$ and negative dominance
give
\[
 0<\beta^M-V_r=\rho^M-\gamma^M<2\rho^M
 <2\cdot2^{-51/40}<1\quad(r\ge1).
\]
The floor identity follows. For $r\ge2$, $\beta^{3^r}>2^9=512$.
\end{proof}

\begin{proof}[Proof of Theorem~\ref{R03:thm:finite}]
The coefficient bounds and exact tests above enumerate the locus
exhaustively for the stated trace range. The associated finite data
provide one recurrent-divisor certificate for each of its $32,232$
triples. The finite-ring cycle argument proves divisibility on the
recorded infinite progression. By Lemma~\ref{R03:lem:uniform-floor},
every certified return with $r\ge2$ is a floor equal to $V_r\ge512$,
and hence is strictly larger than its divisor $p\le59$. It is
therefore composite. Every progression has infinitely many such
indices, proving the theorem.
\end{proof}

\subsection{Reproducibility and the remaining problems}\label{R03:sec:limits}

\subsubsection{Finite evidence and its exact role}

The accompanying supplement contains the following proof artifacts.
They are ordinary Python~3 programs and JSON data; the theorem
statements do not rely on exploratory scans or probabilistic primality.
\begin{itemize}\raggedright
\item \path{minimality_verify.py} checks the fifth-shadow identities,
rational inequalities, and leading cancellation terms used in
Section~\ref{R03:sec:height}. The analytic prime theorem and the infinite
continuation are proved or cited in the text, not inferred from this
finite program.
\item \path{pell_cover_certificate.json},
\path{pell_cover_real_certificate.json}, and
\path{pell_cover_verify.py} contain and verify the seven matrix
certificates for each ray, the Pell invariants, and both covers of
all $15$ coefficient-index residue classes.
\path{discriminant_pell_cover_audit.py} independently checks the
$14$ return identities by polynomial multiplication.
\item \path{totreal_independent_verify.py} performs the integer-only
root isolation and quotient-ring cycle verification in
Section~\ref{R03:sec:real}. Its output
\path{totreal_independent_verify.json} records one recurrent-divisor
certificate for every accepted triple. The corresponding log records
the counts in Tables~\ref{R03:tab:counts} and \ref{R03:tab:prime-counts}.
\end{itemize}
The SHA-256 digest of the independent totally real verifier is
\begin{center}\small\ttfamily
6a514e149ec38f5203daa92bd925a54eaa2955746676bdca013f501cdb7a1bfb
\end{center}
Running the Pell verifier and then the totally real verifier reproduces
the two finite certificate checks. The larger JSON file is supplemental
data rather than a printed code appendix. The mathematical proofs
explain why each finite identity entails an infinite conclusion.

\subsubsection{What remains open after these results}

For the shifted clock, the exact global criterion is least viable
choice after \emph{every} prefix; the exact local criterion must also
keep the truncated boundary child. The first-window theorem turns
a usable numerical prime-gap threshold into a numerical height bound,
but no such value is certified in this report. An independent small
admissible base with a proved infinite continuation could also supply
a useful bound. A long finite all-prime prefix is not a substitute.

The Pell theorem and its real reflection remove two specified
unbounded families without minimality or phase assumptions.
Proposition~\ref{R03:prop:nogo}
shows why one cannot extrapolate its five-prime cover to every
boundary-based ray. It leaves open adaptive recurring-divisor arguments,
phase arguments, and comparisons with smaller admissible bases.

For the classical schedule, Saito's Problem~1.1 asks for infinitely
many composite cubic Pisot floors at powers $3^r$
\cite{SaitoIntervals}. His theorem for iterated reversible recurrences
does not cover the inner recurrence $3^r$. Theorem~\ref{R03:thm:finite}
settles the specified bounded part of the shrinking locus, but every
trace coefficient $a>1000$ remains outside its scope. The earlier
fixed-norm density-one theorem in [031] is stronger asymptotic
background, yet density one does not remove a possibly nonempty
exceptional family. No finite trace bound for a hypothetical
algebraic classical Mills shadow is proved here.

Finally, universal compositeness of cubic Pisot floors is a sufficient
route to transcendence of the least Mills constant, not an established
equivalent formulation. An algebraic prime-generating branch need
not be the least branch. Keeping that distinction is precisely why
the global-minimum argument of Section~\ref{R03:sec:height} and the
unconditional family exclusions are stated separately.



\endgroup

\section{Finite prime forests and explicit continuation criteria}\label{src:R05}
The preceding cells give genuine finite forests, but a large forest alone does not prove infinite survival. This section records the quantitative gap-budget conditions that would provide that continuation and identifies the missing effective estimate.
\medskip
\begingroup



\subsection{The finite and infinite questions}

Put
\begin{equation}\label{R05:eq:schedule}
 n_k=2\cdot3^k-1\quad(k\ge1),\qquad
 n_1=5,\quad n_2=17,\quad n_3=53,\quad n_{k+1}=3n_k+2.
\end{equation}
Write $\xi_2$ for the least member of
\[
 \mathcal W_2=\{A>1:\ \floor{A^{n_k}}\in\PP\text{ for every }k\ge1\}.
\]
The existence of this minimum follows from
\cite[Theorem~1.3]{Saito}, since $n_{k+1}\ge2n_k$.
We seek a small member of $\mathcal W_2$, not merely a long finite
sequence of prime floors.

For a prime $p$, the half-open interval
\[
 [p^{1/n},(p+1)^{1/n})
\]
is the prime cell at exponent $n$. Its closed envelope is denoted
$J_n(p)=[p^{1/n},(p+1)^{1/n}]$.
A cell is \emph{reachable} if its closed envelope is nested in a
chain of cells at the preceding required exponents. In our finite
constructions the containments in earlier half-open cells are strict.

Let $p_1,p_2,\ldots$ be the consecutive primes. Throughout the paper,
\begin{equation}\label{R05:eq:moment}
 M_2(X)=\sum_{p_j\le X}(p_{j+1}-p_j)^2.
\end{equation}
The gap crossing $X$ is included. This convention is essential in the
interval-counting argument below.

\begin{mainthm}[Unconditional finite forests]\label{R05:thm:finite}
The schedule \eqref{R05:eq:schedule} has the following finite forests.
\begin{enumerate}[label=\textup{(\alph*)}]
\item Starting with the first prime floor $31$, there are $1151$
strictly internal full cells at exponent $17$, and at least
\[
 \Nzero=149\,564\,695\,915
\]
distinct reachable full cells at exponent $53$. All their bases lie
in $[31^{1/5},131064^{1/17}]\subset(1,2)$.
\item A forest with $4332$ first prime labels, from $241361$ through
$295703$, has at least
\[
 \Nb=107\,204\,570\,624\,222
\]
reachable full cells at exponent $17$, all with bases in
$[241361^{1/5},295704^{1/5}]\subset(1,25/2)$.
\item Restricting the second forest to its $1865$ first prime labels
at least $271927$ leaves at least
\[
 \No=52\,537\,179\,616\,066
\]
full cells, with base lower endpoint $271927^{1/5}$.
Separately, the second forest contains at least
\[
 \Nu=90\,210\,603\,204\,803
\]
full cells whose prime labels at exponent $17$ lie in
$[2^{61},2^{62})$.
\end{enumerate}
\end{mainthm}

This theorem uses the published finite bound $1476$ for prime gaps
in its verified range \cite{R05-OHP}. It does not assert that any of its
cells has descendants at every subsequent level.

\begin{mainthm}[Explicit conditional continuations]\label{R05:thm:conditional}
Each row below is a sufficient hypothesis for its stated conclusion.
Here the hypothesis is $M_2(X)\le C X^\tau$ for every $X\ge X_0$.
\begin{center}
\begin{tabular}{@{}llll@{}}
\toprule
$C$&$\tau$&$X_0$&Conclusion\\\midrule
$6\cdot10^{15}$&$31/25$&$2^{161}$&$\xi_2<2$\\
$2\cdot10^9$&$4/3$&$2^{161}$&$\xi_2<2$\\
$15\cdot10^{19}$&$31/25$&$(25/2)^{53}$&$\xi_2<25/2$\\
$15\cdot10^{11}$&$4/3$&$(25/2)^{53}$&$\xi_2<25/2$\\
\bottomrule
\end{tabular}
\end{center}
None of these numerical moment hypotheses is established in this paper.
\end{mainthm}

\begin{mainthm}[Conditional long-gap alternatives]\label{R05:thm:tailmain}
Put $X_0=(25/2)^{53}$ and $G(X)=X^{2/3}/64$. Either of the estimates
\begin{align*}
 \sum_{\substack{p_j\le X\\p_{j+1}-p_j\ge G(X)}}(p_{j+1}-p_j)
   &\le10^{13}X^{2/3},\\
 \sum_{\substack{p_j\le X\\p_{j+1}-p_j\ge G(X)}}(p_{j+1}-p_j)^2
   &\le7.5\cdot10^{11}X^{4/3},
\end{align*}
if valid for every $X\ge X_0$, implies $\xi_2<25/2$.
These are alternative sufficient hypotheses, not claimed equivalent
to one another or to the full-moment hypotheses. Neither is proved here.
\end{mainthm}

The finite forests and the implications in
Theorems~\ref{R05:thm:conditional} and \ref{R05:thm:tailmain} are the results proved here.
An unspecified asymptotic constant, a sufficiently-large threshold,
or a finite prime prefix is not a substitute for the displayed
hypotheses. The later sections identify concrete issues in applying
known prime-gap results at these particular scales.

\subsection{Certified finite forests}
\label{R05:sec:finite}

\subsubsection{A finite prime-gap counting principle}

Oliveira e Silva, Herzog and Pardi computed the consecutive-prime
gaps below $4\cdot10^{18}$, with maximum $1476$
\cite[Section~2.2 and Table~8]{R05-OHP}. The upper endpoint of a gap must
lie within the computed range when invoking this datum.
Suppose an integer interval $[A,B]$ lies in a range for which all
enclosing consecutive-prime gaps are covered. Subdivide $(A,B]$
into $\floor{(B-A)/1476}$ disjoint blocks of length $1476$, ignoring
a final remainder. Every such block $(u,u+1476]$ contains a prime:
the last prime at or below $u$ has its successor no farther than
$1476$ away. Hence
\begin{equation}\label{R05:eq:count1476}
 \#([A,B]\cap\PP)\ge\floor{\frac{B-A}{1476}}.
\end{equation}
This is a conservative count, not an enumeration of all primes in
these very large intervals.

If $m>n$ are consecutive required exponents, define
\begin{equation}\label{R05:eq:AB}
 A_{p;n,m}=\floor{p^{m/n}}+1,\qquad
 B_{p;n,m}=\floor{(p+1)^{m/n}}-1.
\end{equation}
A prime $q$ between these endpoints supplies a full child cell.
All endpoint evaluations below use integer roots: for instance
$\floor{p^{m/n}}$ is the unique integer $z$ satisfying
$z^n\le p^m<(z+1)^n$.
The supplied data also verify the strict upper inequality
$(B_{p;n,m}+1)^n<(p+1)^m$ for every retained row.
The intervals for distinct parent primes are pairwise disjoint.

\subsubsection{The seed with first floor 31}

Let $a_0=31^{1/5}$. Exact arithmetic gives
\[
 \ceil{31^{17/5}}=117661.
\]
There are exactly $1152$ primes between $117661$ and $2^{17}-1$;
the first is $117671$ and the last is $131071$.
Their closed exponent-$17$ cells lie in $[a_0,2]$.
The last cell has upper endpoint $2$ and is not contained in the
half-open first cell $[31^{1/5},2)$, so it is discarded.
The remaining $1151$ primes end at $131063$, and all retained closed
cells lie in
\begin{equation}\label{R05:eq:smallwindow}
 [a_0,b_0],\qquad b_0=131064^{1/17}<2.
\end{equation}

For every retained prime $q$, apply \eqref{R05:eq:AB} with $(n,m)=(17,53)$.
The smallest and largest integer endpoints across these rows are
\[
 6435135249309836\quad\hbox{and}\quad9005485420906690.
\]
They are below $2^{53}$. Bertrand's postulate supplies a prime above
any such endpoint and below $2^{54}<4\cdot10^{18}$, so both endpoints
of every enclosing prime gap lie in the verified finite range.
Formula~\eqref{R05:eq:count1476} now gives the exact certified lower count
\begin{equation}\label{R05:eq:N0}
 \sum_{\substack{q\in\PP\\117661\le q\le131063}}
 \floor{\frac{B_{q;17,53}-A_{q;17,53}}{1476}}
 =\Nzero=149564695915.
\end{equation}
These are distinct reachable cells. No individual primality claims
about $149$ billion listed integers are being made; the count is a
consequence of the verified maximum-gap result and exact endpoints.

\subsubsection{The larger and restricted forests}

Put $H=399\cdot10^{16}=3.99\cdot10^{18}$, and let $\mathcal P$ be
the primes in
\begin{equation}\label{R05:eq:first-range}
 \ceil{(2\cdot10^{18})^{5/17}}
 \ \le p\le\ \floor{H^{5/17}}-1.
\end{equation}
The integer limits are $241355$ and $295714$.
There are $4332$ such primes, with first $241361$ and last $295703$.
For $p\in\mathcal P$, set $A_p=A_{p;5,17}$ and $B_p=B_{p;5,17}$.
The endpoints range from
\[
 2000175614492238064\quad\hbox{to}\quad3989463330471287939,
\]
and every row satisfies $2\cdot10^{18}<A_p<B_p<H$.

To justify the rightmost enclosing prime, use the corrected
Cully-Hugill--Lee interval theorem \cite[Theorem~1 and Table~2]{R05-CHL}.
At $x=4\cdot10^{18}$ its parameter $\Delta=39097000$ supplies a prime
in $(x(1-1/\Delta),x]$. The exact inequality
\begin{equation}\label{R05:eq:right-end}
 4\cdot10^{18}(39097000-1)>H\cdot39097000
\end{equation}
places this prime above $H$. Since $4\cdot10^{18}$ is composite,
the prime is strictly below the finite computation's upper limit.
Thus the gap bound $1476$ applies to every block used here.
Consequently
\begin{equation}\label{R05:eq:Nb}
 \sum_{p\in\mathcal P}\floor{\frac{B_p-A_p}{1476}}
 =\Nb=107204570624222.
\end{equation}
All bases lie between $241361^{1/5}$ and $295704^{1/5}$; the inequality
$32\cdot295704<25^5$ proves that the latter is below $25/2$.

For the fixed-window improvement, retain only first primes
\[
 p\ge\ceil{(3\cdot10^{18})^{5/17}}=271924.
\]
This subset $\mathcal P_{\mathrm O}$ has $1865$ members, from
$271927$ through $295703$, and
\begin{equation}\label{R05:eq:No}
 \sum_{p\in\mathcal P_{\mathrm O}}
 \floor{\frac{B_p-A_p}{1476}}=\No=52537179616066.
\end{equation}
Its larger lower base, $271927^{1/5}$, improves the later moment budget
more than the lost nodes diminish it.

For a different restriction, clip every large-forest integer interval
to a dyadic band. Define
\[
 L_{p,j}=\max(A_p,2^j),\qquad
 U_{p,j}=\min(B_p,2^{j+1}-1),\qquad j\in\{60,61\}.
\]
Counting blocks only when $L_{p,j}<U_{p,j}$ gives
\begin{equation}\label{R05:eq:bands}
 \sum_{p\in\mathcal P}\max\left(0,
    \floor{\frac{U_{p,j}-L_{p,j}}{1476}}\right)
 =\begin{cases}
 16993967419419,&j=60,\\
 90210603204803,&j=61.
 \end{cases}
\end{equation}
The latter is $\Nu$. These counts can lose short blocks at the band
boundary, so they are not inferred by partitioning the already
rounded total \eqref{R05:eq:Nb}; they are recomputed from clipped endpoints.
Equations~\eqref{R05:eq:N0}, \eqref{R05:eq:Nb}, \eqref{R05:eq:No} and
\eqref{R05:eq:bands}, with their exact endpoint tests, prove
Theorem~\ref{R05:thm:finite}.

\subsection{A gap budget for a whole generation}
\label{R05:sec:cauchy}

\begin{lemma}[Clipped-gap Cauchy inequality]\label{R05:lem:cauchy}
Let $I_1,\ldots,I_N$ be real intervals in $[2,X]$ with disjoint
interiors. Let $L$ be their total length, and let $R$ count the primes
in their interiors. Then
\begin{equation}\label{R05:eq:cauchy}
 R\ge\frac{L^2}{M_2(X)}-N.
\end{equation}
In particular, if all intervals have length at least $\ell$, then
$M_2(X)\le N\ell^2/2$ implies $R\ge N$.
\end{lemma}
\begin{proof}
Split each interval at its interior primes. There are $R+N$ resulting
segments, with total length $L$. Each segment lies in a gap between
consecutive primes. If one full gap meets several intervals, the
corresponding clipped segments have disjoint interiors; the sum of
their squared lengths is no greater than the square of the full gap.
All gaps used have left endpoint at most $X$, including a possible
gap crossing $X$. Thus the sum of all segment-length squares is at
most $M_2(X)$. Cauchy--Schwarz gives
$L^2\le(R+N)M_2(X)$, proving the result.
\end{proof}

At exponent $n$, write $d=(3n+2)/n$. For each reachable prime parent
$p$ consider the real interval
\begin{equation}\label{R05:eq:childinterval}
 I_p=[p^d,(p+1)^d-1].
\end{equation}
The intervals belonging to distinct parents are disjoint.
We use only primes $q$ in their \emph{interiors}. Then
\[
 p^d<q,\qquad q+1<(p+1)^d,
\]
so $J_{3n+2}(q)$ is strictly inside the parent prime cell.
This endpoint convention gives strict nesting directly, without
needing to rule out integral values of the real endpoints.

\begin{lemma}[Passage to a base]\label{R05:lem:branch}
Suppose a finite reachable frontier consists of full cells in a
fixed interval $[a,b_*]$, where $1<a<b_*<b$. If at each subsequent
level one can retain a positive fixed number of cells, each strictly
inside a retained parent, then $\mathcal W_2$ contains a number below
$b$, and therefore $\xi_2<b$.
\end{lemma}
\begin{proof}
The retained cells form a finitely branching forest with finitely
many roots and nodes at every depth. K\"onig's lemma gives an infinite
branch. Its nested closed cells lie in $[a,b_*]$. Their diameters tend
to zero, since the mean value theorem gives, for $p\ge a^n$,
\[
 (p+1)^{1/n}-p^{1/n}\le\frac1{n a^{n-1}}.
\]
Their unique common point lies in the interior of every earlier
cell, because the next cell is strictly internal. It has all the
required prime floors and belongs to $[a,b_*]\subset(1,b)$.
\end{proof}

\subsection{Continuation in a fixed base window}
\label{R05:sec:fixed}

\begin{theorem}[Fixed-window criterion]\label{R05:thm:window}
Suppose a reachable frontier at exponent $n_*$ has at least $N_*$
full cells in $[a,b_*]$, where $1<a<b_*<b$.
Let $C>0$ and $\tau>0$. The three conditions
\begin{align}
 M_2(X)&\le C X^\tau\qquad(X\ge b^{3n_*+2}),\label{R05:eq:windowH}\\
 a^4&\ge b^{3\tau},\label{R05:eq:windowgrowth}\\
 C&\le\frac{2N_*a^{4n_*+4}}{b^{\tau(3n_*+2)}}\label{R05:eq:windowinitial}
\end{align}
imply $\xi_2<b$.
\end{theorem}
\begin{proof}
Retain $N_*$ parents at the starting level. At any later level $n$,
let $N\ge N_*$ retained parents have bases in the same window.
Their child intervals \eqref{R05:eq:childinterval} lie below
$X=b^{3n+2}$. Convexity and $p\ge a^n$ give
\begin{equation}\label{R05:eq:length}
 |I_p|\ge d p^{d-1}-1
 \ge3a^{2n+2}-1\ge2a^{2n+2}.
\end{equation}
Lemma~\ref{R05:lem:cauchy} therefore gives
\[
 R\ge\frac{4N^2a^{4n+4}}{C b^{\tau(3n+2)}}-N.
\]
This is at least $N$ whenever
$C\le2N_*a^{4n+4}/b^{\tau(3n+2)}$.
Condition~\eqref{R05:eq:windowgrowth} makes that right-hand side
nondecreasing in $n$, and \eqref{R05:eq:windowinitial} verifies it initially.
We may retain $N_*$ children and repeat indefinitely.
Lemma~\ref{R05:lem:branch} completes the proof.
\end{proof}

For the seed $31$, use $(a,b,n_*,N_*)=(31^{1/5},2,53,\Nzero)$ and
$\tau=31/25$. The exact integer checks are
\begin{equation}\label{R05:eq:smallchecks}
 \begin{aligned}
 (6\cdot10^{15})^{25}2^{31\cdot161}
   &<(2\Nzero)^{25}31^{1080},\\
 31^{20}&>2^{93}.
 \end{aligned}
\end{equation}
They prove the first implication in Theorem~\ref{R05:thm:conditional}.
The coefficient ceiling in \eqref{R05:eq:windowinitial} is approximately
$6.3834568534\cdot10^{15}$; this decimal is descriptive, not a proof step.
More generally this forest works for
\[
 \tau<\frac{4\log31}{15\log2}=1.3211190161\ldots,
 \qquad C\le\frac{2\Nzero31^{216/5}}{2^{161\tau}}.
\]

For the larger restricted forest, use
\[
 (a,b,n_*,N_*)=(271927^{1/5},25/2,17,\No).
\]
The checks for $C=15\cdot10^{19}$ and $\tau=31/25$ are
\begin{equation}\label{R05:eq:optimizedchecks}
 \begin{aligned}
 (15\cdot10^{19})^{25}25^{31\cdot53}
   &<(2\No)^{25}271927^{360}2^{31\cdot53},\\
 271927^{20}2^{93}&>25^{93}.
 \end{aligned}
\end{equation}
Thus the third implication in Theorem~\ref{R05:thm:conditional} holds.
The precise coefficient ceiling is approximately
$1.5441314472\cdot10^{20}$, and
$\log_{10}((25/2)^{53})=58.13623\ldots$.
The hypothesis is needed only at the scheduled values
$b^{n_{k+1}}$, but the all-$X$ form is convenient and sufficient.

\subsection{Dyadic concentration at exponent four thirds}
\label{R05:sec:dyadic}

The exponent restriction of the fixed-window criterion is not
intrinsic. Concentrating the children into a dyadic band gives a
different continuation theorem at the critical exponent $4/3$.

\begin{theorem}[Dyadic criterion]\label{R05:thm:dyadic}
Suppose at exponent $n_*\ge17$ a reachable frontier contains at least
$N_*$ full cells with prime labels in $[Y_*,2Y_*)$, where $Y_*\ge4$
is a power of two. Suppose all these cells have bases in
$[a,b_*]\subset(1,b)$ and $N_*\ge60C$.
If
\[
 M_2(X)\le C X^{4/3}\qquad(X\ge X_0)
\]
and $(2Y_*+1)^{(3n_*+2)/n_*}\ge X_0$, then $\xi_2<b$.
\end{theorem}
\begin{proof}
At any stage let $N\ge N_*$ retained parent labels lie in $[Y,2Y)$,
and set $d=(3n+2)/n$ and $X=(2Y+1)^d$.
Since $n\ge17$, we have $3<d\le53/17<13/4$.
Every child interval has length at least $2Y^{d-1}$.
Moreover
\begin{equation}\label{R05:eq:criticalratio}
 \frac{4Y^{2d-2}}{X^{4/3}}
 =\frac{4Y^{2d/3-2}}{(2+1/Y)^{4d/3}}>\frac1{10}.
\end{equation}
Indeed $Y^{2d/3-2}\ge1$, $2+1/Y\le9/4$, and
$4d/3<9/2$, while $(9/4)^{9/2}=(3/2)^9<40$.
The moment hypothesis and Lemma~\ref{R05:lem:cauchy} imply
\begin{equation}\label{R05:eq:Rdyadic}
 R>\frac{N^2}{10C}-N\ge5N.
\end{equation}

The child labels lie between $Y^d$ and $(2Y+1)^d$. Their ratio is
less than $(9/4)^{10/3}<16$; the last inequality is equivalent to
$3^{20}<2^{32}$. This range meets at most five dyadic bands, so one
band contains at least $R/5>N$ child labels. Retain $N_*$ in that
band and call its lower endpoint $Y'$.
Because it contains a label exceeding $Y^d$, $Y'>Y^d/2$.
For the next exponent ratio $d'>3$, its moment endpoint satisfies
\[
 X'=(2Y'+1)^{d'}>Y^{3d}>(2Y+1)^d=X,
\]
since $Y^3>2Y+1$. Thus successive moment endpoints increase, and
all remain above the stated onset. The process can continue forever;
Lemma~\ref{R05:lem:branch} gives the desired base.
\end{proof}

For the seed-$31$ forest, all exponent-$53$ labels lie in
$[2^{52},2^{53})$. Since
\[
 \Nzero>60(2\cdot10^9),\qquad
 (2^{53}+1)^{161/53}>2^{161},
\]
Theorem~\ref{R05:thm:dyadic} proves the second implication in
Theorem~\ref{R05:thm:conditional}.

For the large forest, use the directly counted upper band from
\eqref{R05:eq:bands}, with $Y_*=2^{61}$ and $n_*=17$.
The exact capacity and onset checks are
\begin{equation}\label{R05:eq:upperchecks}
 \Nu=90210603204803>60(15\cdot10^{11}),\qquad
 (2^{62}+1)^{53/17}>(25/2)^{53}.
\end{equation}
The latter follows from the exact integer comparison $25^{17}<2^{79}$.
This proves the fourth implication in Theorem~\ref{R05:thm:conditional}.
The proof does not require a prime between every two consecutive
cubes, nor does it require every parent interval to contain a prime.

\subsection{Only the long gaps need be controlled}
\label{R05:sec:tails}

For $X_0=(25/2)^{53}$ and $G(X)=X^{2/3}/64$, put
\[
 T_j(X)=\sum_{\substack{p_i\le X\\p_{i+1}-p_i\ge G(X)}}
       (p_{i+1}-p_i)^j\qquad(j=1,2).
\]
As in $M_2$, the full gap crossing $X$ is included if it meets the
size cutoff. We prove Theorem~\ref{R05:thm:tailmain} using the upper-band
frontier of size $N=\Nu$; at every subsequent level retain this same
number of nodes in one dyadic band.

Use $\ell=2Y^{d-1}$, $X=(2Y+1)^d$, and let $L$ be the total length
of the child intervals. By \eqref{R05:eq:criticalratio},
\[
 L\ge N\ell,\qquad \ell^2>X^{4/3}/10,\qquad G(X)<\ell/12.
\]
The last comparison follows from $10\cdot12^2<64^2$.
The range of children meets at most five dyadic bands, as before.

\begin{proof}[First-moment alternative]
Assume $T_1(X)\le10^{13}X^{2/3}$ and suppose that fewer than $5N$
prime children occur in the interiors. Splitting the child intervals
at these primes produces fewer than $6N$ clipped segments.
Those contained in full gaps shorter than $G(X)$ have total length
less than $6NG(X)<N\ell/2\le L/2$.
The remaining clipped segments have total length at most $T_1(X)$:
even when a long full gap meets several intervals, the clipped
lengths are disjoint and sum to at most that full gap's length.
Since $N=90210603204803>8\cdot10^{13}$ and
$\ell>X^{2/3}/\sqrt{10}>X^{2/3}/4$,
\[
 T_1(X)\le10^{13}X^{2/3}<N\ell/2\le L/2.
\]
The short and long parts cannot add up to $L$, a contradiction.
Thus there are at least $5N$ children, and one dyadic band retains
at least $N$ of them. The moment endpoints stay above $X_0$ and
increase by the proof of Theorem~\ref{R05:thm:dyadic}; strict nesting and
Lemma~\ref{R05:lem:branch} finish the construction.
\end{proof}

\begin{proof}[Second-moment alternative]
Assume $T_2(X)\le C X^{4/3}$, where $C=7.5\cdot10^{11}$.
In the clipped-gap proof, the squares of the segments lying in
short full gaps sum to at most $G(X)L$. Those from long gaps have
square sum at most $T_2(X)$. The exact capacity $N\ge120C$ gives
\[
 T_2(X)<10C\ell^2\le\frac{N\ell^2}{12}
       \le\frac{L^2}{12N},\qquad
 G(X)L<\frac{L^2}{12N}.
\]
Cauchy--Schwarz now gives
\[
 R+N\ge\frac{L^2}{T_2(X)+G(X)L}\ge6N.
\]
Hence $R\ge5N$, and the same dyadic retention proves the conclusion.
\end{proof}

These hypotheses avoid controlling all the smaller prime gaps.
They are not inferred merely from the existence of numerical
zero-density constants: an actual large-gap estimate at the stated
onset is still required.

\subsection{What the available prime-gap theorems supply}
\label{R05:sec:analytic}

\subsubsection{An asymptotic exponent is not a numerical initialization}

Stadlmann proves, for every $\varepsilon>0$,
\begin{equation}\label{R05:eq:stadlmann}
 M_2(X)\ll_\varepsilon X^{1.23+\varepsilon}
 \qquad\text{\cite[Theorem~1]{R05-Stadlmann}}.
\end{equation}
The convention in that theorem is the same as \eqref{R05:eq:moment}, so
there is no missing crossing-gap term. Taking $\varepsilon=0.01$
supplies the exponent $31/25$, but does not supply the particular
coefficient and onset in Theorem~\ref{R05:thm:conditional}.
A bound with unknown constant and unknown starting point cannot be
applied to the first unknown generation of a prescribed finite forest.

Nor is a direct numerical substitution in the proof automatic.
For example, \cite[Lemma~21, p.~43]{R05-Stadlmann} uses
\[
 D=\floor{\frac{\log(x^\beta/x^{1/\log\log x})}{\log2}},
 \qquad \beta\in[0.01,0.15],
\]
in a decomposition whose intended large-$x$ regime requires
$x^{1/\log\log x}\le x^\beta$.
At the two fixed-window onsets,
\[
 \begin{array}{c|c|c}
 x&\log\log x&1/\log\log x\\\hline
 2^{161}&4.71489\ldots&0.21209\ldots\\
 (25/2)^{53}&4.89682\ldots&0.20421\ldots
 \end{array}
\]
so this ordering fails for every allowed $\beta$.
Increasing the external coefficient $C$ does not repair that
parameter-ordering condition.

A second issue appears in \cite[Lemma~20 and equation~(5.15)]{R05-Stadlmann}.
The sieve application has weights supported on $d<x^{0.0001}$,
and a remainder involves
\[
 \frac{x}{\tau}(\log x)^{2-\log\log\log x/(8K_2)},
 \qquad K_2=10000,
\]
which is subsequently absorbed into an arbitrary negative power of
$\log x$ in the sufficiently-large regime.
At the displayed onsets $x^{0.0001}<1.014$, so that support does not
even contain an integer $d\ge2$. The literal asymptotic absorption
therefore is not a quantitative error estimate at our initialization
scales. Other small-parameter choices, large-value constants and
prime-counting errors also require numerical control.

These observations rule out claiming a direct specialization of
\eqref{R05:eq:stadlmann} from its displayed asymptotic bounds. They do
not disprove any moment hypothesis in this paper, and are not lower
bounds on the true moment. No claim of inherent ineffectivity is
made: extracting or replacing the necessary constants is a separate
analytic problem.

\subsubsection{A corrected pointwise bound is quantitatively too weak}

The corrected version of Cully-Hugill--Lee \cite{R05-CHL} gives intervals
$(x(1-1/\Delta),x]$ containing a prime. At log-thresholds $60,90,105$,
its parameters are respectively
\[
 7.69184\cdot10^{10},\qquad1.84304\cdot10^{11},\qquad
 D_*=1.91886\cdot10^{11}.
\]
At $X=2^{161}$ these rows, together with Bertrand's postulate for the
smaller range, give a uniform bound
$G=X/(D_*-1)$ for gaps whose left endpoints are at most $X$.
Summing $g^2\le Gg$ over consecutive gaps gives
\begin{equation}\label{R05:eq:pointmoment}
 M_2(X)\le G(X+G).
\end{equation}
At this first unverified scale, the right-hand side of
\eqref{R05:eq:pointmoment}, divided by $6\cdot10^{15}X^{31/25}$,
lies strictly between $5.927\cdot10^9$ and $5.928\cdot10^9$.
The portable verification checks this comparison using rational
powers. Thus this particular pointwise-to-moment deduction does not
meet the required budget. It says nothing against a sharper
mean-square estimate.

\subsubsection{Other routes retain a quantitative gap}

Maynard's proof of the earlier exponent $5/4+\varepsilon$
\cite{Maynard} is another possible starting point for an explicit
argument, but does not state a constant and onset that establish a
row of Theorem~\ref{R05:thm:conditional}. The $+\varepsilon$ cannot
simply be discarded. Its Lemma~5.5 also uses a zero-free-region
saving to absorb a power $(\log x_1)^{246}$, without the numerical
threshold needed here.

There are genuinely explicit estimates for tails of large gaps.
For example, \cite[Theorem~3.3]{R05-AFGW} bounds the first moment of gaps
larger than a parameter $N$, with explicit constants, but assumes
$x\ge\exp(4000)$ and $N\ge8x^{1/4+\alpha}\log^3x$.
This onset is far beyond either prescribed frontier scale, and that
first-moment result is not, by itself, one of our second-moment
hypotheses. These distinctions matter: a numerical zero-density
estimate or an asymptotic exceptional-set bound is not automatically
an explicit branching theorem at a chosen seed.

\subsection{Arithmetic scope and the remaining step}

The forest construction is independent of an algebraic-minimum
classification. It proves finite prime-floor realizations and
conditional upper bounds for the least infinite base.
The earlier necessary-form and finite-certificate results
\cite[Corollary~4.1 and Proposition~5.2]{N01} exclude an algebraic
$\xi_2<100$, subject to their stated structural inputs.
Indeed the norm bound gives $|c|\le\xi_2^{3/20}<2$ and hence
$c=\pm1$, while the inverse-trace relation gives $b=\pm1$.
The Pisot property puts the trace coefficient in $0\le a\le101$.
Changing the sign of the algebraic number when $c=-1$ reduces both
norm signs to the signed coefficient grid in the cited proposition;
its recurrent-divisor certificates exclude the schedule $r=2$.
Thus, \emph{with those earlier results as an explicit dependency},
any large-forest hypothesis of Theorem~\ref{R05:thm:conditional} or
\ref{R05:thm:tailmain} would imply transcendence of $\xi_2$.
Neither the structural reduction nor that earlier finite certificate
is reproved here, and neither is needed for the forest theorems.

In particular, none of the arguments above proves unconditionally
that $\xi_2<2$, that $\xi_2<25/2$, or that $\xi_2$ is transcendental.
They also do not resolve the classical Mills constant, whose schedule
is different. The outstanding step is now a quantitative analytic
statement: prove one of the displayed global moment budgets, or a
suitable clipped-gap budget on precisely the child intervals needed
by the forest.

A bound merely counting empty parents would need additional
multiplicity control: a nonempty parent could have only one child,
which might die in the next generation. Lemma~\ref{R05:lem:cauchy}
controls the total number of children, and the two continuation
arguments preserve a whole frontier. This is why the conditional
implications give an actual infinite branch rather than only
arbitrarily long, unrelated prefixes.


\subsection{Portable reproduction data}
\label{R05:app:data}

The accompanying \texttt{certificates} directory contains the files
in Table~\ref{R05:tab:files}. The finite source bound $1476$ is a cited
external computation \cite{R05-OHP}; these programs do not re-sieve
all primes below $4\cdot10^{18}$.
They verify the much smaller first-label lists, every exact integer
endpoint, all block-count sums, the dyadic clipping, and the integer
comparisons used in the continuation criteria.

\begin{table}[htbp]
\centering\small
\caption{Finite reproduction files.}\label{R05:tab:files}
\begin{tabular}{@{}p{0.58\linewidth}p{0.36\linewidth}@{}}
\toprule
File&Contents or role\\\midrule
\texttt{branching\_seed\_verify.py}&Sieve and integer-root computation\\
\texttt{branching\_seed\_verify.json}&1152 exponent-17 label rows\\
\texttt{branching\_seed\_broad\_verify.json}&4332 large-forest endpoint rows\\
\texttt{branching\_seed\_optimized\_verify.json}&1865 restricted-forest rows\\
\texttt{audit\_branching\_seed\_verify.py}&Trial-division and integer-root check\\
\texttt{audit\_branching\_broad\_verify.py}&Large forest and dyadic checks\\
\texttt{audit\_branching\_seed\_verification.json}&Seed and moment-comparison results\\
\texttt{audit\_branching\_broad\_verification.json}&Large-forest results\\
\texttt{audit\_branching\_optimized\_verification.json}&Restricted forest and upper-band counts\\
\bottomrule
\end{tabular}
\end{table}

The first program uses Python's standard library only. The two
checking programs additionally use SymPy's exact
\texttt{integer\_nthroot}; their primality lists use trial division
rather than the first program's sieve. Run without Python's
optimization flag, so assertion checks remain enabled:
\begin{lstlisting}
cd certificates
python3 branching_seed_verify.py
python3 audit_branching_seed_verify.py
python3 audit_branching_broad_verify.py
\end{lstlisting}
The programs record passing terminal results and reproduce the
exact counts in Theorem~\ref{R05:thm:finite}.
The four data sums can also be reconstructed directly from
\eqref{R05:eq:N0}, \eqref{R05:eq:Nb}, \eqref{R05:eq:No} and \eqref{R05:eq:bands};
no floating-point approximation enters those definitions.

\paragraph{Frozen endpoint-data SHA256 digests.}
Each line below is a single 64-character digest; any wrapping is
only typographical.
\begin{lstlisting}
branching_seed_verify.json
aeb167c9de52d80287700cc6bcbcbb168581b2fa9d1bbfe96f360391cffeea51
branching_seed_broad_verify.json
6489814babe0e5514f6e72deca4baf2017b28d0426237153dcdec774bb03cda2
branching_seed_optimized_verify.json
5dbcb247af22d3b86c51ee08a43e9cb7649650c4ab969bfa257c4910455e0378
\end{lstlisting}



\endgroup

\section{Shifted prime patterns and compactness criteria}\label{src:R09}
Two different bridges are useful here: prime-pattern theorems give translated finite configurations, whereas compactness converts uniformly bounded full-prefix feasibility into one infinite base. The uniform bound and the absence of forbidden boundary limits are essential.
\medskip
\begingroup



\subsection{The question and the scope}

A Pisot number is a real algebraic integer $\beta>1$ whose other
conjugates have modulus less than one. If degree one is permitted,
every integer $p+1>1$ is a Pisot number vacuously. Its powers of exponent
at least two are composite, however, and it is not the degree-three
object appearing in cubic algebraic-candidate reductions. The useful
question is whether $p+1$ can be a \emph{coefficient or trace} of a
genuine cubic Pisot number, and what prime-pattern theorems then imply.
Here $p$ always denotes a positive rational prime, not a Pisot number.

For a Pisot number, write $T_n=\Tr(\beta^n)$ and
$S_n=T_n-\beta^n$. The integer $T_n$ is the sum of all conjugate powers;
$S_n\to0$ exponentially. Our two clocks retain their original indices:
\begin{equation}\label{R09:eq:clocks}
 E_k^+=3^k,\qquad E_k^-=2\cdot3^k-1,\qquad k\ge1.
\end{equation}
Both consist of odd exponents. A prime integer part, a finite pattern
among several different Pisot numbers, and prime floors on the whole
power orbit of one fixed number are different conclusions.

The external inputs are used only in their stated forms:
Green--Tao gives finite arithmetic progressions in a positive-relative-density
prime set \cite[Theorems 1.1--1.2]{GreenTao}; Maynard's linear-forms conclusion
gives arbitrarily many, not necessarily all, prime values in a
sufficiently large admissible fixed tuple \cite[p.~2]{Maynard};
L\^e--Wolf permits shifted-prime polynomial steps with a free additive
translation \cite[Theorem 4]{R09-LW}. We prove the algebraic deductions
and the compactness statement below. No historical priority claim is made.

\subsection{Two actual cubic realizations of prime plus one}

\begin{proposition}\label{R09:prop:unit}
For every prime $p\ge2$, the polynomial
\[
 F_p(X)=X^3-(p+1)X^2+X-1
\]
has a cubic Pisot unit root $\alpha_p$ satisfying
\[
 p<\alpha_p<p+1,\qquad \Tr(\alpha_p)=p+1,
 \qquad \floor{\alpha_p}=p.
\]
Furthermore $\floor{\alpha_p^2}=(p+1)^2-2$, which is composite for every
odd prime $p$.
\end{proposition}
\begin{proof}
Put $a=p+1\ge3$. We have $F_p(p)=-p^2+p-1<0$ and $F_p(p+1)=p>0$.
The only possible rational roots are $\pm1$, and neither is a root.
The discriminant is
\[
 D(a)=-4a^3+a^2+18a-31.
\]
It is negative at $a=3$, and $D'(a)=-12a^2+2a+18<0$ for $a\ge3$.
Thus there is one real root $\alpha_p$ and a nonreal conjugate pair.
Their product is $1/\alpha_p$, so each has modulus
$\alpha_p^{-1/2}<1$. The constant term gives norm one.

Set $\epsilon=a-\alpha_p$. The equation gives
$\epsilon=(\alpha_p-1)/\alpha_p^2$, hence $0<\epsilon<1/\alpha_p$ and
\[
 \alpha_p^2-(a^2-2)=\frac2{\alpha_p}-\epsilon^2\in(0,1).
\]
For odd $p$, $a$ is even, so $a^2-2>2$ is even.
\end{proof}

In the entire family $X^3-aX^2+X-1$, $a\ge3$, the same proof gives
\[
 \floor{\alpha_a}\text{ is prime}\quad\Longleftrightarrow\quad a-1
 \text{ is prime}.
\]
This exact equivalence realizes the shifted-prime idea without any
prime-pattern theorem. Its square calculation does not settle either
clock in \eqref{R09:eq:clocks}, which never uses exponent two.

\begin{lemma}\label{R09:lem:real}
Let $C\ge1$ and $a>2C$ be integers. The polynomial
\[
 f_{a,C,-1}(X)=X^3-aX^2-CX+C
\]
is irreducible and has roots $-w,v,\beta$ with
$0<v<w<1$ and $a<\beta<a+1$. Thus $\beta$ is totally real cubic Pisot,
and for every positive odd $n$,
\begin{equation}\label{R09:eq:real-floor}
 0<\beta^n-T_n=w^n-v^n<1,\qquad \floor{\beta^n}=T_n.
\end{equation}
\end{lemma}
\begin{proof}
The values at $-1,0,1$ have signs $-,+,-$, since
$f(-1)=2C-a-1$ and $f(1)=1-a$. Also $f(a)=C(1-a)<0$ and
$f(a+1)=(a+1)^2-Ca>0$. There is one root in each of
$(-1,0),(0,1),(a,a+1)$. None is integral, so a monic integral cubic
cannot have a rational root and is irreducible. The two small roots
sum to $a-\beta<0$, proving $w>v$ and \eqref{R09:eq:real-floor}.
\end{proof}

For the signed polynomial $f_{a,C,s}=X^3-aX^2-CX-sC$, $s\in\{1,-1\}$,
Newton's identities give
\begin{equation}\label{R09:eq:rec}
\begin{split}
 T_0&=3,\quad T_1=a,\quad T_2=a^2+2C,\\
 T_n&=aT_{n-1}+CT_{n-2}+sCT_{n-3}\qquad(n\ge3).
\end{split}
\end{equation}
These identities define integral trace polynomials even at parameters
where the polynomial is reducible.

\begin{proposition}\label{R09:prop:even}
For every prime $p\ge11$, the dominant root $\beta_p$ of
\[
 X^3-(p+1)X^2-4X+4
\]
is a totally real cubic Pisot number of trace $p+1$ and norm $-4$.
Every floor $\floor{\beta_p^n}$ at a positive odd exponent is composite.
In particular this holds at every index of both clocks \eqref{R09:eq:clocks}.
Moreover,
\[
 p+1<\beta_p<p+2,\qquad
 \frac{4(p+1)}{(p+2)^2}<\beta_p-(p+1)<\frac4{p+1}.
\]
\end{proposition}
\begin{proof}
Apply Lemma~\ref{R09:lem:real} with $a=p+1$ and $C=4$.
Since $a$ and $C$ are even, \eqref{R09:eq:rec} makes every $T_n$, $n\ge1$,
even. Equation~\eqref{R09:eq:real-floor} identifies the odd floors with
these traces, which exceed two. Finally,
$\beta_p-a=4(\beta_p-1)/\beta_p^2$; the function
$(x-1)/x^2$ decreases for $x>2$, giving the displayed bounds.
\end{proof}

Euclid's infinitude of primes already makes this an infinite family.
It lies in the elementary even-coefficient region, not in the
remaining odd-$a$ locus. More generally, in Lemma~\ref{R09:lem:real} with
$a$ even, an even $C$ gives the same parity obstruction. If $C$ is odd,
the trace sequence modulo two has period seven and initial residues
$(1,0,0,1,0,1,1)$. Consequently its classical-clock traces are even
when $k$ is even, and its negative-clock traces are even when
$k\equiv0,3,5\pmod6$. This follows directly from \eqref{R09:eq:rec} modulo two.

\subsection{A shifted-prime trace does not determine its floor}

\begin{lemma}\label{R09:lem:sign}
Suppose $|S_n|<1$ and $T_n=q+1$ for an odd prime $q$. Then
\[
 \floor{\beta^n}=
 \begin{cases}q,&S_n>0,\\q+1,&S_n\le0.\end{cases}
\]
\end{lemma}
\begin{proof}
Use $\beta^n=T_n-S_n$ and the strict bound on $|S_n|$.
\end{proof}

The sign distinction occurs at a classical scheduled exponent, not
only in abstract examples. Let $\beta$ be the real root of
\[
 X^3-6X^2-2X-2.
\]
Its discriminant is $-2092$, and its unique real root lies in $(6,7)$.
It is irreducible, since a rational root would be integral and real.
The nonreal conjugates have modulus $\sqrt{2/\beta}<1$. Thus $\beta$ is
cubic Pisot and $T_3=258=257+1$, but
\begin{equation}\label{R09:eq:257}
 \floor{\beta^3}=257\quad\hbox{is prime}.
\end{equation}
For an exact check, set $l=6363/1000$ and $u=1591/250$. Then
\[
 f(l)=-\frac{28939853}{10^9}<0,
 \quad f(u)=\frac{221571}{15625000}>0,
 \quad 257<l^3<u^3<258.
\]
Trial division up to $\sqrt{257}<17$ proves the primality assertion.

There are two immediate limitations in the signed family. If $a$ is
odd, \eqref{R09:eq:rec} makes \emph{every} positive trace odd: the coefficient
sum modulo two is $a+C+sC\equiv1$. Thus $T_n=q+1$ with odd prime $q$
cannot occur anywhere in this odd-$a$ locus. Also
\begin{equation}\label{R09:eq:norm}
 T_n\equiv a^n\pmod C\qquad(n\ge1).
\end{equation}
If $\gcd(a,C)=1$, putting $p+1$ into the norm parameter $C$ does not
make any prime divisor of $C$ divide a trace. Neither a composite norm
nor a shifted-prime trace alone supplies the desired floor conclusion.

\subsection{Prime coefficient patterns with certified composite powers}

We now use a family in which both a prime first floor and composite
scheduled floors are rigorously controlled. For
\begin{equation}\label{R09:eq:class}
 a=43+186t,\quad t\ge0,\qquad f_a(X)=X^3-aX^2-4X+4,
\end{equation}
let $\beta_a$ be its dominant root. Lemma~\ref{R09:lem:real} gives
$\floor{\beta_a}=a$ and exact odd-power traces.

\begin{proposition}\label{R09:prop:31}
Every member of \eqref{R09:eq:class} satisfies
\[
\begin{array}{ll}
31\mid\floor{\beta_a^{3^k}},& k\equiv1\pmod4,\\
31\mid\floor{\beta_a^{2\cdot3^k-1}},& k\equiv3\pmod4.
\end{array}
\]
All these indicated terms, with $k\ge1$, are composite.
\end{proposition}
\begin{proof}
Since $a\equiv12\pmod{31}$,
$f_a(X)=(X+6)^2(X+7)$ modulo $31$. Hence, for odd $n$,
\[
 T_n\equiv 2(-6)^n+(-7)^n
       =-6^n(2-4^n)\pmod{31}.
\]
The element $4$ has order five modulo $31$ and $4^3=2$, so the trace
vanishes exactly when $n\equiv3\pmod5$. The powers of $3$ modulo five
have period four, giving the two index classes. The actual floors
are $T_n$, and $n\ge3$, $\beta_a>43$ make them larger than $31$.
\end{proof}

This is a direct repeated-reduction calculation; it needs no claim
about ramification in the cubic number field. It is also independent
of whether $a$ itself is prime.

\subsubsection{Arithmetic progressions of prime integer parts}

The prime number theorem in the fixed progression $43\pmod{186}$ gives
\begin{equation}\label{R09:eq:density}
 \frac{\#\{p\le x:p\equiv43\pmod{186}\}}{\pi(x)}
 \longrightarrow \frac1{\varphi(186)}=\frac1{60}.
\end{equation}
For the needed fixed-modulus assertion one may also use the
Bombieri--Vinogradov estimate recalled in \cite[\S1]{Maynard}.
Infinitude alone would not justify a relative-density application.

\begin{corollary}\label{R09:cor:gt}
For every fixed $m\ge2$ and every fixed lower cutoff, there are infinitely
many nonconstant arithmetic progressions
\[
 a_j=A+jD\quad(0\le j<m),\qquad D>0,
\]
above that cutoff, consisting of primes in the class $43\pmod{186}$.
The corresponding $m$ cubic Pisot numbers $\beta_{a_j}$ have prime
integer parts in that exact progression and share both composite-index
classes of Proposition~\ref{R09:prop:31}.
\end{corollary}
\begin{proof}
Deleting a finite initial segment preserves the positive density
\eqref{R09:eq:density}. Apply the relative Green--Tao theorem
\cite[Theorem 1.2]{GreenTao}, followed by Proposition~\ref{R09:prop:31}.
\end{proof}

The \emph{Pisot values themselves are not asserted to form an AP}.
Indeed the root equation is
\[
 a=g(\beta_a)=\beta_a-\frac4{\beta_a}+\frac4{\beta_a^2},
 \qquad 0<\beta_a-a<\frac4a.
\]
For $\beta>8$, $g'>0$ and $g''=8(3-\beta)/\beta^4<0$.
Its inverse is strictly convex, so an equally spaced coefficient
triple gives a non-equally-spaced root triple. The roots are, however,
within $4/\min_j a_j$ of their coefficient AP; this error can be made
arbitrarily small by raising the cutoff.

\subsubsection{Bounded clusters from admissible linear forms}

\begin{corollary}\label{R09:cor:clusters}
For each fixed $m\ge1$, there exists a finite $H_m$ such that infinitely
many intervals of length $H_m+1$ contain at least $m$ distinct totally
real cubic Pisot numbers from \eqref{R09:eq:class} with prime integer parts.
Every selected number has the composite-power progressions in
Proposition~\ref{R09:prop:31}.
\end{corollary}
\begin{proof}
Choose $K$ sufficiently large for Maynard's positive-linear-forms
conclusion to give at least $m$ primes. Put
\[
 M=\prod_{\substack{\ell\le K\text{ prime}\\\ell\nmid186}}\ell,
 \qquad L_i(n)=186n+43+186Mi\quad(0\le i<K).
\]
This tuple is admissible. If $\ell\mid186$, all its values are
$43\not\equiv0\pmod\ell$. If $\ell\le K$ and $\ell\nmid186$, then
$\ell\mid M$, so just one residue of $n$ modulo $\ell$ is forbidden.
If $\ell>K$ and $\ell\nmid186$, at most $K<\ell$ residues are forbidden.
Thus no prime divides the product for every integer $n$.

The linear-forms statement in \cite[p.~2]{Maynard} supplies infinitely
many $n$ with at least $m$ prime values. Their coefficient span is at
most $H_m=186M(K-1)$. Since $\beta_a\in(a,a+1)$, the root span is less
than $H_m+1$. Distinct coefficients have distinct roots, and the fixed
index certificates remain valid.
\end{proof}

We have not claimed simultaneous primality of all $K$ forms, nor used
a general assertion that arbitrary positive-density prime subsets
inherit Maynard's theorem. The explicit admissibility argument is why
this application is justified.

\subsection{The shifted-prime polynomial theorem meets the trace recurrence}

\begin{proposition}\label{R09:prop:normalize}
Fix $C\ge1$ and $s\in\{1,-1\}$. For each odd $n\ge1$, define
\[
 P_n(m)=T_n(m-s,C,s)+s.
\]
Then $P_n\in\Z[m]$ and $P_n(0)=0$. For every finite set $I$ of positive
odd exponents, the shifted-prime theorem of L\^e--Wolf gives infinitely
many configurations in which all numbers
\begin{equation}\label{R09:eq:lw}
 x+P_n(p+1)=x+T_n(p+1-s,C,s)+s,\qquad n\in I,
\end{equation}
are prime, with $p$ prime and $x$ an integer. The same holds in a prime
subset of positive relative upper density.
\end{proposition}
\begin{proof}
Integrality follows from \eqref{R09:eq:rec}. At $a=-s$ the polynomial
factors as
\[
 f_{-s,C,s}(X)=(X+s)(X^2-C).
\]
The two square-root contributions cancel at odd powers, leaving
$T_n(-s,C,s)=-s$. Hence $P_n(0)=0$. These are precisely the polynomial
hypotheses of \cite[Theorem 4]{R09-LW}, including its $p+1$ version.
\end{proof}

This is an exact algebraic connection, not an analogy based on the
word ``prime.'' The theorem still includes a free translation:
\[
\begin{array}{c|c|c}
\text{sign}&\text{coefficient }a=p+1-s&\text{required }x\text{ for actual traces}\\
\hline
s=-1 & p+2 & 1\\
s=1 & p & -1
\end{array}
\]
Only on the slice $x=-s$ do the values in \eqref{R09:eq:lw} equal the
untranslated traces. The cited theorem does not prescribe this slice.
For real Pisot parameters satisfying Lemma~\ref{R09:lem:real}, actual odd
traces are the desired floors. In a nonreal family a shadow-sign
condition remains necessary even after the trace values are fixed.

The obstruction to fixing $x$ is concrete. Since $P_1(m)=m$, the real
slice $x=1$ already asks for $p$ and $p+2$ both prime when exponent one
is included. Infinitely many such configurations would imply infinitely
many twin primes. This is a warning about the proposed strengthening,
not an equivalence between the Pisot problem and the twin-prime problem.
For the $s=1$ slice, the first trace is automatically $p$, but the third is
\[
 p^3+3C(p+1),
\]
a prescribed cubic at prime arguments. A linear-pattern theorem does
not supply its infinitely many prime values. Tao--Ziegler's polynomial
progression theorem also retains an additive translation
\cite{TZ2008}; ``polynomial'' does not remove the fixed-slice issue.

\subsubsection{Why the translation is not another choice of Pisot root}

\begin{lemma}\label{R09:lem:rigidity}
Let $\beta,\gamma>1$ be Pisot numbers. If, on an unbounded sequence
$n_j$ of positive integers,
\[
 \Tr(\gamma^{n_j})=\Tr(\beta^{n_j})+c
\]
for one fixed constant $c$, then $\gamma=\beta$ and $c=0$.
Within a fixed $(C,s)$ signed family, even a common translation of
both $T_1$ and $T_3$ at two positive coefficients forces equality of
the coefficients and zero translation.
\end{lemma}
\begin{proof}
Each full trace is its dominant power plus an exponentially decaying
error. If $\gamma>\beta$, the difference of dominant powers tends to
infinity; if $\gamma<\beta$, it tends to minus infinity. Thus the roots
are equal, and their full traces are equal.

For the finite assertion, $T_1(a')-T_1(a)=c$ gives $c=a'-a$. But
\[
 T_3(a')-T_3(a)
 =(a'-a)(a'^2+a'a+a^2+3C).
\]
The factor in parentheses exceeds one, so equality with $c$ forces
$a'=a$ and $c=0$.
\end{proof}

This does not forbid fitting a finite list by unrelated algebraic
data. It forbids silently absorbing a constant shift into one fixed
Pisot trace profile, or into the same signed family at its first two
odd trace coordinates.

\subsection{What linearization and prime witnesses would have to accomplish}

The first trace polynomials are
\[
 T_1=a,\qquad T_3=a^3+3C(a+s),\qquad
 T_5=a^5+5C(a+s)(a^2+C).
\]
For fixed $C,s$, $T_n$ is monic of degree $n$ in $a$, by the recurrence.
Choosing $a$ prime or one above a prime does not linearize its higher
traces. Fixing $a$ instead makes the real third trace
$a^3+3(a-1)C$ affine-linear in $C$, but the Pisot locus $a>2C$ restricts
$C$ to a finite interval. Dirichlet's theorem for an unrestricted
progression does not guarantee a prime in this particular interval.
Its residue is locally admissible exactly when $3\nmid a$, since
$\gcd(a^3,3(a-1))=\gcd(a^3,3)$; and the fifth trace is already quadratic
in $C$. A successful isolated step would not prove a whole prime clock.

Similarly, abundance of primes cannot replace dependence on a fixed
candidate. A criterion requiring a divisor $\ell\mid a$ with
$\ell\equiv2\pmod3$ cannot apply to $a=7^j$, however rich its witness
prime set may be. For $C=4$ and $j\ge2$, these coefficients still lie
in the real Pisot locus. This is only a limitation of that particular
criterion, not evidence that their scheduled floors are prime.

Here is an exact candidate-dependent alternative. Let $\beta$ satisfy
Lemma~\ref{R09:lem:real}, let $\ell\nmid C$ be prime, and let $M$ be its
companion matrix modulo $\ell$. The matrix is invertible, and
$\Tr(M^n)\equiv T_n\pmod\ell$. Put
\[
 E=\ord(M)=3^vU,\quad \gcd(U,3)=1,\qquad
 Z_\ell=\{e\bmod E:\Tr(M^e)=0\}.
\]
For a fixed positive integer $r$, the exponent sequence $r3^k-1$ is
periodic modulo $E$ once $k\ge\max(v,1)$. Consequently
\begin{equation}\label{R09:eq:orbit}
 Z_\ell\text{ intersects that eventual orbit}
 \quad\Longleftrightarrow\quad
 \ell\mid T_{r3^k-1}\text{ for infinitely many }k.
\end{equation}
Indeed if $k_0\ge\max(v,1)$ is a hit and $3^h\equiv1\pmod U$, then
$E\mid r3^{k_0}(3^{th}-1)$ for every $t\ge0$; conversely, infinitely
many hits include an eventual one. For even $r$, these exponents are
odd, so \eqref{R09:eq:real-floor} turns the hits into composite floors once
the trace exceeds $\ell$. The same argument applies to the classical
orbit $3^k$.

Thus one genuine exclusion strategy is to find such a prime $\ell$
\emph{for each candidate}. Prime patterns in an unrelated additive
parameter do not themselves establish \eqref{R09:eq:orbit}. A different
genuine strategy would be to establish actual finite prime-floor
prefixes with a bound independent of prefix length. The next section
shows exactly why that latter input would suffice for our two clocks.

\subsection{A boundary-safe equivalence for returning exponent clocks}

Individual prime-floor constraints are not closed: numbers tending
to $4$ from below may have prime floor $3$, but their limit has floor
$4$. A compactness argument must therefore address integer endpoints.
The following return condition supplies the necessary arithmetic.

\begin{theorem}[Divisibility-return compactness]\label{R09:thm:compact}
Let $(C_k)_{k\ge1}$ be a sequence of positive integers such that
\begin{equation}\label{R09:eq:return}
 \text{for every }m\text{ there are }l>m\text{ and }q\in\mathbb Z_{\ge2}
 \text{ with }C_l=qC_m.
\end{equation}
Fix $1<B<\infty$. The following are equivalent:
\begin{enumerate}[label=(\roman*)]
\item For every $N$, there is $A_N\in(1,B]$ such that
$\floor{A_N^{C_k}}$ is prime for every $1\le k\le N$.
\item There is $A\in(1,B]$ such that $\floor{A^{C_k}}$ is prime for
 every $k\ge1$.
\end{enumerate}
The finite witnesses need not have a common prime prefix.
\end{theorem}
\begin{proof}
Only (i)$\Rightarrow$(ii) needs proof. The first constraint gives
$A_N\ge2^{1/C_1}>1$. Pass to a convergent subsequence with indices
tending to infinity and limit $A\in[2^{1/C_1},B]$. Fix $m$.
If $A^{C_m}$ is not an integer, its floor is locally constant and
hence prime. If it is a prime integer, the desired condition already
holds.

Otherwise $A^{C_m}=t$ is a composite integer, with $t\ge4$.
For all sufficiently late terms, $A_N\ge A$ would force
$\floor{A_N^{C_m}}=t$, contrary to (i). Thus these terms approach $A$
from below. Choose $l,q$ from \eqref{R09:eq:return}. We then have
$A_N^{C_l}\to t^q$ from below, so eventually, with $N\ge l$,
\[
 \floor{A_N^{C_l}}=t^q-1
 =(t-1)(1+t+\cdots+t^{q-1}).
\]
Both factors exceed one. This contradicts primality at coordinate
$l$ and excludes every forbidden integer boundary.
\end{proof}

Both clocks \eqref{R09:eq:clocks} satisfy \eqref{R09:eq:return}. For $3^k$,
take $l=m+1$, $q=3$. For $C_m=2\cdot3^m-1$, we have
$\gcd(C_m,3)=1$. If $L=\ord_{C_m}(3)$, then
$C_{m+L}\equiv C_m\equiv0\pmod{C_m}$ and $C_{m+L}>C_m$.
If one instead starts the negative clock at $k=0$, its first exponent
is $1$, which divides the next exponent $5$; the other returns are
unchanged.

\begin{corollary}[Finite-prefix minima]\label{R09:cor:minima}
Let $W_N=\{A>1:\floor{A^{C_k}}\text{ is prime for }1\le k\le N\}$.
Every nonempty $W_N$ has a least element $a_N$. If all $W_N$ are
nonempty, the sequence $(a_N)$ is nondecreasing. Under
\eqref{R09:eq:return}, it is bounded if and only if an all-prime base
exists; when bounded, its limit is exactly the least all-prime base.
\end{corollary}
\begin{proof}
Choose a witness $U\in W_N$. Below $U$ only finitely many tuples
$(p_1,\ldots,p_N)$ of prime floors occur. The cell of each tuple is
\[
 [2^{1/C_1},U]\cap
 \bigcap_{k=1}^N[p_k^{1/C_k},(p_k+1)^{1/C_k}).
\]
Every nonempty such interval contains its left endpoint. Their finite
union therefore has a minimum, which is also the global minimum of
$W_N$. The inclusion $W_{N+1}\subseteq W_N$ gives monotonicity.
If the minima are bounded, apply the proof of
Theorem~\ref{R09:thm:compact} to their convergent sequence. Its limit is
an all-prime base. Every all-prime base is at least every $a_N$, so
this limit is the least one. Conversely any all-prime base bounds
all $a_N$ above.
\end{proof}

If some $W_N$ is empty, no all-prime base exists. With the convention
$\min\varnothing=+\infty$, the statement can equivalently be written
as convergence of the extended finite-prefix minima to the extended
least all-prime base.

For comparison, a fixed bound $B$ also gives a purely finite argument
if only cubic Pisot candidates are allowed: there are finitely many
such numbers at most $B$. Their integer polynomial coefficients are
bounded, for example by $B+2$, $2B+1$, and $B$ in absolute value for
the three elementary symmetric functions. Decreasing nonempty
finite sets of surviving candidates have nonempty intersection.
Theorem~\ref{R09:thm:compact} is stronger: it treats all real bases and
handles the nonclosed endpoints by an arithmetic return.

\subsection{Conclusions and the remaining input}

There are genuine bridges, but they prove different things.
Proposition~\ref{R09:prop:unit} realizes trace $p+1$ and prime first floor
$p$. Proposition~\ref{R09:prop:even} realizes a nearby real cubic whose
odd floors are all composite. Corollaries~\ref{R09:cor:gt} and
\ref{R09:cor:clusters} organize many distinct Pisot numbers by prime
coefficients without changing their recurring composite powers.
Proposition~\ref{R09:prop:normalize} identifies the exact shifted-prime
polynomial pattern and the missing translation $x=-s$.

For either returning clock, the correct positive existence target is
\emph{every full finite prime-floor prefix with one uniform bound on
the base}. Theorem~\ref{R09:thm:compact} then supplies one fixed infinite
base; compatibility is not an additional assumption under its return
hypothesis. The prime-pattern results used here supply neither those
full power prefixes nor that uniform bound. Conversely, the modular
criterion \eqref{R09:eq:orbit} requires a witness depending on each fixed
candidate, not merely an abundant set of primes.

No assertion here closes all odd-coefficient cubic candidates,
constructs an algebraic base with a full prime power clock, or proves
a new transcendence statement for a Mills-type minimum. In particular
the classical Mills case and the negative-shift minimum with exponents
$2\cdot3^k-1$ are not resolved by these bridges.


\subsection{Portable exact checks}

The accompanying \texttt{certificates/} directory contains standard-library
Python checks and their JSON outputs. Run from the package directory:
\begin{verbatim}
python3 certificates/shifted_prime_trace_verify.py
python3 certificates/prime_pattern_bridge_verify.py
\end{verbatim}
The first checks the rational isolation in \eqref{R09:eq:257}, primality
of $257$, sample parity identities, and exact trace normalization.
The second checks polynomial coefficients, the full residue periods
used in the prime-$31$ proof, and finite instances of the unit and
admissibility calculations. Neither performs a numerical test of an
infinite prime-pattern theorem. The universal algebraic, admissibility,
and endpoint arguments are the proofs in the paper, not extrapolations
from the checks.

\begin{center}
\begin{tabular}{ll}
\toprule
Portable file & Purpose\\
\midrule
\texttt{shifted\_prime\_trace\_verify.py} & Exact trace and root checks\\
\texttt{shifted\_prime\_trace\_verify.json} & Recorded result\\
\texttt{prime\_pattern\_bridge\_verify.py} & Coefficient and period checks\\
\texttt{prime\_pattern\_bridge\_verify.json} & Recorded result\\
\bottomrule
\end{tabular}
\end{center}
The package README specifies scope and execution. External papers are
cited by links and are not redistributed. No third-party proof is
claimed to be reproduced by the finite certificates.



\endgroup

\clearpage\part{Coefficient density and obstructions to fixed congruence covers}

\section{Nonunit cubic families and fixed modulus limitations}\label{src:R06}
The next sections turn from a distinguished least constant to coefficient families of cubic Pisot numbers. This first step supplies pointwise congruence certificates and proves why any fixed finite collection of primes can miss infinitely many members.
\medskip
\begingroup



\subsection{The families and the main results}

For integers $C>1$ and $s\in\{1,-1\}$, put
\begin{equation}\label{R06:eq:family}
 f_{a,C,s}(X)=X^3-aX^2-CX-sC.
\end{equation}
A Pisot number is a real algebraic integer $\beta>1$ all of whose
other conjugates have modulus below one. When \eqref{R06:eq:family}
is its minimal polynomial, let $T_n$ be the sum of the $n$th powers
of its three conjugates. Then
\begin{equation}\label{R06:eq:traces}
 \begin{gathered}
 T_0=3,\quad T_1=a,\quad T_2=a^2+2C,\\
 T_n=aT_{n-1}+CT_{n-2}+sCT_{n-3}\qquad(n\ge3).
 \end{gathered}
\end{equation}
In the ranges considered below, $s=1$ gives a nonreal conjugate pair,
and $s=-1$ gives two small real conjugates. The norm is $sC$, so
these numbers are nonunits.

Composite floors on the classical schedule $3^k$ relate to the
Pisot problem in \cite[Problem~1.1]{SaitoIntervals}.
The negative schedule $2\cdot3^k-1$ also arises in algebraic-minimum
reductions for shifted prime-representing constants
\cite{Saito,N01}. Our theorems are stated directly for explicit
polynomials and do not assume that their roots are least constants.

The author's earlier manuscript \cite[theorem ``Nonreal cubic
conjugates'']{K031} already proves infinitely many composite classical
floors for every non-totally-real cubic Pisot number. Thus the complex
classical families below are explicit recurring-divisor refinements,
not newly resolved open cases. The totally real nonunit families are
restricted progress in the remaining direction; the complex
negative-clock statement concerns a different schedule.

\begin{mainthm}[An adaptive discriminant-prime certificate]
\label{R06:thm:adaptive}
Let $C>1$, let $a>2C$ be odd, and let $s\in\{1,-1\}$.
Suppose a prime $p>3$ satisfies
\[
 p\mid C+27,\qquad 4a\equiv45s\pmod p.
\]
Then $f_{a,C,s}$ is irreducible cubic Pisot. Put $d=\operatorname{ord}_p(2)$.
For every positive odd $n$,
\[
 p\mid T_n\quad\Longleftrightarrow\quad d\mid3n+1.
\]
For $k\ge1$ this yields the following exact trace-divisibility schedules.
\begin{enumerate}[label=\textup{(\roman*)}]
\item If $3\nmid d$, put $D=d/\gcd(d,2)$ and
$h_-=\operatorname{ord}_{D}(3)$. The negative clock has exactly the
hits $k\equiv-1\pmod{h_-}$; at every such index
$\floor{\beta^{2\cdot3^k-1}}$ is composite.
\item Suppose $3\nmid d$ and $-1$ belongs to the subgroup generated
by $3$ modulo $d$. Then $h_+=\operatorname{ord}_{d}(3)$ is even,
and the classical clock has exactly the hits
$k\equiv h_+/2-1\pmod{h_+}$. At each such index
$\floor{\beta^{3^k}}$ is composite.
\end{enumerate}
If $3\mid d$, neither clock has a hit from this certificate.
For the classical clock, failure of the subgroup condition likewise
means there are no hits from this prime.
\end{mainthm}

\begin{corollary}[The prime-$31$ families]\label{R06:thm:31main}
Let $C>1$ satisfy $C\equiv4\pmod{31}$, and let $a>2C$ be odd.
Suppose either
\[
 s=1,\ a\equiv19\pmod{31},\qquad\text{or}\qquad
 s=-1,\ a\equiv12\pmod{31}.
\]
Then \eqref{R06:eq:family} is an irreducible cubic Pisot polynomial.
Its dominant root $\beta$ satisfies
\begin{align*}
 \floor{\beta^{3^k}}&\text{ is composite whenever }k\ge1,
       \ k\equiv1\pmod4,\\
 \floor{\beta^{2\cdot3^k-1}}&\text{ is composite whenever }k\ge1,
       \ k\equiv3\pmod4.
\end{align*}
Both conclusions hold at every stated index, not merely eventually.
\end{corollary}

In particular, at $C=4$ the classical theorem covers the entire
positive progressions $a=19+62t$ for $s=1$ and $a=43+62t$ for $s=-1$,
$t\ge0$. If one also imposes the negative-shift coefficient condition
$a\equiv-s\pmod3$, the progressions become $143+186t$ and $43+186t$.
For general norms these additional conditions can be combined with
$C\equiv4\pmod{93}$; coprimality conditions, if desired for another
application, remain separate restrictions.

\begin{mainthm}[No fixed modulus excludes the real shadow locus]
\label{R06:thm:nogo}
Fix $C>1$ with $C\equiv1\pmod3$, and an integer $m\ge1$.
There are infinitely many odd integers $a\equiv1\pmod3$ with
$\gcd(a,C)=1$ such that $f_{a,C,-1}$ is irreducible and totally real
Pisot, with dominant root $\beta$, and
\begin{gather*}
 \Nm(\beta)=-C,\qquad \Tr(\beta^{-1})=1,
 \qquad C<\beta^{3/20},\\
 0<\beta^n-T_n<2^{-n}\beta^{-17n/40},\qquad
 \floor{\beta^n}=T_n\equiv1\pmod m
 \quad\text{for every positive odd }n.
\end{gather*}
The numbers $\beta$ depend on $m$. No assertion is made that their
traces are prime or that they are least prime-representing constants.
\end{mainthm}

Thus a finite set of fixed primes can exclude useful infinite
subfamilies, as in Corollary~\ref{R06:thm:31main}, but cannot exclude the
entire unbounded real coefficient and shadow locus of
Theorem~\ref{R06:thm:nogo}. Sections~\ref{R06:sec:geometry}--\ref{R06:sec:norm}
explain the nonunit discriminant structure, and
Section~\ref{R06:sec:phase} gives a separate limitation of a purely
phase-based approach.

\subsection{Root configuration and prime certificates}
\label{R06:sec:families}

\begin{lemma}\label{R06:lem:roots}
If $C\ge2$ and $a>2C$, both signs in \eqref{R06:eq:family} give an
irreducible cubic Pisot polynomial with $\beta\in(a,a+1)$.
For $s=-1$, write the other roots as $-w,v$, where $0<v<w<1$.
Then $\floor{\beta^n}=T_n$ for every positive odd $n$.
For $s=1$, the other roots are nonreal, of common modulus
$\sqrt{C/\beta}$, and $|T_n-\beta^n|<1$ for every $n\ge3$.
If $a$ is odd, $T_n$ is odd for every $n\ge1$.
\end{lemma}
\begin{proof}
For $s=-1$, the signs at $-1,0,1,a,a+1$ are respectively
\[
 2C-a-1<0,\quad C>0,\quad1-a<0,\quad C(1-a)<0,
 \quad(a+1)^2-Ca>0.
\]
Thus the three roots lie in $(-1,0),(0,1),(a,a+1)$.
They are all nonintegral, so the monic integral cubic is irreducible.
Vieta gives $v-w=a-\beta<0$, and therefore
\[
 -1<v^n-w^n<0\quad(n\text{ positive and odd}).
\]
The trace identity $T_n=\beta^n+v^n-w^n$ proves the floor assertion.

For $s=1$, the discriminant is
\begin{equation}\label{R06:eq:discriminant}
 \Delta=C H_{C,s}(a),\qquad
 H_{C,s}(a)=-4sa^3+C(a^2-18sa-27)+4C^2.
\end{equation}
For $s=1$ and $a\ge C\ge2$ this is negative:
$-4a^3+Ca^2\le-3a^3$ and $-3a^3+4C^2<0$.
Also $f(a)=-C(a+1)<0$ and
$f(a+1)=(a+1)^2-C(a+2)>0$.
There is one real root, in $(a,a+1)$, hence no rational root.
The other-root product is $C/\beta<1$, so the pair has modulus
$\sqrt{C/\beta}$. Since $a>2C$,
\[
 |T_n-\beta^n|\le2(C/\beta)^{n/2}<2(1/2)^{3/2}<1
 \qquad(n\ge3).
\]
Finally, if $a$ is odd, the first three traces are odd and the sum
of the recurrence coefficients in \eqref{R06:eq:traces} is odd modulo
two. Induction gives oddness of every trace.
\end{proof}

\begin{proof}[Proof of Theorem~\ref{R06:thm:adaptive}]
Lemma~\ref{R06:lem:roots} gives the root configuration. In $\F_p[X]$,
\begin{equation}\label{R06:eq:adaptivefactor}
 f_{a,C,s}(X)=(X-6s)^2(X+3s/4).
\end{equation}
The roots are nonzero and distinct apart from the displayed double
root, since $p>3$. Their ratio is $-1/8$, so at odd exponents
\[
 T_n\equiv(6s)^n(2-8^{-n})\pmod p.
\]
This vanishes exactly when $2^{3n+1}=1$, or $d\mid3n+1$.
For $n=2\cdot3^k-1$ this is
$D\mid3^{k+1}-1$. It is impossible if $3\mid d$; otherwise it is
exactly $k+1\equiv0\pmod{h_-}$.
For $n=3^k$ it is $3^{k+1}\equiv-1\pmod d$.
Here $d>2$, because orders one or two would force $p\le3$.
Within the cyclic group generated by $3$, its element of order two,
when present, is $3^{h_+/2}$. Thus the stated subgroup condition
and index class are necessary and sufficient. Merely knowing that
$h_+$ is even would not suffice in a noncyclic ambient unit group.

At every certified $k\ge1$, $n\ge3$ and the floor is $T_n$ or
$T_n-1$, by Lemma~\ref{R06:lem:roots}. Since $a$ is odd, $T_n$ is odd.
Also
\[
 T_n>\beta^n-1>8C^3-1>C+27\ge p\qquad(C\ge2).
\]
Thus $T_n$ is a composite multiple of $p$, and $T_n-1$ is an even
integer greater than two. This proves the floor assertions.
\end{proof}

The test requires factoring the comparatively small integer $C+27$
and checking orders and a coefficient residue. It is not universal:
$C=37$ gives $C+27=64$, with no eligible odd prime, while $C=22$
gives $49$ and $\operatorname{ord}_7(2)=3$. Even when an eligible
prime exists, its coefficient congruence need not hold.
Failure of this test is not evidence that any floor is prime.
The certificate uses repeated reduction of the polynomial, not a claim
that $p$ ramifies in the cubic number field: an index factor can absorb
a polynomial-discriminant factor. No field-ramification assertion is
needed in the argument.

\begin{proof}[Proof of Corollary~\ref{R06:thm:31main}]
For $p=31$ one has $d=\operatorname{ord}_{31}(2)=5$ and
$\operatorname{ord}_5(3)=4$, with $3^2\equiv-1\pmod5$.
Also $45s/4\equiv19s\pmod{31}$, giving residues $19$ and $12$.
Theorem~\ref{R06:thm:adaptive} therefore gives precisely the stated
index classes. The direct factorizations are
\begin{equation}\label{R06:eq:31factors}
 \begin{aligned}
 f_{19,4,1}(X)&=(X-6)^2(X-7),\\
 f_{12,4,-1}(X)&=(X+6)^2(X+7)
 \end{aligned}\qquad\text{over }\F_{31}.
\end{equation}
Equivalently, $31\mid T_n$ at odd exponents exactly when
$n\equiv3\pmod5$.
\end{proof}

This direct factorization proof is important: when $C=4$, the prime
$31$ divides $C+27$, so the generic local criterion in
Proposition~\ref{R06:prop:local} below does not automatically cover it.

\begin{proposition}[An additional prime-$53$ family]\label{R06:prop:53}
For $C=4$, every dominant root of
\[
 f_{239+318t,4,1}\quad\text{or}\quad f_{79+318t,4,-1},\qquad t\ge0,
\]
has a composite floor at exponent $2\cdot3^k-1$ whenever
$k\ge1$ and $k\equiv1\pmod3$.
\end{proposition}
\begin{proof}
Modulo $53$, the complex coefficient is $a=27$ and
\[
 f=(X-19)^2(X-42),\qquad 42/19=-1/21.
\]
The element $21$ has order $52$ and $21^5=27=1/2$.
The exponent $2\cdot3^k-1$ is respectively $1,5,17$ modulo $52$
when $k\equiv0,1,2\pmod3$, because $3^3\equiv27\pmod{52}$ and
$2\cdot3^3\equiv2\pmod{52}$.
Thus the trace vanishes modulo $53$ exactly when $k\equiv1\pmod3$.
Reflection gives the real coefficient $a=26$ modulo $53$.
The displayed progressions impose these residues, oddness, and
$a\equiv-s\pmod3$. Lemma~\ref{R06:lem:roots} applies, and the same
odd-trace/even-predecessor argument gives composite floors.
All relevant traces already exceed $53$ at the first exponent $5$.
\end{proof}

For each fixed $C$, arbitrarily large members of the prime-$31$
families also satisfy $C\le\beta^{3/20}$ and small-conjugate modulus
below $\beta^{-17/40}$. In the complex case this follows from
$\rho=\sqrt{C/\beta}$; in the real case it follows from the estimate
in Section~\ref{R06:sec:nogo}. These additional inequalities are not
needed for the composite-floor proofs.

\subsubsection{Extension to every admissible even multiplier}
% Optional integration fragment. Requires theorem/corollary/proof environments.
% Label prefix audit06 is unique. No inherited macros are required.
\begin{corollary}[The prime-$31$ families on every even multiplier]
\label{R06:audit06:even-multipliers}
Let $C>1$ and $a>2C$ be integers, with $a$ odd and
$C\equiv4\pmod{31}$. Let $s\in\{1,-1\}$ and suppose
$a\equiv19s\pmod{31}$. Let $\beta$ be the dominant root of
\[
                 X^3-aX^2-CX-sC.
\]
For every positive even integer $r$ with $5\nmid r$, there is exactly
one residue class $k_r\pmod4$ such that
\[
 3^{k_r+1}\equiv 2r^{-1}\pmod5.
\]
For every $k\ge1$ with $k\equiv k_r\pmod4$, the integer
$\lfloor\beta^{r3^k-1}\rfloor$ is composite. In particular, each
such $\beta$ is excluded from the prime-representing set of every
one of these clocks. No assertion about the minimum for any clock
is needed.
\end{corollary}
\begin{proof}
The root-location argument for this polynomial gives a cubic Pisot
number $\beta\in(a,a+1)$. If $s=-1$, its small roots are $-w,v$
with $0<v<w<1$, so at every positive odd exponent $n$ its floor is
exactly its integral power trace $T_n$. If $s=1$, the small roots
are nonreal of modulus $\sqrt{C/\beta}<1/\sqrt2$, and hence
$|T_n-\beta^n|<1$ for $n\ge3$. In either case the floor at such an
odd exponent is $T_n$ or $T_n-1$. The Newton recurrence
\[
 T_0=3,\quad T_1=a,\quad T_2=a^2+2C,\qquad
 T_n=aT_{n-1}+CT_{n-2}+sCT_{n-3}
\]
shows that every $T_n$ is odd.

Modulo $31$, the polynomial factors as
\[
 (X-6s)^2(X+3s/4).
\]
Thus for odd $n$,
\[
 T_n\equiv(6s)^n(2-8^{-n})\pmod{31}.
\]
Since $2$ has order $5$ modulo $31$, this trace vanishes exactly
when $5\mid3n+1$. For $n=r3^k-1$ this is
$r3^{k+1}\equiv2\pmod5$. The number $3$ generates the nonzero
residues modulo $5$, giving the claimed unique class modulo $4$.
As $r$ is even, every prescribed $n$ is odd and at least $5$.
Furthermore
\[
 T_n>\beta^n-1>(2C)^3-1>C+27\ge31.
\]
At every indicated index, $T_n$ is therefore an odd composite
multiple of $31$, while $T_n-1$ is an even integer greater than two.
Either possible floor is composite.
\end{proof}

\begin{remark}
The condition $5\nmid r$ is exact for this particular witness:
if $5\mid r$, the prime $31$ has no zero trace on this orbit. This
failure says nothing about other possible divisors. The corollary
is a strengthened family exclusion, not a solution of $\xi_2$ or
of the general even-multiplier minimum problem.
\end{remark}

\subsection{Nonunit discriminants and the rational-square locus}
\label{R06:sec:geometry}

The factorization at $a=-s$ is
\begin{equation}\label{R06:eq:unitcongruences}
 \begin{gathered}
 f_{-s,C,s}(X)=(X+s)(X^2-C),\\
 T_n\equiv-s\pmod{a+s}\quad(n\text{ odd}),\qquad
 \Delta\equiv4C(C-1)^2\pmod{a+s}.
 \end{gathered}
\end{equation}
The trace congruence follows because the odd power sums of the two
roots of $X^2-C$ cancel. Equivalently, it is a polynomial identity
in the Newton sums, so is valid also modulo nonprime integers.

\begin{proposition}\label{R06:prop:irredH}
For every integer $C>1$, $H_{C,s}$ is irreducible over $\Q$, and
\begin{equation}\label{R06:eq:discH}
 \disc_a H_{C,s}=-16C^2(C+27)^3.
\end{equation}
Its root field is generated by a root $u$ of
$u^3+Cu+2sC=0$, and its quadratic resolvent is
$\Q(\sqrt{-(C+27)})$.
\end{proposition}
\begin{proof}
Eliminating $a$ from $f(u)=f'(u)=0$ gives
\begin{equation}\label{R06:eq:rootmap}
 u^3+Cu+2sC=0,\qquad a=\frac{3u^2-C}{2u}.
\end{equation}
With $t=-su$, the cubic becomes $t^3+Ct-2C$.
Its derivative is positive and its unique real root lies in $(1,2)$,
since its values at $1,2$ are $1-C<0$ and $8>0$.
A rational root of a monic integral cubic is integral, so this cubic
is irreducible. The inverse rational formula is
\begin{equation}\label{R06:eq:inversemap}
 u=-\frac{C(a+9s)}{2(a^2+3C)}.
\end{equation}
Substitution verifies both maps. If the denominator in
\eqref{R06:eq:inversemap} vanished at a root of $H$, then reducing $H$
under $a^2=-3C$ would give $C(C-27-6sa)=0$, forcing $C=-27$;
this is impossible here. Thus the two generated fields coincide
and have degree three, proving irreducibility of $H$.
Expansion of its cubic discriminant gives \eqref{R06:eq:discH}.
Removing its rational square factor gives the stated resolvent.
\end{proof}

For comparison, at $C=1$ the exceptional factorization is
\[
 H_{1,s}(a)=-s(a+s)(4a^2-5sa+23).
\]
This linear factor is responsible for the quadratic, Pell-type
reduction used in the unit case \cite{R03}.

\begin{proposition}[No integral cancellation under the norm bound]
\label{R06:prop:nocancel}
Suppose $a$ is odd, \eqref{R06:eq:family} is cubic Pisot, $C>1$, and
$C\le\beta^{3/20}$. Then neither $(a+s)/2$ nor $a+s$ divides $\Delta$.
\end{proposition}
\begin{proof}
The two small conjugates give $\beta<a+2$, so
\[
 a-1>C^{20/3}-3\ge C^6-3>8C(C-1)^2.
\]
The last inequality follows from $C^6\ge8C^3$ for $C\ge2$, and
$8C^3-8C(C-1)^2=8C(2C-1)>3$.
Therefore $0<4C(C-1)^2<(a+s)/2$.
The discriminant congruence in \eqref{R06:eq:unitcongruences} excludes
both divisibilities.
\end{proof}

This does not exclude rational squares whose denominators cancel
square factors of $a+s$. The next theorem treats that exact condition.

\begin{theorem}[The direct nonunit square relation]\label{R06:thm:genus}
Fix $C>1$ and $s\in\{1,-1\}$. For integer $a\ne-s$, the condition
\begin{equation}\label{R06:eq:square}
 \frac{\Delta}{-2s(a+s)}\in\Q^2
\end{equation}
is equivalent to the existence of an integer $Y$ satisfying
\begin{equation}\label{R06:eq:quartic}
 Y^2=-2sC(a+s)H_{C,s}(a).
\end{equation}
The smooth projective model of this quartic has genus one.
Consequently only finitely many integer coefficients $a\ne-s$
satisfy \eqref{R06:eq:square} for each fixed $C,s$.
\end{theorem}
\begin{proof}
Multiplying \eqref{R06:eq:square} by $[-2s(a+s)]^2$ gives the right side
of \eqref{R06:eq:quartic}. A rational square which is an integer is an
integer square; division by the same nonzero square proves the
converse. Proposition~\ref{R06:prop:irredH} shows that $H$ has three
simple roots, and $H(-s)=4(C-1)^2\ne0$.
Thus the quartic has four distinct roots. Its double cover of the
$a$-line has four simple branch points, so Riemann--Hurwitz gives
$2g-2=-4+4=0$. The affine curve is smooth and of genus one.
Siegel's integral-point theorem \cite{Siegel,R06-Lang} gives finitely
many integral pairs $(a,Y)$, hence finitely many such $a$.
\end{proof}

For the positive range $M=a+s>0$, put $B=-2s\Delta$ and suppose
$B>0$. If $g=\gcd(M,B)$, the discriminant congruence gives
$g\mid8C(C-1)^2$. The ratio in \eqref{R06:eq:square} is $B/(4M)$;
it is a rational square exactly when the coprime integers $M/g,B/g$
are both squares. Hence the condition is equivalent to
\begin{equation}\label{R06:eq:genus2}
 \begin{gathered}
 a=gu^2-s,\qquad v^2=F_{C,s,g}(u),\qquad\gcd(u,v)=1,\\
 F_{C,s,g}(u)=\frac{-2sC}{g}H_{C,s}(gu^2-s),\qquad u,v>0.
 \end{gathered}
\end{equation}
For every divisor $g$ of $8C(C-1)^2$, this $F$ has integral
coefficients: its constant term is $-8sC(C-1)^2/g$, and the
nonconstant terms of the substituted $H$ contain $g$.
Its degree is six, its leading coefficient is $8Cg^2$, and it has
six distinct roots because $H$ is squarefree and $H(-s)\ne0$.
Thus each displayed hyperelliptic curve has genus two.
The coprimality requirement ensures the stated converse, including
the identification of $g$ as the gcd.

At $C=1$, the repeated factor $a+s$ in \eqref{R06:eq:quartic} leaves a
quadratic after cancellation. For $C>1$ this degeneration disappears.
Theorem~\ref{R06:thm:genus} does not provide an effective bound, a bound
uniform in $C$, or a classification of all discriminant-character
survivors. It excludes only an infinite fixed-$C$ family obtained
by imposing this particular direct rational-square relation.

\subsection{A cubic-field norm for trace resultants}
\label{R06:sec:norm}

Set
\begin{equation}\label{R06:eq:h}
 w=su+2,\qquad h_C(W)=(W-2)^3+CW
 =W^3-6W^2+(12+C)W-8.
\end{equation}
Then $\Nm(w)=8$. At a root $a$ of $H_{C,s}$ given by
\eqref{R06:eq:rootmap}, the polynomial has roots $u,u,v$, where
$v=sC/u^2$ and $v/u=-1/w$.
For odd $n$ it follows that
\begin{equation}\label{R06:eq:ramtrace}
 T_n(a)=u^n(2-w^{-n}).
\end{equation}

\begin{proposition}\label{R06:prop:resultant}
Use the ordered resultant convention for which $\Res_a(a,H)=H(0)$.
For every positive odd $n$,
\begin{equation}\label{R06:eq:resultant}
 \Res_a(T_n,H_{C,s})
 =C^n\Nm_{\Q(w)/\Q}(1-2w^n)
 =C^n\bigl(1-2L_n+4M_n-8^{n+1}\bigr),
\end{equation}
where $L_n=\sum_iw_i^n$ and $M_n=\sum_{i<j}(w_iw_j)^n$ for the
three roots of $h_C$. In particular this resultant is independent
of $s$ at odd exponents.
\end{proposition}
\begin{proof}
As a polynomial in $a$, $T_n$ is monic of degree $n$.
The leading coefficient of $H$ is $-4s$, and
$\prod_i u_i=-2sC$. Since $n$ is odd, switching the resultant order
contributes $(-1)^{3n}=-1$. Thus \eqref{R06:eq:ramtrace} gives
\[
 \Res(T_n,H)=-(-4s)^n(-2sC)^n
       \prod_i(2-w_i^{-n})
 =-C^n\prod_i(2w_i^n-1).
\]
There are three factors, so this is $C^n\prod_i(1-2w_i^n)$.
Expanding the product and using $\prod_iw_i=8$ gives the last equality.
\end{proof}

The integers in \eqref{R06:eq:resultant} can be computed without algebraic
approximations, using
\[
 \begin{aligned}
 L_0&=3,\quad L_1=6,\quad L_2=12-2C,\\
 L_n&=6L_{n-1}-(12+C)L_{n-2}+8L_{n-3},\\
 M_0&=3,\quad M_1=12+C,\quad M_2=(12+C)^2-96,\\
 M_n&=(12+C)M_{n-1}-48M_{n-2}+64M_{n-3}.
 \end{aligned}
\]
Writing $E_C(n)=\Nm(1-2w^n)$, examples are
\[
 E_C(1)=4C-27,
\]
\[
 E_C(5)=4C^5+240C^4+4800C^3+35780C^2+77600C-250047.
\]
In particular $E_4(3)=31\cdot127$ and $E_4(5)=53\cdot18973$.

\begin{proposition}[A local scalar criterion]\label{R06:prop:local}
Let $p$ be odd with $p\nmid C(C+27)$ and $H_{C,s}(a)\equiv0\pmod p$.
Define $u$ by \eqref{R06:eq:inversemap} modulo $p$ and $w=su+2$.
These are defined and nonzero, and for every positive odd $n$,
\[
 p\mid T_n(a)\quad\Longleftrightarrow\quad w^n=1/2\text{ in }\F_p.
\]
\end{proposition}
\begin{proof}
The denominator calculation in Proposition~\ref{R06:prop:irredH}
works modulo $p$ under the excluded factors. The identities
\eqref{R06:eq:rootmap} and \eqref{R06:eq:h} imply $u,w\ne0$, and factor
$f$ as $(X-u)^2(X-v)$ over $\F_p$.
Equation~\eqref{R06:eq:ramtrace} gives the equivalence.
\end{proof}

For the negative clock the scalar equation is
$(w^2)^{3^k}=w/2$. If it has a hit at an index
$k\ge v_3(\operatorname{ord}_p(w))$, that hit recurs infinitely often:
write $\operatorname{ord}_p(w)=3^eU$ with $3\nmid U$, and choose
$h>0$ with $3^h\equiv1\pmod U$.
Then $\operatorname{ord}_p(w)$ divides $2\cdot3^k(3^{th}-1)$ for
all $t\ge0$. This is a recurrence implication, not an existence
theorem for a suitable prime or a suitable initial hit.
Similarly, divisibility of $E_C(n)$ is only a necessary condition
for a common trace/discriminant prime outside $C$; a coefficient
specialization is still required for a converse.

\subsection{The fixed-modulus obstruction and its quantitative form}
\label{R06:sec:nogo}

\begin{proof}[Proof of Theorem~\ref{R06:thm:nogo}]
Let $A_0=2^{27}C^7$ and $L=\lcm(6C,m)$.
Choose any $a\ge A_0$ with $a\equiv1\pmod L$.
Then $a$ is odd, $a\equiv1\pmod3$, $\gcd(a,C)=1$, and $a>2C$.
Lemma~\ref{R06:lem:roots} gives roots $-w,v,\beta$ with
$0<v<w<1$ and $\beta\in(a,a+1)$.
The negative-root equation gives
\[
 aw^2=C(1+w)-w^3<2C.
\]
Since $\beta<a+1<2a$, this yields $w<2\sqrt{C/\beta}$.
Moreover
\[
 a^3\ge A_0^3=2^{81}C^{21}>2^{80}C^{20}=(4\sqrt C)^{40},
\]
so $\beta^{3/40}>4\sqrt C$ and
\[
 w<\tfrac12\beta^{-17/40},\qquad C<\beta^{3/20}.
\]
For positive odd $n$,
\[
 0<\beta^n-T_n=w^n-v^n<w^n
 <2^{-n}\beta^{-17n/40}.
\]
In particular its floor is $T_n$; also $T_n>0$ and
$\beta^{-17n/40}<T_n^{-17/40}$, if the numerical shadow bound is
expressed in terms of the integer trace instead.

At $a=1$, the factorization $f=(X-1)(X^2-C)$ gives the polynomial
identity $T_n(1,C,-1)=1$ for every odd $n$.
Because \eqref{R06:eq:traces} consists of polynomials with integral
coefficients, $a\equiv1\pmod m$ implies $T_n\equiv1\pmod m$ for
all those exponents. This includes primes dividing $C$, the prime
two, and arbitrary prime powers; no division in a residue ring is
involved. The norm and inverse trace follow from the coefficients.
Finally there are infinitely many such $a$, and their distinct
intervals $(a,a+1)$ give distinct Pisot numbers.
\end{proof}

\begin{corollary}[Necessary product size for a finite-prime cover]
\label{R06:cor:product}
Fix $C>1$ with $C\equiv1\pmod3$, let $A_0=2^{27}C^7$, and let
$H>A_0$. Suppose a finite set $\mathcal P$ of primes excludes every
coefficient in the real locus of Theorem~\ref{R06:thm:nogo} with $a\le H$
by divisibility of some odd scheduled trace. Put
$Q=\prod_{p\in\mathcal P}p$. Then
\[
 6CQ>H-A_0,\qquad
 \sum_{p\in\mathcal P}\log p>
       \log\!\left(\frac{H-A_0}{6C}\right).
\]
For $C=4$, this reads $24Q>H-2^{41}$.
\end{corollary}
\begin{proof}
The first integer $a\ge A_0$ congruent to one modulo $6CQ$ lies
strictly below $A_0+6CQ$. If $6CQ\le H-A_0$, this coefficient lies
below $H$ and satisfies every conclusion of Theorem~\ref{R06:thm:nogo},
while all its odd traces are one modulo every $p\in\mathcal P$.
That contradicts the proposed cover. Taking logarithms proves the
second inequality.
\end{proof}

This is not a lower bound for arbitrary algorithms. It does not
exclude coefficient-dependent divisors, unbounded prime sets,
additional discriminant restrictions, or arguments using minimality.
Nor does it construct primes: the trace could have a prime divisor
outside any specified finite set.

\subsection{Irrational phase survivors near zero}
\label{R06:sec:phase}

For the complex branch, write its small conjugates as
$\rho e^{\pm i\theta}$. Vieta gives
\[
 2\rho\cos\theta=-\frac C\beta\left(1+\frac1\beta\right)<0,
\]
so $\theta=\pi/2+\delta$ for $0<\delta<\pi/2$.
For $n_k=2\cdot3^k-1\equiv1\pmod4$,
\[
 T_{n_k}-\beta^{n_k}=-2\rho^{n_k}\sin(n_k\delta).
\]
Once this shadow has magnitude below one, the floor equals $T_{n_k}$
if and only if $\sin(n_k\delta)\ge0$.
Writing $t=\delta/\pi\in(0,1/2)$ gives the exact phase condition
\begin{equation}\label{R06:eq:survivor}
 \{3^kt\}\in[t/2,(1+t)/2].
\end{equation}
Here $n_k\delta/(2\pi)=3^kt-t/2$, and the translated sine
half-circle does not wrap around the unit interval.

\begin{theorem}\label{R06:thm:phase}
For every integer $m\ge1$, there is a set of Hausdorff dimension
$\log2/(2\log3)$ inside
$[3^{-(m+1)},\tfrac12 3^{-m}]$ such that every member satisfies
\eqref{R06:eq:survivor} strictly for every $k\ge0$.
All but countably many members are irrational.
\end{theorem}
\begin{proof}
Use the ternary expansions
\[
 t=0.\underbrace{00\cdots0}_{m\ \mathrm{digits}}
       1b_1,1b_2,1b_3\cdots\quad\text{in base three},
 \qquad b_j\in\{0,1\}.
\]
All digits are zero or one, and every zero run has length at most
$m$. Every shifted tail therefore satisfies
\[
 3^{-(m+1)}\le\{3^kt\}\le\tfrac12.
\]
The initial prefix gives
$3^{-(m+1)}\le t\le\tfrac12 3^{-m}$, and hence
\[
 t/2\le\tfrac14 3^{-m}<3^{-(m+1)},\qquad
 (1+t)/2>\tfrac12.
\]
Thus all survivor inequalities are strict.

After the fixed prefix, the two allowed digit blocks are $10$ and
$11$. They give two separated similarity maps of contraction $1/9$.
At depth $j$ there are $2^j$ disjoint cylinder sets of diameter
proportional to $9^{-j}$. These covers give dimension at most
$\log2/\log9$. Assigning mass $2^{-j}$ to each cylinder gives the
reverse inequality by the mass distribution principle, since an
interval at scale $9^{-j}$ meets a bounded number of cylinders.
A fixed finite prefix does not alter the dimension.
Rational numbers have eventually periodic ternary expansions and
are countable, whereas this separated binary construction is
uncountable. Removing its rational members leaves the stated
irrational survivors.
\end{proof}

Taking $m$ arbitrarily large places these sets arbitrarily close
to zero. No claim is made that their angles arise from conjugates
of cubic algebraic integers. The theorem shows precisely that
smallness of the angle, irrationality, and the ternary recurrence
alone do not eliminate all phase survivors. A complete argument
must use additional arithmetic or minimality; equidistribution of
an algebraic-angle orbit must not be assumed without proof.

\subsection{Scope}

The explicit prime certificates prove composite floors for proper
infinite nonunit subfamilies on both schedules. The fixed-modulus
construction proves that a finite prime-divisor cover cannot settle
the whole unbounded real coefficient/shadow locus. The genus-one
finiteness theorem concerns one specified rational-square relation
for each fixed norm, not all possible local or character conditions.
The phase construction concerns real angles, not a realized family
of algebraic Pisot conjugates.

None of these results proves transcendence of the least constant
for $2\cdot3^k-1$, transcendence of the classical Mills constant,
or the full assertion that every cubic Pisot number has infinitely
many composite floors along $3^k$. Their roles are complementary:
explicit excluded families, an exact discriminant-norm calculation,
and limitations of several otherwise natural uniform methods.


\subsection{Portable exact verification}

The universal proofs are given above. The supplementary files in
\texttt{certificates} check symbolic identities, conventional
Sylvester-determinant signs, finite-field factorizations, exponent
cycles, and exact sample root/shadow inequalities.
\begin{center}\small
\begin{tabular}{@{}p{0.60\linewidth}p{0.34\linewidth}@{}}
\toprule
File&Role\\\midrule
\texttt{adaptive\_ramification\_verify.py}&Adaptive order and clock certificates\\
\texttt{adaptive\_ramification\_verify.json}&Adaptive certificate results\\
\texttt{audit\_nonunit\_discriminant\_verify.py}&Independent norm and local checks\\
\texttt{audit\_nonunit\_discriminant\_verify.json}&Independent terminal results\\
\texttt{nonunit\_discriminant\_verify.py}&Discriminant, norm and prime certificates\\
\texttt{nonunit\_discriminant\_verify.json}&Exact terminal results\\
\texttt{audit\_fixed\_nonunit\_obstruction.py}&Real-locus identities and modulus checks\\
\texttt{audit\_fixed\_nonunit\_obstruction.json}&Real-locus terminal results\\
\texttt{phase\_survivors\_verify.py}&Exact rational survivor samples\\
\texttt{phase\_survivors\_verify.json}&Phase sample results\\
\bottomrule
\end{tabular}
\end{center}
The algebraic checking programs use SymPy for exact
polynomial arithmetic. The phase program uses the Python standard
library. Run without Python's optimization flag:
\begin{lstlisting}
cd certificates
python3 adaptive_ramification_verify.py
python3 nonunit_discriminant_verify.py
python3 audit_nonunit_discriminant_verify.py
python3 audit_fixed_nonunit_obstruction.py
python3 phase_survivors_verify.py
\end{lstlisting}
The phase file checks $1008$ periodic rational examples. Those checks
are illustrations only; irrationality, Hausdorff dimension and the
all-index conclusion are proved symbolically in
Theorem~\ref{R06:thm:phase}. Likewise finite samples do not establish the
universal modulus theorem or the genus statements.
For sign-sensitive resultants the supplied code constructs the
ordered Sylvester determinant directly, with
$\Res(X,X^3+2)=2$ as a convention check.



\endgroup

\section{Density on fixed norm lines}\label{src:R08}
Failure of a finite universal cover is compatible with a density-one theorem. Varying the auxiliary prime with the coefficient gives the following result on every fixed norm line, without deciding any prescribed exceptional coefficient.
\medskip
\begingroup




\subsection{The coefficient line and the main theorem}

For $C\ge1$, $s\in\{1,-1\}$ and an integer $a>2C$, set
\begin{equation}\label{R08:eq:family}
               f_{a,C,s}(X)=X^3-aX^2-CX-sC.
\end{equation}
Its dominant real root will be denoted by $\beta_{a,C,s}$.
Lemma~\ref{R08:lem:roots} shows that it is cubic Pisot: an algebraic integer
larger than $1$ whose other conjugates have modulus below $1$.
Its norm is $sC$, so $C>1$ gives nonunits, while $C=1$ is included too.

Let $M$ be a fixed positive even integer and let $\mathcal R$ be a
nonempty collection of odd residue classes modulo $M$. Write
\[
 \mathcal A(H)=\{a\in\N:2C<a\le H,\ a\bmod M\in\mathcal R\}.
\]
All density statements below concern coefficients in this explicitly
specified one-dimensional set.

\begin{theorem}[Simultaneous fixed-line density one]\label{R08:thm:main}
Fix $C\ge1$, $M$ and $\mathcal R$ as above. For relative natural density
one of the coefficients $a\in\mathcal A(H)$, there is an integer
$L=L(a)\ge1$ such that, for both $s=1$ and $s=-1$, all four integers
\[
 \floor{\beta_{a,C,s}^{3^k}},\qquad
 \floor{\beta_{a,C,s}^{2\cdot3^k-1}}
\]
are composite whenever $k\ge1$ and $L\mid k$.
More precisely, if $\mathcal G(H)$ denotes the coefficients with this
property, then
\[
                         \frac{\#\mathcal G(H)}{\#\mathcal A(H)}\longrightarrow1.
\]
A single suitable prime divisor of $a$ provides the common progression.
\end{theorem}

The progression has density $1/L$ in the index $k$, but $L$ may depend
on $a$. Neither a uniform lower bound on this index density nor uniformity
in $C$ is asserted.

\paragraph{Relation to the earlier corpus.}
The small-prime return mechanism for the classical clock already occurs
in \cite[lemma ``A uniform small-prime return condition'']{K031}.
The same manuscript's theorem ``Density one at every fixed nonzero norm''
varies both coefficients $(a,b)$ in $X^3-aX^2+bX-c$ and gives
$o(B^2)$ exceptional pairs. Its theorem ``Density one within Saito's
shrinking conjugate locus'' still varies $b$ over a growing interval
and gives $o(B^{63/40})$ exceptions in a family of that order.
Neither estimate controls the line $b=-C$, $c=sC$, containing only
$O(B)$ points: the entire line could lie in either exceptional set.
Theorem~\ref{R08:thm:main} supplies a separate line-specific argument and
a common return for two clocks.

The nonreal cubic Pisot case of infinitely many composite classical
floors was also established in \cite[theorem ``Nonreal cubic conjugates'']{K031}.
We retain that antecedent. Our simultaneous certificate is not presented
as a new resolution of that known case, and no literature-wide priority
claim is made.

\subsection{Root locations, traces, and actual floors}

Let $T_n$ be the sum of the $n$th powers of the three roots of
\eqref{R08:eq:family}. Newton's identities give integers
\begin{equation}\label{R08:eq:newton}
 \begin{gathered}
 T_0=3,\qquad T_1=a,\qquad T_2=a^2+2C,\\
 T_n=aT_{n-1}+CT_{n-2}+sCT_{n-3}\qquad(n\ge3).
 \end{gathered}
\end{equation}

\begin{lemma}\label{R08:lem:roots}
For $C\ge1$ and $a>2C$, both polynomials in~\eqref{R08:eq:family}
are irreducible cubic Pisot, with $\beta\in(a,a+1)$.
For $s=-1$, their other roots are $-\rho,\sigma$ with
$0<\sigma<\rho<1$, and
\[
                     \floor{\beta^n}=T_n\quad(n\ge1\text{ odd}).
\]
For $s=1$, the other roots are nonreal, of modulus $\sqrt{C/\beta}$,
and $|T_n-\beta^n|<1$ for every $n\ge3$.
If $a$ is odd, every $T_n$ is odd.
\end{lemma}
\begin{proof}
For $s=-1$, the signs at $-1,0,1,a,a+1$ are respectively
\[
 2C-a-1<0,\quad C>0,\quad1-a<0,\quad C(1-a)<0,
 \quad(a+1)^2-Ca>0.
\]
Thus the roots lie in $(-1,0),(0,1),(a,a+1)$.
None is an integer, so the monic cubic has no rational root and is
irreducible. Vieta gives $\sigma-\rho=a-\beta<0$. For positive odd $n$,
\[
                  0<\beta^n-T_n=\rho^n-\sigma^n<1,
\]
proving the floor identity.

In general, direct expansion gives
\begin{equation}\label{R08:eq:discriminant}
 \Delta=C H_{C,s}(a),\qquad
 H_{C,s}(a)=-4sa^3+C(a^2-18sa-27)+4C^2.
\end{equation}
For $s=1$ this is negative when $a>2C$:
$-4a^3+Ca^2+4C^2<-3a^3+a^2<0$, and the remaining terms are negative.
Also $f(a)=-C(a+1)<0$ and
$f(a+1)=(a+1)^2-C(a+2)>0$.
Hence the sole real root lies in $(a,a+1)$, giving irreducibility.
The other roots have product $C/\beta$ and common modulus
$\sqrt{C/\beta}<1/\sqrt2$. Consequently
\[
             |T_n-\beta^n|\le2(C/\beta)^{n/2}<1\qquad(n\ge3).
\]
Finally, if $a$ is odd, the first three traces are odd and the sum of
the recurrence coefficients is $1+C+C=1\pmod2$.
Induction proves the parity assertion.
\end{proof}

\begin{lemma}[The coefficient-divisor floor bridge]\label{R08:lem:floor}
Let $a>2C$ be odd and let a prime $p\mid a$. At any odd $n\ge3$
for which $p\mid T_n$, the integer $\floor{\beta_{a,C,s}^n}$ is composite.
\end{lemma}
\begin{proof}
We have $p\le a$ and
\[
                    T_n>\beta^n-1>a^3-1>a\ge p.
\]
For $s=-1$, the floor is this proper multiple of $p$.
For $s=1$, it is either the odd integer $T_n$ or the even integer
$T_n-1>2$, by Lemma~\ref{R08:lem:roots}. Both are composite.
No assumption on the sign of the complex conjugate sum is needed.
\end{proof}

\subsection{A common modular return for both clocks}

Put
\[
 h_{C,s}(X)=X^3-CX-sC,
\]
which is the reduction of $f_{a,C,s}$ whenever $p\mid a$.

\begin{lemma}\label{R08:lem:return}
Suppose $p\nmid C$ and there is an integer $N\ge1$, coprime to $3$,
such that
\[
                    X^N=1\quad\text{in }\F_p[X]/(h_{C,s}).
\]
If $p\mid a$ and $L=\ord_N(3)$, with $L=1$ when $N=1$, then
\[
                     T_{3^k}\equiv T_{2\cdot3^k-1}\equiv0\pmod p
                             \qquad(L\mid k).
\]
A common $N$ for the two signs gives a common progression for all four floors.
\end{lemma}
\begin{proof}
At $L\mid k$, both exponents are $1\pmod N$, so every root power
returns to its first power. Their sum is $T_1=a\equiv0\pmod p$.
Equivalently, apply the trace functional in the quotient ring to
$X^{3^k}=X^{2\cdot3^k-1}=X$.
\end{proof}

For an elementary source of these identities, define
\begin{equation}\label{R08:eq:goodprime}
 D_C=4C-27,\qquad
 \mathcal P_C=\{p>3:p\equiv2\pmod3,\ p\nmid C D_C,
                                      \leg{D_C}{p}=1\}.
\end{equation}
Here and below $p$ denotes a prime.

\begin{proposition}\label{R08:prop:elementaryreturn}
For every $p\in\mathcal P_C$ and both signs,
\[
                         X^{p^3-1}=1\quad\text{in }\F_p[X]/(h_{C,s}).
\]
In particular $N=p^3-1$ is a common prime-to-$3$ order bound.
\end{proposition}
\begin{proof}
The discriminant of $h_{C,s}$ is $C^2D_C$, a nonzero square modulo $p$.
Its Frobenius permutation is therefore even. To see this directly,
Frobenius acts on a Vandermonde square root of the discriminant by the
sign of its permutation; that square root lies in $\F_p$ precisely
when the discriminant is a square. A separable cubic with even
Frobenius has splitting type $1+1+1$ or $3$.
All roots are nonzero and lie in $\F_{p^3}$, so their orders divide
$p^3-1$. Separability gives the quotient-ring identity.
Since $p\equiv-1\pmod3$, $3\nmid p^3-1$.
The argument is independent of $s$; alternatively the two root triples
are negatives of one another.
\end{proof}

\subsection{Divergent good-prime sets, including the exceptional norm}

\subsubsection{Explicit prime progressions when \texorpdfstring{$C\ne6$}{C is not 6}}

\begin{proposition}\label{R08:prop:primes}
For every positive integer $C\ne6$,
\[
                             \sum_{p\in\mathcal P_C}\frac1p=\infty.
\]
\end{proposition}
\begin{proof}
The odd integer $D=D_C$ is always $1\pmod4$ and never zero.
Quadratic reciprocity, including its sign convention for negative $D$,
gives $\leg{D}{p}=\leg{p}{|D|}$; on the right is the Jacobi symbol,
with the denominator-one value defined as $1$.
The following reduced residue classes supply usable primes, apart
from the finitely many divisors of $2CD$:
\begin{enumerate}[label=\textup{(\roman*)}]
\item If $C\ge7$, then $D>0$ and take $p\equiv-1\pmod{3D}$.
Here $p\equiv2\pmod3$ and $\leg{p}{D}=\leg{-1}{D}=1$.
\item If $C\in\{1,2,4,5\}$, impose $p\equiv1\pmod{|D|}$
and $p\equiv2\pmod3$. The moduli are coprime, and CRT gives a
reduced class with Jacobi symbol $1$.
\item If $C=3$, take $p\equiv2\pmod{15}$. Then $D=-15$ and
$\leg{p}{15}=\leg{2}{3}\leg{2}{5}=1$.
\end{enumerate}
The reciprocal sum of primes in a reduced arithmetic progression
diverges, by the harmonic form of Dirichlet's theorem
\cite[Theorem~4.11 and its proof]{Kedlaya}; see also
\cite[Equation~27.11.9]{R08-DLMF}. Removing finitely many primes does not
change divergence.
\end{proof}

At $C=6$, $D_C=-3$ and every $p\equiv2\pmod3$ has
$\leg{D_C}{p}=-1$. Thus this particular prime selection is empty.
The next argument supplies orders directly rather than inferring
them from a splitting condition.

\subsubsection{The case \texorpdfstring{$C=6$}{C=6}: prescribed reduction orders}

We use the following precise specialization of a published theorem.
For a point $R$ on an algebraic group, write $G_R$ for the Zariski
closure of the cyclic subgroup it generates.

\begin{theorem}[Perucca, Theorem~3.1 \cite{R08-Perucca}]\label{R08:thm:perucca}
Let $G$ be a torus over a number field $K$, let $F/K$ be finite,
and let $R\in G(F)$ with $G_R$ connected.
For any nonzero integer $m$, there is a positive-Dirichlet-density set
of primes $\mathfrak p$ of $K$ such that some prime $\mathfrak q$ of
$F$ above $\mathfrak p$ gives a reduction of $R$ whose order is coprime to $m$.
\end{theorem}

\begin{proposition}\label{R08:prop:C6}
There is a set $\mathcal P_6$ of rational primes with divergent
reciprocal sum such that, for each $p\in\mathcal P_6$, a common integer
$N_p$ coprime to $3$ satisfies
\[
       X^{N_p}=1\quad\text{in }\F_p[X]/(X^3-6X-6s)
                              \quad(s=1,-1).
\]
\end{proposition}
\begin{proof}
Let $\alpha_1,\alpha_2,\alpha_3$ be the roots of $X^3-6X-6$,
and let $F$ be its splitting field. The polynomial is Eisenstein at $3$
and has discriminant $-108$. Its root absolute values are $b,\rho,\rho$,
where $b>0$ is the real root and $\rho^2=6/b$.
Since $b^3=6b+6>6$, we have $b>\rho$.

We verify multiplicative independence, which is the essential hypothesis.
Suppose $\prod_i\alpha_i^{e_i}=1$ with $e_i\in\Z$, and put
$S=e_1+e_2+e_3$. Transitivity of the Galois action lets us send each
coordinate in turn to the real root. Taking absolute values gives
\[
                 e_i\log b+(S-e_i)\log\rho=0\qquad(i=1,2,3).
\]
Subtracting shows $e_1=e_2=e_3=e$, because $b\ne\rho$.
The product of the root absolute values is $6$, so $6^e=1$ and $e=0$.
Thus there is no nontrivial multiplicative relation. Every proper
algebraic subgroup of the split torus is annihilated by a nontrivial
integer character, possibly after multiplying the character to remove
finite components. Hence the point $R=(\alpha_1,\alpha_2,\alpha_3)$
has $G_R=\mathbb G_m^3$, which is connected.

Apply Theorem~\ref{R08:thm:perucca} with
\[
                    K=\Q,\qquad G=\mathbb G_m^3,\qquad m=3.
\]
It supplies rational primes directly, although $R$ is only $F$-rational.
Discard finitely many bad reductions, including $2$ and $3$.
At a prime $\mathfrak q$ above each remaining $p$, the reduced roots
are distinct nonzero elements of the common finite field
$\kappa(\mathfrak q)$. The order of their tuple is the least common
multiple of their orders, and is coprime to $3$.
Consequently their separable polynomial divides $X^N-1$ over that
field, and hence over $\F_p$.
Negating the roots changes orders only by a possible factor $2$;
replacing $N$ by $\lcm(N,2)$ gives a common bound for both signs.

Positive Dirichlet density implies reciprocal divergence: otherwise
$\sum p^{-z}$ would remain bounded as real $z$ decreases to $1$,
contradicting the definition of positive Dirichlet density.
\end{proof}

Theorem~\ref{R08:thm:perucca} is the only additional input needed for $C=6$.
No residue-degree-one restriction or descent of prime-ideal density
is used. In particular, complete splitting alone would not justify
the claimed prime-to-$3$ orders.

\subsection{From good primes to density-one coefficients}

\begin{proof}[Proof of Theorem~\ref{R08:thm:main}]
For each fixed $C$, Propositions~\ref{R08:prop:primes} and~\ref{R08:prop:C6}
give a set of good primes with divergent reciprocal sum.
Delete the finitely many primes dividing $M$.
For a fixed finite set $\mathcal F$ of remaining good primes,
the Chinese remainder theorem gives
\begin{equation}\label{R08:eq:CRT}
 \lim_{H\to\infty}
 \frac{\#\{a\in\mathcal A(H):p\nmid a\text{ for all }p\in\mathcal F\}}
      {\#\mathcal A(H)}
                 =\prod_{p\in\mathcal F}\left(1-\frac1p\right).
\end{equation}
Indeed each allowed residue modulo $M$ has $\prod_{p\in\mathcal F}p$
extensions, of which exactly $\prod_{p\in\mathcal F}(p-1)$ avoid all
those divisibilities.

The coefficients with no good prime divisor avoid every such finite set.
First take their relative upper density as $H\to\infty$, and then
increase $\mathcal F$. Since $\sum1/p=\infty$, the products
in~\eqref{R08:eq:CRT} tend to zero. Thus relative density one have a good
prime divisor $p\mid a$.
For each such coefficient, Lemma~\ref{R08:lem:return} supplies a common
progression $L\mid k$ of trace divisors for both signs and clocks.
Lemma~\ref{R08:lem:floor} turns every positive index in that progression
into four composite floors, proving the theorem.
\end{proof}

\begin{corollary}[The earlier residual congruence line]\label{R08:cor:residual}
Fix $C>1$ with $C\equiv1\pmod3$, and fix $s\in\{1,-1\}$.
Theorem~\ref{R08:thm:main} holds relative to the denominator
\[
 \mathcal A_{C,s}(H)=\{2C<a\le H:a\text{ odd},\ a\equiv-s\pmod3,
                                                  \gcd(a,C)=1\}.
\]
Its cardinality is $\delta_C H+O_C(1)$, where
\[
          \delta_C=\frac{\varphi(6C)}{12C}
                =\frac16\prod_{\substack{q\mid C\\q>3}}\left(1-\frac1q\right)>0.
\]
\end{corollary}
\begin{proof}
These conditions form a nonempty union of odd classes modulo $6C$.
There are $\varphi(6C)/2$ such classes, since exactly half of the
units modulo $6C$ have the chosen nonzero residue modulo $3$.
Apply Theorem~\ref{R08:thm:main}. The factor at $2$ is already included
by oddness; it must not be counted again.
\end{proof}

Because $a<\beta_{a,C,s}<a+1$, ordering either sign's roots by
$\beta<B$ rather than $a\le B$ changes the cutoff by at most one
coefficient. The same density conclusion follows.
These statements do not permit conditioning on an arbitrary thin nonlinear
subset, such as a selected Pell family or a phase-survivor set.

\subsection{Every odd norm seed supplies a linear residue specialization}
\label{R08:sec:norm}

The main proof uses unramified coefficient divisors and is distinct from
the preceding discriminant-prime method \cite{R06}.
We next give a positive result about that other method, with its precise
existential quantifiers.

For $C>1$ define
\begin{equation}\label{R08:eq:norm}
 g_C(W)=(W-2)^3+CW,\qquad
 E_C(n)=\Nm_{\Q(w)/\Q}(1-2w^n),\qquad g_C(w)=0.
\end{equation}
The derivative $3(W-2)^2+C$ is positive, and the unique real root
lies in $(0,1)$ since $g_C(0)=-8$ and $g_C(1)=C-1>0$.
Thus the monic cubic has no integer root and is irreducible;
the norm in~\eqref{R08:eq:norm} is well defined.

\begin{theorem}[A genuine degree-one prime at every odd seed]\label{R08:thm:seed}
Fix $C>1$ with $C\equiv1\pmod3$, and let $j\ge1$ be odd.
Put $n=2\cdot3^j-1$. There are a prime $p\equiv2\pmod3$,
$p\nmid C$, and $w_0\in\F_p$ such that
\[
                       g_C(w_0)=0,\qquad 2w_0^n=1.
\]
For either sign, the formulas
\begin{equation}\label{R08:eq:reconstruct}
 u=s(w_0-2),\qquad a_0=\frac{3u^2-C}{2u}\quad\text{in }\F_p
\end{equation}
give a repeated root of $f_{a_0,C,s}$ and a zero trace at exponent $n$.
The trace zero recurs on an arithmetic progression of negative-clock
indices containing $j$.
\end{theorem}
\begin{proof}
Here $C\ge4$. The real root $r$ of $g_C$ lies in $(0,2/3)$, since
$g_C(2/3)=-64/27+2C/3>0$.
If $z,\overline z$ are the other roots, their product has modulus
$|z|^2=8/r>12$. As $n\ge5$,
\[
                 E_C(n)=(1-2r^n)|1-2z^n|^2>0.
\]
This integer is odd: the multiplication matrix of $1-2w^n$ is the
identity modulo $2$.

Modulo $3$,
\[
                     g_C(W)=(W-1)(W^2+W+2).
\]
For a root $\zeta$ of the irreducible quadratic in $\F_9$,
$\zeta^4=-1$ and $\Nm_{\F_9/\F_3}(1-\zeta)=1$.
Since $j$ is odd, $n\equiv5\pmod8$, and reduction of the norm gives
\[
 E_C(n)\equiv(1-2)\Nm_{\F_9/\F_3}(1+2\zeta)
                     =-1\pmod3.
\]
Consequently some prime $p\equiv2\pmod3$ divides $E_C(n)$ to an
odd valuation. It is not $2$, because $E_C(n)$ is odd, and hence $p>3$.

If $p\mid C$, then $g_C(W)=(W-2)^3\pmod p$, so norm zero would give
$2^{n+1}=1\pmod p$. The order of $2$ would divide both
$n+1=2\cdot3^j$ and $p-1$, forcing it to divide $2$.
This would imply $p\mid3$, a contradiction. Thus $p\nmid C$.

If $p\mid C+27$, then
\[
       g_C(W)=(W-8)(W+1)^2,\qquad
       E_C(n)\equiv9(1-2\cdot8^n)\pmod p.
\]
Therefore $w_0=8$ supplies the asserted root.
Otherwise
\[
                         \disc(g_C)=-4C^2(C+27)
\]
is nonzero modulo $p$. The order $\Z[w]$ is maximal at $p$, and the
norm valuation is the sum of local valuations weighted by residue degrees.
A degree-two factor contributes an even integer to that valuation.
A degree-three factor cannot contribute: a root $\xi\in\F_{p^3}$
with $\xi^n=1/2$ would have field norm $8$, yielding
\[
                         8^n=2^{-3},\qquad 2^{3(n+1)}=1\pmod p.
\]
Again the order of $2$ would divide $2$, since $3\nmid p-1$, which
is impossible. The total norm valuation is odd, so a degree-one
factor must contribute. Its root $w_0\in\F_p$ satisfies both equations.

The values $w_0=0,2$ are impossible, since $g_C(0)=-8$ and
$g_C(2)=2C$. Thus~\eqref{R08:eq:reconstruct} is defined, and substitution
gives
\[
                 u^3+Cu+2sC=0,\qquad f(u)=f'(u)=0.
\]
The remaining root $v$ satisfies $v/u=-1/w_0$.
At odd exponents $m$, the trace is consequently
\begin{equation}\label{R08:eq:repeatedtrace}
                         T_m=u^m(2-w_0^{-m})\pmod p.
\end{equation}
Let $d=\ord_p(w_0)$ and $L=\ord_d(3)$.
Because $d\mid p-1$, $3\nmid d$; at every $k\equiv j\pmod L$,
$2\cdot3^k-1\equiv n\pmod d$, so~\eqref{R08:eq:repeatedtrace} vanishes.
\end{proof}

For the real sign $s=-1$, CRT combines $a\equiv a_0\pmod p$ with
$a\equiv1\pmod{6C}$, because $p\nmid6C$.
This gives arbitrarily large admissible real Pisot coefficients.
Every odd floor equals its trace, so the recurring divisible floors
are composite once they exceed $p$. Unlike Theorem~\ref{R08:thm:main},
this construction need not have $p\mid a$, and an eventual size
qualification is appropriate.

\begin{table}[htbp]
\centering
\begin{tabular}{rrrr}
\toprule
$C$ & $p$ & $w_0\pmod p$ & $a_0\pmod p$ for $s=1$\\
\midrule
4 & 53 & 21 & 27\\
7 & 17 & 8 & 7\\
22 & 557 & 295 & 45\\
37 & 1021912841 & 691653895 & 197820572\\
\bottomrule
\end{tabular}
\caption{Exact examples at the seed $j=1$, $n=5$.
The real coefficient residue is the negative of the displayed one.}
\label{R08:tab:seeds}
\end{table}

The $C=37$ example has $C+27=64$, so it lies beyond the special method
using only odd primes dividing $C+27$.
The theorem does not place an arbitrary prescribed coefficient in the
constructed residue class. It also does not assert that primes obtained
from different seeds are distinct or have a divergent reciprocal sum.

\subsection{The entire discriminant-prime method is not complete}

By a discriminant-prime trace certificate we mean a prime divisor of
$\Delta$ that divides the trace at a relevant odd exponent.
The next example excludes all such primes, not just a selected finite list.

\begin{theorem}[An exact method-completeness counterexample]\label{R08:thm:counter}
Let $\beta$ be the largest root of
\[
                      f(X)=X^3-1000003X^2-4X+4.
\]
It is an irreducible totally real cubic Pisot number, and its floors at
all positive odd exponents equal its traces.
No prime dividing the polynomial discriminant divides any of these
odd traces. Nevertheless
\[
                 \floor{\beta^{2\cdot3^k-1}}\text{ is composite}
                          \quad\text{whenever }k\ge1,\quad k\equiv10\pmod{30}.
\]
Thus the example disproves completeness of the discriminant-prime-only
certificate method; it is not an example of prime floors on the full clock.
\end{theorem}
\begin{proof}
Here $C=4$, $a=1000003$ and $s=-1$.
Lemma~\ref{R08:lem:roots} gives the asserted root and floor properties.
The complete discriminant factorization is
\begin{equation}\label{R08:eq:counterfactor}
 \Delta=4H,\qquad
 H=4000040000204000316=2^2\cdot3\cdot3373\cdot98824982710841.
\end{equation}
Both factors above $3$ are prime; the accompanying exact recursive
Lucas certificates prove their primality and are described in
Section~\ref{R08:app:certs}.

At the two large primes the reductions are $(X-u)^2(X-v)$, with
\[
\begin{array}{c|rrr}
p&u&v&w=2-u\pmod p\\\hline
3373&2495&3351&880\\
98824982710841&70198375859310&57253214703065&28626606851533
\end{array}
\]
All displayed roots are nonzero, and direct multiplication verifies
the two factorizations. Formula~\eqref{R08:eq:repeatedtrace} applies.
The exact quadratic-character data are
\[
\begin{array}{c|rr}
p&\leg{w}{p}&\leg{2}{p}\\\hline
3373&1&-1\\
98824982710841&-1&1
\end{array}
\]
For odd $n$, the character of $w^n$ equals that of $w$, whereas
$1/2$ has the character of $2$. These mismatches rule out
$w^n=1/2$ and hence $p\mid T_n$ for every odd $n$.
Parity rules out $p=2$. Modulo $3$ the polynomial is
$(X-1)^2(X+1)$, with odd traces $2-1=1$.
Equation~\eqref{R08:eq:counterfactor} now excludes every discriminant prime.

A separate unramified prime gives actual composite floors.
In $\F_5[X]/(f)$ put $U_k=X^{2\cdot3^k-1}$.
The recurrence is $U_{k+1}=X^2U_k^3$; exact multiplication gives
\[
                U_0=X,\qquad U_{10}=4+2X+X^2,\qquad U_{30}=X.
\]
The trace of $U_{10}$ is $0\pmod5$, so $5\mid T_{2\cdot3^k-1}$
for every $k\equiv10\pmod{30}$. These real floors equal their traces
and exceed $5$. Also $5\nmid\Delta$, as required.
\end{proof}

This example satisfies the usual strict real conjugate bounds as well.
Writing its small roots as $-\rho,\sigma$, the exact checks
\[
 f(-1/400)=-\frac{143361201}{64000000}<0,\quad
 400^{40}>(1000004)^{17},\quad 4^{20}<(1000003)^3
\]
give
\[
                \rho<\beta^{-17/40},\qquad 4<\beta^{3/20}.
\]
Since $\rho\sigma=4/\beta$, they also imply
$\rho<\sigma^{17/23}$.
Thus the counterexample to that method does not disappear by imposing
these strict shrinking-conjugate inequalities.

\subsection{What the density theorem does and does not cover}

For every fixed $C$, Theorem~\ref{R08:thm:main} gives density one on each
fixed odd coefficient progression. It does not prove that the exceptional
set is empty, that a specified coefficient is certified, or that the
coefficient of a hypothetical algebraic least constant is typical.
Neither a uniform error term in $C$ nor a uniform positive lower bound
on the composite-index density is supplied.

Theorem~\ref{R08:thm:seed} concerns existence of a prime and a coefficient
residue at each seed, not coverage of every coefficient.
If $\mathcal B_{C,p}$ denotes the coefficient residues certified by a
recurrent discriminant prime $p$, a sufficient condition for the same
CRT strategy to yield density one for that restricted method would be
\[
                          \sum_p\frac{\#\mathcal B_{C,p}}p=\infty.
\]
No such divergence is proved here. Infinitely many norm seeds, or even
infinitely many distinct prime divisors, would not establish it.
Theorem~\ref{R08:thm:counter} already shows that completeness is false.

Finite unions of fixed residue classes are covered; arbitrary thin
nonlinear conditioning is not. In particular, one cannot transfer this
density statement to a selected Pell ray or a complex phase-survivor set
without a separate distribution theorem. These distinctions are essential
when using coefficient statistics in questions about individual Pisot
numbers or global prime-representing minima.


\subsection{Portable exact certificates}\label{R08:app:certs}

The accompanying \texttt{certificates/} directory contains the mathematical
verifiers and their outputs.
Run the following commands from the corresponding extracted source archive's directory:
\begin{lstlisting}[language=bash]
python3 certificates/fixed_norm_adaptive_verify.py
python3 certificates/C6_kummer_verify.py
python3 certificates/audit_fixed_line_density_verify.py
python3 certificates/audit_norm_seed_verify.py
python3 certificates/audit_discriminant_counterexample_verify.py
(cd certificates && sha256sum -c SHA256SUMS.txt)
\end{lstlisting}
The first and fourth programs require SymPy. The other three use only
the Python standard library. The \texttt{audit\_} prefix distinguishes
separate implementations of the arithmetic, not an additional mathematical
hypothesis. Each program checks exact assertions and prints a PASS result;
JSON records contain the finite data.

The principal verifier checks $984$ local returns, $932$ explicit
prime-progression selections, $10$ complete-period CRT counts, $99$ norm
sign/residue checks and $46$ degree-one norm specializations.
The independent line verifier checks $1188$ local cases, $11$ progressions,
$24$ CRT periods and $990$ root/parity cases.
The independent norm verifier computes the determinant of $I-2W^n$,
where
\[
                        W=\begin{pmatrix}0&0&8\\1&0&-C-12\\0&1&6\end{pmatrix},
\]
and checks $99$ norms and $43$ linear-root specializations.
These finite ranges supplement the proofs; none is used as evidence
for prime reciprocal divergence.

For $C=6$, the exact checker uses the fact that every factor of a cubic
has degree $1$, $2$ or $3$, so all roots lie in $\F_{p^6}$.
Removing the full $3$-part from $p^6-1$ and testing $X^N=1$ in the
quotient ring is an exact good-prime test. The sample $p=11$ gives
$X^{40}=1$ for both signs and $\ord_{40}(3)=4$.
The prime-density conclusion comes from Perucca's theorem, not the search.

\subsubsection{Primality and the counterexample}
The file \texttt{audit\_discriminant\_counterexample\_certificates.json}
contains complete recursive Lucas certificates for the two prime factors
in~\eqref{R08:eq:counterfactor}. The standalone verifier checks $15$ primes
in the recursion using only integer products, modular powers and gcds.
For every non-base entry $n$, it verifies a complete factorization of
$n-1$ into recursively certified primes and a base $b$ satisfying
\[
 b^{n-1}\equiv1\pmod n,\qquad
 \gcd\bigl(b^{(n-1)/q}-1,n\bigr)=1\quad(q\mid n-1).
\]
These conditions prove primality: for every prime divisor $\ell$ of $n$,
the order of $b\pmod\ell$ contains the full power of every $q\mid n-1$,
so $n-1\mid\ell-1$, forcing $\ell=n$.
Thus the factorization is not accepted on a probable-prime label.
The same program independently verifies both local factorizations,
Euler-criterion mismatches, the strict endpoint inequalities, and the
unramified prime-$5$ cycle.



\endgroup

\section{Prime coefficients and simultaneous composite floors}\label{src:R10}
The coefficient itself can be restricted to primes if the chosen local residue conditions are compatible with prime distribution. The Frobenius construction below verifies that compatibility rather than assuming it from a splitting condition.
\medskip
\begingroup




\subsection{Prime-relative density and the main theorem}

Fix $C\ge1$ with $C\equiv1\pmod3$.
For $s\in\{1,-1\}$ and an integer $a>2C$, put
\begin{equation}\label{R10:eq:family}
             f_{a,C,s}(X)=X^3-aX^2-CX-sC.
\end{equation}
Its dominant real root is denoted by $\beta_{a,C,s}$.
Section~\ref{R10:sec:floors} proves that these are irreducible cubic Pisot
polynomials, with $\beta_{a,C,s}\in(a,a+1)$.
The sign $s=-1$ gives the totally real branch; $s=1$ gives two nonreal
small conjugates.

Fix an even integer $M\ge2$ and a reduced residue class $b\pmod M$.
Write
\[
 \pi(H;M,b)=\#\{a\le H:a\text{ prime},\ a\equiv b\pmod M\}.
\]
The reduced class is odd because $M$ is even.

\begin{theorem}[Density one among prime coefficients]\label{R10:thm:main}
Relative density one of the prime coefficients $a\equiv b\pmod M$,
$a>2C$, have an integer $L=L(a)\ge1$ such that, for both signs,
\[
 \floor{\beta_{a,C,s}^{3^k}},\qquad
 \floor{\beta_{a,C,s}^{\,2\cdot3^k-1}}
\]
are composite at every index
\[
                         k=2+tL,\qquad t\ge1.
\]
Equivalently, the exceptional prime coefficients up to $H$ are
$o(\pi(H;M,b))$ as $H\to\infty$.
The same conclusion holds in any fixed nonempty finite union of reduced
odd residue classes.
\end{theorem}

This density is relative to \emph{primes}, not to all integers.
The preceding fixed-line theorem \cite{R08} proved integer-relative
density one using witness primes dividing $a$.
It cannot simply be restricted to prime coefficients, which form a
zero-density subset of the integers.
Here the witness primes supply \emph{nonzero} trace-zero residues.

The common seed index is $k=2$, with exponents $9$ and $17$.
We do not assert compositeness at the seed: a seed trace may equal its
witness prime. The theorem concerns all subsequent returns.
No uniform bound on $L$, uniformity in $C$ or $M$, or exceptional-set
rate is asserted.

The nonreal cubic Pisot result for infinitely many composite classical
floors was already established in \cite[theorem
\emph{Nonreal cubic conjugates}]{K031}.
Its inclusion here serves the simultaneous certificate, not a claim
to resolve that known case anew.

\subsection{Two trace polynomials and exact residue identities}

Write $T_{n,C,s}(A)$ for the Newton trace polynomial of
$X^3-AX^2-CX-sC$, and set $P_n(A)=T_{n,C,-1}(A)$.
Its recurrence is
\begin{equation}\label{R10:eq:recurrence}
 \begin{gathered}
 P_0=3,\qquad P_1=A,\qquad P_2=A^2+2C,\\
 P_n=AP_{n-1}+CP_{n-2}-CP_{n-3}\qquad(n\ge3).
 \end{gathered}
\end{equation}
Induction on the leading term shows that every $P_n$, $n\ge1$,
is monic of degree $n$ in $\Z[A]$.

\begin{lemma}\label{R10:lem:identities}
For every positive odd $n$,
\begin{equation}\label{R10:eq:boundary}
                             P_n(1)=1.
\end{equation}
Moreover
\begin{equation}\label{R10:eq:constants}
 P_9(0)=-3C^3(3C+1),\qquad
 P_{17}(0)=-17C^6(C^2+5C+1).
\end{equation}
If $C\equiv1\pmod3$ and $A\equiv5\pmod9$, then
\begin{equation}\label{R10:eq:mod9}
                             P_9(A)\equiv2,\qquad
                             P_{17}(A)\equiv4\pmod9.
\end{equation}
\end{lemma}
\begin{proof}
The factorization $f_{1,C,-1}(X)=(X-1)(X^2-C)$ proves
\eqref{R10:eq:boundary}, since the latter two roots have cancelling odd powers.
Substitution of $A=0$ into~\eqref{R10:eq:recurrence} gives~\eqref{R10:eq:constants}.

The recurrence modulo $9$ gives the following complete table.
Each row lists values at $A=0,1,\ldots,8$:
\[
\begin{array}{c|c|c}
 C\bmod9&P_9(A)\bmod9&P_{17}(A)\bmod9\\\hline
 1&(6,1,2,6,1,2,6,1,2)&(7,1,7,4,1,4,1,1,1)\\
 4&(6,1,2,6,1,2,6,1,2)&(1,1,7,7,1,4,4,1,1)\\
 7&(6,1,2,6,1,2,6,1,2)&(4,1,7,1,1,4,7,1,1)
\end{array}
\]
These are all possible $C\pmod9$ under the hypothesis.
The entry at $A=5$ proves~\eqref{R10:eq:mod9}.
Section~\ref{R10:app:certificates} supplies exact implementations of these
finite identities.
\end{proof}

Both constant terms in~\eqref{R10:eq:constants} are nonzero for $C>0$.
This will keep the trace-zero residues accessible to prime coefficients.

\subsection{An integer anchor and a realized Frobenius class}
\label{R10:sec:anchor}

We use Chebotarev in the following standard form: if $E/\Q$ is finite
Galois and $\mathcal D$ is a conjugacy class in $\Gal(E/\Q)$, the
unramified rational primes with that Frobenius class have natural
prime density $|\mathcal D|/[E:\Q]$.
Away from the discriminant of a monic squarefree integral polynomial,
its factor degrees modulo a prime are the cycle lengths of Frobenius
on its roots \cite[Fact~2.1 and Theorem~3.1]{Lenstra}.
The cyclotomic specialization also gives equidistribution of primes
among fixed reduced residue classes.

\begin{proposition}\label{R10:prop:witnessprimes}
For each $n\in\{9,17\}$ there is a set $\mathcal S_n$ of rational primes
of positive natural prime density such that every $p\in\mathcal S_n$
satisfies
\[
             p\nmid6C,\qquad p\bmod9\in\{2,4,5,7\},
\]
and $P_n$ has a nonzero root modulo $p$.
In particular, $\sum_{p\in\mathcal S_n}1/p=\infty$.
\end{proposition}
\begin{proof}
Fix $n$ and write $P=P_n$.
Factor it over $\Q$ as
\[
                    P=\prod_i G_i^{e_i},\qquad F=\prod_i G_i,
\]
where the distinct $G_i$ are monic irreducible integer polynomials,
by Gauss's lemma. Thus $F\in\Z[A]$ is monic and squarefree, with
$\disc(F)\ne0$ and $F(0)\ne0$.
For \emph{every} prime $p$ and residue $r$,
\begin{equation}\label{R10:eq:root-equivalence}
                         P(r)=0\pmod p\quad\Longleftrightarrow\quad
                         F(r)=0\pmod p.
\end{equation}
This follows from the product of integral factors, not an assumption
that taking squarefree parts commutes with reduction.

Let $\mathcal B$ be the finite set of prime divisors of
$6C\,\disc(F)\,F(0)$.
Choose an arbitrarily large positive \emph{integer} $A_0$ satisfying
\begin{equation}\label{R10:eq:anchor}
 A_0\equiv5\pmod9,\qquad
 A_0\equiv1\pmod q\quad(q\in\mathcal B,\ q\ne3).
\end{equation}
CRT makes these compatible and gives an odd $A_0$.
Since $P$ is monic, we may also require $P(A_0)>0$.
For every $q\in\mathcal B\setminus\{3\}$,
equation~\eqref{R10:eq:boundary} gives $P(A_0)\equiv1\pmod q$;
equation~\eqref{R10:eq:mod9} excludes $3$.
Thus $P(A_0)$ is coprime to every prime in $\mathcal B$.

The positive integer $P(A_0)$ is $2$ or $4\pmod9$.
If all its prime divisors were $1$ or $-1\pmod9$, their product,
with multiplicities, would also be $1$ or $-1\pmod9$.
Hence some prime divisor $p_0$ satisfies
\begin{equation}\label{R10:eq:p0}
                            p_0\bmod9\in\{2,4,5,7\}.
\end{equation}
It lies outside $\mathcal B$.
By~\eqref{R10:eq:root-equivalence}, $A_0\pmod{p_0}$ is a root of $F$.
It is nonzero since $p_0\nmid F(0)$, and simple since
$p_0\nmid\disc(F)$. The corresponding root of $P$ need not be simple.

Let $K$ be the splitting field of $F$, and put $E=K\Q(\zeta_9)$.
The prime $p_0$ is unramified in both fields and hence in $E$.
Take its \emph{actual} Frobenius conjugacy class
$\mathcal D\subset\Gal(E/\Q)$.
Its restriction to $K$ fixes a root of $F$, and its cyclotomic
restriction sends $\zeta_9$ to $\zeta_9^{p_0}$.
The fixed-root property is invariant under conjugation; the cyclotomic
restriction is constant on the class because the cyclotomic group is abelian.

Chebotarev supplies a prime set of density $|\mathcal D|/[E:\Q]>0$
with both properties. Delete the finitely many primes in $\mathcal B$.
Each remaining prime has a nonzero root of $F$, hence of $P$, and
satisfies~\eqref{R10:eq:p0}. This is $\mathcal S_n$.
Positive prime density implies reciprocal divergence, for example by
partial summation and the prime number theorem.
\end{proof}

The anchor is not a final candidate coefficient: it need not satisfy
$a\equiv b\pmod M$ or any other residual condition.
It realizes a compatible Frobenius class.
No disjointness between $K$ and $\Q(\zeta_9)$ is asserted.
One could choose the anchor prime by Dirichlet's theorem, since
\eqref{R10:eq:anchor} is reduced, but an integer anchor already suffices.

\subsection{A universal exponent for cubic quotient rings}
\label{R10:sec:orders}

The cubic at a witness prime may have repeated factors.
The next bound handles this without another separability requirement.

\begin{lemma}\label{R10:lem:unit-exponent}
Let $p\ge5$, $p\nmid C$, $a\in\F_p$, and $s\in\{1,-1\}$.
In
\[
                 R=\F_p[X]/(X^3-aX^2-CX-sC),
\]
the class $x$ of $X$ satisfies
\begin{equation}\label{R10:eq:Ep}
              x^{E_p}=1,\qquad E_p=p(p^2-1)(p^2+p+1).
\end{equation}
For $p\bmod9\in\{2,4,5,7\}$, write $E_p=3^e m_p$, $3\nmid m_p$.
Then
\begin{equation}\label{R10:eq:v3}
 e=\begin{cases}1,&p\equiv2,5\pmod9,\\2,&p\equiv4,7\pmod9.\end{cases}
\end{equation}
With $L_p=\ord_{m_p}(3)$, taking $L_p=1$ if $m_p=1$, every
$k\ge2$ with $k\equiv2\pmod{L_p}$ satisfies
\begin{equation}\label{R10:eq:returnpowers}
                  3^k\equiv9,\qquad 2\cdot3^k-1\equiv17\pmod{E_p}.
\end{equation}
\end{lemma}
\begin{proof}
The constant term is nonzero, so $x$ is a unit.
Modulo the nilradical, each field factor has degree $d\in\{1,2,3\}$,
and $x$ has order dividing a common multiple of $p^d-1$.
The integer $B_p=(p^2-1)(p^2+p+1)$ is such a multiple.
Therefore $x^{B_p}=1+z$ with $z$ nilpotent.
Every irreducible-factor multiplicity is at most $3$, so $z^3=0$.
In characteristic $p\ge5$, $(1+z)^p=1+z^p=1$,
proving~\eqref{R10:eq:Ep} and explaining its extra factor $p$.

The four residue classes modulo $9$ give~\eqref{R10:eq:v3} by direct
valuation of the two factors of $B_p$.
Modulo $m_p$ the order definition gives $3^k\equiv3^2$.
Modulo $3^e$ both $3^k$ and $9$ vanish, since $k\ge2$ and $e\le2$.
CRT gives the first congruence in~\eqref{R10:eq:returnpowers}, and the second follows.
\end{proof}

The multiplication-map trace of $x^n$ equals $T_{n,C,s}(a)$ also
in a nonreduced algebra: use the companion matrix and Newton's identities.
Consequently
\begin{equation}\label{R10:eq:returntraces}
 \begin{aligned}
 T_{3^k,C,s}(a)&\equiv T_{9,C,s}(a)\pmod p,\\
 T_{2\cdot3^k-1,C,s}(a)&\equiv T_{17,C,s}(a)\pmod p
 \end{aligned}
 \quad(k\ge2,\ k\equiv2\pmod{L_p}).
\end{equation}

\subsection{The prime-coefficient sieve and its denominator}\label{R10:sec:sieve}

For $p\in\mathcal S_n$ put
\[
 \mathcal R_{n,p}=\{r\in\F_p^\times:P_n(r)=0\},\qquad
 r_{n,p}=\#\mathcal R_{n,p}\ge1.
\]
Delete the finitely many witness primes dividing $M$.

\begin{lemma}[Prime-relative CRT]\label{R10:lem:primeCRT}
For each fixed finite set $\mathcal Q\subset\mathcal S_n$ of remaining primes,
\begin{equation}\label{R10:eq:primeCRT}
 \lim_{H\to\infty}
 \frac{\#\{a\le H:a\text{ prime},\ a\equiv b\pmod M,\
                   a\bmod p\notin\mathcal R_{n,p}\ (p\in\mathcal Q)\}}
      {\pi(H;M,b)}
       =\prod_{p\in\mathcal Q}\left(1-\frac{r_{n,p}}{p-1}\right).
\end{equation}
\end{lemma}
\begin{proof}
Apart from the finitely many coefficients equal to a prime in $\mathcal Q$,
prime coefficients occupy only nonzero residues modulo each $p$.
There are $p-1$ reduced choices, of which $r_{n,p}$ are forbidden.
CRT combines these with the fixed reduced class modulo $M$.
All resulting reduced classes modulo $M\prod_{p\in\mathcal Q}p$
have the same limiting prime density, by cyclotomic Chebotarev.
Counting them gives~\eqref{R10:eq:primeCRT}.
\end{proof}

The denominator is $p-1$, \emph{not} $p$.
Since $r_{n,p}\ge1$ and $\sum_{p\in\mathcal S_n}1/p=\infty$, the products
tend to zero as $\mathcal Q$ increases.
Coefficients missing every trace-zero residue avoid each fixed
$\mathcal Q$. Take their relative upper density as $H\to\infty$
\emph{before} enlarging $\mathcal Q$.
Thus relative density one of the prime coefficients have a witness
prime for this seed and the real sign.

For the complex sign, reflection of the cubic gives, for odd $n$,
\begin{equation}\label{R10:eq:reflection}
                           T_{n,C,1}(a)=-P_n(-a).
\end{equation}
The same proof therefore applies with root sets $-\mathcal R_{n,p}$.
Apply it to the two seeds and the two signs.
The intersection of the resulting four sets still has relative density
one: their four exceptional sets have a union of relative density zero.
No independence among these four conditions is assumed.

For a coefficient in this intersection choose one witness prime for
each seed/sign pair, and let $L(a)$ be the least common multiple of
their four periods $L_p$.
Equation~\eqref{R10:eq:returntraces} makes all four traces divisible by their
witnesses whenever $k\ge2$ and $k\equiv2\pmod{L(a)}$.
The common offset $2$ matters: four arbitrary progressions need not
intersect, whereas these ones do.

\subsection{Returned traces and actual composite floors}\label{R10:sec:floors}

\begin{lemma}\label{R10:lem:floor}
For odd $a>2C$, both signs give irreducible cubic Pisot numbers
$\beta\in(a,a+1)$, and all their traces are odd.
For $s=-1$, $\floor{\beta^n}=T_{n,C,-1}(a)$ at every positive odd $n$.
For $s=1$ and $n\ge3$, the floor is either $T_{n,C,1}(a)$ or
$T_{n,C,1}(a)-1$.
If a prime divides the seed trace at exponent $n_0=9$ or $17$, every
later return on the corresponding clock gives a composite floor.
\end{lemma}
\begin{proof}
For $s=-1$ the signs at $-1,0,1,a,a+1$ are
\[
 2C-a-1<0,\quad C>0,\quad1-a<0,\quad C(1-a)<0,\quad(a+1)^2-Ca>0.
\]
The roots lie in $(-1,0),(0,1),(a,a+1)$.
Writing the small roots as $-\rho,\sigma$,
their sum $\sigma-\rho=a-\beta<0$ gives $\rho>\sigma$.
For odd $n$,
\[
                         0<\beta^n-T_n=\rho^n-\sigma^n<1.
\]

For $s=1$ there is no nonpositive root: at $X=-y$ the value is
$-y^3-ay^2+Cy-C<0$ for $y\ge0$.
For $y\le1$ use $C(y-1)\le0$; for $y\ge1$ use $a>2C$.
Descartes's rule gives exactly one positive root, and
$f(a)<0<f(a+1)$ places it in $(a,a+1)$.
The other roots are nonreal, with common modulus
$\sqrt{C/\beta}<1/\sqrt2$. Thus
\[
                        |T_n-\beta^n|\le2(C/\beta)^{n/2}<1
                                      \qquad(n\ge3).
\]
In either sign no root is an integer: the large root lies between
consecutive integers, while any other real root is nonzero with modulus
below $1$. The monic cubics are irreducible.

The initial traces $3,a,a^2+2C$ are odd, and their recurrence preserves
parity because its coefficient sum is $1+C+C=1\pmod2$.
This proves the floor alternatives, including the even predecessor
in the complex case.

Let $p$ be a seed witness. Its seed trace $T_{n_0}$ is positive and
divisible by $p$, hence $T_{n_0}\ge p$.
At a later index $k\ge3$, the clock exponent $n$ exceeds $n_0$.
Since $\beta>a\ge3$,
\[
                     T_n>\beta^n-1>\beta^{n_0}+1>T_{n_0}\ge p.
\]
At a return, $T_n$ is a proper multiple of $p$.
The real floor is this composite trace.
The complex floor is either that trace or its even predecessor,
which exceeds $2$; both are composite.
\end{proof}

\begin{proof}[Completion of Theorem~\ref{R10:thm:main}]
Section~\ref{R10:sec:sieve} supplies a relative-density-one set of prime
coefficients and a common period $L(a)$ for their four trace witnesses.
Every $k=2+tL(a)$, $t\ge1$, is strictly later than the seed.
Lemma~\ref{R10:lem:floor} proves all four floors composite.

For a finite union of reduced classes, sum the finitely many exceptional
estimates. Different fixed moduli may first be replaced by a common
modulus; nonreduced lifts contain only finitely many primes and can be discarded.
\end{proof}

One may in particular impose either residual condition $a\equiv-s\pmod3$
along with oddness. For prime coefficients, $\gcd(a,C)=1$ deletes only
finitely many values.
In the real branch, writing the small roots as $-\rho,\sigma$, the
equation at $-\rho$ gives
$a\rho^2=C(1+\rho)-\rho^3<2C$.
Since $\beta\in(a,a+1)$, for fixed $C$ this yields
$\rho<\beta^{-17/40}$ and $C<\beta^{3/20}$ for all sufficiently large
$a$. The norm relation $\rho\sigma=C/\beta$ also gives
$\rho<\sigma^{17/23}$ in that range.
These cofinite real restrictions preserve the prime-relative conclusion.
No such transfer to an arbitrary thin phase-defined subset is asserted.

\subsection{Arithmetic progressions of certified prime coefficients}

We use Green--Tao's relative-prime theorem: every subset of the primes
of positive relative upper density contains arbitrarily long nonconstant
arithmetic progressions \cite[Theorem~1.2]{GreenTao}.
This is an imported theorem, not a new transference principle.

\begin{corollary}\label{R10:cor:GT}
For every length $\ell\ge2$ and bound $B$, there is an arithmetic
progression of prime coefficients
\[
             a_j=u+jd>B,\qquad j=0,\ldots,\ell-1,\qquad d>0,
\]
all congruent to $b\pmod M$, and an integer $L^*\ge1$ such that
at every $k=2+tL^*$, $t\ge1$, all $4\ell$ associated floors,
from both signs and both clocks, are composite.
\end{corollary}
\begin{proof}
Let $\mathcal G$ be the prime set certified in Theorem~\ref{R10:thm:main}.
Since $\pi(H;M,b)/\pi(H)\to1/\varphi(M)$, its density relative to all
primes is $1/\varphi(M)>0$.
Remove the finite segment up to $\max(B,2C)$ and apply Green--Tao.

For the resulting finite coefficient progression put
\[
                              L^*=\lcm_{0\le j<\ell}L(a_j).
\]
Every $k=2+tL^*$, $t\ge1$, is a later common return for every coefficient,
proving the assertion.
\end{proof}

The exact arithmetic progression is in the coefficients, equivalently
in the integer parts of the large roots, not necessarily in the algebraic
roots themselves. No bounded-gap or bounded-cluster conclusion is drawn.
Prime-relative density is not substituted for the distribution input of
a Maynard-type argument.

\subsection{Scope and relation to earlier work}

The preceding integer-density proof used primes dividing $a$.
That mechanism does not establish density among prime coefficients.
Here nonzero seed-root classes and reduced-residue equidistribution
supply a separate result.

The local corpus \cite{K031} already treats projective returns and gives
a counterexample showing that an arbitrary observed prime trace need not
recur after a $3$-primary transient.
We do not assume automatic recurrence of such a seed.
The chosen residues modulo $9$ force the relevant transient to have
ended by index $2$, as quantified in~\eqref{R10:eq:v3}.
The inspected antecedent passages do not provide the prime-coefficient
density argument used here; no literature-wide novelty claim is made.

The anchor in~\eqref{R10:eq:anchor} certifies compatibility through an actual
Frobenius class. There is no hidden splitting-field disjointness,
primitive-root hypothesis, or separability assumption on the later
cubic quotient algebra.

The result concerns each fixed $C\equiv1\pmod3$ and fixed reduced
coefficient progression. It does not exclude every coefficient,
give an effective exceptional-set rate, or identify a coefficient attached
to a hypothetical least prime-representing constant.
An arbitrary further thin conditioning requires another theorem.
Neither the density result nor Corollary~\ref{R10:cor:GT} gives a universal
cubic Pisot or Mills-transcendence closure.


\subsection{Portable exact checks}\label{R10:app:certificates}

The original source archive includes two separately implemented mathematical verifiers
with JSON/log outputs. Both require Python~3 and SymPy.
Run from the corresponding extracted source archive's directory:
\begin{lstlisting}[language=bash]
python3 certificates/prime_coefficient_mod9_verify.py
python3 certificates/audit_prime_coefficient_density_verify.py
(cd certificates && sha256sum -c SHA256SUMS.txt)
\end{lstlisting}
The second filename identifies the independent implementation.

The first program checks:
\begin{itemize}
\item monicity and degrees through $n=17$, odd boundary and reflection
identities, and the two symbolic constant terms;
\item all $81$ pairs $(A,C)\pmod9$ for the integer-valued congruence
$P_9(A)\equiv P_3(A)\pmod9$, and the complete seed table;
\item $42400$ residue-root equivalence checks for sampled squarefree
parts, including primes of bad specialization;
\item $5152$ cubic unit-exponent checks, including $210$ nonreduced cases;
\item $68$ exponent returns, $282$ seed-trace returns, and $24$
exact reduced-CRT counts.
\end{itemize}
The congruence $P_9(A)\equiv P_3(A)\pmod9$ is integer-valued,
not a coefficientwise assertion in $\Z[A,C]$.

The independent program checks the symbolic identities, $10$ exact
squarefree/CRT anchors, $12044$ quotient algebras (including $370$
nonreduced cases), and $24$ reduced-CRT counts.
All acceptance arithmetic is exact.

These computations validate finite algebraic inputs and sample
specializations. They do not prove Chebotarev, prime equidistribution,
or Green--Tao, and do not compute a numerical prime density or an
exceptional-set rate. The analytic conclusions use the general arguments
and the precisely identified imported theorems.




\endgroup

\section{Signature controlled density and quantitative exceptional sets}\label{src:R11}
The signature argument generalizes the fixed-norm construction to fixed pairs of coefficients and quantifies the exceptional set. The projective synthesis incorporated here improves its original logarithmic exponent from $1/6$ to $2/3$ and retains the full sieve proof.
\medskip
\begingroup
\label{R11:sec:synthesis}
% Standalone section fragment. Requires amsmath, amssymb, amsthm.
% Use the existing theorem/lemma/corollary environments; no new macros required.
% Bibliographic keys: AKK (Arango-Pineros--Keliher--Keyes, Theorem A),
% Lenstra (Chebotarev, Fact 2.1 and Theorem 3.1).
\label{R11:syn:projective-density-improvement}

The full-state return in the fixed-line density argument can be replaced
by a scalar return. This admits four, rather than one, reduced residue
classes modulo $9$. The resulting improvement is a consequence of the
signature and projective mechanisms developed in this article, not an assumption on
primitive roots or on the distribution of a distinguished coefficient.

Fix positive integers $B,q$ and consider
\[
 f_{a,B,q,s}(X)=X^3-aX^2-BX-sq,\qquad s\in\{-1,1\}.
\]
Write $\beta_{a,B,q,s}$ for its dominant root and put
\[
 A_-=B+q+1,\qquad
 A_+=\max\{B+q+1,B^2/(4q),2q\}.
\]
Fix an even modulus $M$ and a reduced residue class $b\pmod M$.
Let $E_{\mathbb Z}(X)$ count the integer coefficients
$a\equiv b\pmod M$, $A_-<a\le X$, for which there is no positive
integer $L$ such that both real-branch floors
\[
 \left\lfloor\beta_{a,B,q,-1}^{3^k}\right\rfloor,
 \qquad
 \left\lfloor\beta_{a,B,q,-1}^{2\cdot3^k-1}\right\rfloor
\]
are composite for every $k=1+tL$, $t\ge1$.
Let $E_{\mathbb P}(X)$ count the prime coefficients among these exceptions.
For the simultaneous four-floor conclusion use $a>A_+$ and include
both signs; this version will require $B\equiv q\pmod2$.

\begin{theorem}[Projective improvement of the fixed-line rates]
\label{R11:syn:projective-density-rate}
For each fixed $B,q,M$ as above,
\[
 E_{\mathbb Z}(X)\ll_{B,q,M}\frac{X}{(\log X)^{2/3}},
 \qquad
 E_{\mathbb P}(X)\ll_{B,q,M}\frac{X}{(\log X)^{5/3}}.
\]
The same estimates hold for exceptions to the simultaneous four-floor
conclusion when $B\equiv q\pmod2$.
Thus its exceptional proportion among prime coefficients is
$O_{B,q,M}((\log X)^{-2/3})$.
Finite unions of fixed reduced odd classes are allowed.
On the diagonal $B=q=C$, the sufficient root threshold may be replaced
by $a>2C$, for every fixed positive integer $C$.
No bound on $L$, uniformity in the fixed parameters, or exclusion of
an individual coefficient is asserted.
\end{theorem}

\subsubsection{A universal projective exponent}

\begin{lemma}[Scalar rather than full-state returns]
\label{R11:syn:projective-universal}
Let $p\ge5$ be prime, let $f\in\mathbb F_p[X]$ be a monic cubic,
and let $u$ be any unit in $R=\mathbb F_p[X]/(f)$. Set
\[
 D_p=p\,\operatorname{lcm}(p^2-1,p^2+p+1)
     =\frac{p(p^2-1)(p^2+p+1)}{\gcd(p-1,3)}.
\]
Then $u^{D_p}\in\mathbb F_p^\times$, viewed as the scalar units of $R$.
Moreover
\[
 v_3(D_p)=
 \begin{cases}
  v_3(p+1),&p\equiv2\pmod3,\\
  v_3(p-1),&p\equiv1\pmod3.
 \end{cases}
\]
In particular $v_3(D_p)=1$ for
$p\pmod9\in\{2,4,5,7\}$.
\end{lemma}
\begin{proof}
If $f$ is irreducible, $R$ is the field $\mathbb F_{p^3}$ and
$u^{p^2+p+1}$ is its nonzero field norm, hence is scalar.
Since $p^2+p+1$ divides $D_p$, the assertion follows.
If $f$ is reducible, every field factor of the reduced algebra has
degree at most two. Thus $u^{p^2-1}=1+z$ with $z$ in the nilradical.
The nilradical of a cubic quotient has cube zero. Consequently
$(1+z)^p=1+z^p=1$, and $u^{p(p^2-1)}=1$.
This exponent divides $D_p$, proving the result even at repeated
reductions. Finally
$\gcd(p^2-1,p^2+p+1)=\gcd(p-1,3)$.
If $p\equiv1\pmod3$, then $v_3(p^2+p+1)=1$; if
$p\equiv2\pmod3$, that factor is prime to $3$.
These facts give the displayed valuations.
\end{proof}

For the cubic $f_{a,B,q,s}$, exclude the primes dividing $q$ so that
$x=X\bmod f$ is a unit. At any remaining prime in the four indicated
classes put
\[
 m_p=D_p/3,\qquad L_p=\operatorname{ord}_{m_p}(3),
\]
with order one if the modulus is one. For $k\ge1$ and
$k\equiv1\pmod{L_p}$, the differences
\[
 3^k-3,\qquad (2\cdot3^k-1)-5=2(3^k-3)
\]
are multiples of $D_p$. They are divisible by $m_p$ by definition
of the order and by $3$ directly. Therefore each returned ring
state is a nonzero scalar multiple of its seed state.
The trace of multiplication is linear over $\mathbb F_p$, so
\begin{equation}\label{R11:syn:projective-trace-return}
 \begin{split}
 T_{3,s}(a)\equiv0\pmod p
 &\ \Longrightarrow\ T_{3^k,s}(a)\equiv0\pmod p,\\
 T_{5,s}(a)\equiv0\pmod p
 &\ \Longrightarrow\ T_{2\cdot3^k-1,s}(a)\equiv0\pmod p.
 \end{split}
\end{equation}
The equalities of the trace values themselves are not required and
are not asserted. The multiplication trace agrees with the Newton
trace even in a nonreduced cubic algebra, by the companion matrix.

\subsubsection{Four compatible Frobenius fibers}

For the real branch the two seed polynomials are
\[
 P_3(A)=A^3+3BA-3q,
 \qquad
 P_5(A)=A^5+5BA^3-5qA^2+5B^2A-5Bq.
\]
Each has exactly one real root, which is positive and simple.
For $P_3$ this follows from its positive derivative and negative value
at zero. For $P_5$, write
\[
 P_5(A)=5(A^2+B)\bigl(Q_B(A)-q\bigr),\qquad
 Q_B(A)=\frac{A(A^4+5BA^2+5B^2)}{5(A^2+B)}.
\]
The derivative
\[
 Q_B'(A)=\frac{3A^6+10BA^4+10B^2A^2+5B^3}{5(A^2+B)^2}>0
\]
and $Q_B(A)\sim A^3/5$ at both ends prove the assertion.
Choose the monic irreducible factor $h\in\mathbb Z[A]$ containing
this unique real root $\alpha$. Its degree $d$ is $1$, $3$, or $5$,
its constant term is nonzero, and $K=\mathbb Q(\alpha)$ has exactly
one real embedding.

Let $Z=\mathbb Q(\zeta_9)$. Then $K\cap Z=\mathbb Q$.
Indeed, degrees one and five are immediate from the degrees;
in degree three a nontrivial intersection would identify $K$ with
the unique cubic subfield of $Z$, which is totally real, contradicting
its signature. Let $E$ be the compositum of the splitting field of
$h$ with $Z$, and write
\[
 G=\operatorname{Gal}(E/\mathbb Q),\qquad
 \pi:G\longrightarrow\operatorname{Gal}(Z/\mathbb Q).
\]
The stabilizer $H$ of $\alpha$ surjects onto this cyclic quotient of
order six: the fixed field in $Z$ of $\pi(H)$ is
$Z\cap E^H=Z\cap K=\mathbb Q$.
Every other root stabilizer also surjects, by conjugacy.
This establishes compatibility, without a linear-disjointness
assumption on the splitting field.

For each $r\in\{2,4,5,7\}$ let $\tau_r(\zeta_9)=\zeta_9^r$.
Since each root stabilizer has order $|G|/d$, its intersection with
each fiber $\pi^{-1}(\tau_r)$ has size $|G|/(6d)$.
Counting pairs consisting of an element and one of its fixed roots
therefore gives
\begin{equation}\label{R11:syn:projective-weighted-fibers}
 \frac1{|G|}
 \sum_{g\in\pi^{-1}(\{\tau_2,\tau_4,\tau_5,\tau_7\})}
      |\operatorname{Fix}(g)|=\frac46=\frac23.
\end{equation}
Delete the finitely many primes at most seven, the primes dividing
$qM h(0)\operatorname{disc}(h)$, and any additional ramified primes
of $E$. Let $r(p)$ be the number of roots of $h\bmod p$ when
$p\bmod9\in\{2,4,5,7\}$ is not deleted, and zero otherwise.
All these roots are nonzero coefficient residues; $0\le r(p)\le5$.
Factorization--Frobenius compatibility and Chebotarev
\cite[Fact~2.1 and Theorem~3.1]{Lenstra}, together with the
Chebotarev Mertens theorem \cite[Theorem~A]{R11-AKK}, imply
\begin{equation}\label{R11:syn:projective-weighted-mertens}
 \sum_{p\le y}\frac{r(p)}p
       =\frac23\log\log y+O_{B,q,M}(1).
\end{equation}
For clarity, the Mertens input is applied separately to the finitely
many conjugacy classes and weighted by their fixed-root counts;
\eqref{R11:syn:projective-weighted-fibers} is exactly the resulting
coefficient. Thus the value $2/3$ does not come from an unproved
independence heuristic.

A root $a\bmod p$ is a seed-trace zero because $h$ divides $P_n$.
For the partner sign, the identity
$T_{n,1}(a)=-P_n(-a)$ for odd $n$ gives the reflected nonzero root
classes, with the same cardinalities and the same Mertens estimate.

\subsubsection{Quantitative sieve and simultaneous floor returns}

Here is a self-contained upper-sieve calculation recording the
logarithmic exponent and its error term. Fix one seed/sign pair.
For integer coefficients forbid its $r(p)$ root classes, writing
$\rho(p)=r(p)$. For prime coefficients additionally forbid the zero
class at every undeleted prime, writing $\rho(p)=1+r(p)$ there.
At deleted primes impose no condition. Omitting primes with zero
local count, we have $0<\rho(p)\le6<p$ and
\[
 \sum_{p\le y}\frac{\rho(p)}p
  =\kappa\log\log y+O(1),\qquad
 \kappa=\frac23\ \hbox{or}\ \frac53,
\]
respectively. Nonzero seed roots are essential for the extra class
to be distinct in the prime case.

Let $N$ be the number of integer coefficients in the fixed progression
and relevant interval up to $X$. On squarefree integers supported on
the sieve primes extend $\rho$ multiplicatively. CRT gives
$N\rho(v)/v+R_v$ coefficients forbidden at every prime dividing $v$,
where $|R_v|\le\rho(v)$. Put
\[
 u(p)=\rho(p)/p,\qquad g(p)=u(p)/(1-u(p)),\qquad
 G(z)=\sum_{\substack{d\le z\\d\text{ squarefree, supported}}}g(d).
\]
The Selberg quadratic weights, supported on $d\le z$, are
\[
 \lambda_d=\mu(d)\prod_{p\mid d}(1-u(p))^{-1}
       \frac{G_d(z/d)}{G(z)},\qquad
 G_d(v)=\sum_{\substack{m\le v\\(m,d)=1}}g(m).
\]
All sums here use squarefree supported integers. They satisfy
$\lambda_1=1$, $|\lambda_d|\le1$, and
\[
 \sum_{d,e\le z}\lambda_d\lambda_e u([d,e])=G(z)^{-1}.
\]
For verification, set $Y_r=\sum_{r\mid d}\lambda_du(d)$.
Then the quadratic form is $\sum_rY_r^2/g(r)$ and
$\sum_r\mu(r)Y_r=1$; its minimizer is
$Y_r=\mu(r)g(r)/G(z)$. Divisor inversion gives the displayed weights.
Their absolute bound follows by including in $G(z)$ all distinct
terms $\ell m$ with $\ell\mid d$, $m\le z/d$, and $(m,d)=1$.
Squaring the divisor sum majorizes the indicator of survival, so
\[
 S(X,z)\le \frac{N}{G(z)}+
        \sum_{d,e\le z}\rho([d,e])
 \le\frac{N}{G(z)}+z^2(1+\log z)^{10}.
\]
For the last inequality use
$\rho([d,e])\le\rho(d)\rho(e)$,
$\rho(d)\le6^{\omega(d)}\le\tau_6(d)$, and
$\sum_{d\le z}\tau_6(d)\le z(1+\log z)^5$.

Finally $G(z)\gg(\log z)^\kappa$.
Indeed set $y=z^{1/24}$ and give a squarefree product $D$ of primes
at most $y$ probability proportional to $g(D)$.
Its expected logarithm is
\[
 \sum_{p\le y}\frac{\rho(p)\log p}{p}
 \le6\log y+O(1)=\tfrac14\log z+O(1).
\]
Markov's inequality puts at least half the product weight on $D\le z$
for all sufficiently large $z$. This weight is included in $G(z)$,
while
\[
 \prod_{p\le y}(1+g(p))
   =\prod_{p\le y}(1-\rho(p)/p)^{-1}
   \asymp(\log y)^\kappa
\]
follows by taking logarithms and using the convergence of
$\sum_p\rho(p)^2/p^2$.
Choose $z=X^{1/4}$. The result is
\[
 S(X,z)\ll\frac{X}{(\log X)^\kappa}.
\]
In the prime case the coefficients at most $z$ contribute at most
$z$, and every larger prime avoids the zero classes at sieve primes.
Thus the same estimate applies to prime coefficients lacking this
seed/sign witness. This does not require prime equidistribution
in a growing modulus.

It remains to check that avoiding these exceptional sets supplies
actual composite floors at common indices. For $a>A_-$, the real
cubic has roots $-\rho,\sigma,\beta$, with
$0<\sigma<\rho<1$ and $a<\beta<a+1$. This follows from its signs
at $-1,0,1,a,a+1$, and the ordering of the small roots from their sum
$a-\beta<0$. At every positive odd exponent,
\[
 0<\beta^n-T_{n,-1}(a)=\rho^n-\sigma^n<1,
\]
so its floor is its trace. For $a>A_+$ the partner has no nonpositive
real root, since for $y\ge0$
\[
 f_{a,B,q,1}(-y)=-y^3-ay^2+By-q
       \le-y^3+B^2/(4a)-q<0.
\]
Descartes's rule and the signs at $a,a+1$ give its sole real root
in $(a,a+1)$. Its other two roots have modulus $\sqrt{q/\beta}<1/\sqrt2$,
so $|T_{n,1}(a)-\beta^n|<1$ for $n\ge3$.
These cubics are irreducible since they have no integer root.
When $a$ is odd and $B\equiv q\pmod2$, Newton's recurrence preserves
oddness of every trace. Thus the partner floor is either its odd
trace or its even predecessor.

Outside the union of the two one-pair exceptional sets, choose a
witness prime for each real seed; for four floors choose witnesses
for both seeds and signs. Take $L$ to be the least common multiple
of their $L_p$. All seed trace zeros recur at every
$k=1+tL$, $t\ge1$, by \eqref{R11:syn:projective-trace-return}.
The common offset $1$ is crucial; arbitrary progressions need not
intersect. If $p$ divides a positive seed trace $T_{n_0,s}(a)$, then
$p\le T_{n_0,s}(a)$. At a later clock index, since $\beta>a\ge3$,
\[
 T_{n,s}(a)>\beta^n-1>\beta^{n_0}+1>T_{n_0,s}(a)\ge p.
\]
The returned trace is a proper multiple of $p$. The real floor is
that trace. In the partner case the alternative is an even integer
larger than two, also composite. Taking a finite union of the
exceptional estimates proves the theorem. Division by
$\pi(X;M,b)\sim X/(\varphi(M)\log X)$ gives its relative prime form.
The sharper diagonal threshold follows from the same root signs
and the bound $C/\beta<1/2$; it changes only a finite set of coefficients.

\begin{corollary}[A finite catalogue with one index progression]
\label{R11:syn:projective-finite-catalogue}
Fix finitely many positive pairs $(B_i,q_i)$ and one fixed reduced odd
coefficient progression. For all but
$O(X/(\log X)^{2/3})$ integer coefficients up to $X$, and all but
$O(X/(\log X)^{5/3})$ prime coefficients up to $X$, there is one
progression $k=1+tL(a)$, $t\ge1$, on which the two real floors are
composite simultaneously for every pair in the catalogue.
The partner floors may simultaneously be included for every pair
with $B_i\equiv q_i\pmod2$. All coefficients are taken above the
largest applicable fixed threshold.
\end{corollary}
\begin{proof}
Take a finite union of the preceding exceptional sets and the least
common multiple of the finitely many resulting periods. The common
offset remains one. Constants may depend on the entire fixed catalogue;
no assertion is made for a catalogue growing with $X$.
\end{proof}

\paragraph{Relation to the earlier rate.} The original full-state exponent used only the residue class $2\pmod9$ and gave logarithmic exponents $1/6$ and $7/6$. The scalar-return argument above replaces that proof and improves those exponents to $2/3$ and $5/3$ with the same fixed-parameter scope.
\subsection{Prime coefficient progressions and scope}

\begin{corollary}[Finite coefficient progressions with common returns]
\label{R11:cor:GT}
For each fixed $B,q,M,b$, every length $\ell\ge2$ and every lower bound
admit a nonconstant arithmetic progression of certified prime
coefficients above that bound.
All its coefficients have a common progression
$k=1+tL_*$, $t\ge1$, of composite floors for both real clocks.
If $B\equiv q\pmod2$, all four sign and clock combinations can
be included simultaneously.
\end{corollary}

\begin{proof}
The certified prime set has relative density $1/\varphi(M)>0$
among all primes. Removing an initial segment preserves that density.
The relative-prime form of the Green--Tao theorem
\cite[Theorem~1.2]{GreenTao} gives an arithmetic progression of any
prescribed finite length.
For its finitely many coefficients choose their periods $L(a)$
and take $L_*=\lcm_a L(a)$. Their common offset is $1$,
so this least common multiple synchronizes every asserted return.
\end{proof}

The arithmetic progression is in the coefficients, equivalently
in the integer parts of the dominant roots, not necessarily in
the roots themselves.
No bounded-gap or Maynard bounded-cluster conclusion is inferred.

The mechanism differs from the earlier integer-line result
\cite{R08} and from the congruence-restricted prime-line argument
\cite{R10}. Here the real signature of a factor of $P_3$ or $P_5$
replaces the earlier integer CRT anchor.
It proves precisely the needed compatibility even if the
splitting field has nontrivial cyclotomic intersection.
It also avoids assuming that every quintic is irreducible.

A close methodological precedent is Br\"udern--Soundararajan
\cite[Section~2, Lemma~3]{R11-BS}.
They characterize containment of the $m$th roots of unity in an
irreducible polynomial's root field by the condition that all but
finitely many prime divisors of its values lie in $1\pmod m$,
using Chebotarev. The same lemma is numbered~2 in their
arXiv version~1.
This relates polynomial prime divisors to root-field cyclotomic
information, as does our construction.
The restriction-image argument in the root-stabilizer argument
is standard Galois correspondence, not a new general Galois principle.

For a neighboring Pisot-floor construction, Hahn--Ismailescu--Kim--Kim
\cite[Section~4, Theorem~4]{R11-HIKK} prove compositeness at every positive
power for the dominant root of $X^3-aX^2+cX-1$ when
$c\ge0$, $a\ge2c+5$, and $a\equiv c\equiv3\pmod6$.
Those conditions force $a>3$ to be divisible by $3$, so the coefficient
$a$ is composite. Their all-power construction concerns a different
linear-coefficient sign; it is not the fixed negative-linear,
prime-coefficient density assertion proved here.

The existing corpus \cite{K031} already discusses modular and
projective returns and warns that a completely split prime seed
need not recur along a prescribed clock.
The original choice $p\equiv2\pmod9$ controlled the $3$-primary part
of a universal cubic exponent, which makes recurrence at the
first seed valid. No literature-wide priority claim is made.

These statements concern each fixed positive pair $(B,q)$.
They do not cover arbitrary coefficient signs or positive-norm
totally real cubics. With $B\not\equiv q\pmod2$, the partner trace
congruences still hold, but this proof gives no corresponding
partner-floor conclusion.
The estimates are not uniform in the parameters, the periods,
or an index-density parameter. They supply neither an explicit
onset nor emptiness of the exceptional set.
Powering a Pisot number generally changes $B,q$, so the fixed-parameter
estimates do not control such a nonlinear orbit.
An arbitrary thin phase condition, a specified individual survivor,
or a candidate for a least prime-representing constant needs a
separate argument.
In particular no universal Mills or Pisot-transcendence closure follows.


\subsection{Portable exact certificates}

The accompanying programs use exact integer, rational, and symbolic
arithmetic. In the original source archive, run from its extracted directory:
\begin{lstlisting}[language=bash]
python3 certificates/signature_prime_density_verify.py
python3 certificates/fixed_B_q_signature_verify.py
python3 certificates/quantitative_sieve_verify.py
(cd certificates && sha256sum -c SHA256SUMS.txt)
\end{lstlisting}
The first two programs require Python~3 and SymPy.
The third uses only the Python standard library.
Their JSON outputs and concise logs are included.

The diagonal checker verifies the symbolic trace and derivative
identities, $252$ exact Sturm and factor samples, the order-six
group modulo $9$, $6168$ cubic quotient algebras including $162$
nonreduced cases, $44$ exponent returns, $40$ usable sample
prime witnesses, $260$ reflected trace returns, and $48$
exact reduced-CRT counts.

The general $(B,q)$ checker verifies the displayed derivative identity,
$362$ real-root and factor samples including three linear-factor
cases, $63$ two-sign root-location and parity samples,
$8916$ cubic quotient algebras including $418$ repeated reductions,
and $17832$ seed-trace equalities.

The original quantitative checker verifies three exact finite group models
of the single-fiber identity underlying \eqref{R11:syn:projective-weighted-fibers}, including a non-disjoint subdirect product. It uses the original
restriction $p\equiv2\pmod9$ and weighted density $1/6$, not the new
four-fiber density $2/3$. It also verifies
four rational Selberg quadratic-form and residue-sieve calculations.
It checks the minimizing weights, their absolute bound by $1$,
the exact quadratic minimum, the squared-sum majorant,
and the remainder estimate.

These finite tests supplement the general proofs.
They do not prove the imported Chebotarev, Mertens, or Green--Tao
theorems; do not turn the asymptotic constants into numerical ones;
and do not establish exclusion of every coefficient.



\endgroup

\clearpage\part{Thin families and adaptive projective witnesses}

\section{The thin prime power family}\label{src:R12}
Density statements need not resolve a thin sequence such as $a=7^m$. We therefore keep this family separate and prove explicit residue-class returns, together with an obstruction to covering it by finitely many fixed primes.
\medskip
\begingroup



\subsection{The thin family and the two conclusions}

For $m\ge2$ put $a=7^m$ and
\begin{equation}\label{R12:eq:family}
 f_m(X)=X^3-aX^2-4X+4.
\end{equation}
We will show that its dominant root $\beta_m$ is totally real cubic
Pisot. Recall that a Pisot number is a real algebraic integer greater
than one whose other conjugates have modulus less than one. Let
$T_n(a)$ be the sum of the $n$th powers of the three roots, and retain
the original clock indices
\begin{equation}\label{R12:eq:clocks}
 E_k^+=3^k,\qquad E_k^-=2\cdot3^k-1,\qquad k\ge1.
\end{equation}
Here $\floor{\beta_m}=7^m$ is composite: the family is not already a
classical Mills power $\xi^{3^j}$ at a required prime index, even though
it satisfies the root, norm, and small-conjugate inequalities proved below.
Our clocks start at $k=1$, not at exponent one.

This exponential coefficient family is thin among integer coefficients.
A density-one result in the ambient coefficient line does not imply
that any particular $7^m$ belongs to its successful set. The arguments
below work directly on this family.

\begin{mainthm}\label{R12:thm:A}
For every $m\ge2$ with $m\not\equiv0,4\pmod{12}$, each sequence
\[
 \floor{\beta_m^{E_k^+}},\qquad \floor{\beta_m^{E_k^-}}
\]
has an explicit infinite arithmetic progression of composite terms
in its original index $k$. The negative clock is certified for every
$12\nmid m$. The classical clock is certified whenever $4\nmid m$ or
$m\equiv2\pmod3$.

Every indicated term is composite, including the first index in its
certificate. The two clocks may use \emph{different} progressions;
no common simultaneous progression is asserted.
\end{mainthm}

\begin{mainthm}\label{R12:thm:B}
For every finite set $P$ of primes, there are infinitely many $m\ge2$
such that, for every $p\in P$ and every $k\ge1$,
\[
 p\nmid\floor{\beta_m^{E_k^+}},\qquad
 p\nmid\floor{\beta_m^{E_k^-}}.
\]
These $m$ can be arbitrarily large and their Pisot numbers satisfy
\[
 4<\beta_m^{3/20},\qquad
 |\gamma|<\tfrac12\beta_m^{-17/40}
 \quad\text{for both other conjugates }\gamma.
\]
\end{mainthm}

Theorem~\ref{R12:thm:B} does not say that the floors are prime. Its
quantifiers are $\forall P$ finite, $\exists$ infinitely many $m$,
not the reverse. It excludes a fixed finite-prime cover of the whole
family, but does not exclude a witness prime depending on $m$.
All results needed here are proved below; no prime-distribution
theorem or ambient density assertion is used.

\subsection{Actual Pisot roots and exact odd floors}

\begin{lemma}\label{R12:lem:roots}
For every integer $a>8$, the polynomial $X^3-aX^2-4X+4$ is irreducible
and has roots $-w,v,\beta$ with
\[
 0<v<w<1,\qquad a<\beta<a+1.
\]
For every positive odd $n$,
\begin{equation}\label{R12:eq:floor}
 0<\beta^n-T_n(a)=w^n-v^n<1,
 \qquad \floor{\beta^n}=T_n(a).
\end{equation}
If $a=7^m$ and $m\ge15$, both small conjugates and the norm satisfy
the bounds in Theorem~\ref{R12:thm:B}.
\end{lemma}
\begin{proof}
The polynomial values at $-1,0,1$ have signs $-,+,-$, since
$f(-1)=7-a$ and $f(1)=1-a$. Also
$f(a)=4(1-a)<0$ and $f(a+1)=(a-1)^2>0$.
There is one root in each of $(-1,0),(0,1),(a,a+1)$.
These intervals contain no integer, so the monic integral cubic has
no rational root and is irreducible. Vieta gives $v-w=a-\beta<0$,
which proves \eqref{R12:eq:floor}.

The equation at $-w$ gives $aw^2=4+4w-w^3<8$.
For $m\ge15$ we have $a>2^{41}$ and $\beta<a+1<2a$, hence
\[
 (2w)^{40}\beta^{17}
 <(32/a)^{20}(2a)^{17}=\frac{2^{117}}{a^3}<1.
\]
This gives the asserted strict small-conjugate bound, also for $v<w$.
The norm has absolute value four and $4^{20}<\beta^3$.
\end{proof}

Newton's identities give the integral trace recurrence
\begin{equation}\label{R12:eq:trace}
\begin{split}
 T_0(a)&=3,\quad T_1(a)=a,\quad T_2(a)=a^2+8,\\
 T_n(a)&=aT_{n-1}(a)+4T_{n-2}(a)-4T_{n-3}(a),\qquad n\ge3.
\end{split}
\end{equation}
It also defines trace polynomials at arbitrary integer parameters,
including reducible specializations. When $a=7^m$ all positive traces
are odd. At every clock index $k\ge1$, the exponent is odd and at
least three, so
\begin{equation}\label{R12:eq:large}
 T_{E_k^\pm}(a)>\beta^{E_k^\pm}-1>a^3-1
 \ge49^3-1>19.
\end{equation}
Thus a divisor in $\{5,11,13,19\}$ proves actual floor compositeness
at every certified index, not merely eventually.

\subsection{Exact returns in a cubic finite algebra}

Fix a prime $p$ and a coefficient residue $a\in\F_p$. Work in
\[
 R_{p,a}=\F_p[X]/(X^3-aX^2-4X+4),\qquad x=X\bmod f.
\]
The trace of multiplication by an element of this three-dimensional
algebra is the linear functional
\begin{equation}\label{R12:eq:tau}
 \tau(c_0+c_1x+c_2x^2)=3c_0+ac_1+(a^2+8)c_2.
\end{equation}
In particular $\tau(x^n)=T_n(a)\pmod p$. This follows from
\eqref{R12:eq:trace}, or directly from the multiplication matrix
\begin{equation}\label{R12:eq:matrix}
 M_{p,a}=\begin{pmatrix}0&0&-4\\1&0&4\\0&1&a\end{pmatrix}.
\end{equation}
These identities do not require a separable reduction.

Starting from $u_0=x$, define
\[
 F_+(u)=u^3,\qquad F_-(u)=x^2u^3.
\]
Their respective iterates are $u_k=x^{E_k^+}$ and $u_k=x^{E_k^-}$,
now allowing $k=0$ just for initialization.

\begin{lemma}[Exact return certificate]\label{R12:lem:return}
For either map, suppose $k_0,h$ are positive integers and
\[
 u_{k_0+h}=u_{k_0},\qquad \tau(u_{k_0})=0.
\]
Then $p\mid T_{E_k^\pm}(a)$ at every original index
$k=k_0+th$, $t\ge0$.
\end{lemma}
\begin{proof}
Equality of two states under a deterministic map makes their later
iterates equal. Repeating the return gives
$u_{k_0+th}=u_{k_0}$, and taking \eqref{R12:eq:tau} proves the claim.
No purely periodic behavior from $k=0$ is assumed.
\end{proof}

\subsubsection{The finite certificate table}

The orders of $7$ modulo $5,11,13,19$ are $4,10,12,3$, respectively.
Table~\ref{R12:tab:returns} gives exact return certificates, with the
coefficient residue forced by its condition on $m$. The state column
lists $(c_0,c_1,c_2)$ for $u_{k_0}$; in every row,
$u_{k_0+h}=u_{k_0}$ and \eqref{R12:eq:tau} is zero. The displayed $h$
need not be the minimal period.

\begin{table}[htbp]
\centering
\small
\setlength{\tabcolsep}{5pt}
\begin{tabular}{llrrrrrl}
\toprule
Label & Clock & $p$ & $m$ class & $a\bmod p$ & $k_0$ & $h$ & $u_{k_0}$\\
\midrule
C1 & $+$ &5&$1\bmod4$&2&1&2&$(1,4,2)$\\
C2 & $+$ &5&$2\bmod4$&4&1&30&$(1,4,4)$\\
C3 & $+$ &5&$3\bmod4$&3&6&30&$(0,1,1)$\\
N2 & $-$ &5&$2\bmod4$&4&5&30&$(4,3,4)$\\
N3 & $-$ &5&$3\bmod4$&3&10&30&$(4,2,1)$\\
N11a & $-$ &11&$1\bmod10$&7&2&4&$(7,8,0)$\\
N11b & $-$ &11&$5\bmod10$&10&13&36&$(6,9,10)$\\
N11c & $-$ &11&$7\bmod10$&6&2&4&$(2,10,8)$\\
N11d & $-$ &11&$9\bmod10$&8&8&36&$(10,7,4)$\\
N19a & $-$ &19&$1\bmod3$&7&5&4&$(11,17,5)$\\
N19b & $-$ &19&$2\bmod3$&11&40&126&$(13,2,1)$\\
N13a & $-$ &13&$9\bmod12$&8&2&6&$(3,11,1)$\\
N13b & $-$ &13&$4\bmod12$&9&9&10&$(9,4,12)$\\
N13c & $-$ &13&$8\bmod12$&3&4&10&$(5,5,12)$\\
C19 & $+$ &19&$2\bmod3$&11&19&126&$(5,1,16)$\\
\bottomrule
\end{tabular}
\caption{Each row certifies all $k=k_0+th$, $t\ge0$, on its stated clock.}
\label{R12:tab:returns}
\end{table}

For a transparent reproduction of the finite identities, multiply
two coefficient triples by convolution to obtain $z_0,\ldots,z_4$.
For $j=4,3$, replace $z_jx^j$ using $x^3=ax^2+4x-4$:
\begin{equation}\label{R12:eq:reduction}
 z_{j-1}\leftarrow z_{j-1}+az_j,\qquad
 z_{j-2}\leftarrow z_{j-2}+4z_j,\qquad
 z_{j-3}\leftarrow z_{j-3}-4z_j.
\end{equation}
Reduce the first three coefficients modulo $p$, iterate the indicated
map from $x$, and compare its states at $k_0$ and $k_0+h$.
This requires at most $166$ cubing steps for any row. The appendix
provides this exact algorithm and a separate matrix-exponentiation
checker for the same identities. Thus the infinite conclusions rest
on finite algebraic equalities and Lemma~\ref{R12:lem:return}, not on
extrapolating a sample of clock terms.

The polynomial discriminant is
\begin{equation}\label{R12:eq:disc}
 \Delta(a)=16(a^3+a^2+18a-11).
\end{equation}
It is nonzero modulo $p$ at every table row, and $p\nmid a$.
These primes are therefore unramified in the cubic number field:
the field discriminant divides the polynomial discriminant.
In particular the table does not use coefficient-divisor witnesses
or primes of repeated reduction.

\subsubsection{Why the prime five rows have a simple explanation}

For $a=2\pmod5$,
\[
 f=(X+1)(X^2+2X+4).
\]
The quadratic roots are $2\zeta,2\zeta^2$, with
$\zeta^3=1$ and $\zeta\ne1$. For $n=3^k$, $k\ge1$, the trace is
$-1+2^{n+1}$ modulo five, which vanishes when $k$ is odd. This is C1.

For $a=3,4\pmod5$, the cubics have no root in $\F_5$ and are
irreducible. Their roots lie in $\F_{125}$ and have norm $-4=1$,
so belong to the cyclic norm-one group of order $31$. Each root is
nonidentity, hence has order $31$. This group contains a trace-zero
element: any root of the irreducible polynomial $X^3+2X+4$ supplies
one. Since $3$ generates $(\Z/31\Z)^*$, the classical exponent orbit
reaches every nonzero exponent residue. The negative exponent orbit
reaches every residue except $-1$. The traces at exponents $0$ and
$-1$ are $3$ and $1$, respectively, so neither can be the trace-zero
exponent. Both clocks therefore reach a trace-zero group element.
The table records explicit indices and returns for this mechanism.

\subsubsection{The coefficient classes are covered pointwise}

\begin{proof}[Proof of Theorem~\ref{R12:thm:A}]
C1--C3 cover every $4\nmid m$ on the classical clock, and C19 adds
$m\equiv2\pmod3$. Their union is precisely the classes
$m\not\equiv0,4\pmod{12}$.

For the negative clock, N2 and N3 cover $m\equiv2,3\pmod4$.
Suppose $m\equiv1\pmod4$. If $m\bmod10$ is in $\{1,5,7,9\}$,
one of N11a--N11d applies. Otherwise $m\equiv3\pmod{10}$, so
$m\equiv13\pmod{20}$. If $3\nmid m$, use N19a or N19b.
If $3\mid m$, then $m\equiv33\pmod{60}$ and hence
$m\equiv9\pmod{12}$, covered by N13a. This covers every odd $m$.
Finally, $4\mid m$ and $12\nmid m$ give $m\equiv4,8\pmod{12}$,
covered by N13b or N13c.

Lemma~\ref{R12:lem:return} supplies the index progression for each row.
Equations~\eqref{R12:eq:floor} and \eqref{R12:eq:large} turn its divisor into
actual compositeness at every indicated index. No matching of the
two clocks' offsets is needed or asserted.
\end{proof}

This is a complete argument on the stated coefficient-exponent
classes, not a density transfer. Coefficients outside the table's
guarantee may still have infinitely many composite floors.

\subsection{Every fixed finite set of primes misses infinitely many members}

The prime $7$ divides the coefficient, but never a scheduled trace.
At $a=0\pmod7$, the two finite-algebra orbits have the following
exact states:
\[
\begin{array}{c|c|c|c}
 \text{clock}&u_1&u_2&(\tau(u_1),\tau(u_2))\\\hline
 +&(3,4,0)&(2,0,4)&(2,3)\\
 -&(5,2,3)&(3,3,6)&(4,1)
\end{array}
\]
For each clock, its map sends $u_2$ back to $u_1$.
Thus the indicated nonzero trace residues alternate for all $k\ge1$.
This two-cycle is checked directly by \eqref{R12:eq:reduction}.

\begin{proof}[Proof of Theorem~\ref{R12:thm:B}]
Given finite $P$, set
\begin{equation}\label{R12:eq:L}
 L=\lcm\{\ord_p(7):p\in P,\ p\ne7\},
\end{equation}
with the empty least common multiple equal to one. For every positive
multiple $m$ of $L$ and every $p\in P\setminus\{7\}$, we have
$7^m\equiv1\pmod p$. The integral factorization
\[
 X^3-X^2-4X+4=(X-1)(X^2-4)
\]
gives, at every positive odd exponent,
\begin{equation}\label{R12:eq:a1}
 T_n(1)=1+2^n+(-2)^n=1.
\end{equation}
The trace polynomial therefore specializes to $1\pmod p$.
This also holds at primes where the factors meet; it is an integral
identity and does not assume separability. The separate two-cycle
excludes $p=7$. Equation~\eqref{R12:eq:floor} identifies all scheduled
floors with these traces. Taking arbitrarily large multiples $m$ of
$L$, in particular $m\ge15$, proves the conclusion and the strict
bounds through Lemma~\ref{R12:lem:roots}.
\end{proof}

\begin{corollary}[A quantitative fixed-prime limitation]\label{R12:cor:bound}
Suppose a finite prime set $P$ supplies a scheduled trace divisor
for every integer $m\in[M_0,M]$, where $M\ge M_0\ge2$.
Then the $L$ in \eqref{R12:eq:L} satisfies
\[
 L>M-M_0+1,\qquad
 \prod_{\substack{p\in P\\p\ne7}}(p-1)>M-M_0+1.
\]
\end{corollary}
\begin{proof}
If $L\le M-M_0+1$, the first multiple of $L$ at least $M_0$ is
at most $M_0+L-1\le M$. The proof of Theorem~\ref{R12:thm:B} shows that
this member has no such divisor. Finally
$L\le\prod_{p\in P,\,p\ne7}(p-1)$.
\end{proof}

The limitation concerns fixed-prime covers, not every possible
algorithm or Diophantine proof. In particular it is compatible with
a prime witness that changes with $m$.

\subsection{Exact classical three adic persistence}

Let $v_3$ denote the usual $3$-adic valuation of a nonzero rational
number, with $v_3(0)=+\infty$ when needed for coefficient bounds.

\begin{theorem}\label{R12:thm:valuation}
For every $m\ge1$ and $k\ge1$,
\begin{equation}\label{R12:eq:valuation}
 v_3\bigl(T_{3^k}(7^m)-1\bigr)=k+1+v_3(m).
\end{equation}
More generally, for every integer $a\equiv1\pmod3$, $a\ne1$,
\[
 v_3\bigl(T_{3^k}(a)-1\bigr)=k+v_3(a-1).
\]
\end{theorem}
\begin{proof}
Put $a=1+h$ and $D(z)=(1-z)(1-4z^2)$. The formal trace identity
\[
 \sum_{n\ge1}\frac{T_n(a)}n z^n
 =-\log(1-az-4z^2+4z^3)
\]
follows by factoring into the three root factors and expanding the
formal logarithms. Subtracting the identity at $a=1$ gives
\[
 \sum_{n\ge1}\frac{T_n(1+h)-T_n(1)}n z^n
 =\sum_{j\ge1}\frac{h^jz^j}{jD(z)^j}.
\]
Consequently the coefficient $A_{n,j}$ of $h^j$ is exactly
\begin{equation}\label{R12:eq:coeff}
 A_{n,j}=\frac nj[z^{n-j}]D(z)^{-j}.
\end{equation}
The bracket is an integer. When $n=3^k$, we have
$v_3(A_{n,j})\ge k-v_3(j)$. For $j=1$,
\[
 A_{n,1}=n\sum_{r=0}^{(n-1)/2}4^r
 =n\frac{4^{(n+1)/2}-1}{3}.
\]
The geometric sum is a $3$-adic unit: modulo three it is
$(n+1)/2\not\equiv0$. Thus $v_3(A_{n,1})=k$.

Writing $b=v_3(h)\ge1$, every term with $j\ge2$ satisfies
\[
 v_3(A_{n,j}h^j)\ge k-v_3(j)+jb>k+b,
\]
since $v_3(j)<j-1$. The linear term is therefore the unique term
of least valuation and cannot cancel. Equation~\eqref{R12:eq:a1} gives
$T_n(1)=1$, proving the general assertion.

Finally $v_3(7^m-1)=1+v_3(m)$. For completeness, if $3\nmid r$,
the binomial expansion gives $v_3((1+6)^r-1)=1$; each cubing of a
number $u\equiv1\pmod3$ adds exactly one, since
$u^2+u+1=3+3(u-1)+(u-1)^2$ has valuation one. Writing
$m=3^tr$ proves the formula and hence \eqref{R12:eq:valuation}.
\end{proof}

This describes congruence persistence near trace one, not
compositeness. Its negative-clock analogue is false, even though
the first derivative alone has the suggestive valuation.
Indeed
\[
 T_5(1+h)-1=105h+50h^2+30h^3+5h^4+h^5.
\]
At $m=4$, $a=2401$, one obtains exactly
\[
 T_5(2401)-1=79792543008252000\equiv54\pmod{81}.
\]
The valuation is three, whereas $k+1+v_3(m)=2$ for $k=1$.
Higher Taylor terms can cancel the leading term. No general
negative-clock valuation formula is claimed.

\subsection{The precise remaining adaptive problem}

A coefficient-divisor criterion requiring a prime $p\equiv2\pmod3$
with $p\mid a$ cannot apply to $a=7^m$, whose only prime divisor is
$7\equiv1\pmod3$. More generally, constructing useful coefficient
residue classes does not automatically place a power of seven in them.
Two concrete specializations illustrate this obstruction:
\[
 \ord_{31}(7)=15,\quad12^{15}=-1\pmod{31};\qquad
 \ord_{53}(7)=26,\quad26^{26}=-1\pmod{53}.
\]
Thus $a\equiv12\pmod{31}$ and $a\equiv26\pmod{53}$ are both disjoint
from powers of seven. This only describes those residue classes;
it does not say those primes are useless at all other coefficients.

For a fixed candidate and an odd prime $p$, the matrix in
\eqref{R12:eq:matrix} is invertible. If its order is $E$, a recurring
divisor certificate asks whether
\[
 \{e\bmod E:\Tr(M_{p,a}^e)=0\}
\]
intersects the eventual orbit of the desired exponent clock modulo
$E$. The finite table proves particular intersections. To cover all
remaining $m$ by this strategy would require a suitable prime
\emph{for each specified $m$}; no such universal statement is proved.
Infinitely many composite terms need not logically come from one
prime recurring, so this is a sufficient strategy, not an equivalent
reformulation of the entire compositeness problem.

If a splitting or Chebotarev argument instead supplies a useful
nonzero residue $r\pmod p$, transferring it to the thin family first
requires
\[
 r\in\langle7\rangle\subseteq\F_p^*.
\]
With $d=\ord_p(7)$, this is exactly $r^d=1\pmod p$.
When true, a discrete logarithm supplies one exponent class
$m\equiv m_0\pmod d$, not coverage of every preassigned $m$.
A splitting or quadratic-character condition alone does not imply
this membership. An unproved primitive-root assertion for seven
would be an additional assumption; even full subgroup membership
would not force the residue to equal $7^m$ for a given $m$.

There is also a direct obstruction to simply taking many prime
factors of $7^m-1$: \eqref{R12:eq:a1} makes every odd trace congruent to
one modulo each such prime. Theorem~\ref{R12:thm:B} similarly explains
why enlarging a fixed finite table cannot close the entire family.
The classes not covered by this initial table include $12\mid m$. Adaptive certificates in Sections~\ref{src:R14}--\ref{src:R16} handle some of them, including $m=12,60$, but not all. The results here make no claim of all-$m$
closure, of a distinguished least prime-representing base, or of
a new Mills-type transcendence theorem.


\subsection{Portable verification}

Two standard-library Python programs accompany the paper:
\begin{verbatim}
python3 certificates/prime_power_coefficient_verify.py
python3 certificates/audit_prime_power_coefficient_verify.py
\end{verbatim}
The first uses the quotient-algebra arithmetic \eqref{R12:eq:reduction}.
The second uses binary exponentiation of the $3\times3$ matrices
\eqref{R12:eq:matrix}, with no import of the first program's arithmetic.
Both verify all fifteen returns, the coefficient classes, nonzero
discriminants, and the nonzero prime-seven cycles. The complete
coefficient-class comparison is performed on $\Z/60\Z$, since every
table modulus divides sixty. It is not a sample of large $m$.

The matrix program also verifies exact matrix orders. In table order
they are
\[
\begin{gathered}
12,31,31,31,31,10,1330,120,1330,\\
90,6858,168,549,549,6858.
\end{gathered}
\]
The starting indices satisfy $k_0\ge v_3(E)$ in every row, and the
two original exponents differ by a multiple of $E$. These checks
confirm the handling of possible initial transients. Supplemental
finite valuation tests support the examples; the universal valuation
statement is established by Theorem~\ref{R12:thm:valuation}.

Each program has a corresponding \texttt{.json} result file. The README
specifies execution and scope. The finite identities and the mathematical
proofs together establish the stated infinite conclusions.

\endgroup

\section{A prime seed outlier and certified returns}\label{src:R13}
The following example shows that prime leading coefficients and several prime seeds can coexist with later recurring composites. It supplies exact prime certificates rather than treating probable-prime output as proof.
\medskip
\begingroup




\subsection{One coefficient and the two exponent clocks}

For an integer $a>8$, put
\[
 f_a(X)=X^3-aX^2-4X+4
\]
and let $\beta_a$ be its largest real root.
The two clocks have exponents $3^k$ and $2\cdot3^k-1$, with
$k\ge1$; their first exponents are $3$ and $5$.
We call $a=\floor{\beta_a}$ the leading coefficient or first floor.

\begin{theorem}\label{R13:thm:point}
Let $a=100981$ and $\beta=\beta_a$.
The polynomial $f_a$ is irreducible and totally real, with roots
$-\rho,\sigma,\beta$ satisfying
\begin{equation}\label{R13:eq:rootsbounds}
 0<\sigma<\rho<\beta^{-17/40}<1,\qquad
 a<\beta<a+1,\qquad 4<\beta^{3/20},\qquad
 \rho<\sigma^{17/23}.
\end{equation}
Moreover, $\Norm(\beta)=-4$ and $\Tr(\beta^{-1})=1$.
The three integers
\begin{equation}\label{R13:eq:threeprimes}
 \begin{aligned}
 \floor\beta&=100981,\\
 \floor{\beta^3}&=1029719653587901,\\
 \floor{\beta^5}&=10500218502186178331982901
 \end{aligned}
\end{equation}
are primes congruent to $1$ modulo $9$.
Both $\floor{\beta^9}$ and $\floor{\beta^{17}}$ are composite.
There are also the separate infinite divisibility certificates
\begin{equation}\label{R13:eq:separate}
 \begin{array}{rlll}
 13&\mid\floor{\beta^{3^k}},
     &k=6t,&t\ge1,\\
 41&\mid\floor{\beta^{\,2\cdot3^k-1}},
     &k=23+3444t,&t\ge0.
 \end{array}
\end{equation}
Every multiple in \eqref{R13:eq:separate} is proper, so every indicated
floor is composite.
\end{theorem}

The fixed-seed criterion under discussion requires a prime
$p\equiv2\pmod9$ dividing the trace at exponent $3$ or $5$,
possibly with further finite prime exclusions.
Neither seed in \eqref{R13:eq:threeprimes} has such a divisor.
This is a point outside that sufficient criterion, not a counterexample
to a density-one conclusion.
In particular, the obstruction is not merely that the first floor
was composite: all three displayed initial values are proved prime.

\subsection{Roots and exact floors}

For every integer $a>8$, direct substitution gives
\[
 \begin{gathered}
 f_a(-1)=7-a<0,\qquad f_a(0)=4>0,\qquad f_a(1)=1-a<0,\\
 f_a(a)=4(1-a)<0,\qquad f_a(a+1)=(a+1)^2-4a>0.
 \end{gathered}
\]
There is one root in each of $(-1,0)$, $(0,1)$ and $(a,a+1)$,
and these exhaust the three roots. Write them as $-\rho,\sigma,\beta_a$.
Their sum gives $\sigma-\rho=a-\beta_a<0$, hence
$0<\sigma<\rho<1$.
None is an integer. A reducible monic cubic over $\Q$ would have an
integer root, so $f_a$ is irreducible.
Thus $\beta_a$ is a totally real cubic Pisot number.
The elementary symmetric functions give
\begin{equation}\label{R13:eq:norm}
 \Norm(\beta_a)=-4,\qquad
 \Tr(\beta_a^{-1})=\frac{-4}{-4}=1,\qquad
 \rho\sigma=\frac4{\beta_a}.
\end{equation}

Define the integral trace polynomials by
\begin{equation}\label{R13:eq:recurrence}
 \begin{gathered}
 T_0(a)=3,\qquad T_1(a)=a,\qquad T_2(a)=a^2+8,\\
 T_n(a)=aT_{n-1}(a)+4T_{n-2}(a)-4T_{n-3}(a)\quad(n\ge3).
 \end{gathered}
\end{equation}
Newton's identities and the equation for each root show that
$T_n(a)=\beta_a^n+(-\rho)^n+\sigma^n$.
For every positive odd $n$,
\begin{equation}\label{R13:eq:floor}
 0<\beta_a^n-T_n(a)=\rho^n-\sigma^n<1,
 \qquad \floor{\beta_a^n}=T_n(a).
\end{equation}
This identity holds from the first odd exponent, not just eventually.

For $a=100981$, exact rational evaluation gives
\[
 f_a(-1/150)=-\frac{1557151}{3375000}<0.
\]
The unique negative root therefore lies in $(-1/150,0)$, so
$\rho<1/150$. The exact integer inequalities
\begin{equation}\label{R13:eq:integerbounds}
 150^{40}>100982^{17},\qquad 4^{20}<100981^3
\end{equation}
imply $\rho<\beta^{-17/40}$ and $4<\beta^{3/20}$.
The same small-conjugate bound holds for $\sigma<\rho$.
Finally, \eqref{R13:eq:norm} gives
\[
 \rho^{40}<\beta^{-17}=(\rho\sigma/4)^{17},
 \qquad \rho^{23}<\sigma^{17}/4^{17}<\sigma^{17},
\]
which proves the last inequality in \eqref{R13:eq:rootsbounds}.
The coefficient is odd, is $1$ modulo $3$, and is coprime to $4$.

\subsection{Certified primality of the first seeds}

We use the following complete-factor form of the Lucas criterion.
Its proof is included to specify exactly what the certificates establish.

\begin{lemma}\label{R13:lem:lucas}
Let $N>2$ and suppose the complete prime factorization of $N-1$ is
known. If an integer $b$ satisfies
\begin{equation}\label{R13:eq:lucas}
 b^{N-1}\equiv1\pmod N,\qquad
 \gcd\bigl(b^{(N-1)/r}-1,N\bigr)=1
 \quad\text{for every prime }r\mid N-1,
\end{equation}
then $N$ is prime.
\end{lemma}

\begin{proof}
Let $\ell$ be any prime divisor of $N$.
The first congruence makes $b$ invertible modulo $\ell$ and gives
$d=\ord_\ell(b)\mid N-1$.
If $r^e$ is the exact power of $r$ dividing $N-1$, the gcd condition
implies $b^{(N-1)/r}\not\equiv1\pmod\ell$.
Consequently $r^e\mid d$. This holds for every $r$, so $d=N-1$.
As $d\mid\ell-1$, we have $\ell\ge N$.
Since $\ell\mid N$, it follows that $\ell=N$ and $N$ is prime.
\end{proof}

Write $N_1,N_3,N_5$ for the three integers in \eqref{R13:eq:threeprimes}.
Their complete factorizations and certificate bases are
\begin{equation}\label{R13:eq:mainfactors}
 \begin{aligned}
 N_1-1&=2^2\cdot3^3\cdot5\cdot11\cdot17,
       &&b_1=7,\\
 N_3-1&=2^2\cdot3^4\cdot5^2\cdot11\cdot17
          \cdot131\cdot227\cdot22861,
       &&b_3=6,\\
 N_5-1&=2^2\cdot3^4\cdot5^2\cdot11\cdot17
          \cdot809\cdot2154553\cdot3977095391,
       &&b_5=6.
 \end{aligned}
\end{equation}
For each row, exact modular powers and gcds satisfy \eqref{R13:eq:lucas}.
All prime factors on the right are themselves recursively certified
by the same criterion, terminating at the elementary prime $2$.
The complete certificate graph has $36$ nodes and also certifies
$13$, $41$ and $256603$.
The accompanying data and dependency-free acceptance programs are
described in the appendix; no probable-prime test is an acceptance step.

At $a=100981$, evaluation of \eqref{R13:eq:recurrence} proves that
$T_1(a)=N_1$, $T_3(a)=N_3$ and $T_5(a)=N_5$.
Each is $1$ modulo $9$.
Since $T_3(a)$ and $T_5(a)$ are themselves prime, neither has a
prime divisor congruent to $2$ modulo $9$.
Adding further exclusions to that witness criterion cannot detect
this point.

\subsection{The next floor on each clock}

In this section $a=100981$, and we abbreviate $T_n=T_n(a)$.
The recurrence gives
\begin{equation}\label{R13:eq:T9}
 T_9=1091834984307981235989043476507177498066859621.
\end{equation}
Its exact base-$2$ Fermat remainder is
\begin{equation}\label{R13:eq:fermat}
 2^{T_9-1}\bmod T_9
 =247023223208684492599063887120591198040833796\ne1.
\end{equation}
As $T_9>1$ is odd, Fermat's necessary condition for primality
proves it composite, without requiring its factorization.

For longer integers, a bracket $[\cdots]_{10}$ below means the
single decimal integer obtained by concatenating its digit rows
from top to bottom; no digit is omitted.
Exact recurrence evaluation gives
\begin{equation}\label{R13:eq:T17}
 T_{17}=\decrows{1180522704757521266282466391894926537172}
                    {5694973373598468435215524879359555088042987461}.
\end{equation}
There are the proper factorizations
\begin{equation}\label{R13:eq:factor17a}
 T_{17}=229247\cdot
 \decrows{5149566645397851515101468686154787356748}
         {7011709525526913919115734903224709976763}
\end{equation}
and
\begin{equation}\label{R13:eq:factor17b}
 T_{17}=256603\cdot
 \decrows{4600580292348574515038664364387503408660}
         {7307682971744166807151611163390744021087}.
\end{equation}
Each divisor and quotient exceeds $1$.
The first factorization alone proves compositeness; no primality
assertion about $229247$ is needed.
The second uses the certified prime $256603\equiv4\pmod9$,
which will also give a later-seed return.
By \eqref{R13:eq:floor}, these computations prove compositeness of
both clock floors at the common index $k=2$.

\subsection{Exact modular returns}

For a prime $p\nmid4$, let
\[
 R_p=\F_p[X]/(f_a),\qquad x=X\bmod f_a.
\]
The element $x$ is a unit. Multiplication by $x$ in the basis
$1,x,x^2$ has matrix
\[
 M_a=\begin{pmatrix}0&0&-4\\1&0&4\\0&1&a\end{pmatrix}.
\]
Consequently the multiplication trace of $x^n$ is $T_n(a)$ modulo $p$,
even when the quotient is not reduced. More generally,
\begin{equation}\label{R13:eq:ringtrace}
 \Tr_{R_p/\F_p}(v_0+v_1x+v_2x^2)
       =3v_0+av_1+(a^2+8)v_2.
\end{equation}

\subsubsection{The classical clock modulo 13}

Here $a\equiv10\pmod{13}$ and
\[
 f_a(X)=(X+5)(X^2-2X+6)\quad\text{in }\F_{13}[X].
\]
Exact reduction by $x^3=ax^2+4x-4$ gives
\begin{equation}\label{R13:eq:13data}
 x^{168}=1,\qquad \ord_{56}(3)=6,\qquad
 x^{3^6}=3+8x+7x^2.
\end{equation}
The last state's trace is $9+80+4\cdot7=117\equiv0\pmod{13}$,
since $a^2+8\equiv4\pmod{13}$.
For $k=6t$, $t\ge1$, the exponents $3^k$ and $3^6$ agree
modulo $56$ and are both $0$ modulo $3$.
They therefore agree modulo $168$, proving
$13\mid T_{3^k}(a)$ for every such $k$.
The range starts at $k=6$, not at $k=0$.

\subsubsection{The negative clock modulo 41}

Here $a\equiv39\pmod{41}$ and
$f_a=X^3+2X^2-4X+4$ modulo $41$.
Evaluation at all $41$ residues gives no root, so this cubic is
irreducible and $R_{41}=\F_{41^3}$.
In particular,
\begin{equation}\label{R13:eq:41data}
 x^{68920}=1,\qquad \ord_{68920}(3)=3444.
\end{equation}
Starting at $u_0=x$, the exact recurrence
\[
 u_{k+1}=x^2u_k^3
\]
represents $u_k=x^{2\cdot3^k-1}$.
Reduction in $R_{41}$ gives
\[
 u_{23}=18x+3x^2.
\]
As $a^2+8\equiv12\pmod{41}$, its trace is
$39\cdot18+12\cdot3=738=18\cdot41$.
Equation~\eqref{R13:eq:41data} returns this state for every
$k=23+3444t$, $t\ge0$.
Thus $41\mid T_{2\cdot3^k-1}(a)$ on the asserted progression.
The additional identity $x^{8615}=1$ yields the same clock period,
but is not needed for this argument.

On both displayed progressions, the odd exponent $n$ is at least $3$,
and
\begin{equation}\label{R13:eq:proper}
 T_n(a)>\beta_a^n-1>a^3-1>41.
\end{equation}
Equations~\eqref{R13:eq:floor} and \eqref{R13:eq:proper} turn the trace
divisibilities into actual composite floors, completing the proof
of Theorem~\ref{R13:thm:point}.

The classical progression is $0$ modulo $6$.
Since $3444\equiv0\pmod6$, the displayed negative progression is
$23\equiv5\pmod6$. These two \emph{displayed progressions} are disjoint;
the certificates do not supply a common-index progression.
The full sets of composite indices are not disjoint: The earlier direct floor calculation
already proves that both contain $k=2$.

\subsubsection{A complementary later negative seed}

For completeness, the prime $p=256603$ in \eqref{R13:eq:factor17b}
gives another negative-clock certificate.
We justify the universal exponent used for this return.
If $p>2$, $g$ is a monic cubic with nonzero constant term over $\F_p$,
and $x=X\bmod g$, then
\begin{equation}\label{R13:eq:universal}
 x^{E_p}=1,\qquad E_p=p(p^2-1)(p^2+p+1).
\end{equation}
Indeed, factor $g=\prod h_i^{e_i}$ into distinct monic irreducibles.
In the reduced quotient the factors are fields of degrees at most $3$.
Each possible multiplicative order divides
$H=(p^2-1)(p^2+p+1)$, since $p^d-1\mid H$ for $d=1,2,3$.
Hence $x^H=1+z$ with $z^3=0$ in the original quotient.
In characteristic $p\ge3$, $(1+z)^p=1$, proving \eqref{R13:eq:universal}
also for nonreduced quotients.

For our prime $p\equiv4\pmod9$,
$v_3(p^2-1)=v_3(p^2+p+1)=1$, so $v_3(E_p)=2$.
Set $m=E_p/9$ and $L=\ord_m(3)$.
If $k\ge2$ and $k\equiv2\pmod L$, then
$3^k\equiv9\pmod m$, while both sides are $0$ modulo $9$.
Thus $2\cdot3^k-1\equiv17\pmod{E_p}$.
The trace at every such exponent is divisible by $256603$;
\eqref{R13:eq:factor17b} proves properness at $k=2$, and larger
exponents give still larger floors.

This is the offset-two, exponent-$17$ negative-seed mechanism.
It complements the fixed exponent-$3/5$ test at this point.
It makes no assertion about membership in a complete simultaneous
four-way certificate set involving other signs or clocks.
The much smaller prime $41$ supplies the displayed negative
progression in Theorem~\ref{R13:thm:point}; no period optimization
is needed for the $256603$ certificate.

\subsection{Coefficient families and finite prime filters}

\begin{proposition}\label{R13:prop:AP}
Every coefficient
\begin{equation}\label{R13:eq:coefficientAP}
 a=1843+3198t,\qquad t\ge0,
\end{equation}
has the exact all-odd floor identity \eqref{R13:eq:floor} and both
separate composite-floor progressions in \eqref{R13:eq:separate}, with
$\beta$ replaced by $\beta_a$.
There are infinitely many prime coefficients in this progression.
\end{proposition}

\begin{proof}
Here $3198=6\cdot13\cdot41=\lcm(6,13,41)$.
Every coefficient in \eqref{R13:eq:coefficientAP} is $1$ modulo $6$,
$10$ modulo $13$, $39$ modulo $41$, and larger than $8$.
Thus the root proof, both quotient-ring calculations and the
properness estimate remain valid.
As $\gcd(1843,3198)=1$, the prime-progression theorem supplies
infinitely many primes in this class; equivalently, use the
cyclotomic case of Chebotarev \cite[Theorem~3.1 and p.~4]{Lenstra}.
\end{proof}

Our explicit point is the member $t=31$.
Primality of both first clock seeds is proved here only for
$a=100981$, not throughout \eqref{R13:eq:coefficientAP}.
All sufficiently large members satisfy the same strict inequalities,
as follows from the estimate below.

\begin{lemma}[Prime-leading escape from each finite filter]
\label{R13:lem:finite}
For every fixed finite set $P$ of primes, there are infinitely many
prime integers $a>8$ such that $\beta_a$ satisfies
\eqref{R13:eq:rootsbounds} and
\begin{equation}\label{R13:eq:filter}
 \floor{\beta_a^n}\equiv1\pmod p
 \quad\text{for every }p\in P\text{ and every positive odd }n.
\end{equation}
These coefficients can be chosen arbitrarily large.
\end{lemma}

\begin{proof}
Put $M=\lcm(24,\prod_{p\in P}p)$, with empty product $1$.
Take primes $a\equiv1\pmod M$ above the threshold specified below.
There are infinitely many by the same reduced-residue-class theorem.
At the boundary coefficient $1$,
\[
 f_1(X)=(X-1)(X^2-4),\qquad
 T_n(1)=1+2^n+(-2)^n=1\quad(n\text{ odd}).
\]
Each $T_n(a)$ is an integer polynomial in $a$, by
\eqref{R13:eq:recurrence}. Therefore $T_n(a)\equiv T_n(1)=1\pmod p$
for every $p\in P$, proving \eqref{R13:eq:filter} by \eqref{R13:eq:floor}.
This polynomial congruence does not require separability modulo $p$.

For an explicit sufficient cutoff, the negative root equation gives
\[
 a\rho^2=4(\rho+1)-\rho^3<8,
 \qquad \rho<\sqrt{8/a},\qquad \beta_a<a+1<2a.
\]
Consequently, if $a\ge2^{27}$,
\[
 \rho\beta_a^{17/40}
 <2^{77/40}a^{-3/40}\le2^{-1/10}<1,
 \qquad 4^{20}<a^3<\beta_a^3.
\]
These prove the strict bounds for $\rho$ and the norm magnitude;
$\sigma<\rho$ and the argument following \eqref{R13:eq:integerbounds}
give the remaining inequalities in \eqref{R13:eq:rootsbounds}.
The congruence $a\equiv1\pmod{24}$ also gives oddness and
$a\equiv1\pmod3$.
\end{proof}

\begin{corollary}\label{R13:cor:GT}
Fix a finite prime set $P$, a length $\ell\ge2$, and a lower bound $A$.
There is a nonconstant arithmetic progression of $\ell$ prime
coefficients $a>A$, each satisfying the conclusion of
Lemma~\ref{R13:lem:finite} for this same $P$.
\end{corollary}

\begin{proof}
The prime coefficients in $1\pmod M$ above
$\max\{A,2^{27}\}$ have relative density $1/\varphi(M)>0$
among all primes, by the cyclotomic prime-progression theorem
\cite{Lenstra}. Apply the relative-prime Green--Tao theorem
\cite[Theorem~1.2]{GreenTao} to this set.
\end{proof}

These finite arithmetic progressions are in the integer coefficients,
not necessarily in their algebraic roots.
The quantifiers are
\[
 \text{for every finite }P,\quad
 \text{there exist infinitely many prime coefficients escaping }P.
\]
They do not produce one coefficient escaping every prime.
Nor is simultaneous primality of $T_3(a)$ and $T_5(a)$ asserted
for these families.

\subsubsection{A common progression for the prime-seed outlier}
\label{R13:syn:outlier-common}

The coefficient with prime leading floor and prime first seeds on both
clocks also admits a small-prime simultaneous certificate. The following
result does not claim that the earlier displayed separate progressions
intersect; it supplies a different negative-clock witness.

\begin{theorem}[Synchronization of the outlier]
\label{R13:syn:outlier-synchronization}
For every integer $a=504+767u$, $u\ge0$, let $\beta_a$ be the largest
root of $X^3-aX^2-4X+4$. Both floors
\[
 \left\lfloor\beta_a^{3^k}\right\rfloor,
 \qquad \left\lfloor\beta_a^{2\cdot3^k-1}\right\rfloor
\]
are composite at every $k=30+1770t$, $t\ge0$.
The first is divisible by $13$, and the second by $59$.
In particular this holds for the prime-seed outlier $a=100981$.
\end{theorem}
\begin{proof}
Every such $a$ satisfies $a\equiv10\pmod{13}$ and
$a\equiv32\pmod{59}$; also $a\ge504>59$.
In each cubic quotient write $x$ for the class of $X$ and use
\[
 \tau(c_0+c_1x+c_2x^2)=3c_0+ac_1+(a^2+8)c_2.
\]
This multiplication trace equals the reduced Newton trace at every
power. The relation is $x^3=ax^2+4x-4$.
Exact binary powering supplies the finite identities
\[
\begin{array}{c|c|c}
 \text{modulus}&\text{scalar identity}&\text{zero-trace seed}\\ \hline
 13&x^{168}=1&x^{3^6}=3+8x+7x^2\\
 59&x^{3541}=55&x^{2\cdot3^{30}-1}=41+45x+58x^2.
\end{array}
\]
Their seed traces are respectively
\[
 3\cdot3+10\cdot8+4\cdot7=117=9\cdot13,
 \qquad
 3\cdot41+32\cdot45+1032\cdot58=61419=1041\cdot59.
\]
The elementary congruences
$3^6\equiv1\pmod{56}$ and $3^{1770}\equiv1\pmod{3541}$
complete the certificates. If $6\mid k$ and $k\ge6$, then
$3^k\equiv3^6\pmod{168}$: both are zero modulo $3$ and one
modulo $56$. Thus $13$ divides every indicated classical trace.
Replacing $k=30$ by $30+1770t$ changes the negative-clock exponent
by a multiple of $3541$. The second scalar identity then preserves
zero trace modulo $59$. Since $6\mid30$ and $6\mid1770$, both
certificates hold on the stated common progression.

The real cubic root calculation gives $a<\beta_a<a+1$ and
$\lfloor\beta_a^n\rfloor=T_n(a)$ at every positive odd $n$.
All indicated exponents are odd and at least three, so these floors
exceed $a^3-1>59$. Both divisors are therefore proper, including at
the first displayed index. Finally
$100981=504+767\cdot131$, proving the specialization.
\end{proof}

The coefficient progression is reduced, since $\gcd(504,767)=1$,
so it contains infinitely many prime coefficients. This assertion
concerns the coefficient, not simultaneous primality of the initial
seed traces throughout the progression. For odd coefficients it can
equivalently be written $a=1271+1534v$, $v\ge0$.

\subsection{Scope of the example}

The explicit point separates a sufficient fixed-seed test from a
universal pointwise test while retaining prime first floor and
prime first clock seeds.
Its later composite floors exclude it from an all-prime floor
sequence. The separate return certificates also persist on a
reduced coefficient progression, but their displayed index offsets
do not synchronize. The new witness in Section~\ref{R13:syn:outlier-common} instead yields a common progression for the same outlier and an entire coefficient class.
The fixed-filter lemma shows that requiring a prime leading
coefficient does not make any fixed finite divisibility filter universal.
None of these conclusions excludes every remaining candidate,
establishes a global minimum, or resolves a universal Mills-type
or Pisot prime-floor problem.


\subsection{Portable exact certificates}

The directory \texttt{certificates/} contains
\path{prime_seed_lucas_certificates.json}, the full $36$-node recursive
primality graph, and two independent acceptance programs.
Run from the paper directory, using Python~3 with no third-party dependency:
\begin{lstlisting}[language=bash]
python3 certificates/prime_seed_outlier_verify.py
python3 certificates/audit_prime_seed_outlier_verify.py
(cd certificates && sha256sum -c SHA256SUMS.txt)
\end{lstlisting}
Each program has a corresponding \texttt{.json} report and
\texttt{.log} file.
The primality checks multiply each complete factorization to $N-1$,
verify every recursive dependency, and check the modular powers and
gcd conditions of Lemma~\ref{R13:lem:lucas}.
Every dependency is smaller than its parent, so induction from $2$
proves all $36$ asserted prime nodes.

The first verifier recomputes the traces, rational endpoint,
strict integer inequalities, Fermat remainder, proper factorizations,
quotient-ring states, multiplicative orders and coefficient congruences.
The second uses direct $3$-by-$3$ companion-matrix powers for the
modular return checks, rather than importing the first program's
quotient-ring implementation. It also checks the exact matrix orders
$168$ modulo $13$ and $8615$ modulo $41$.
The data at one return state, together with the proved exponent
congruences, establish infinitely many indices; no finite sampling
is extrapolated to an infinite claim.
These exact checks do not replace the imported prime-density and
Green--Tao theorems.



\endgroup

\section{Adaptive returns and ancestral minimality}\label{src:R14}
Adaptive certificates extend the thin-family coverage, while the second half of this section analyzes the actual minimality condition. Skipping several levels enlarges a raw interval but does not preserve every ancestral prime cell automatically.
\medskip
\begingroup
\renewcommand{\PP}{\mathcal P}




\subsection{Scope and the two exponent clocks}

For a real number $B>1$, the classical and negative clocks in this
paper are respectively
\[
 \floor{B^{3^k}}\quad\text{and}\quad
 \floor{B^{\,2\cdot3^k-1}},\qquad k\ge1.
\]
The word ``negative'' only names the second exponent sequence; both
sequences have positive odd exponents. Their first exponents are $3$
and $5$.

The arithmetic part studies
\begin{equation}\label{R14:eq:family}
 f_a(X)=X^3-aX^2-4X+4,\qquad a=7^m,\quad m\ge2,
\end{equation}
and its dominant real root $\beta_a$. We prove a polynomial
factorization obstruction on the entire classical clock, a
consecutive-coprimality identity, and separate numerical return
certificates for two coefficients. In particular, polynomial
irreducibility is not being confused with primality of polynomial
values. The finite-ring certificates concern actual values and do
not supply a divisor-existence theorem for all $m$.

The minimum problem has a different quantifier. For a set $S$ of
positive integers, write
\[
 \mathcal W(S)=
 \{B>1:\floor{B^{3^k}}\in S\text{ for every }k\ge1\}.
\]
For the full set $\PP$ of primes, its least element is the classical
Mills constant, denoted here by $\xi$. We use Saito's conditional
Pisot reduction for this least element, explicitly as an external
input \cite[Theorem 1.7]{Saito}. We then give the necessary elementary
proofs of the strict gap exponent and the eventual least-prime
property before analyzing skipped levels.

The conclusions about altered sets are deliberately narrower than a
conclusion about $\mathcal W(\PP)$. Our thinning theorem changes the
allowed set $S$; its prime specialization is conditional on an
admissible prime tower with a sufficiently small gap exponent. The
unconditional example uses integers coprime to $6$, not primes.
No transfer of an ambient coefficient-density statement to the thin
sequence $7^m$ is used, and no all-$m$ exclusion or transcendence
conclusion is obtained.

\subsection{Exact floors and universal trace identities}

Let $P_n(A)$ be the sum of the $n$th powers of the three formal roots
of $X^3-AX^2-4X+4$. Newton's identities give
\begin{equation}\label{R14:eq:trace-recurrence}
 \begin{gathered}
 P_0=3,\qquad P_1=A,\qquad P_2=A^2+8,\\
 P_n=A P_{n-1}+4P_{n-2}-4P_{n-3}\qquad(n\ge3).
 \end{gathered}
\end{equation}
Thus $P_n\in\Z[A]$, and $P_n$ is monic of degree $n$ for every
$n\ge1$: the first two cases are explicit, and the first summand
in the recurrence supplies the unique term of degree $n$.

\begin{lemma}[Roots and odd floors]\label{R14:lem:odd-floors}
For every integer $a>8$, the polynomial $f_a$ is irreducible over
$\Q$ and has roots $-\rho,\sigma,\beta_a$ satisfying
\[
 0<\sigma<\rho<1,\qquad a<\beta_a<a+1.
\]
In particular, $\beta_a$ is a totally real cubic Pisot number, and
for every positive odd integer $n$,
\begin{equation}\label{R14:eq:exact-floor}
 \floor{\beta_a^n}=P_n(a),\qquad
 0<\beta_a^n-P_n(a)=\rho^n-\sigma^n<1.
\end{equation}
The floors in \eqref{R14:eq:exact-floor} strictly increase along the
positive odd exponents.
\end{lemma}

\begin{proof}
Direct substitution gives
\[
 \begin{gathered}
 f_a(-1)=7-a<0,\qquad f_a(0)=4>0,\qquad f_a(1)=1-a<0,\\
 f_a(a)=4(1-a)<0,\qquad
 f_a(a+1)=(a+1)^2-4a=(a-1)^2>0.
 \end{gathered}
\]
There is a root in each of $(-1,0)$, $(0,1)$, and $(a,a+1)$.
These exhaust the roots of the cubic. None is an integer, so the
rational-root test proves irreducibility. Their sum is $a$, whence
$\sigma-\rho=a-\beta_a<0$. For odd $n$ their power sum is
$P_n(a)=\beta_a^n-\rho^n+\sigma^n$. The displayed bounds on
$\rho$ and $\sigma$ prove \eqref{R14:eq:exact-floor}. Finally,
$\beta_a^{n+2}-\beta_a^n>1$ because $\beta_a>a>8$; taking
floors proves strict increase.
\end{proof}

The next identity will be used for both irreducibility and
coprimality. If the three formal roots are $\alpha_1,\alpha_2,\alpha_3$,
put
\[
 Q_n(A)=\sum_{i<j}(\alpha_i\alpha_j)^n.
\]
This symmetric expression lies in $\Z[A]$, by the fundamental
theorem on symmetric polynomials. Applying the cubic power-sum
identity to $\alpha_1^n,\alpha_2^n,\alpha_3^n$ gives
\begin{equation}\label{R14:eq:cubing}
 P_{3n}=P_n^3-3P_nQ_n+3(-4)^n,
\end{equation}
because $\alpha_1\alpha_2\alpha_3=-4$.

\begin{theorem}[All monomial substitutions on the classical clock]
\label{R14:thm:eisenstein}
For every $k,r\ge1$, the polynomial $P_{3^k}(X^r)$ is Eisenstein
at $3$, hence irreducible over $\Q$. Also $P_5(X^r)$ is Eisenstein
at $5$ for every $r\ge1$.
\end{theorem}

\begin{proof}
Reducing \eqref{R14:eq:cubing} modulo $3$ and starting from $P_1(A)=A$
shows that
\begin{equation}\label{R14:eq:frobenius}
 P_{3^k}(A)\equiv A^{3^k}\pmod3\qquad(k\ge0).
\end{equation}
Let $c_k=P_{3^k}(0)$. Then $c_0=0$, and $3\mid c_k$ for every
$k$. Specializing \eqref{R14:eq:cubing} at $A=0$, with $n=3^k$, yields
\[
 c_{k+1}=c_k^3-3c_kQ_{3^k}(0)+3(-4)^{3^k}
 \equiv6\pmod9.
\]
Indeed, the first term is divisible by $27$, the second by $9$,
and $(-4)^{3^k}\equiv-1\pmod3$. Thus every nonleading coefficient
of $P_{3^k}$ is divisible by $3$, whereas its constant coefficient
is not divisible by $9$. Substitution $A=X^r$ preserves the leading
and constant coefficients and only inserts zero coefficients.
Eisenstein's criterion applies.

For completeness, the irreducibility implication can be seen
without an additional factorization theorem. By Gauss's lemma, a
nontrivial rational factorization of this monic integer polynomial
would give a product of two positive-degree monic integer
polynomials. Modulo $3$ their product is a power of $X$, so both
reductions are positive-degree powers of $X$. Both constant
coefficients would be divisible by $3$, making their product
divisible by $9$, a contradiction.

The recurrence gives
\[
 P_5(A)=A^5+20A^3-20A^2+80A-80.
\]
Every nonleading coefficient is divisible by $5$ and the constant
is not divisible by $25$. The same substitution and Eisenstein
argument, at $5$, proves the last assertion.
\end{proof}

\begin{remark}[Exact scope of the factorization obstruction]
If $m=rt$, writing $7^m=(7^t)^r$ cannot lead to a nonconstant
factorization of $P_{3^k}(X^r)$ in $\Q[X]$. This holds for every
positive monomial substitution, not merely for the first classical
trace. It says nothing about whether a particular integer value is
prime, and it does not exclude congruence arguments, nonpolynomial
factorizations, or other Diophantine mechanisms. The same proof
works for the power sums of $X^3-AX^2-bX+c$, with fixed integers
$b,c$ and $3\nmid c$: the constant coefficient at a classical
exponent is $-3c\pmod9$.
\end{remark}

\begin{theorem}[Consecutive classical coprimality]\label{R14:thm:gcd}
For every $m\ge1$ and $k\ge0$,
\[
 \gcd\bigl(P_{3^k}(7^m),P_{3^{k+1}}(7^m)\bigr)=1.
\]
\end{theorem}

\begin{proof}
At an odd integer $a$, \eqref{R14:eq:trace-recurrence} modulo $2$ shows
that every positive trace is odd. At $a=7^m$, equation
\eqref{R14:eq:frobenius} gives $P_{3^k}(a)\equiv1\pmod3$.
For $n=3^k$, the identity \eqref{R14:eq:cubing} implies
\[
 \gcd(P_n(a),P_{3n}(a))\mid3\cdot4^n.
\]
The gcd is neither even nor divisible by $3$, and is therefore $1$.
\end{proof}

This theorem concerns consecutive scheduled traces, not all pairs
of traces. It prevents automatic propagation of a divisor to the
very next cubing step. The nonconsecutive returns below are fully
consistent with it.

\subsection{Exact adaptive returns for two coefficients}

\subsubsection{The finite-ring mechanism}

Let $d>1$ be an integer and set
\[
 R_d=(\Z/d\Z)[X]/(f_a),\qquad x=X\pmod{f_a}.
\]
A triple $(c_0,c_1,c_2)$ denotes $c_0+c_1x+c_2x^2$ in this ring.
The relation is $x^3=ax^2+4x-4$. Define
\begin{equation}\label{R14:eq:tau}
 \tau(c_0+c_1x+c_2x^2)
 =3c_0+ac_1+(a^2+8)c_2\pmod d.
\end{equation}
The sequences $\tau(x^n)$ and $P_n(a)\pmod d$ obey the same
recurrence and have the same first three values. Hence
\begin{equation}\label{R14:eq:trace-ring}
 \tau(x^n)=P_n(a)\pmod d\qquad(n\ge0).
\end{equation}
Equivalently, multiplication by $x$ has the companion matrix
\begin{equation}\label{R14:eq:companion}
 M=\begin{pmatrix}0&0&-4\\1&0&4\\0&1&a\end{pmatrix}
 \pmod d,
\end{equation}
and $\tau$ is the matrix trace of multiplication by a ring element.

\begin{lemma}[Two return criteria]\label{R14:lem:returns}
Suppose $E,h$ are positive integers with $x^E=1$ in $R_d$ and
$3^h\equiv1\pmod E$. If $d\mid P_{3^{k_0}}(a)$, then
\[
 d\mid P_{3^{k_0+ht}}(a)\qquad(t\ge0).
\]
For the negative clock, put $u_0=x$ and
$u_{k+1}=x^2u_k^3$. If $u_{k_0+h}=u_{k_0}$ and
$\tau(u_{k_0})=0$, then
\[
 d\mid P_{2\cdot3^{k_0+ht}-1}(a)\qquad(t\ge0).
\]
\end{lemma}

\begin{proof}
The congruence $3^{k_0+ht}\equiv3^{k_0}\pmod E$ implies equality
of the corresponding powers of $x$, because $x^E=1$. Apply
\eqref{R14:eq:trace-ring}. For the second assertion, induction gives
$u_k=x^{2\cdot3^k-1}$. The deterministic map $u\mapsto x^2u^3$
repeats every subsequent state after a return, so
$u_{k_0+ht}=u_{k_0}$ for every $t\ge0$. Apply
\eqref{R14:eq:trace-ring} again.
\end{proof}

No primality of $d$ is needed in this lemma, and $h$ need not be a
least period. If $a>8$ and the seed trace is larger than $d$,
Lemma~\ref{R14:lem:odd-floors} makes every displayed divisible floor a
proper multiple, including the seed itself.

\subsubsection{A common progression for the coefficient with exponent twelve}

Set
\[
 a=7^{12}=13\,841\,287\,201.
\]
The exact first classical trace factors as
\begin{equation}\label{R14:eq:m12factor}
 \begin{aligned}
 P_3(a)&=a^3+12a-12\\
 &=2\,651\,730\,845\,859\,653\,471\,945\,118\,828\,001\\
 &=3\,037\,679\,595\,827\cdot872\,946\,195\,346\,756\,763.
 \end{aligned}
\end{equation}
Define
\begin{equation}\label{R14:eq:m12data}
 \begin{aligned}
 p&=3\,037\,679\,595\,827,\\
 E&=p^3-1
   =28\,030\,180\,350\,483\,512\,012\,134\,943\,524\,776\,874\,282,\\
 h&=91\,332\,334\,257\,806\,616\,700\,607\,778\,662\,532.
 \end{aligned}
\end{equation}
The packaged exact verifiers establish
\begin{equation}\label{R14:eq:m12certificate}
 x^E=1\ \text{in }R_p,\qquad
 3^h\equiv1\pmod E,\qquad
 x^3=(-4,4,a),\qquad\tau(x^3)=0\pmod p.
\end{equation}
These are finite integer identities, checked by binary powering;
the verification method and filenames are specified in
Section~\ref{R14:sec:verification}.

For the negative clock, reduce the same coefficient modulo $11$.
Then $a\equiv5\pmod{11}$, and the exact states are
\begin{equation}\label{R14:eq:m12negative}
 u_1=(5,8,7)=u_{37},\qquad\tau(u_1)=0\pmod{11}.
\end{equation}
The integer $h$ in \eqref{R14:eq:m12data} is divisible by $36$.

\begin{theorem}[Simultaneous composite floors for one coefficient]
\label{R14:thm:m12}
For every $t\ge0$, both floors
\[
 \floor{\beta_{7^{12}}^{\,3^k}}\quad\text{and}\quad
 \floor{\beta_{7^{12}}^{\,2\cdot3^k-1}}
 \quad\text{are composite at }k=1+ht,
\]
where $h$ is given in \eqref{R14:eq:m12data}. The first is divisible
by $p$ and the second by $11$.
\end{theorem}

\begin{proof}
Equations \eqref{R14:eq:m12certificate} and
Lemma~\ref{R14:lem:returns} give the classical divisibility at
$k=1+ht$. The negative return \eqref{R14:eq:m12negative} gives
divisibility by $11$ at $k=1+36t$. Since $36\mid h$, the same
indices $1+ht$ satisfy both certificates. Their offsets are both
$1$; no assertion about intersecting incompatible progressions is
being made. The factorization \eqref{R14:eq:m12factor} proves that $p$
is a proper divisor at the first classical seed. The negative seed
has exponent $5$ and exceeds $a^3-1>11$. Strict increase from
Lemma~\ref{R14:lem:odd-floors} makes all later multiples proper.
\end{proof}

The verifiers also prove that $p$ is prime by exhaustive trial
division through $\floor{\sqrt p}=1\,742\,894$. This is not a
logical prerequisite for the return proof: a proper integer divisor
and \eqref{R14:eq:m12certificate} already suffice. No iteration through
all $h$ states, unverified order computation, or prime-distribution
hypothesis is used.

\begin{remark}[Field interpretation of this certificate]
Here $p\equiv2\pmod3$. The independent matrix computation also gives
\[
 \det(M^p-M)\equiv2\,129\,913\,761\,719\not\equiv0\pmod p.
\]
Thus $x^p-x$ is a unit in the cubic quotient. If $f_a$ had a root
in $\F_p$, evaluation at that root would send this unit to zero,
which is impossible. The cubic is therefore irreducible modulo
$p$, and $R_p\cong\F_{p^3}$. Since $\gcd(3,p^3-1)=1$, this
explains the absence of a $3$-primary transient in the discovered
return. The proof of Theorem~\ref{R14:thm:m12} uses only the displayed
ring identities, not this interpretation.
\end{remark}

\subsubsection{A prime first trace and later returns at exponent four}

For $a=7^4=2401$,
\begin{equation}\label{R14:eq:m4prime}
 P_3(a)=13\,841\,316\,001
\end{equation}
is prime. Both the quotient-ring verifier and the independent matrix
verifier prove this by exhaustive trial division through
$\floor{\sqrt{P_3(a)}}=117\,649$. In particular, universal
compositeness at the first classical scheduled term is false even
within the subfamily where $4\mid m$.

For the later classical return, the exact data are
\begin{equation}\label{R14:eq:m4certificate}
 \begin{gathered}
 d=59,\qquad a\equiv41\pmod{59},\qquad
 E=59^3-1=205\,378,\qquad h_4=24\,780,\\
 x^E=1,\qquad 3^{h_4}\equiv1\pmod E,\qquad
 x^9=(36,16,24),\qquad\tau(x^9)=0\pmod{59}.
 \end{gathered}
\end{equation}
For the negative clock, an independent matrix calculation modulo
$13$ establishes the state return
\begin{equation}\label{R14:eq:m4negative}
 u_9=u_{19}=(9,4,12),\qquad\tau(u_9)=0\pmod{13}.
\end{equation}
Here the negative-clock state can equivalently be read as the first
column of $M^{2\cdot3^k-1}$.

\begin{corollary}\label{R14:cor:m4}
For every $t\ge0$, the classical floor for $a=7^4$ is composite at
\[
 k=2+24\,780t,
\]
and the negative-clock floor is composite at
\[
 k=9+10t.
\]
\end{corollary}

\begin{proof}
Apply Lemma~\ref{R14:lem:returns} to
\eqref{R14:eq:m4certificate} and \eqref{R14:eq:m4negative}. The first
corresponding exponents are $9$ and $2\cdot3^9-1$. At every odd
exponent $n\ge3$, Lemma~\ref{R14:lem:odd-floors} gives
$P_n(a)>\beta_a^n-1>a^3-1>59>13$. Thus every displayed divisor
is proper.
\end{proof}

The two progressions in Corollary~\ref{R14:cor:m4} are separate; no
common-index assertion is needed or claimed for them. The
certificates for $m=4$ and $m=12$ do not show that the same adaptive
procedure must succeed for each $m$. In both examples the initial
floor $\floor{\beta_{7^m}}=7^m$ is itself composite. Neither is an
all-prime tower or a candidate for a prime leading floor.

\subsection{The conditional classical Pisot setting}

\subsubsection{External prime and Pisot inputs}

The prime-distribution input is the following every-interval lower
bound. There exist constants $c>0$ and $x_0$ such that, with
$\theta_0=21/40$,
\begin{equation}\label{R14:eq:bhp}
 \#\bigl(\PP\cap[x,x+x^{\theta_0}]\bigr)
 \ge c\frac{x^{\theta_0}}{\log x}\qquad(x\ge x_0).
\end{equation}
This is the Baker--Harman--Pintz input \cite{BHP}, in the precise
cardinality formulation stated by Saito
\cite[Theorem 1.6]{Saito}. No numerical value for $x_0$ is assumed.

Saito's reduction \cite[Theorem 1.7]{Saito} says that, if the
classical minimum $\xi$ is algebraic, some power
$\beta=\xi^{3^m}$ is a totally real cubic Pisot number. Its
conjugates can be written as
\begin{equation}\label{R14:eq:saito-shape}
 \beta,\ -\rho,\ \gamma,\qquad
 0<\sigma=|\gamma|<\rho<1,\qquad
 \rho\le\min(\sigma^{17/23},\beta^{-17/40}).
\end{equation}
This is an inherited reduction, not a result established in this
paper. It does not assert that any particular cubic with these
root bounds has a prime tower. The subsequent elementary lemmas
make explicit the further consequences needed here.

Put
\begin{equation}\label{R14:eq:pisot-notation}
 \begin{gathered}
 M=3^r,\qquad p_r=\Tr(\beta^{3^r}),\qquad
 e_r=\beta^{3^r}-p_r=\rho^{3^r}-\gamma^{3^r},\\
 u=-\frac{\log\rho}{\log\beta},\qquad
 \kappa=\frac{2-u}{3}.
 \end{gathered}
\end{equation}
Unlike the family in \eqref{R14:eq:family}, the sign of $\gamma$ in
this conditional setting is not prescribed. Since
$|\gamma|<\rho$, however, $e_r>0$, $e_r\sim\rho^{3^r}$, and
$e_r<1$ for all sufficiently large $r$. Therefore on a sufficiently
late tail
\begin{equation}\label{R14:eq:prime-tail}
 p_r=\floor{\beta^{3^r}}=\floor{\xi^{3^{m+r}}}\in\PP.
\end{equation}

\begin{lemma}[The strict gap exponent]\label{R14:lem:strict-gap}
Under \eqref{R14:eq:saito-shape}, the number $u$ is irrational and
\begin{equation}\label{R14:eq:u-range}
 \frac{17}{40}<u<\frac12,\qquad
 \frac12<\kappa<\frac{21}{40}.
\end{equation}
Furthermore,
\begin{equation}\label{R14:eq:gap-asymptotic}
 G_r:=p_{r+1}-p_r^3
 \sim3(\beta^2\rho)^{3^r}
 \sim3(p_r^3)^\kappa.
\end{equation}
In particular, eventually
$p_r^3<p_{r+1}<(p_r+1)^3$.
\end{lemma}

\begin{proof}
The absolute norm $\beta\rho\sigma$ is a nonzero positive integer.
Since $\sigma<\rho$, we have $\beta\rho^2>1$, so $u<1/2$.
The bound $\rho\le\beta^{-17/40}$ gives $u\ge17/40$.

We prove irrationality to exclude equality at this endpoint.
If $u=s/t$ with positive integers $s,t$, then
$\beta^s\rho^t=1$, or
\[
 \beta^s(-\rho)^t=(-1)^t.
\]
The Galois group of the splitting field of an irreducible cubic is
either $S_3$ or the cyclic group of order $3$. In the first case,
apply the transposition interchanging $\beta$ and $-\rho$, and
divide the two relations in absolute value. This gives
$(\beta/\rho)^{s-t}=1$, so $s=t$. It follows that
$\beta\rho=1$, contradicting
$\beta\rho\sigma=\sigma<1$ and the integral nonzero norm.
In the cyclic case, let $\ell_i$ be the logarithms of the absolute
values of the three conjugates in cyclic order. The relation and
its two conjugates give
\[
 s\ell_1+t\ell_2=0,\qquad
 s\ell_2+t\ell_3=0,\qquad
 s\ell_3+t\ell_1=0.
\]
The coefficient matrix has determinant $s^3+t^3>0$, forcing all
$\ell_i=0$, in contradiction with $\log\beta>0$.
This proves irrationality, and therefore \eqref{R14:eq:u-range}.

With $M=3^r$, expand exactly:
\[
 \begin{aligned}
 G_r
 &=\beta^{3M}-e_{r+1}-(\beta^M-e_r)^3\\
 &=3\beta^{2M}e_r-3\beta^Me_r^2+e_r^3-e_{r+1}.
 \end{aligned}
\]
We have $e_r\sim\rho^M$ and $e_{r+1}\sim\rho^{3M}$.
Dividing the last three terms by $\beta^{2M}\rho^M$ shows that
each tends to zero, since $\rho/\beta<1$. Thus
$G_r\sim3(\beta^2\rho)^M$. Since $p_r\sim\beta^M$ and
$\rho=\beta^{-u}$, the second equivalence in
\eqref{R14:eq:gap-asymptotic} follows. The gap is eventually positive.
On the tail where $e_r\in(0,1)$,
$p_{r+1}<\beta^{3M}<(p_r+1)^3$, proving the upper nesting bound.
\end{proof}

\subsubsection{Integer cells and the eventual least child}

\begin{lemma}[Exact ancestor of an integer cell]\label{R14:lem:ancestor}
Let $q\ge1$ and $D\ge1$ be integers. Throughout the half-open
interval
\[
 [q^{1/D},(q+1)^{1/D}),
\]
the integer part is constant and equals $\floor{q^{1/D}}$.
Consequently a terminal label $q=\floor{B^{LD}}$, where $L>0$,
determines its earlier label exactly by
\[
 \floor{B^L}=\floor{q^{1/D}}.
\]
\end{lemma}

\begin{proof}
Put $j=\floor{q^{1/D}}$. Then $j^D\le q<(j+1)^D$.
Both $q$ and $(j+1)^D$ are integers, so $q+1\le(j+1)^D$.
Every point of the half-open interval is at least $j$ and strictly
less than $j+1$. This proves both assertions, including the case
where its excluded right endpoint is an integer.
\end{proof}

\begin{lemma}[Continuation and least-prime choice]\label{R14:lem:least-prime}
The estimate \eqref{R14:eq:bhp} implies that every sufficiently large
prime has an infinite strict cubic prime continuation. If
\eqref{R14:eq:prime-tail} comes from the actual global minimum $\xi$,
then, for every sufficiently large $r$, $p_{r+1}$ is the least
prime strictly greater than $p_r^3$.
\end{lemma}

\begin{proof}
At $x=q^3$, estimate \eqref{R14:eq:bhp} supplies a prime $q'$ with
\[
 q^3<q'\le q^3+q^{63/40}<(q+1)^3-1
\]
for all sufficiently large primes $q$. The left endpoint is a
composite cube. The last inequality holds eventually because
$63/40<2$. Repeating this choice produces primes $q_j$ with
$q_j^3<q_{j+1}$ and $q_{j+1}+1<(q_j+1)^3$.
The corresponding closed inverse-power cells are nested, with
strict containment at both endpoints at each successive step.
Their lengths tend to zero. Their unique common point lies below
every right endpoint, since the next cell already does, and hence
realizes all the prescribed prime floors. This proves continuation.

Now suppose, at a sufficiently late index, that a prime $q$ obeys
$p_r^3<q<p_{r+1}$. In the original $\xi$ clock, replace the label
at level $m+r+1$ by $q$ and continue from it as just constructed.
The cell of this replacement label lies within the cell of $p_r$:
its lower endpoint is above the $p_r$ lower endpoint, while
$q+1\le p_{r+1}<(p_r+1)^3$. Lemma~\ref{R14:lem:ancestor} then
preserves every earlier label as well. The resulting base $B$ is
strictly smaller than $\xi$, because
\[
 B\le(q+1)^{1/3^{m+r+1}}
 \le p_{r+1}^{1/3^{m+r+1}}<\xi.
\]
It has prime floors at all original levels, contradicting global
minimality. Thus no such prime $q$ exists. Since $p_{r+1}$ is
itself prime and exceeds $p_r^3$, it is the least prime above that
cube.
\end{proof}

The strict endpoint and the least-prime property are preparatory
consequences of the cited reduction, the prime input, and actual
global minimality. They are not consequences of a long
skipped-level interval.

\subsection{Ancestral conditions remove the apparent exponent gain}

Work on the sufficiently late tail of the conditional minimum
setting above. Fix a depth $d\ge2$, put $M=3^r$, and consider
terminal labels at level $r+d$. A label $q$ has the half-open cell
\[
 \left[q^{1/(M3^d)},(q+1)^{1/(M3^d)}\right)
\]
in the $\beta$ coordinate. With $D=3^{d-1}$,
Lemma~\ref{R14:lem:ancestor} makes its first descendant of level $r$
exactly
\begin{equation}\label{R14:eq:first-ancestor}
 \floor{q^{1/D}}.
\end{equation}

\begin{proposition}[The terminal search corridor]\label{R14:prop:corridor}
Every prime terminal label in
\begin{equation}\label{R14:eq:forbidden-raw}
 (p_r^{3^d},\ p_{r+1}^{3^{d-1}})
\end{equation}
has a composite ancestor at level $r+1$. More generally, for the
search among earlier terminal labels $q<p_{r+d}$, preserving the
first $j$ descendants leaves the geometric corridor
\begin{equation}\label{R14:eq:corridors}
 I_{r,d}^{(j)}=(p_{r+j}^{3^{d-j}},\ p_{r+d}),
 \qquad 0\le j<d.
\end{equation}
For $j<d-1$, requiring the next ancestor to be prime reduces this
corridor to the one with $j$ replaced by $j+1$. After all intermediate
conditions are imposed, the remaining corridor is
\begin{equation}\label{R14:eq:final-corridor}
 I_{r,d}^{(d-1)}=(p_{r+d-1}^3,\ p_{r+d}).
\end{equation}
Here a corridor is a real search interval, not an assertion that
it contains a prime.
\end{proposition}

\begin{proof}
For a prime $q$ in \eqref{R14:eq:forbidden-raw}, formula
\eqref{R14:eq:first-ancestor} gives an integer between $p_r^3$ and
$p_{r+1}-1$. The first is a composite cube, and every larger
integer in this range is composite by
Lemma~\ref{R14:lem:least-prime}. Thus the ancestor is not prime.

For the general step, suppose the label $p_{r+j}$ has been
preserved. Its cell gives the lower terminal bound
$p_{r+j}^{3^{d-j}}$. Its upper terminal bound causes no additional
restriction on $q<p_{r+d}$, because the candidate itself satisfies
\[
 p_{r+d}<\beta^{M3^d}
 <(p_{r+j}+1)^{3^{d-j}}.
\]
For prime terminal labels the lower endpoint, a proper integer
power, is already excluded. This gives the corridor
\eqref{R14:eq:corridors}.

If $j<d-1$, every terminal prime below
$p_{r+j+1}^{3^{d-j-1}}$ has its next ancestor in
$[p_{r+j}^3,p_{r+j+1})$. That ancestor is composite by the same
least-prime argument. Conversely, within the remaining corridor,
the next ancestor must equal $p_{r+j+1}$: its lower bound is that
integer, and its upper bound follows from
$p_{r+d}<(p_{r+j+1}+1)^{3^{d-j-1}}$. Thus imposing that next
condition removes exactly the earlier forbidden portion of the
geometric search and gives the next corridor. Iterate to obtain
\eqref{R14:eq:final-corridor}.
\end{proof}

A continuation after a prime terminal label cannot repair a
composite earlier label. Also, this proposition concerns a first
change by the chosen terminal depth. A smaller real base can still
have terminal label $q=p_{r+d}$ and first differ at a deeper level;
such a base is not being ruled out merely by the terminal-label
search at depth $d$.

\begin{theorem}[Fixed-depth width loss]\label{R14:thm:width}
Fix $d\ge2$. As $r\to\infty$, put $X=\beta^{M3^d}$ and
\[
 \eta_s=1-\frac{1+u}{3^s}\qquad(s\ge1).
\]
For each $0\le j<d$,
\begin{equation}\label{R14:eq:width}
 |I_{r,d}^{(j)}|
 \sim3^{d-j}\beta^{(3^d-3^j)M}\rho^{3^jM}
 =3^{d-j}X^{\eta_{d-j}}.
\end{equation}
In particular,
\begin{equation}\label{R14:eq:eta-ranges}
 \eta_1=\kappa\in\left(\frac12,\frac{21}{40}\right),
 \qquad
 \eta_2\in\left(\frac56,\frac{101}{120}\right),
 \qquad\eta_s\longrightarrow1.
\end{equation}
The ratio of the final to the raw corridor width is
\begin{equation}\label{R14:eq:width-ratio}
 \frac{|I_{r,d}^{(d-1)}|}{|I_{r,d}^{(0)}|}
 \sim3^{1-d}
 \left(\frac\rho\beta\right)^{(3^{d-1}-1)M}
 \longrightarrow0.
\end{equation}
\end{theorem}

\begin{proof}
Write $L=3^{d-j}$ and $Y=\beta^{M3^j}$. Since
$p_{r+j}=Y-e_{r+j}$ and $p_{r+d}=Y^L-e_{r+d}$,
\[
 |I_{r,d}^{(j)}|
 =Y^L-e_{r+d}-(Y-e_{r+j})^L.
\]
For fixed $d$ and $j$, the binomial expansion gives the leading
term $LY^{L-1}e_{r+j}$. The next binomial term relative to this
term is $O(e_{r+j}/Y)$, which tends to zero like
$O((\rho/\beta)^{M3^j})$. Also
\[
 \frac{e_{r+d}}{Y^{L-1}e_{r+j}}
 \sim\left(\frac\rho\beta\right)^{M(3^d-3^j)}
 \longrightarrow0.
\]
Because $e_{r+j}\sim\rho^{M3^j}$, this proves the first
asymptotic in \eqref{R14:eq:width}. Substituting $\rho=\beta^{-u}$
proves the power of $X$. The bounds in \eqref{R14:eq:eta-ranges}
follow from \eqref{R14:eq:u-range}. Dividing the expressions in
\eqref{R14:eq:width} for $j=d-1$ and $j=0$ yields
\eqref{R14:eq:width-ratio}.
\end{proof}

A raw two-level corridor therefore has an exponent greater than
$5/6$, well above the BHP exponent. Its part compatible with every
ancestor has exponent $\kappa<21/40$, exactly the original
one-step exponent at a later index. The gain is lost through the
ancestral conditions, not through an imprecise estimate of the
terminal interval. All these asymptotics have fixed depth; no
uniform assertion for depth increasing with $r$ is made.

\subsection{Thinning while preserving short-interval richness}

The following theorem isolates the information retained by a
short-interval lower bound. It does not assume that $S$ consists
of primes, that the specified base is algebraic, or that it was
already a minimum for $S$.

\begin{theorem}[Altered-set minimum]\label{R14:thm:thinning}
Let $0<\kappa<\theta<1$, and suppose $S\subseteq\N$ satisfies
\begin{equation}\label{R14:eq:richness}
 \#(S\cap[x,x+x^\theta])
 \ge c\frac{x^\theta}{\log x}\qquad(x\ge x_0)
\end{equation}
for some $c>0$ and $x_0>1$. Let $A>1$, and suppose that for every
$k\ge1$,
\begin{equation}\label{R14:eq:thinning-labels}
 P_k=\floor{A^{3^k}}\in S,\qquad P_1\ge2,\qquad
 P_k^3<P_{k+1}<(P_k+1)^3.
\end{equation}
Assume also that, for all sufficiently large $k$,
\begin{equation}\label{R14:eq:small-gap}
 P_{k+1}-P_k^3=O((P_k^3)^\kappa).
\end{equation}
Define
\begin{equation}\label{R14:eq:altered-set}
 S_A=(S\cap[P_1,\infty))
 \setminus\bigcup_{k\ge1}[P_k^3,P_{k+1}).
\end{equation}
Then every $P_k$ belongs to $S_A$, the number $A$ is the attained
global minimum of $\mathcal W(S_A)$, and for all sufficiently large
$x$,
\begin{equation}\label{R14:eq:retained-richness}
 \#(S_A\cap[x,x+x^\theta])
 \ge\frac c2\frac{x^\theta}{\log x}.
\end{equation}
More precisely, the number of integers removed from such a window
is $O(x^\kappa)$.
\end{theorem}

\begin{proof}
Because $P_1\ge2$ and $P_{k+1}>P_k^3$, the labels strictly
increase and $P_k<P_k^3$ for every $k$. The deletion intervals are
disjoint: the right endpoint of the $k$th is $P_{k+1}$, smaller
than the left endpoint $P_{k+1}^3$ of the next. Their right endpoints
are excluded. Thus no $P_k$ is deleted, and all labels survive the
initial cutoff. Hence $A\in\mathcal W(S_A)$.

Suppose $1<B<A$ also belonged to $\mathcal W(S_A)$, and write
$Q_k=\floor{B^{3^k}}$. Monotonicity gives $Q_k\le P_k$.
The cutoff gives $Q_1\ge P_1$, so $Q_1=P_1$. The two label
sequences cannot agree forever. If they did, both $A$ and $B$
would belong to the same inverse-power cell at every $k$, so
\[
 0<A-B\le(P_k+1)^{1/3^k}-P_k^{1/3^k}\le3^{-k}.
\]
The last inequality is the mean value theorem, since the derivative
of $y^{1/3^k}$ is at most $3^{-k}$ for $y\ge1$.
This is impossible as $k\to\infty$.

At the first differing index $k+1$, the preceding label is
$Q_k=P_k$, whereas $Q_{k+1}<P_{k+1}$. The preceding equality
implies $B^{3^k}\ge P_k$, and therefore
\[
 P_k^3\le Q_{k+1}<P_{k+1}.
\]
This is precisely a deleted interval. It contradicts
$Q_{k+1}\in S_A$, proving global leastness.

For the quantitative statement put $X_k=P_k^3$ and
$\Delta_k=P_{k+1}-X_k$. Since $\kappa<1$,
\eqref{R14:eq:small-gap} implies $P_{k+1}/X_k\to1$, and hence
\begin{equation}\label{R14:eq:deletion-spacing}
 X_{k+1}=P_{k+1}^3\sim X_k^3.
\end{equation}
The fixed cutoff and the finitely many early deletion intervals
are irrelevant for sufficiently large $x$. If the window
$[x,x+x^\theta]$ meets a late interval $[X_k,P_{k+1})$, then
$X_k\le x+x^\theta$ and $x<P_{k+1}=X_k+O(X_k^\kappa)$.
Since both $\theta$ and $\kappa$ are below $1$, these inequalities
give $X_k\sim x$. Two late deletion intervals cannot both meet
such a window: their starting points would have ratio tending to
$1$, whereas \eqref{R14:eq:deletion-spacing} makes the ratios of
successive starting points tend to infinity. Thus at most one late
deletion interval is met. It contains $O(x^\kappa)$ integers.
Subtracting this loss from \eqref{R14:eq:richness} gives
\[
 \#(S_A\cap[x,x+x^\theta])
 \ge c\frac{x^\theta}{\log x}-O(x^\kappa).
\]
Finally $x^\kappa=o(x^\theta/\log x)$ because $\kappa<\theta$.
This proves \eqref{R14:eq:retained-richness} for every sufficiently
large real $x$, not merely for an averaged collection of windows.
\end{proof}

\begin{remark}[Why the first-label hypothesis matters]
The hypothesis $P_1\ge2$ is essential to the stated general
construction. If $P_1=1<P_2$, its first deletion interval
$[P_1^3,P_2)=[1,P_2)$ removes $P_1$ itself. Merely assuming
$A>1$ would not prevent this defect. In the prime applications and
the explicit example below, the stronger hypothesis holds.
\end{remark}

\subsubsection{Conditional specialization to the actual primes}

Take $S=\PP$ and $\theta=21/40$. If an all-prime classical
candidate $A$ has the gap bound \eqref{R14:eq:small-gap} with
$\kappa<21/40$, its labels automatically have $P_1\ge2$ and
strict nesting: a prime cannot equal the cube of an earlier prime,
and the upper nesting bound follows from its floor cell.
Theorem~\ref{R14:thm:thinning} then gives a subset $\PP_A\subseteq\PP$
which retains the entire candidate tower, still satisfies the BHP
lower-bound scale, and has $A$ as its exact global minimum.
No composite has been added to the allowed outputs.

This is a conditional stability statement, not an unconditional
construction of an algebraic prime tower. The existence of the
required algebraic all-prime candidate has not been established.
Moreover, even if such a candidate were supplied, the theorem
changes the allowed set from all primes to a subset; it is not a
counterexample about the full prime set. The conclusion is that
the particular short-interval information
\eqref{R14:eq:richness} survives deletion of the competing-child
intervals when their exponent is strictly smaller.

If only an eventually all-prime Pisot tail is given, its indexing
must first be normalized. Choose $R$ so that
$\floor{\beta^{3^{R+j}}}$ is prime for every $j\ge0$, and put
\begin{equation}\label{R14:eq:tail-normalization}
 \beta'=\beta^{3^R},\qquad A_{\mathrm{tail}}=(\beta')^{1/3},
 \qquad
 P_k=\floor{(\beta')^{3^{k-1}}}\quad(k\ge1).
\end{equation}
Choose $R$ late enough also for strict nesting. Then
$P_k=\floor{A_{\mathrm{tail}}^{3^k}}$ is a complete clock in the
convention of Theorem~\ref{R14:thm:thinning}. In the Pisot setting,
\eqref{R14:eq:gap-asymptotic} supplies its small-gap bound. The altered
set has minimum $A_{\mathrm{tail}}$, which need not be the original
Mills minimum. An already complete classical candidate requires no
such normalization.

The definition of primes and further arithmetic information about
the full prime set can distinguish it from a thinned subset. Thus
this theorem does not prohibit an argument using BHP together with
additional arithmetic. It identifies the missing information in an
argument that uses only the displayed richness bound, primality of
the retained elements, and existence of an admissible tower.

\subsection{An explicit algebraic minimum for integers coprime to six}

The conditional prime specialization has an unconditional analogue
for a larger integer set. It illustrates the scope of the thinning
lemma without asserting that the outputs are primes.

\begin{theorem}[An explicit nonprime-output witness]\label{R14:thm:wheel}
Let
\begin{equation}\label{R14:eq:wheel-poly}
 a=2^{40}+1,\qquad b=2^{22},\qquad
 f(X)=X^3-aX^2-bX+2,
\end{equation}
and let $\beta$ be its largest real root. There is a set
$S_\ast\subseteq\{n\ge1:\gcd(n,6)=1\}$ with the following
properties:
\begin{enumerate}[label=\textup{(\roman*)}]
 \item $\beta^{1/3}$ is the attained global minimum of
       $\mathcal W(S_\ast)$.
 \item For every sufficiently large real $x$,
       \[
       \#\bigl(S_\ast\cap[x,x+x^{21/40}]\bigr)
       =\left(\frac13+o(1)\right)x^{21/40}.
       \]
 \item The roots are $\beta,-\rho,\gamma$ with
       $0<\gamma<\rho<1$ and
       \[
       \rho<\min(\gamma^{17/23},\beta^{-17/40}).
       \]
       Their trace gaps have an exponent
       $\kappa\in(1/2,21/40)$.
\end{enumerate}
The minimizing base is algebraic, but its first classical label is
$2^{40}+1$, which is composite.
\end{theorem}

\begin{proof}
The only possible rational roots of the monic cubic are
$\pm1,\pm2$. Their values are
\[
 \begin{array}{ll}
 f(1)=3-a-b,& f(-1)=1-a+b,\\
 f(2)=10-4a-2b,& f(-2)=-6-4a+2b.
 \end{array}
\]
All are negative and nonzero, so the cubic is irreducible.
Exact rational substitutions give the signs
\[
 \begin{array}{c|rrrrrr}
 z&-5\cdot10^{-6}&-4\cdot10^{-6}&3\cdot10^{-7}&10^{-6}&a&a+1\\
 \hline
 \operatorname{sgn}f(z)&-&+&+&-&-&+
 \end{array}
\]
For the last two, $f(a)=2-ba<0$ and
$f(a+1)=(a+1)(a+1-b)+2>0$. The other four are direct rational
calculations, also checked exactly by the packaged portable
verifier. Hence the three disjoint root brackets are
\begin{equation}\label{R14:eq:wheel-brackets}
 \begin{gathered}
 a<\beta<a+1,\\
 4\cdot10^{-6}<\rho<5\cdot10^{-6},\qquad
 3\cdot10^{-7}<\gamma<10^{-6}.
 \end{gathered}
\end{equation}
They exhaust the roots and give $0<\gamma<\rho<1$.

We record elementary bounds proving the required root-size
inequalities. Set $t=5\cdot10^{-6}$. Since
$25\cdot2^{35}<10^{12}$, we have $t^2<2^{-35}$.
Also $a+1=2^{40}+2<2^{41}$. Therefore
\[
 t^{40}(a+1)^{17}<2^{-700}2^{697}<1,
\]
which, with \eqref{R14:eq:wheel-brackets}, gives
$\rho<\beta^{-17/40}$. For the other inequality, use
$t<2^{-17}$ and $3\cdot10^{-7}>2^{-22}$ to get
\[
 t^{23}<2^{-391}<2^{-374}<(3\cdot10^{-7})^{17}.
\]
Thus $\rho^{23}<\gamma^{17}$, as required.
The norm is $-2$, so $\beta\rho\gamma=2$ and
$\beta\rho^2>2$. Consequently
\[
 \frac{17}{40}<-\frac{\log\rho}{\log\beta}<\frac12.
\]
The gap-exponent assertion now follows by the same expansion as in
Lemma~\ref{R14:lem:strict-gap}; no prime hypothesis was used in that
expansion.

Let $T_n$ be the power sum of the three roots. It satisfies
\begin{equation}\label{R14:eq:wheel-trace}
 \begin{gathered}
 T_0=3,\qquad T_1=a,\qquad T_2=a^2+2b,\\
 T_n=aT_{n-1}+bT_{n-2}-2T_{n-3}\qquad(n\ge3).
 \end{gathered}
\end{equation}
For every odd positive $n$,
$T_n=\beta^n-\rho^n+\gamma^n$ and
$0<\rho^n-\gamma^n<1$, so $T_n=\floor{\beta^n}$.
Put $p_r=T_{3^r}$ for $r\ge0$. These are positive integers.
Modulo $2$, recurrence \eqref{R14:eq:wheel-trace} and its initial
values show that every positive $T_n$ is odd. Modulo $3$, the
cubing identity for three variables gives
$T_{3n}\equiv T_n^3\pmod3$. Thus
$p_r\equiv a\equiv-1\pmod3$. We have proved
\begin{equation}\label{R14:eq:wheel-residues}
 p_r\equiv5\pmod6\qquad(r\ge0).
\end{equation}

Strict nesting holds at every index here, not only eventually.
For $n=3^r$, the exact algebraic identity is
\begin{equation}\label{R14:eq:wheel-gap}
 p_{r+1}-p_r^3
 =3(\beta^n-\rho^n)(\beta^n+\gamma^n)(\rho^n-\gamma^n)>0.
\end{equation}
The upper inequality follows from
$p_{r+1}<\beta^{3n}<(p_r+1)^3$. The asymptotic proved above also
gives $p_{r+1}-p_r^3=O((p_r^3)^\kappa)$ with
$\kappa<21/40$.

Apply Theorem~\ref{R14:thm:thinning} to
\[
 S=\{n\ge1:\gcd(n,6)=1\},\qquad
 A=\beta^{1/3},\qquad P_k=p_{k-1}.
\]
This indexing gives
$P_k=\floor{A^{3^k}}$ for every $k\ge1$, and
$P_1=p_0=a\ge2$. Define explicitly
\begin{equation}\label{R14:eq:wheel-set}
 S_\ast=(S\cap[a,\infty))
 \setminus\bigcup_{r\ge0}[p_r^3,p_{r+1}).
\end{equation}
The base $A$ is algebraic because $\beta$ is algebraic, and the
thinning theorem proves that $A$ is the exact minimum for this
set. For an interval of length $L$, the original set $S$ has
$L/3+O(1)$ elements, by counting the residue classes $1$ and $5$
modulo $6$. The proof of the thinning theorem removes only
$O(x^\kappa)$ elements from $[x,x+x^{21/40}]$. Its count is
therefore
\[
 \frac{x^{21/40}}3+O(1)+O(x^\kappa)
 =\left(\frac13+o(1)\right)x^{21/40}.
\]
Finally $2^8\equiv-1\pmod{257}$, so
$2^{40}+1\equiv0\pmod{257}$ and $a>257$. The first classical
label is composite.
\end{proof}

This example combines the specified cubic root pattern, exact
odd trace floors, a sub-BHP gap exponent, elementary coprimality
to $6$, an algebraic attained minimum, and richness on every
BHP-scale window. It does not combine these properties with prime
outputs. Its explicitly composite first label prevents an
interpretation as a Mills counterexample.

\subsection{Exact verification and reproducibility}\label{R14:sec:verification}

The numerical certificates are retained in the original source archive in the
\path{certificates/} directory. All four listed programs use only
Python 3 and its standard library.

\begin{enumerate}
 \item \path{adaptive_case_and_eisenstein_verify.py} uses the
       quotient ring $R_d$ and integer binary powering. It checks
       \eqref{R14:eq:m12factor}--\eqref{R14:eq:m12negative}, the divisibility
       $36\mid h$, the numerical primality assertions, and
       \eqref{R14:eq:m4certificate}. Its finite polynomial and gcd
       tests are sanity checks, not proofs of the universal
       theorems.
 \item \path{audit_adaptive_matrix_verify.py} independently uses
       the $3$ by $3$ matrix \eqref{R14:eq:companion}. It checks the
       same returns, the field-interpretation determinant, and the
       negative-clock return \eqref{R14:eq:m4negative}. Its exhaustive
       primality routine tests the possible divisors of the form
       $6j\pm1$, separately from the quotient-ring program.
 \item \path{independent_audit.py} supplies a further independent
       exact arithmetic check of the adaptive witnesses.
 \item \path{classical_minimality_portable_verify.py} checks the
       rational root brackets and inequalities for
       \eqref{R14:eq:wheel-poly}, the initial trace residues and
       nesting, the polynomial identity \eqref{R14:eq:wheel-gap},
       and the algebra of the corridor exponents and width ratio.
       It includes a regression check for the indispensable
       hypothesis $P_1\ge2$. These are not primality or
       transcendence tests.
\end{enumerate}

For the ring arithmetic, multiply triples as polynomials of degree
at most $4$, then reduce the terms of degrees $4$ and $3$, in that
order, using $x^3=ax^2+4x-4$. All coefficients are reduced modulo
$d$. Binary powering starts with $v=1$ and successively squares the
base, multiplying it into $v$ whenever the current binary digit of
the exponent is $1$. The invariant is the usual exact product of
the already processed powers. This computes the claimed identities
in logarithmically many ring multiplications. The matrix verifier
performs ordinary modular matrix products instead, so the two
implementations do not share the quotient-ring multiplication
routine.

Exhaustive trial division is a deterministic primality proof: a
composite integer $N$ has a prime divisor at most $\sqrt N$.
After $2$ and $3$ have been checked, every possible prime divisor
larger than $3$ is of the form $6j\pm1$. No probable-prime result
is accepted as a proof. The proper-divisor conclusions themselves
do not require that the recurring divisors be prime.

From the artifact directory the checks can be repeated with
\begin{lstlisting}
python3 certificates/adaptive_case_and_eisenstein_verify.py
python3 certificates/audit_adaptive_matrix_verify.py
python3 certificates/independent_audit.py
python3 certificates/classical_minimality_portable_verify.py
sha256sum -c SHA256SUMS.txt
\end{lstlisting}
The accompanying successful logs and available JSON reports record
these finite checks. The all-$k,r$ Eisenstein theorem,
all-$m,k$ consecutive-coprimality theorem, fixed-depth asymptotics,
and general thinning theorem are proved in the text, rather than
inferred from finite samples.


\subsection{Remaining questions}

The arithmetic certificates settle the two specified coefficients,
not the whole thin family $a=7^m$. A universal adaptive approach
would still need a theorem guaranteeing a suitable divisor and
return, or another compositeness mechanism, for each remaining
coefficient. The irreducibility theorem rules out a particular
polynomial-composition shortcut; it does not prove that such a
universal argument is impossible.

For the actual classical minimum, a skipped-level prime is useful
only when every prescribed earlier floor is prime. The width
calculation identifies the resulting one-step scale exactly.
Possible additional inputs include a prime inside a competing-child
interval, a correlation theorem excluding an entire candidate
sequence of exceptional prime-free intervals, or a direct arithmetic
obstruction to the candidate trace orbit. None of those inputs is
proved here.

The altered-set theorem establishes the stability of a specified
short-interval lower bound under thinning. It does not establish
an algebraic all-prime tower, does not replace the full prime set
by an equivalent one, and does not preclude a proof using further
arithmetic information. No Riemann Hypothesis, Density Hypothesis,
Artin-type conjecture, or presumed prime pattern on a prescribed
orbit is used in the results proved here.



\endgroup

\section{Projective return criteria and finite horizon limitations}\label{src:R15}
The finite-ring examples admit a common projective formulation. This section gives recurrence criteria, finite-horizon obstructions, and the precise behavior of witness orders under powering; its projective exponent also strengthens the density theorem.
\medskip
\begingroup




\subsection{The clocks and the exact floor bridge}

For an integer $a>8$, put
\[
 f_a(X)=X^3-aX^2-4X+4,
 \qquad n_+(k)=3^k,\quad n_-(k)=2\cdot3^k-1\quad(k\ge1).
\]
The names classical and negative refer to these two exponent sequences;
both have positive odd exponents. Write $P_n(A)$ for the formal Newton
power-sum polynomial of $X^3-AX^2-4X+4$. Its recurrence is
\begin{equation}\label{R15:eq:rec}
 P_0=3,\quad P_1=A,\quad P_2=A^2+8,\qquad
 P_n=AP_{n-1}+4P_{n-2}-4P_{n-3}\quad(n\ge3).
\end{equation}
In particular $P_n\in\Z[A]$, monic of degree $n$ for $n\ge1$.

\begin{lemma}\label{R15:lem:floor}
The roots of $f_a$ are $-\rho,\sigma,\beta_a$, where
$0<\sigma<\rho<1$ and $a<\beta_a<a+1$. The polynomial is
irreducible, and at every positive odd exponent
\[
 \floor{\beta_a^n}=P_n(a),\qquad
 0<\beta_a^n-P_n(a)=\rho^n-\sigma^n<1.
\]
These floors strictly increase with the positive odd exponent. If
$a>2^{41}$, then additionally
\begin{equation}\label{R15:eq:shape}
 \rho<\tfrac12\beta_a^{-17/40},\qquad
 \rho<\sigma^{17/23},\qquad 4<\beta_a^{3/20}.
\end{equation}
\end{lemma}
\begin{proof}
The values at $-1,0,1,a,a+1$ are respectively
$7-a,4,1-a,4(1-a),(a-1)^2$. Their signs put one root in each
of $(-1,0),(0,1),(a,a+1)$. No root is integral, so the rational-root
test proves irreducibility. Their sum gives $\rho-\sigma=\beta_a-a>0$.
The odd power sum proves the floor identity. Since $\beta_a>8$,
$\beta_a^{n+2}-\beta_a^n>1$, which proves strict increase.
The equation at $-\rho$ gives $a\rho^2=4+4\rho-\rho^3<8$.
With $\beta_a<2a$,
\[
 (2\rho)^{40}\beta_a^{17}
 <(32/a)^{20}(2a)^{17}=2^{117}/a^3<1.
\]
The norm identity $\beta_a\rho\sigma=4$ shows that
$\rho<\sigma^{17/23}$ is equivalent to
$\rho^{40}\beta_a^{17}<4^{17}$, implied by the preceding bound.
Finally $4^{20}<a^3<\beta_a^3$.
\end{proof}

The bounds \eqref{R15:eq:shape} are of the kind arising in the conditional
Pisot reduction for the classical Mills minimum
\cite[Theorem 1.7]{Saito}. Satisfying them does not make a cubic an
all-prime tower or a least constant. Our main arithmetic family is
$a=7^m$, $m\ge2$, whose integer part is already composite.
Prime-leading examples will occur separately in the bounded-witness
and finite-prime limitations.

For a modulus $D\ge2$, let
\[
 R_D=(\Z/D\Z)[X]/(f_a),\qquad x=X\bmod f_a.
\]
Multiplication trace is the linear function
\begin{equation}\label{R15:eq:tau}
 \tau(c_0+c_1x+c_2x^2)=3c_0+ac_1+(a^2+8)c_2\pmod D.
\end{equation}
It satisfies $\tau(x^n)=P_n(a)\pmod D$, by the common recurrence
and initial values. Multiplication by $x$ has matrix
\begin{equation}\label{R15:eq:matrix}
 M_a=\begin{pmatrix}0&0&-4\\1&0&4\\0&1&a\end{pmatrix}\pmod D.
\end{equation}
This interpretation includes repeated factors and composite moduli.

\subsection{A positive projective recurrence mechanism}

\subsubsection{Every zero seed recurs at an irreducible reduction}

If $D$ is odd, $x$ is a unit because the constant term is $4$.
Let $Q_D$ be its order modulo scalar units:
\[
 Q_D=\ord\bigl([x]\in R_D^*/(\Z/D\Z)^*\bigr).
\]
Equivalently, $Q_D$ is the least positive integer such that $x^{Q_D}$
is a scalar unit. Multiplication by such a scalar preserves trace zero
in both directions. Thus trace-zero status depends only on the exponent
modulo $Q_D$, although the trace value itself need not.

\begin{theorem}[Irreducible reduction]\label{R15:thm:irreducible}
Let $p$ be an odd prime and suppose $f_a$ is irreducible modulo $p$.
For either clock, if $p\mid P_{n_\pm(k_0)}(a)$ at any original
index $k_0\ge1$, then the same prime divides infinitely many terms
on that clock. For the real root, these yield infinitely many composite
floors, with at most a finite initial exception for properness.
\end{theorem}
\begin{proof}
The quotient is $\F_{p^3}$. Put $S=p^2+p+1$.
Its field norm gives $x^S=-4\in\F_p^*$, so $Q_p\mid S$.
If $p\equiv1\pmod3$, then $v_3(S)=1$; otherwise, including
$p=3$, $v_3(S)=0$. Write $S=3^v B$ with $3\nmid B$, and use
$\ord_1(3)=1$. For $h=\ord_B(3)$ and $k_0\ge1$, both
\[
 3^{k_0}(3^h-1),\qquad 2\cdot3^{k_0}(3^h-1)
\]
are divisible by $S$. Hence the corresponding clock states at
$k_0+h$ and $k_0$ differ by a scalar unit. Their traces vanish
together. Repetition gives every $k_0+th$, $t\ge0$.
Lemma~\ref{R15:lem:floor} makes the growing floors eventually larger than $p$.
\end{proof}

This proof needs a scalar return, not equality of the full ring states.
For example, if $O=\ord(x)$ and $d=\ord_p(-4)$, the norm gives
$d=O/\gcd(O,S)$. When $p\equiv1\pmod3$ and $v_3(d)>0$,
this forces $v_3(O)=v_3(d)+1$, which can exceed one. Those extra
powers of $3$ lie in scalar information and disappear projectively.
Irreducibility cannot simply be omitted; an exact reducible transient
example is given in Section~\ref{R15:sec:bounded}.

\subsubsection{Hits without factoring a scheduled integer trace}

The following trace-zero geometry is classical. Trace-zero cosets in
$\F_{p^3}^*/\F_p^*$ form the planar Singer difference set; its
norm-one polynomial description is given explicitly in
\cite[Section 1.2 and Corollary 1.2]{MRS}. We prove the needed
root count directly and make no novelty claim for that ingredient.

\begin{theorem}[A finite auxiliary-prime test]\label{R15:thm:singer}
Suppose $p\ge5$ is prime, $S=p^2+p+1$ is prime,
$\ord_S(3)=S-1$, and $f_a$ is irreducible modulo $p$.
Each clock has exactly $p+1$ trace-zero phases modulo $S-1$.
Consequently each clock has infinitely many composite floors for
every integer $a>8$ in that coefficient residue class. The two sets
of phases need not intersect.
\end{theorem}
\begin{proof}
The projective order of $x$ divides the prime $S$ and is not $1$,
since an irreducible cubic root is not scalar. Thus it equals $S$.
The homomorphism $z\mapsto z^{p-1}$ has kernel $\F_p^*$,
so $y=x^{p-1}$ generates the norm-one subgroup of order $S$.
For every integer $n$,
\begin{equation}\label{R15:eq:zeropoly}
 \frac{\Tr(x^n)}{x^n}=1+y^n+y^{(p+1)n}.
\end{equation}
The polynomial $F(Z)=Z^{p+1}+Z+1$ has exactly $p+1$ distinct
roots, all in that subgroup. Indeed, a root $z$ is nonzero; raising
its equation to the $p$th power and multiplying by $z$ gives
$z^{p^2+p+1}=-z^{p+1}-z=1$. These roots belong to $\F_{p^3}$
because $S\mid p^3-1$. Also $F'(Z)=Z^p+1$; a common root of
$F$ and $F'$ would give $1=0$. The roots are therefore distinct.

As $y$ generates the group, exactly $p+1$ classes $n\pmod S$
give trace zero. Class $0$ does not, because $\Tr(1)=3\ne0$.
Class $-1$ does not either: in this cubic
$\Tr(x^{-1})=(-4)/(-4)=1$.
Over one index period $S-1$, $3^k$ runs through every nonzero
residue modulo $S$. The values $2\cdot3^k-1$ run through every
residue except $-1$, since $S$ is odd. Each orbit therefore visits
all $p+1$ zero classes. Projective periodicity proves the asserted
phase counts and recurrence. Only a finite prefix can fail to exceed $p$.
\end{proof}

For a general irreducible cubic with nonzero norm, the same argument
still guarantees hits on both clocks under the two assumptions on $S$:
the negative clock omits at most one of the $p+1>1$ zero classes.
The exact negative-clock count above uses the nonzero inverse trace of
our specific cubic. The hypotheses are finite, checkable conditions;
we do not assume or prove infinitely many primes satisfying them.
In particular no primitive-root conjecture is used.

\subsubsection{The exact same-prime simultaneous count}

\begin{theorem}[Common phases]\label{R15:thm:common}
Under the hypotheses of Theorem~\ref{R15:thm:singer}, the number of
phases modulo $S-1$ where both clocks have zero trace modulo $p$ is
$2$ if $a$ is a nonzero square modulo $p$, $0$ if it is a nonsquare,
and $1$ if $a\equiv0\pmod p$.
\end{theorem}
\begin{proof}
Put $z=x^{3^k}$. The conditions are exactly
$\Tr(z)=0$ and $\Tr(z^2/x)=0$, with no reindexing.
The trace-zero subspace has basis
\[
 v=3x-a,\qquad w=3x^2-(a^2+8),
\]
since $p\ne3$ and the $x,x^2$ coefficients show independence.
For $z=uv+tw$, its second condition is
\[
 A u^2+2But+C t^2=0,
\]
where direct trace calculation gives
\[
 A=a(a-9),\quad B=a(a^2-3a+8),\quad
 C=a^4+3a^3+16a^2+60a-44.
\]
These formulas follow from $\Tr(x^{-1})=1$ and the first four
power sums. Their quadratic discriminant, up to the square factor $4$, is
\begin{equation}\label{R15:eq:quadratic-discriminant}
 B^2-AC=36a(a^3+a^2+18a-11)=\frac94 a\Delta(f_a).
\end{equation}
For an irreducible cubic over an odd finite field, $\Delta(f_a)$
is a nonzero square: Frobenius cyclically permutes the roots
$x,x^p,x^{p^2}$, an even permutation, so it fixes their nonzero
Vandermonde product. Its square is the discriminant.
For $a\ne0$ the form is nondegenerate. The elementary binary
quadratic criterion, including a possible root at infinity, gives
two projective zeros exactly for square $a$, and none otherwise.
Indeed, when $A\ne0$, multiplying the form by $A$ gives
$(Au+Bt)^2-(B^2-AC)t^2$; when $A=0$, nondegeneracy forces
$B\ne0$ and the form factors as $t(2Bu+Ct)$ with two distinct lines.
For $a=0$ the form is $-44t^2$, nonzero because $p\ge5$ and
$p=11$ would give $\Delta(f_0)=-176=0$, contradicting
irreducibility. It has one projective zero.

Finally, $[x]$ generates the projective group of prime order $S$,
and primitive $3$ makes $[x^{3^k}]$ visit every projective point
except the scalar identity exactly once per period. The omitted
point has trace $3\ne0$, so every projective zero of the form is
visited. The two equations scale by a nonzero scalar and its square,
respectively, so this projective count is legitimate.
\end{proof}
The zero count concerns this particular auxiliary prime, not all
possible divisors or the complete sets of composite indices.

\subsection{The prime 173 and complete coefficient residue cells}

Take $a=7^{12}$. Modulo $p=173$, its value is $81$.
Exhaustive evaluation shows that $f_{81}$ has no root in $\F_{173}$,
hence is irreducible. The integers $173$ and $S=30103$ are prime by
trial division. Moreover
\[
 \begin{gathered}
 S-1=30102=2\cdot3\cdot29\cdot173,
 \qquad 3^{30102}\equiv1\pmod{30103},\\
 3^{30102/r}\not\equiv1\pmod{30103}
 \qquad(r=2,3,29,173).
 \end{gathered}
\]
Thus Theorem~\ref{R15:thm:singer} applies without factoring any integer
trace. Exact ring or companion-matrix calculations give the states
in Table~\ref{R15:tab:173}. Each tuple denotes coefficients in $1,x,x^2$.
\begin{table}[ht]
\centering
\begin{tabular}{rlll}
\toprule
Original index $k$ & Classical state & Negative state & Verified zero traces\\
\midrule
318 & $(143,48,33)$ & --- & Classical\\
421 & --- & $(99,114,107)$ & Negative\\
14733 & $(95,89,11)$ & $(20,140,31)$ & Both\\
25911 & $(18,43,50)$ & $(3,142,157)$ & Both\\
\bottomrule
\end{tabular}
\caption{Exact states modulo $173$ at their original clock indices.}
\label{R15:tab:173}
\end{table}
The complete scan has $174$ zero phases for each clock; their intersection
in $1\le k\le30102$ is exactly $\{14733,25911\}$.

\begin{corollary}\label{R15:cor:173}
For $a=7^{12}$, both actual floors are composite at every
\[
 k=14733+30102t\quad(t\ge0),
\]
and also at every $k=25911+30102t$.
\end{corollary}
\begin{proof}
The exact identities $x^{30103}=-4$ and
$3^{30102}\equiv1\pmod{30103}$ propagate each displayed zero
by nonzero scalar multiplication. Every scheduled floor exceeds
$a^3-1>173$, so properness holds already at the displayed seeds.
\end{proof}
The scalar $-4$ has order $43$ modulo $173$. Full-state order therefore
retains information unnecessary for divisibility of its trace.
The much larger previously available return period is not invalidated;
this certificate simply supplies a shorter one.

The order of $7$ modulo $173$ is $172$. The cubic is irreducible at
$a=7^m\pmod{173}$ precisely for the following $60$ residue classes:
\begin{equation}\label{R15:eq:cells}
\begin{gathered}
2,3,4,5,6,11,12,13,16,17,18,20,24,27,28,34,36,39,46,50,\\
53,54,57,58,64,65,67,70,71,74,81,83,87,90,94,101,102,103,\\
104,106,107,110,113,115,121,124,126,129,136,137,138,140,142,\\
144,152,153,154,160,166,170 \pmod{172}.
\end{gathered}
\end{equation}
For every $m\ge2$ in these classes, both clocks have infinitely many
composite floors. The phases can depend on the coefficient cell.
The common phases in Corollary~\ref{R15:cor:173} extend specifically to
$m\equiv12\pmod{172}$, not to all $60$ cells. The list is a complete
finite residue calculation, not a sample extrapolated to all $m$.
Since $7$ has order $172$, it is a nonsquare modulo $173$.
Exactly $36$ listed classes are even:
\[
\begin{gathered}
2,4,6,12,16,18,20,24,28,34,36,46,50,54,58,64,70,74,\\
90,94,102,104,106,110,124,126,136,138,140,142,144,152,\\
154,160,166,170\pmod{172}.
\end{gathered}
\]
By Theorem~\ref{R15:thm:common}, each of these classes has two same-prime
simultaneous phases modulo $30102$. All indicated floors are already
proper multiples, since $7^m\ge49$ for $m\ge2$.
The other $24$ irreducible classes have no common zero phase for
this prime, but still have $174$ zero phases on each clock.
This does not rule out simultaneous composite floors using other primes.

\subsection{Exact criteria for an arbitrary fixed modulus}

The positive results above are sufficient conditions. We now give a complete
local decision criterion for a fixed divisor, distinguishing transient
states, recurrent states, and trace outputs.

\subsubsection{Nilpotent and unit components}

\begin{lemma}\label{R15:lem:decompose}
For $R_D$ and $x$ as above, there are an exact ordinary-power preperiod
$\nu\ge0$ and an exact eventual period $O\ge1$ such that
$x^{n+O}=x^n$ for $n\ge\nu$. They arise from a decomposition into a
nilpotent component and a unit component, and satisfy
\[
 \nu\le3\max_{p^e\parallel D}e,\qquad O\le D^3.
\]
Use $\nu=0$ if the nilpotent component is absent, and $O=1$ if the
unit component is the zero ring.
\end{lemma}
\begin{proof}
First work modulo $p^e$. The descending images
$R\supseteq xR\supseteq x^2R\supseteq\cdots$ are additive
$p$-groups inside a group of order $p^{3e}$. Each strict inclusion
reduces cardinality by at least $p$. Once two consecutive images
agree, all subsequent ones agree. Stabilization therefore occurs
by $N=3e$. Let $I=x^N R$ be a stable image. Multiplication by $x$
is bijective on $I$, so multiplication by $x^N$ is also bijective there.
For every $r\in R$, choose $i\in I$ with $x^N i=x^N r$.
This proves
\[
 R=\ker(x^N\cdot)\oplus I.
\]
The intersection is zero by bijectivity on $I$. Both summands are
ideals, and their products vanish, since elements of $I$ have the
form $x^Nz$. Splitting $1$ gives complementary idempotents and
thus a direct product of rings. On the first factor $x$ is nilpotent;
on the second it is a unit. Combine these decompositions by CRT.
Take $\nu$ to be the nilpotence index on the combined nilpotent
factor and $O$ the order on the unit factor. The bounds follow from
$N\le3\max e$ and the total ring size $D^3$. These are exact
preperiod and period: on a nilpotent factor an equality
$x^{n+h}=x^n$ implies $x^n=0$, because $1-x^h$ is a unit; on
the unit factor equality is equivalent to $O\mid h$.
\end{proof}

\begin{theorem}[Complete eventual zero criterion]\label{R15:thm:local}
Write $O=3^sM$, with $3\nmid M$, and use $\ord_1(3)=1$.
The exact eventual state periods in the original clock index are
\begin{equation}\label{R15:eq:periods}
 L_+=\ord_M(3),\qquad
 L_-=\ord_{M/\gcd(M,2)}(3).
\end{equation}
A clock state at $k\ge1$ is recurrent exactly when
$n_\pm(k)\ge\nu$ and $k\ge s$.
Choose any $K\ge1$ satisfying both conditions, and set, for its clock,
\[
 e_j=\nu+\bigl((n_\pm(K+j)-\nu)\bmod O\bigr),\qquad
 Z_D=\{0\le j<L_\pm:\tau(x^{e_j})=0\pmod D\}.
\]
Then $D$ divides infinitely many scheduled traces if and only if
$Z_D\ne\varnothing$. More precisely, after $K$ the zero indices
are exactly the union of the classes $k\equiv K+j\pmod{L_\pm}$
for $j\in Z_D$.
\end{theorem}
\begin{proof}
Once the nilpotent component has vanished, equality after an index
increment $h$ is equivalent to divisibility by $O$ of, respectively,
\[
 3^k(3^h-1),\qquad2\cdot3^k(3^h-1).
\]
Their $3$-adic valuation is exactly $k$, since $3\nmid3^h-1$.
Thus no unit-state return is possible for $k<s$. For $k\ge s$,
the remaining divisibility is exactly $M\mid3^h-1$ in the
classical case and $M/\gcd(M,2)\mid3^h-1$ in the negative case.
This proves \eqref{R15:eq:periods} and the recurrent-state threshold,
using the last observation of Lemma~\ref{R15:lem:decompose} for the
nilpotent part. Reduction of every exponent to the interval
$[\nu,\nu+O)$ gives the stated states $x^{e_j}$. A finite eventual
cycle has infinitely many zero outputs exactly when it has one.
\end{proof}

The trace-output period can be smaller than the state period.
A transient zero need not recur, and two isolated zero traces do not
by themselves establish an infinite return. The theorem checks the
whole eventual cycle. Powers such as $3^{K+j}$ are evaluated modulo
$O$; no enormous integer trace is required. By Lemma~\ref{R15:lem:floor},
an infinite zero set gives infinitely many proper composite floors
for this fixed divisor. Finitely many zero indices do not.

For odd coefficients all positive traces are odd by \eqref{R15:eq:rec},
so useful divisors are odd and $x$ is a unit. For even $a$, all
positive traces are even, and the scheduled floors are already
composite. The nonunit discussion ensures completeness but is not
a missing obstacle in the odd-coefficient case.

\subsubsection{The sharper projective criterion and uniform exponents}

For odd $D$, write $Q_D=3^qN$, with $3\nmid N$.
The same cyclic-group proof applied modulo scalar units gives exact
projective preperiod $q$ and periods
\begin{equation}\label{R15:eq:projective-periods}
 J_+=\ord_N(3),\qquad J_-=\ord_{N/\gcd(N,2)}(3).
\end{equation}
For $K\ge\max(1,q)$, replace $Z_D$ by
\[
 W_D=\{0\le j<J_\pm:
 \tau(x^{n_\pm(K+j)\bmod Q_D})=0\pmod D\}.
\]
It gives the same necessary-and-sufficient infinite-divisibility
criterion and exactly all eventual zero phases. This can be much
shorter than a full-state computation. The zero pattern may still
have a smaller period than its projective state.
For composite $D$, nonempty zero sets at its prime-power factors
are not enough: their index phases must be simultaneously compatible.
The generalized CRT requires pairwise agreement modulo the gcds of
the phase periods, after the finite thresholds have passed.

\begin{lemma}[A uniform unit exponent]\label{R15:lem:universal}
For $p\ge3$ and $e\ge1$, every unit in a cubic algebra
$(\Z/p^e\Z)[X]/(f)$, with $f$ monic, is killed by
\[
 E_{p,e}=p^e(p^2-1)(p^2+p+1).
\]
For $p\ne3$,
\[
 v_3(E_{p,e})=
 \begin{cases}v_3(p+1),&p\equiv2\pmod3,\\
 v_3(p-1)+1,&p\equiv1\pmod3,
 \end{cases}
\]
and $v_3(E_{3,e})=e$.
\end{lemma}
\begin{proof}
Modulo $p$, the reduced algebra is a product of fields whose degrees
sum to at most $3$. Their unit orders divide
$H=(p^2-1)(p^2+p+1)$. Thus a unit $u$ has $u^H=1+z$, with
$z$ in the nilradical. That ideal has cube zero in a cubic algebra,
so $(1+z)^p=1$ for $p\ge3$. On lifting, the remaining factor
$1+pY$ is killed by $p^{e-1}$, by repeated binomial expansion.
This proves the exponent, without separability. The valuations follow
by reducing the factors modulo $3$; if $p\equiv1\pmod3$,
$p^2+p+1=3(1+3t+3t^2)$ for $p=1+3t$.
\end{proof}

A seed divisible by $p^e$ at $k_0\ge v_3(E_{p,e})$ therefore
has a full-state return. This is sufficient, not necessary.
It explains the safe seeds $k_0\ge1$ for $p\equiv2,5\pmod9$
and $k_0\ge2$ for $p\equiv4,7\pmod9$, while the projective
theorem can return earlier. For a fixed odd composite divisor,
take the lcm of these exponents over its prime-power factors.

\subsubsection{Exact coefficient-cell encoding}

For a fixed $p\ge5$, one can compute all recurrent-witness coefficient
cells without factoring a trace. Work in
\[
 B=\F_p[A]/(A^p-A),\qquad
 R_B=B[X]/(X^3-AX^2-4X+4).
\]
The algebra $B$ is the product of its $p$ coefficient cells.
Using the uniform exponent, start at $K\ge\max(1,v_3(E_{p,1}))$
and take a valid eventual period $L$ for the chosen clock.
Let $F_j(A)$ be the trace function at index $K+j$, represented
modulo $A^p-A$, for $0\le j<L$. Then
\begin{equation}\label{R15:eq:witnesspoly}
 H_{p,\pm}(A)=\gcd\left(A^p-A,\prod_{j=0}^{L-1}F_j(A)\right)
\end{equation}
has a residue $a$ as a root exactly when $p$ divides infinitely many
scheduled traces there. Indeed, evaluation in each field cell makes
the product zero exactly when some recurrent phase is zero.
The separate gcds with $F_j$ record its phases. Reducing after every
operation keeps representatives of degree below $p$.
For $a=7^m$ and $p\ne7$, evaluation on the finite subgroup generated
by $7$ gives a complete exponent-residue table. Backwards iteration
on the prime-to-$3$ cyclic part is also valid, using the inverse of
$3$ modulo its order, but does not place an arbitrary guessed zero
state on the desired orbit. Theorem~\ref{R15:thm:singer} supplies that
coverage in its special setting. None of these finite local decisions
proves the existence of a useful modulus for every coefficient.

\subsection{Finite horizons and unbounded witness primes}\label{R15:sec:bounded}

The next obstruction permits the witness prime to be arbitrarily large.
It is a standard resultant-specialization and congruence-avoidance
argument, applied to these particular trace polynomials.
Put $C_k(A)=P_{3^k}(A)$, including $C_0(A)=A$.

\begin{lemma}\label{R15:lem:eisenstein}
For $k\ge1$, $C_k$ is Eisenstein at $3$. Hence the polynomials
$C_0,C_1,\ldots$ are pairwise relatively prime over $\Q[A]$.
\end{lemma}
\begin{proof}
For the formal roots $\alpha_i$, put
$Q_n(A)=\sum_{i<j}(\alpha_i\alpha_j)^n\in\Z[A]$.
Newton's cubing identity is
\[
 P_{3n}=P_n^3-3P_nQ_n+3(-4)^n.
\]
It gives $C_k(A)\equiv A^{3^k}\pmod3$.
At $A=0$, $C_0(0)=0$ and every $C_k(0)$ is divisible by $3$;
the same identity gives $C_{k+1}(0)\equiv6\pmod9$.
Together with monicity this proves Eisenstein's criterion.
The irreducibles $C_k$ have distinct degrees $3^k$ (including
$C_0=A$), so they are pairwise coprime.
\end{proof}

\begin{theorem}[Pairwise coprime finite prefixes]\label{R15:thm:horizon}
For every $N\ge1$ there is an integer $L_N$ divisible by $12$
such that, for every positive multiple $m$ of $L_N$,
\[
 \gcd(C_i(7^m),C_j(7^m))=1\qquad(0\le i<j\le N).
\]
There are also infinitely many prime coefficients $a>2^{41}$ with
this same pairwise-coprimality property. These prime coefficients
are a separate family, not powers of $7$.
\end{theorem}
\begin{proof}
Let $S_N$ be the finite set of prime divisors of
\[
 R_N=\prod_{0\le i<j\le N}|\Res_A(C_i,C_j)|.
\]
Every factor is a nonzero integer by Lemma~\ref{R15:lem:eisenstein}.
If a prime divides both specialized values $C_i(a),C_j(a)$,
their reductions share the root $a\pmod p$, so their Sylvester
determinant vanishes modulo $p$. Thus that prime belongs to $S_N$.
At $A=1$ there is the integer factorization
\[
 X^3-X^2-4X+4=(X-1)(X-2)(X+2),
 \qquad P_n(1)=1\quad(n\text{ positive odd}).
\]
Choose
\[
 L_N=\lcm\bigl(12,\{\ord_p(7):p\in S_N,\ p\ne7\}\bigr).
\]
For every positive multiple $m$ of $L_N$, all selected traces are $1\pmod p$
for $p\in S_N\setminus\{7\}$.
Prime $7$ requires care: it divides $C_0(7^m)$, but no later
classical trace. In the coefficient-$0$ quotient modulo $7$,
\[
 u_1=x^3=(3,4,0),\quad u_2=u_1^3=(2,0,4),\quad u_2^3=u_1.
\]
Their traces under $\tau(c)=3c_0+8c_2$ are $2$ and $3$.
Determinism gives these nonzero values for every $k\ge1$.
Thus $7$ cannot divide two distinct selected traces. All possible
common prime divisors have been excluded.

For the prime-coefficient assertion set $Q_N=2\prod_{p\in S_N}p$.
Every $a\equiv1\pmod{Q_N}$ avoids all possible common primes.
There are infinitely many primes in this particular residue class,
by the elementary cyclotomic lemma proved next. Take them larger
than $2^{41}$ and apply Lemma~\ref{R15:lem:floor}.
\end{proof}

\begin{lemma}[Prime supply in the residue class one]\label{R15:lem:cyclotomic}
For every integer $Q\ge2$, infinitely many primes satisfy
$q\equiv1\pmod Q$.
\end{lemma}
\begin{proof}
Given a finite set of primes $S$, choose a sufficiently large positive
integer $M$ divisible by $Q\prod_{p\in S}p$. The monic cyclotomic
polynomial $\Phi_Q$ has positive degree and constant term $1$,
so $\Phi_Q(M)>1$. A prime divisor $q$ of this value divides
neither $M$ nor $Q$, and is not in $S$, since the value is $1$
modulo every prime dividing $M$.
Over $\F_q$, the factorization
$X^Q-1=\prod_{d\mid Q}\Phi_d(X)$ is squarefree because $q\nmid Q$.
The residue $M$ is a root of $\Phi_Q$, so its multiplicative order
is exactly $Q$. A smaller order $d\mid Q$ would make it a root
of another factor of this squarefree product. Thus $Q\mid q-1$.
Since the finite forbidden set was arbitrary, the supply is infinite.
\end{proof}
This standard argument is only for residue class $1$; it is not a
proof of the full theorem for arbitrary reduced progressions.

\begin{corollary}[No simultaneous uniform seed and gap bound]
\label{R15:cor:bounded}
Given $B\ge0$ and $H\ge1$, infinitely many positive $m$ divisible
by $12$ admit no prime $p$, seed $0\le k\le B$, and gap
$1\le h\le H$ satisfying
$p\mid C_k(7^m)$ and $p\mid C_{k+h}(7^m)$.
\end{corollary}
\begin{proof}
Apply Theorem~\ref{R15:thm:horizon} with $N=B+H$.
\end{proof}
There is no bound on $p$ here. Even an adaptive recurring-prime
certificate cannot cover every such coefficient while bounding both
its first zero index and its return gap uniformly. This does not
bound all compositeness algorithms. Pairwise coprime values may all
be composite with unrelated factors. Two trace zeros are only a
necessary prerequisite for a two-hit recurring certificate, not a
sufficient state-return proof.

For a reproducible finite illustration, the six absolute resultants
at $N=3$ factor as
\[
\begin{array}{c|l}
(i,j)&|\Res(C_i,C_j)|\\\hline
(0,1)&2^2\cdot3\\
(0,2)&2^6\cdot3\cdot13\\
(0,3)&2^{18}\cdot3\cdot5\cdot53\cdot73\\
(1,2)&2^{18}\cdot3^3\\
(1,3)&2^{56}\cdot3^3\cdot5\cdot11\cdot727\cdot13789\\
(2,3)&2^{162}\cdot3^9.
\end{array}
\]
They give $L_3=216885240$ and $Q_3=332776921680060$.
The certificate checks determinants, factorizations, and orders, without
constructing $7^{L_3}$. The all-horizon result is the proof above,
not an extrapolation from these six computations.

\subsubsection{A divisor that occurs only in a transient}

For $m\equiv3\pmod{16}$, $a=7^m\equiv3\pmod{17}$, and
\[
 P_3(7^3)=40357711=17\cdot2373983.
\]
Modulo $17$,
$f_3=(X-5)(X^2+2X+6)$, the quadratic discriminant $14$ is a
nonsquare, and the cubic discriminant is $6\ne0$.
The classical states under cubing are
\[
\begin{array}{r|rrrrrrrrrrr}
k&0&1&2&3&4&5&6&7&8&9&10\\\hline
\tau(u_k)&3&0&12&11&5&6&12&11&5&6&12.
\end{array}
\]
The exact states at $k=2,\ldots,9$ are, in order,
\[
\begin{gathered}
(10,11,4),(15,0,8),(8,5,5),(12,7,14),\\
(13,8,14),(16,16,0),(3,10,11),(8,11,12),
\end{gathered}
\]
and $u_{10}=u_2$; also $u_1=(13,4,3)$.
Thus $17$ divides the classical clock exactly at original index $1$
and never again. This is a finite exact cycle proof, not a scan
extrapolated past its end. It shows why merely factoring a scheduled
term does not suffice. Its reducible cubic does not contradict
Theorem~\ref{R15:thm:irreducible}, and this example is not confined to
$12\mid m$.

\subsection{Projective orders along the actual powering orbit}

We now allow the coefficients to move, but only along one prescribed
cubic powering orbit. No coefficient-density statement is transferred
to that orbit. Write
\[
 f_0(X)=X^3-aX^2+bX-c,\qquad c\ne0,
\]
and let $f_r$ have root multiset $\alpha_i^{3^r}$, where the
$\alpha_i$ are the roots of $f_0$, with multiplicities. For a totally
real cubic Pisot number, these are still its minimal polynomials:
odd powers preserve distinct real conjugates; a power cannot be
rational since their absolute values differ.
With $f_r=X^3-a_rX^2+b_rX-c_r$, Newton's identities give
\begin{equation}\label{R15:eq:cubemap}
 (a,b,c)\longmapsto
 (a^3-3ab+3c,\ b^3-3abc+3c^2,\ c^3).
\end{equation}
The second expression follows by applying the same cubing identity
to the three pair products, whose sum, pair sum, and product are
$b,ac,c^2$. In the real negative-norm notation $b=-B,c=-q$, this is
\begin{equation}\label{R15:eq:negative-map}
 (a,B,q)\longmapsto
 (a^3+3aB-3q,\ B^3+3aBq-3q^2,\ q^3).
\end{equation}

\subsubsection{New polynomial discriminant divisors and nilpotent kernels}

Let $\Delta_r$ be the polynomial discriminant. Factoring the
pairwise differences of cubes gives
\begin{equation}\label{R15:eq:discriminant-map}
 \Delta_{r+1}=\Delta_r J_r^2,\qquad
 J_r=a_r^2b_r^2-b_r^3-a_r^3c_r.
\end{equation}
Indeed, elementary symmetric expansion identifies $J_r$ with
\[
 \prod_{i<j}(\alpha_i^{2\cdot3^r}
 +\alpha_i^{3^r}\alpha_j^{3^r}+\alpha_j^{2\cdot3^r}).
\]
A new polynomial-discriminant prime need not ramify in the fixed
number field: powered roots can collide modulo that prime or the
power basis can define a nonmaximal order. We include these collisions
rather than excluding every prime dividing $\Delta_r$.

Fix a prime $\ell\nmid c$ and put
\[
 R_0=\F_\ell[X]/(f_0),\quad x=X\bmod f_0,
 \qquad R_r=\F_\ell[Y]/(f_r),\quad y=Y\bmod f_r.
\]
There is a homomorphism
\begin{equation}\label{R15:eq:psi}
 \psi_r:R_r\longrightarrow R_0,\qquad y\longmapsto x^{3^r},
\end{equation}
with image $A_r=\F_\ell[x^{3^r}]$. To justify it without
separability, multiplication by $x$ has characteristic polynomial
$f_0$; its $3^r$th power has characteristic polynomial $f_r$,
including multiplicities. Cayley--Hamilton gives the required relation.
The image need not equal $R_0$.

\begin{lemma}\label{R15:lem:nilkernel}
The kernel of \eqref{R15:eq:psi} is a nilpotent ideal. The induced map
$R_r^*\to A_r^*$ is surjective, with an $\ell$-group kernel.
The same is true after quotienting by the embedded scalar group
$\F_\ell^*$.
\end{lemma}
\begin{proof}
Extend scalars to a finite splitting field. If a polynomial $h(y)$
is in the kernel, evaluation at each original root gives
$h(\alpha_i^{3^r})=0$. These powered roots are precisely the
geometric root support of $f_r$. Thus $h$ is divisible by every
distinct irreducible factor of $f_r$, and its class is in the
nilradical. In a cubic quotient this ideal has cube zero, so the
kernel is nilpotent, even with original repeated roots or new collisions.
To lift a unit from $A_r$, lift it and its inverse arbitrarily;
their product is $1+n$ for a kernel element $n$. This is invertible
by a finite geometric series, proving surjectivity. The kernel on
units is $1+\ker\psi_r$, an $\ell$-group since the ideal has
cardinality a power of $\ell$. If a unit maps to a scalar, dividing
by that scalar puts it in this same kernel. This proves the
projective assertion as well.
\end{proof}

Let $Q_{\ell,r}$ be the order of $y$ in $R_r^*/\F_\ell^*$,
with $x$ for $r=0$.
\begin{theorem}[Exact projective transfer]\label{R15:thm:transfer}
For every $r\ge0$, the quotient
\begin{equation}\label{R15:eq:orderquotient}
 \frac{Q_{\ell,r}}{Q_{\ell,0}/\gcd(Q_{\ell,0},3^r)}
\end{equation}
is a nonnegative integral power of $\ell$. It is at most $\ell$
if $\ell\ge3$ and at most $4$ if $\ell=2$.
In particular, for $\ell\ne3$,
\begin{equation}\label{R15:eq:valuation-transfer}
 v_3(Q_{\ell,r})=\max\{v_3(Q_{\ell,0})-r,0\}.
\end{equation}
Moreover, for every polynomial $h$,
\begin{equation}\label{R15:eq:trace-transfer}
 \tau_r(h(y))=\tau_0(h(x^{3^r})),
\end{equation}
where both traces are regular multiplication traces in their cubic
algebras.
\end{theorem}
\begin{proof}
Equality of an element of $A_r$ to a scalar is the same equality
inside $R_0$. The projective order of $x^{3^r}$ in $A_r$ is
therefore $Q_{\ell,0}/\gcd(Q_{\ell,0},3^r)$.
Restrict the surjection in Lemma~\ref{R15:lem:nilkernel} to the cyclic
subgroup generated by $[y]$. Its order divided by its image order
is the order of a subgroup of an $\ell$-group, proving
\eqref{R15:eq:orderquotient}. If $n^3=0$, then $(1+n)^\ell=1$
for $\ell\ge3$, and $(1+n)^4=1$ for $\ell=2$.
These bounds apply to that cyclic kernel. Taking $3$-adic valuations
proves \eqref{R15:eq:valuation-transfer}. Finally both sides of
\eqref{R15:eq:trace-transfer} sum $h(\alpha_i^{3^r})$ over the same
three-root multiset, including multiplicities. No general invariance
of regular trace under arbitrary quotient maps is asserted.
\end{proof}

Over $\F_7$, $f_0=X^3-1$ has $Q_{7,0}=3$, whereas after cubing
$f_1=(Y-1)^3$ has $Q_{7,1}=7$. The image of $y$ is $1$, but
its nilpotent thickening contributes the extra factor $7$.
Over $\F_3$ the same construction has $Q_{3,0}=Q_{3,1}=3$,
so \eqref{R15:eq:valuation-transfer} fails in characteristic $3$.
The exclusion is essential.

\subsubsection{The original seed index does not reset}

Suppose $\ell\ne3$ and a classical seed in the powered algebra
has trace zero at relative index $s\ge0$. Trace transfer says
\[
 \tau_r(y^{3^s})=\Tr(\beta^{3^{r+s}})\pmod\ell.
\]
A scalar return after some positive index increment $h$ exists if
and only if
\[
 Q_{\ell,r}\mid3^s(3^h-1)\ \text{for some }h\ge1
 \quad\Longleftrightarrow\quad v_3(Q_{\ell,r})\le s.
\]
The reverse implication uses an order of $3$ modulo the prime-to-$3$
part of $Q_{\ell,r}$. The transfer theorem gives the exact equivalence
\begin{equation}\label{R15:eq:threshold}
 v_3(Q_{\ell,r})\le s
 \quad\Longleftrightarrow\quad v_3(Q_{\ell,0})\le r+s.
\end{equation}
Thus the powered-state scalar certificate is the original certificate
at its original total index, not a newly reset threshold. Actual
periods may acquire prime-to-$3$ factors, as the $\F_7$ example
shows. Non-scalar trace-zero returns can also exist; the displayed
criterion characterizes scalar returns only. More general finite-state
zero sequences still transfer exactly to the same original classical
tail. For the negative clock, powering gives exponent
$3^r(2\cdot3^s-1)$, not in general $2\cdot3^{r+s}-1$;
no identical negative-clock reindexing is claimed.

\subsubsection{First collisions and a factorization-free gcd test}

Assume $\Delta_0\ne0$. Let $\ell\nmid6c\Delta_0$ be prime,
and let $r\ge1$ be its first collision stage, so
$\ell\mid\Delta_r$ but $\ell\nmid\Delta_{r-1}$.
Some ratio of two distinct original reduced roots then has exact
order $3^r$. The original reduction is separable.

\begin{proposition}\label{R15:prop:collision}
The original splitting type imposes the following restrictions:
\begin{enumerate}[label=\textup{(\roman*)}]
\item In type $1+1+1$, $\ell\equiv1\pmod{3^r}$.
\item In type $1+2$, $\ell\equiv-1\pmod{3^r}$ and
      $v_3(Q_{\ell,0})=r$.
\item In type $3$, necessarily $r=1$ and $\ell\equiv1\pmod3$;
      no classical trace $T_{3^j}$ with $j\ge1$ vanishes modulo $\ell$.
\end{enumerate}
\end{proposition}
\begin{proof}
In the split case the ratio belongs to $\F_\ell^*$, proving (i).
In the irreducible case its norm is $1$, so its order divides
$\ell^2+\ell+1$, whose $3$-part is at most $3$. A collision is
therefore possible only as in (iii). Frobenius cyclically permutes
all three roots, so a collision of two powered roots makes all
three coincide. Every subsequent classical trace is $3t$ with
$t\ne0$, hence is nonzero. Since the collision is at stage $1$,
this proves the assertion for every $j\ge1$.

For mixed type, let $\alpha\in\F_\ell^*$ be the linear root
and $\lambda,\lambda^\ell$ the quadratic pair. If the colliding
ratio is $\lambda/\lambda^\ell$, its order divides $\ell+1$.
If it is $\lambda/\alpha$, a pure $3$-power-order element outside
$\F_\ell$ is possible only for $\ell\equiv-1\pmod3$:
when $\ell\equiv1\pmod3$, the entire $3$-primary subgroup of
$\F_{\ell^2}^*$ lies in $\F_\ell^*$, a contradiction.
For $\ell\equiv-1\pmod3$, the $3$-part lies in $\ell+1$;
in either case $3^r\mid\ell+1$.
The projective group of $\F_\ell\times\F_{\ell^2}$ identifies
with $\F_{\ell^2}^*$ by dividing the second component by the
first. Thus $Q_{\ell,0}=\ord(y)$ for $y=\lambda/\alpha$.
A cross-pair collision has $\ord(y)=3^r$. A quadratic-pair
collision has $\ord(y^{1-\ell})=3^r$; since
$3\nmid1-\ell$, this also gives $v_3(\ord(y))=r$.
\end{proof}

\begin{corollary}[An adaptive gcd certificate]\label{R15:cor:gcdcert}
If for some $j\ge0$ a prime $\ell$ satisfies
\[
 \ell\nmid6c\Delta_0,\quad \ell\equiv2\pmod3,\quad
 \ell\mid\gcd(T_{3^j},\Delta_j),
\]
then it divides infinitely many classical traces. For a growing
Pisot trace sequence whose late terms equal its floors, this gives
infinitely many composite floors.
\end{corollary}
\begin{proof}
Its first collision stage $r\le j$ must have mixed type by
Proposition~\ref{R15:prop:collision}. Hence $v_3(Q_{\ell,0})=r\le j$,
which is exactly the scalar-return threshold \eqref{R15:eq:threshold}
in the original algebra. The seed trace is zero, so the returns
are zero and eventually exceed the fixed prime.
\end{proof}

This admits a factorization-free sufficient test. Starting with the
positive integer $\gcd(T_{3^j},\Delta_j)$, repeatedly divide by
its gcd with $|6c\Delta_0|$ until it is coprime to that integer.
Call the remaining integer $H_j$. If $H_j\equiv2\pmod3$,
then some prime factor of $H_j$ is $2\pmod3$ and the corollary
applies. This procedure needs gcd computations, not a factorization
of the large trace. It is not necessary: an even number of such
prime factors can give $H_j\equiv1\pmod3$.
If a late classical trace sequence consists of increasing primes
$p_j\equiv-1\pmod3$, the same argument implies
\[
 \gcd(p_j,\Delta_j)=1\qquad\text{for every sufficiently large }j.
\]
Indeed, these prime values eventually avoid the fixed factors of
$6c\Delta_0$; a violation would force the same prime to divide
infinitely many strictly larger prime values. Neither this
restriction nor the gcd test proves that a witness ever exists for
a remaining candidate.

\subsection{The invariant boundary and finite witness sets}

The surface $a=1,B=q$ is invariant under \eqref{R15:eq:negative-map}:
\[
 (1,q,q)\longmapsto(1,q^3,q^3).
\]
Its polynomial factors as $(X-1)(X^2-q)$, so every positive odd
trace equals $1$ as an integer identity, including when reduced
factors happen to coincide.

\begin{proposition}\label{R15:prop:boundary}
For any fixed finite prime set $\mathcal P$, infinitely many prime
coefficients $a>2^{41}$ give cubics $X^3-aX^2-4X+4$ satisfying
\eqref{R15:eq:shape} whose actual powering coefficient triples obey
\[
 (a_r,B_r,q_r)\equiv(1,4^{3^r},4^{3^r})\pmod\ell
 \qquad(\ell\in\mathcal P,\ r\ge0).
\]
No prime in $\mathcal P$ divides any positive odd trace at any
powered stage.
\end{proposition}
\begin{proof}
Put $L=\lcm(24,\prod_{\ell\in\mathcal P}\ell)$.
Lemma~\ref{R15:lem:cyclotomic} supplies infinitely many prime
$a\equiv1\pmod L$; take them above the stated threshold.
Lemma~\ref{R15:lem:floor} supplies the irreducible real Pisot roots,
prime initial floor, and strict shape. Induction under the invariant
coefficient map proves the displayed congruence. At each stage the
odd traces are therefore $1\pmod\ell$ by the factorization.
\end{proof}

Letting the lower coefficients change along a prescribed orbit
therefore does not make a fixed finite divisor list universal.
The quantifiers are ``for every finite list, infinitely many prime-leading
parameters escape it'', not the existence of one parameter escaping
all primes. Nothing here says that the later traces are prime.
The scalar transfer theorem also does not eliminate any candidate by
itself. A useful remaining task is to find, for a given candidate,
an actual $\ell\nmid3c$ and index $j$ with
\[
 \ell\mid T_{3^j},\qquad v_3(Q_{\ell,0})\le j,
\]
or a recurrent trace-zero state not explained by scalar return.
An ambient coefficient-density theorem does not supply this statement
on one prescribed moving triple.

\subsection{Exact verification and scope}

All numerical certificates are in \path{certificates/}; the packaged
programs use only Python's standard library.
\begin{enumerate}
\item \path{fixed_modulus_projective_verify.py} checks full unit and
nonunit state criteria on finite grids, projective criteria, coefficient
witness polynomials, auxiliary primalities and orders, and the complete
$p=173$ phase and coefficient-cell tables.
\item \path{positive_projective_independent_verify.py} independently
uses companion matrices for the $p=173$ certificate, including scalar
returns, all $174$ phases on each clock, the two common phases, and
all $60$ irreducible coefficient cells.
\item \path{simultaneous_projective_quadratic_verify.py} checks
\eqref{R15:eq:quadratic-discriminant} coefficientwise, reconstructs the
quadratic form in irreducible cells, and compares projective zero counts
with complete small-prime clock scans.
\item \path{bounded_witness_verify.py} computes the six $N=3$
Sylvester determinants by exact fraction-free elimination, verifies
their factorizations and orders, and checks the $7$-cycle and the
transient-only $17$-orbit.
\item \path{powering_orbit_portable_verify.py} uses exact sparse
polynomial arithmetic for the coefficient and discriminant identities,
and finite quotient algebras for projective transfer, repeated-factor
examples, matched traces, first collisions, and the invariant boundary.
It is a standard-library transcription of the symbolic checks,
not a probabilistic replacement.
\item \path{powering_orbit_independent_verify.py} independently
checks powering transfer and first collisions with companion matrices,
including a late mixed example whose return period changes.
\end{enumerate}

Triple multiplication uses the defining monic cubic to reduce terms
of degree $4$ and $3$, in that order. Binary powering then evaluates
huge exponents in logarithmically many ring multiplications.
The independent matrix implementation uses ordinary $3$ by $3$
modular products instead. Primality assertions use exhaustive trial
division; order certificates check the exponent and its quotients by
every distinct prime factor. No probable-prime screen is accepted
as a proof. Return identities, not a finite numerical extrapolation,
prove the infinite progressions.

From this directory run
\begin{lstlisting}
python3 certificates/fixed_modulus_projective_verify.py
python3 certificates/positive_projective_independent_verify.py
python3 certificates/simultaneous_projective_quadratic_verify.py
python3 certificates/bounded_witness_verify.py
python3 certificates/powering_orbit_portable_verify.py
python3 certificates/powering_orbit_independent_verify.py
sha256sum -c SHA256SUMS.txt
\end{lstlisting}
The logs and JSON files record the finite calculations. The universal
local criteria, projective theorems, resultant obstruction, and
powering transfer rest on the proofs in this paper, not those grids.


The positive mechanism does not require factoring a scheduled integer
trace once a suitable auxiliary prime and irreducible cell have been
found. It does not prove an infinite supply of such primes, nor that
every coefficient meets one. The bounded-horizon obstruction leaves
unbounded adaptive indices and periods entirely possible. The
same-prime simultaneous count is not a statement about all actual
composite indices. No all-$m$ closure, prime pattern on a prescribed
orbit, global minimum result, or Mills transcendence theorem is proved.
The finite-field trace geometry, finite-module decomposition, resultant
avoidance, cyclotomic prime supply, and Cayley--Hamilton methods are
standard tools; their use here is not a historical-priority claim.



\endgroup

\section{Coset coverage and output prime order bounds}\label{src:R16}
Coset coverage removes an unnecessary primitive-root hypothesis in several explicit families. The subsequent order bounds concern actual hypothetical output primes and must not be inferred from arbitrary local models.
\medskip
\begingroup




\subsection{Two clocks and the exact floor bridge}

For an integer $a>8$, define
\[
 f_a(X)=X^3-aX^2-4X+4,\qquad
 n_+(k)=3^k,\quad n_-(k)=2\cdot3^k-1\quad(k\ge1).
\]
The signs label the classical and negative clocks; both exponents are
positive and odd. Their Newton power sums are the integer polynomials
\begin{equation}\label{R16:eq:newton}
 T_0=3,\quad T_1=a,\quad T_2=a^2+8,\qquad
 T_n=aT_{n-1}+4T_{n-2}-4T_{n-3}\quad(n\ge3).
\end{equation}
We write $T_n(a)$ when the coefficient must be displayed.

\begin{lemma}\label{R16:lem:floor}
The roots of $f_a$ are $-w,v,\beta$ with
$0<v<w<1$ and $a<\beta<a+1$. The polynomial is irreducible over $\Q$.
For every positive odd $n$,
\[
 \floor{\beta^n}=T_n(a),\qquad
 0<\beta^n-T_n(a)=w^n-v^n<1.
\]
These floors strictly increase with the odd exponent and tend to infinity.
\end{lemma}
\begin{proof}
The values of $f_a$ at $-1,0,1,a,a+1$ are respectively
$7-a,4,1-a,4(1-a),(a-1)^2$. The signs place a root in each of
$(-1,0),(0,1),(a,a+1)$. There are no integer roots, so the rational-root
test proves irreducibility. The sum of the roots gives
$w-v=\beta-a>0$. Taking odd power sums proves the floor identity.
Since $\beta>8$, one has $\beta^{n+2}-\beta^n>1$, proving strict increase.
\end{proof}

A recurring congruence $p\mid T_{n_\pm(k)}(a)$ therefore gives infinitely
many composite \emph{actual floors}, once the trace exceeds $p$. No
assumption of prime initial floors is needed for that implication. The thin
family $a=7^m$, $m\ge2$, is used below for explicit tests; its leading
coefficient is not prime, and we do not present it as an all-prime Mills
candidate.

\subsection{Complete coset coverage beyond a primitive-root hypothesis}

\subsubsection{Projective trace-zero classes}

Let $p\ge5$ be prime and put $S=p^2+p+1$. The quotient
\[
 \F_{p^3}^{*}/\F_p^{*}
\]
is cyclic of order $S$. Choose $x\in\F_{p^3}^{*}$ whose projective class
has order $S$, and define
\[
 D=\{n\in\Z/S\Z:\Tr_{\F_{p^3}/\F_p}(x^n)=0\}.
\]
A scalar multiple of $x^n$ has the same zero/nonzero trace status.
The trace is a nonzero linear map, since $\Tr(1)=3\ne0$. Its kernel is a
plane, so
\begin{equation}\label{R16:eq:Dproperties}
 |D|=(p^2-1)/(p-1)=p+1,\qquad 0\notin D,\qquad pD=D.
\end{equation}
Every nonscalar $\gamma$ is $c x^u$ for some $c\in\F_p^*$ and
$u\ne0$ in $\Z/S\Z$, and
\begin{equation}\label{R16:eq:traceexponent}
 \Tr(\gamma^n)=0\quad\Longleftrightarrow\quad un\in D.
\end{equation}
Because $[\F_{p^3}:\F_p]=3$ is prime, nonscalar elements are exactly
those having irreducible cubic minimal polynomial. These elementary
projective trace facts are the standard Singer construction; see
\cite[\S1.2 and Corollary 1.2]{MRS}.

Suppose now that $S$ is prime. Necessarily $p\equiv2\pmod3$, and every
nonscalar element has projective order $S$. Set
\[
 G=(\Z/S\Z)^*,\quad H=\langle3\rangle\subseteq G,
 \quad h=\ord_S(3),\quad A_-=(2H-1)\setminus\{0\}.
\]
For subsets of $G$, products and inverses below are taken in $G$.

\begin{theorem}[Exact coverage criterion]\label{R16:thm:coverage}
Define $W_+=DH^{-1}$ and $W_-=DA_-^{-1}$. For
$\gamma=cx^u$ with $u\in G$:
\begin{align*}
 \exists k\ge1:\Tr(\gamma^{3^k})=0
   &\quad\Longleftrightarrow\quad u\in W_+,\\
 \exists k\ge1:\Tr(\gamma^{2\cdot3^k-1})=0
   &\quad\Longleftrightarrow\quad u\in W_-.
\end{align*}
Every hit recurs with index gap $h$. Thus universal classical coverage is
exactly the condition that $D$ meet every multiplicative $H$-coset;
universal negative-clock coverage is exactly $W_-=G$.
\end{theorem}
\begin{proof}
The phases $k=1,\ldots,h$ run through all $3^k\in H$. Apply
\eqref{R16:eq:traceexponent}. An exponent zero modulo $S$ has scalar nonzero
trace, so it is correctly omitted from $A_-$. Replacing $k$ by $k+h$
changes each exponent by a multiple of $S$, hence changes its field value
only by a nonzero scalar. This preserves trace zero. Finally,
$DH^{-1}=G$ is equivalent to meeting every $H$-coset.
\end{proof}

Changing the chosen projective generator rescales $D$ and relabels $u$;
coverage is unaffected. Importantly, $W_-$ is a separate finite set: it is
not determined merely by the multiplicative coset counts for $D$.

\subsubsection{Exact certificates at 17, 59, and 89}

\begin{theorem}[Three complete finite certificates]\label{R16:thm:three}
At $p=17$, every irreducible cubic element has a recurring classical trace
zero. At each of $p=59$ and $p=89$, every irreducible cubic element has
recurring trace zeros on \emph{both clocks separately}. In all three cases
$3$ is not primitive modulo $S=p^2+p+1$.
At $p=17$, exactly $39$ of the $306$ projective generators have no
negative-clock zero.
\end{theorem}
\begin{proof}[Exact finite verification and its algebraic reduction]
Use a root of $f_{a_0}$ as reference element, with the data in
Table~\ref{R16:tab:coverage}. Trial division proves that both $p$ and $S$ are
prime. Testing the $p$ possible roots proves each reference cubic
irreducible. Exact quotient-ring powering gives $x^S=-4$, while $x$ is
nonscalar. Thus it has projective order $S$.

Enumerating $x^n$ for $0\le n<S$ gives all of $D$, not a sample.
Enumerating $H$ gives all coset counts and all elements of $W_\pm$.
The indicated cardinalities are exact; Theorem~\ref{R16:thm:coverage} then
proves every assertion. The portable verifier
\texttt{frobenius\_coset\_coverage\_verify.py} contains these complete loops;
its JSON retains $D$, all cosets, and every negative miss at $17$.
For example, the class $u=3$ is one of those misses.
\end{proof}

\begin{table}[ht]
\centering\small
\begin{tabular}{rrrrlrr}
\toprule
$p$&$S$&$a_0$&$h$&$H$-coset counts&$|W_+|$&$|W_-|$\\
\midrule
17&307&5&34&$2,2,2,2,2,2,2,2,2$&306&267\\
59&3541&0&1770&$33,27$&3540&3540\\
89&8011&0&2670&$39,24,27$&8010&8010\\
\bottomrule
\end{tabular}
\caption{Complete coverage, despite a proper subgroup $H=\langle3\rangle$.}
\label{R16:tab:coverage}
\end{table}

For precision, coset representatives in increasing least-unused order are
\[
\begin{array}{c|l}
p=17&1,2,4,5,7,10,14,20,28\\
p=59&1,2\\
p=89&1,7,13.
\end{array}
\]
Here is the finite algorithm underlying the table; all operations in its
first step are in the displayed cubic quotient:
\begin{lstlisting}
D = {n in range(S) : trace(x**n) == 0}
H = {3**k mod S : 0 <= k < ord_S(3)}
Aminus = {2*t-1 mod S : t in H} - {0}
Wplus  = {d * inverse(t) mod S : d in D, t in H}
Wminus = {d * inverse(t) mod S : d in D, t in Aminus}
\end{lstlisting}
Quotient multiplication is reduced using $x^3=a_0x^2+4x-4$, and
\begin{equation}\label{R16:eq:tracefunctional}
 \tau(c_0+c_1x+c_2x^2)=3c_0+a_0c_1+(a_0^2+8)c_2\pmod p.
\end{equation}
Thus the certificate requires neither an enormous integer trace nor its
factorization. Universal assertions here exhaust \emph{all nonidentity projective
classes}, not only coefficients in $f_a$.

\begin{corollary}\label{R16:cor:floors}
Let $a>8$ be an integer. If $f_a$ is irreducible modulo $59$ or $89$,
then both sequences $\floor{\beta^{n_+(k)}}$ and
$\floor{\beta^{n_-(k)}}$ contain infinitely many composite terms.
If it is irreducible modulo $17$, the same holds for the classical clock.
\end{corollary}
\begin{proof}
Apply Theorem~\ref{R16:thm:three} to the residue class of $X$, then
Lemma~\ref{R16:lem:floor}. A fixed divisor is proper for all sufficiently
large indices in its recurring phase.
\end{proof}

The four irreducible cells of $f_a$ modulo $17$ are
$a\in\{5,6,9,14\}$. A separate complete check shows that all four also
have negative-clock hits. This restricted-family fact does not remove the
$39$ generic projective misses in Theorem~\ref{R16:thm:three}.

The complete irreducible coefficient cells of the thin family are
\begin{align*}
 p=17:\quad&m\equiv6,11,13,15\pmod{16},\\
 p=59:\quad&m\equiv4,5,6,12,17,18,22,23,24,26\pmod{29},\\
 p=89:\quad&m\equiv6,11,12,15,16,17,19,22,25,30,31,32,\\
 &\hspace{25mm}36,42,46,52,55,58,59,60,62,72,74,80,81\pmod{88}.
\end{align*}
The moduli are the exact respective orders of $7$. For every $m\ge2$
in one of these cells, both clocks have infinitely many composite floors
separately. At $17$ this uses the additional fixed-family checks, rather
than the false generic negative-coverage assertion.

\begin{corollary}[An explicit simultaneous powering class]\label{R16:cor:m60}
For every $m\ge2$ with $m\equiv60\pmod{88}$, both floors associated with
$a=7^m$ are composite at every
\[
 k=49+2670t\quad\text{and every }k=179+2670t,
 \qquad t\in\Z_{\ge0}.
\]
\end{corollary}
\begin{proof}
The order of $7$ modulo $89$ is $88$, so $a\equiv73$ throughout this
class. The reduction $f_{73}$ is irreducible. At $k=49$, direct quotient
powering gives the respective states
\[
 x^{3^{49}}=(22,48,33),\qquad
 x^{2\cdot3^{49}-1}=(20,35,41).
\]
Both have trace zero by \eqref{R16:eq:tracefunctional}. The complete common
zero-phase set is $\{49,179\}$ modulo $2670$, verified by enumerating
one full projective period; $x^{8011}=-4$ and
$\ord_{8011}(3)=2670$ prove recurrence. Every positive $m$ in the
class is at least $60$, so $a>89$. Since $n_\pm(k)\ge3$, the floors
are already larger than $a>89$ at both displayed seeds. Hence the divisor
$89$ is proper from $t=0$ onward.
\end{proof}

The coefficient $m=60$ illustrates a concrete gain over two earlier
finite certificates in this research sequence: it is divisible by $12$,
and it is not an irreducible cell for the earlier auxiliary prime $173$.
Indeed $7^{60}\equiv22\pmod{173}$ and $f_{22}(161)\equiv0\pmod{173}$.
This comparison concerns only those certificates, not a historical claim
that this coefficient was previously unknown. Corollary~\ref{R16:cor:m60}
is a separately verified simultaneous class, not a theorem that every
irreducible coefficient at $89$ has a common zero phase.

\subsubsection{Short certificates for \texorpdfstring{$a=7^{12}$}{a = 7 to the twelfth power}; no common-phase inference}

At $p=59$, the coefficient is $a\equiv9$, and $f_9$ is irreducible.
The complete $1770$-phase cycle contains $33$ classical zero phases and
$31$ negative zero phases. The first positive phases and their states in
$(1,x,x^2)$ are
\[
 \begin{array}{c|c|c}
 \text{clock}&k&x^{n_\pm(k)}\\\hline
 +&36&(26,42,32)\\
 -&19&(31,42,2).
 \end{array}
\]
The trace functional \eqref{R16:eq:tracefunctional} vanishes in each case.
Consequently each clock has infinitely many composite floors, respectively
on $k\equiv36\pmod{1770}$ and $k\equiv19\pmod{1770}$.

The complete phase sets have \emph{empty intersection}. This is stronger
than noticing that the two selected progressions are disjoint: there is no
same-prime common zero phase at $59$ for this coefficient. Although $9$ is
a square, a common-phase theorem requiring a primitive-$3$ orbit cannot
be applied here. Other primes may still witness common composite indices.

An even smaller separate certificate occurs at $p=11$, where
$a\equiv5$, $x^{133}=7$, and $\ord_{133}(3)=18$:
\[
 \begin{array}{c|c|c}
 +&k\equiv16\pmod{18}&(10,5,6)\text{ at }k=16\\
 -&k\equiv1\pmod{18}&(5,8,7)\text{ at }k=1.
 \end{array}
\]
These are the complete zero phases at this prime. They too are separate,
not a simultaneous certificate. The positive assertions above do not
claim that finitely many auxiliary primes cover all $a$, or all $7^m$.

\subsection{What Frobenius enlargement and irreducibility do not supply}

\subsubsection{The full-group shortcut is exactly the primitive-root case}

Frobenius invariance $pD=D$ immediately gives a sufficient classical
criterion:
\begin{equation}\label{R16:eq:fullgroup}
 \langle3,p\rangle=G\quad\Longrightarrow\quad W_+=G.
\end{equation}
Indeed, for $d\in D$ and $u\in G$, write $d/u=3^kp^j$ and apply an
inverse Frobenius power. Then $u3^k\in D$; adding a period makes $k$
positive. This reasoning does not apply to the affine negative exponent:
$p(2t-1)=2pt-p$, not $2pt-1$.

\begin{proposition}[Cubic reciprocity limitation]\label{R16:prop:frob}
If $p\ge5$ and $S=p^2+p+1$ are prime, then
\[
 \langle3,p\rangle=(\Z/S\Z)^*
 \quad\Longleftrightarrow\quad
 3\text{ is primitive modulo }S.
\]
\end{proposition}
\begin{proof}
We first record the cubic supplement in the needed convention:
\begin{equation}\label{R16:eq:cubiclaw}
 3^{(S-1)/3}\equiv p^{(p+1)/3}\pmod S.
\end{equation}
This is a specialization of the standard supplements stated in
\cite[p.~291, equations (1.1)--(1.2)]{R16-Sun}. To see the specialization, let
$\omega^2+\omega+1=0$ and take the primary Eisenstein prime
$\pi=-1-(p+1)\omega$, of norm $S$. Modulo $\pi$ one has $\omega=p$.
The supplements give
\[
 \left(\frac{1-\omega}{\pi}\right)_3=1,
 \qquad
 \left(\frac{\omega}{\pi}\right)_3
      =\omega^{-(p+1)/3}.
\]
Use $3=-\omega^2(1-\omega)^2$ and the fact that $-1$ is a cube.
The cubic symbol of $3$ is $\omega^{(p+1)/3}$; its defining Euler
criterion gives \eqref{R16:eq:cubiclaw}.

Set $H=\langle3\rangle$ and $N=S-1=p(p+1)$. The residue $p$ has
order exactly $3$ modulo $S$. If $\langle3,p\rangle=G$ but $H\ne G$,
then $[G:H]=3$ and $p\notin H$. In the cyclic group $G$, an order
of $H$ divisible by $3$ would force $H$ to contain the unique subgroup
of order $3$, including $p$. Thus $|H|=N/3$ is prime to $3$, and
$v_3(N)=1$. The left side of \eqref{R16:eq:cubiclaw} is then $1$, whereas
its right side is not $1$, since $3\nmid(p+1)/3$. This contradiction
proves the nontrivial direction; the reverse is immediate.
\end{proof}

Frobenius can enlarge a proper subgroup without filling $G$; $p=17$
is an example. Proposition~\ref{R16:prop:frob} does not restrict the complete
$D$-coset coverage in the preceding coset-coverage subsection. That coverage is the reason the positive
certificates go beyond the full-group shortcut.

\subsubsection{Composite projective group order and missing strata}

For an irreducible cubic element $\gamma$, let $Q$ be its projective
order and put $y=\gamma^{p-1}$. Then $\ord(y)=Q\mid S$, and
\begin{equation}\label{R16:eq:tracenormone}
 \frac{\Tr(\gamma^n)}{\gamma^n}
 =1+y^n+(y^n)^{p+1}.
\end{equation}
Write $Q=3^eN$, $3\nmid N$. One has $e\le1$: when $p\equiv1\pmod3$,
$S\equiv3\pmod9$, and otherwise $3\nmid S$.
For every classical index $k\ge1$, $y^{3^k}$ has exact order $N$.
It follows that
\begin{equation}\label{R16:eq:gcdstratum}
 \gcd(Z^{p+1}+Z+1,\Phi_N(Z))=1\text{ in }\F_p[Z]
 \quad\Longrightarrow\quad\text{no classical zero in this stratum}.
\end{equation}
If this gcd is nonconstant and $\langle3,p\rangle=(\Z/N\Z)^*$,
there is conversely a classical zero: the group acts transitively on the
primitive $N$th roots, and trace zero is invariant under Frobenius.
This converse is sufficient, not necessary. Smaller exact coset coverage
may suffice. For $N=1$, the gcd in \eqref{R16:eq:gcdstratum} is $1$, since
$1+1+1\ne0$.

The general exact criterion remains \eqref{R16:eq:traceexponent}, now allowing
nonunit $u$. If $3\nmid S$, enumerate an exponent period
$\ord_S(3)$; if $3\mid S$, start at $k=1$ and use
$\ord_{S/3}(3)$. Both clocks are periodic modulo $S$ from that point.
One must not invert a nonunit exponent when $S$ is composite.

At $p=11$, $S=133=7\cdot19$. The reference $f_5$ has full projective
order and
\[
 D=\{2,5,16,18,22,43,50,55,65,73,74,109\}.
\]
Every member is coprime to $133$. Thus every trace-zero projective point
has full order $133$; the order-$7$ and order-$19$ strata contain none.
Both missing strata actually occur in the family:
\[
 \begin{array}{c|c|c}
 a\bmod11&\text{scalar identity}&\text{projective order}\\\hline
 2&x^7=8&7\\
 4&x^{19}=2&19.
 \end{array}
\]
Both cubics are irreducible. Enumerating one full projective period shows
that $11$ divides \emph{no ordinary power trace} in either cell.
For $a=7^m$, these are respectively $m\equiv3$ and $6\pmod{10}$.
This diagnoses an existence gap, not an obstruction to using other primes.

For contrast, irreducibility does guarantee recurrence of a clock zero
that already exists. Indeed, if $Q=3^eN$ with $e\le1$, take
$L=\ord_N(3)$, with $L=1$ for $N=1$. For $k\ge1$,
\[
 n_+(k+L)-n_+(k)=3^k(3^L-1),\qquad
 n_-(k+L)-n_-(k)=2\cdot3^k(3^L-1)
\]
are multiples of $Q$. The two corresponding states are scalar multiples.
Existence and recurrence are different assertions.

\subsection{Actual output primes: norm and cyclotomic lower bounds}

This section concerns a \emph{different, conditional} setting. Let $\beta$
be a fixed cubic Pisot number, with minimal polynomial
\[
 f(X)=X^3-aX^2+bX-c,\qquad
 c\ne0,\quad C=|c|,\quad\Delta=\operatorname{disc}(f)\ne0.
\]
Write its other conjugates as $\beta_2,\beta_3$. Then
\begin{equation}\label{R16:eq:Csize}
 1\le C=|\beta\beta_2\beta_3|<\beta.
\end{equation}
Assume that, from some index onward,
\begin{equation}\label{R16:eq:primetail}
 p_k=\Tr(\beta^{3^k})=\floor{\beta^{3^k}}
\end{equation}
are strictly increasing primes. We consider only the residue branch
$p_k\equiv-1\pmod3$, and primes not dividing $6c\Delta$. The latter
exclusion is automatic for sufficiently late increasing outputs.
We do not assert that all hypothetical candidates lie in this branch.

\subsubsection{Two independent components of projective order}

For a good output prime, let $Q_k$ be the order of the class of $X$ in
\[
 (\F_{p_k}[X]/(f))^*/\F_{p_k}^*.
\]
Write $Q_k=3^{w_k}q_k$, with $3\nmid q_k$.
If $w_k\le k$, a positive index gap $L$ with $3^L\equiv1\pmod{q_k}$
would make $X^{3^{k+L}}$ a nonzero scalar multiple of $X^{3^k}$.
The trace is zero modulo $p_k$ at the first exponent because its integer
value is $p_k$. It would then be zero at every such later exponent,
contradicting the increasing prime tail. Consequently
\begin{equation}\label{R16:eq:3height}
 w_k\ge k+1.
\end{equation}
An irreducible reduction has projective order dividing $p_k^2+p_k+1$,
whose $3$-part is at most $3$. A completely split reduction has order
dividing $p_k-1$, prime to $3$ in the stated branch. Thus every such
sufficiently late reduction is of mixed type $(1,2)$.

Let $p\equiv-1\pmod3$ be any good mixed prime, with linear root
$\alpha\in\F_p^*$ and quadratic roots $\lambda,\lambda^p$.
Dividing all components by $\alpha$ identifies the projective order with
\[
 Q=\ord(\lambda/\alpha)=3^wq,\qquad3\nmid q.
\]
Its norm component is
\[
 \Norm_{\F_{p^2}/\F_p}(\lambda/\alpha)=c/\alpha^3,
 \qquad g=\ord_{\F_p^*}(c/\alpha^3).
\]
Since $g\mid Q$ and $g\mid p-1$, one has
\begin{equation}\label{R16:eq:gdiv}
 3\nmid g,\qquad g\mid q.
\end{equation}
The large $3$-part in \eqref{R16:eq:3height} therefore cannot account for this
separate norm-component order.

\begin{theorem}[Elementary output-prime bound]\label{R16:thm:norm}
For any good mixed prime $p\equiv-1\pmod3$ in the fixed setting above,
\[
 g>\frac{\log p-\log8}{3\log\beta+2\log C}.
\]
For an actual output in \eqref{R16:eq:primetail}, once $\beta^{3^k}>2$,
\begin{equation}\label{R16:eq:growth}
 g_k>\frac{3^k\log\beta-\log16}{3\log\beta+2\log C}.
\end{equation}
If the indices in the stated residue branch are infinite, then along
those indices
\begin{align}
 \liminf\frac{q_k}{3^k}&\ge
 \frac{\log\beta}{3\log\beta+2\log C},\label{R16:eq:qliminf}\\
 \liminf\frac{Q_k}{9^k}&\ge
 \frac{3\log\beta}{3\log\beta+2\log C}.\label{R16:eq:Qliminf}
\end{align}
Limits here are along the relevant indices tending to infinity.
\end{theorem}
\begin{proof}
The order relation gives $\alpha^{3g}=c^g$ modulo $p$, so $p$ divides
\[
 R_g=\Norm_{\Q(\beta)/\Q}(\beta^{3g}-c^g)\in\Z.
\]
This integer is nonzero: $\beta^3>C$ and
$|\beta_i|^3<1\le C$ for the two small conjugates. Triangle inequalities
and \eqref{R16:eq:Csize} give
\[
 0<|R_g|\le(\beta^{3g}+C^g)
 \prod_{i=2}^3(|\beta_i|^{3g}+C^g)
 <8(\beta^3C^2)^g.
\]
Taking logarithms proves the first bound. For an actual output,
\eqref{R16:eq:primetail} implies
$p_k>\beta^{3^k}-1>\beta^{3^k}/2$, proving \eqref{R16:eq:growth}.
Finally $q_k\ge g_k$ and $Q_k\ge3^{k+1}g_k$ by
\eqref{R16:eq:3height}--\eqref{R16:eq:gdiv}. Divide and take lower limits.
\end{proof}

This is the standard elementary norm-height argument for an order of size
at least a constant times $\log p$. The substitution of the growing
\emph{output equality} $p_k=T_{3^k}$ is essential: it cannot be made for
an arbitrary fixed prime divisor of a large trace.
For checking the norm, if $U_n$ denotes the pair-product power sum, then
\begin{equation}\label{R16:eq:normidentity}
 R_g=c^g(c^gT_{3g}-U_{3g}),\qquad R_1=c(a^3c-b^3).
\end{equation}
Expansion of the three factors proves this: the first and last terms
cancel because $\Norm(\beta^{3g})=c^{3g}$.

If the small conjugates have moduli $\sigma<\rho\le\beta^{-17/40}$,
then $C=\beta\rho\sigma<\beta^{3/20}$. Thus the candidate-specific
constants in \eqref{R16:eq:qliminf} and \eqref{R16:eq:Qliminf} are strictly larger
than $10/33$ and $10/11$, respectively. These root-size assumptions are
compatible with the cubic totally real shape appearing in the algebraic
reduction of Saito \cite[Theorem 1.7]{Saito}; here they are explicit
hypotheses, not an additional reduction theorem.

\subsubsection{Exact order gives a totient bound}

Let $\Phi_j(X,Y)=Y^{\varphi(j)}\Phi_j(X/Y)$ be the homogeneous
cyclotomic polynomial, with $\Phi_1(X,Y)=X-Y$.
Since $\alpha^3/c$ has exact order $g$ and $p\nmid cg$, separability of
$Z^g-1$ modulo $p$ shows that
\begin{equation}\label{R16:eq:cyclonorm}
 p\mid N_g,\qquad
 N_g=\Norm_{\Q(\beta)/\Q}\Phi_g(\beta^3,c)\ne0.
\end{equation}
Indeed, $f(X)$ and $\Phi_g(X^3,c)$ share the root $\alpha$ modulo $p$.
No conjugate ratio $\beta_i^3/c$ has modulus $1$, so the norm cannot
vanish.

\begin{lemma}\label{R16:lem:cyclobound}
For $0<r<1$, put $F(r)=r/(1-r)^2$. If $|z|=r$, then, uniformly for
$j\ge1$,
\[
 \bigl|\log|\Phi_j(z)|\bigr|\le F(r).
\]
\end{lemma}
\begin{proof}
The M\"obius product expresses the logarithm as a signed sum of
$\log|1-z^d|$ over $d\mid j$. Each absolute summand is at most
\[
 -\log(1-r^d)\le\frac{r^d}{1-r^d}\le\frac{r^d}{1-r}.
\]
Summing over all positive integers instead of divisors gives $F(r)$.
For $j=1$, the sign in $\Phi_1(z)=-(1-z)$ has no effect.
\end{proof}

Define the fixed constant
\[
 E=F(C/\beta^3)+F(|\beta_2|^3/C)+F(|\beta_3|^3/C).
\]
Reciprocity in absolute value,
$|\Phi_j(z)|=|z|^{\varphi(j)}|\Phi_j(1/z)|$ for $|z|>1$,
and Lemma~\ref{R16:lem:cyclobound} yield
\[
 0<|N_g|\le e^E(\beta^3C^2)^{\varphi(g)}.
\]
Homogenization contributes $C^{3\varphi(g)}$, while the dominant ratio
contributes $(\beta^3/C)^{\varphi(g)}$; this explains the exact base.
For an actual output with $M=3^k$, we obtain
\begin{equation}\label{R16:eq:totientgeneral}
 \varphi(g)>\frac{M\log\beta-\log2-E}{3\log\beta+2\log C}.
\end{equation}
This improves an order bound to a bound on its totient. It is still only a
necessary condition.

\subsubsection{Cube norms and a unit-norm integer refinement}

Suppose $c=d^3$ for a nonzero integer $d$, and put $D=|d|=C^{1/3}$.
Cubing is an automorphism of $\F_p^*$ because $p\equiv-1\pmod3$.
Consequently
\[
 g=\ord((d/\alpha)^3)=\ord(d/\alpha)=\ord(\alpha/d).
\]
The smaller norm certificates are
\[
 p\mid\Norm(\beta^g-d^g),\qquad
 p\mid\Norm\Phi_g(\beta,d).
\]
They are nonzero since $\beta>D$ and $|\beta_i|<1\le D$. The same
triangle and cyclotomic estimates, now with
\[
 E_d=F(D/\beta)+F(|\beta_2|/D)+F(|\beta_3|/D),
\]
give at an actual output
\begin{equation}\label{R16:eq:cubenorm}
 \begin{split}
 g&>\frac{M\log\beta-\log16}{\log\beta+(2/3)\log C},\\
 \varphi(g)&>\frac{M\log\beta-\log2-E_d}
                     {\log\beta+(2/3)\log C}.
 \end{split}
\end{equation}
The relevant norm bounds are $8(\beta D^2)^g$ and
$e^{E_d}(\beta D^2)^{\varphi(g)}$.

\begin{corollary}[Unit norms, under explicit small-root bounds]
Assume $c=\varepsilon\in\{1,-1\}$, $\beta\ge16$, and
$|\beta_2|,|\beta_3|\le\beta^{-17/40}$. At each good mixed output of
\eqref{R16:eq:primetail} in the stated residue branch, with $k\ge1$,
\[
 \varphi(g_k)\ge3^k+1,\qquad g_k\ge3^k+2,
 \qquad Q_k\ge3\cdot9^k+6\cdot3^k.
\]
\end{corollary}
\begin{proof}
Take $d=\varepsilon$ in \eqref{R16:eq:cubenorm}. The small moduli are below
$1/3$, since $2^{17}>3^{10}$. As $F$ is increasing,
\[
 E_\varepsilon\le F(1/16)+2F(1/3)
 =\frac{16}{225}+\frac32=\frac{707}{450}<2\le\log8.
\]
For the last inequality, $\log2\ge2/3$ follows, for example, by
integrating or differentiating
$\log x-2(x-1)/(x+1)$ for $x\ge1$.
Hence $E_\varepsilon+\log2<\log\beta$.
With $M=3^k$, \eqref{R16:eq:cubenorm} gives $\varphi(g)>M-1$, so
$\varphi(g)\ge M$. The bound excludes $g\le2$; for $g>2$, its totient
is even, whereas $M$ is odd. Thus $\varphi(g)\ge M+1$ and
$g\ge\varphi(g)+1\ge M+2$. Finally
$Q\ge3^{k+1}g=3Mg$ by \eqref{R16:eq:3height}--\eqref{R16:eq:gdiv}.
\end{proof}

All indices are relative to the \emph{fixed chosen} Pisot number, and the
prime trace-floor tail must already have begun. Choosing a later cubic
power may arrange $\beta\ge16$, but shifts the clock. In particular,
replacing $\beta_0$ by $\beta_0^3$ makes its norm a cube, yet the new
certificate $\Phi_g(\beta_0^3,\Norm(\beta_0))$ is exactly the old
general certificate. Forgetting the index shift does not create an extra
improvement. Unit norms already are cubes at the original chosen power.

\subsubsection{Why a local model is not an actual output model}

For fixed $k\ge1$, choose $p\equiv\pm1\pmod{3^{k+1}}$ and an element
$\zeta$ of order $3^{k+1}$ in $\F_p^*$ in the plus case, or in the
norm-one subgroup of $\F_{p^2}^*$ in the minus case. The root multiset
\[
 1,\zeta,\zeta^{-1}
\]
has norm $1$, trace zero at exponent $3^k$, and trace $3$ at every later
classical exponent. Its projective $3$-height is exactly $k+1$.
The plus case is split; the minus case is mixed. These are legitimate
finite-field examples where a seed zero and the required height do not
force a later zero.

In the mixed case the norm-component order is only $g=1$. For a
\emph{fixed} irreducible Pisot candidate, however, this forces
$p\mid R_1=c(a^3c-b^3)$, a fixed nonzero integer. The displayed linear
root $1$ also forces $p\mid f(1)$. It cannot model arbitrarily late
increasing output primes of that candidate. Likewise, lifting coefficients
modulo a fixed prime usually produces a large trace divisible by that
prime, not an equality between the trace and that prime.

None of the lower bounds above contradicts $Q\mid p^2-1$. An actual
output is exponential in $3^k$, leaving far more room than the required
order scale. No inert output, recurring divisor, or transcendence
conclusion has been deduced from these necessary conditions.

\subsection{Ordinary divisors and sparse clocks are different conditions}

For a fixed coefficient and one clock, a sufficient inert-witness strategy
would require
\begin{equation}\label{R16:eq:IW}
 \exists p\ \exists k\ge1:
 \quad f_a\text{ is irreducible modulo }p,
 \quad p\mid T_{n_\pm(k)}(a).
\end{equation}
As proved in the preceding projective-return argument, such a zero recurs. The two clocks can
use different witnesses; \eqref{R16:eq:IW} for both does not itself give a
common index progression. Chebotarev can provide the irreducibility
condition, while a prime-factor theorem can provide divisors. Their
intersection on a fixed sparse clock is additional arithmetic input.

\subsubsection{The even powering family is not algebraically degenerate}

\begin{proposition}\label{R16:prop:S3}
For every even $m\ge2$, the polynomial $f_{7^m}$ has Galois group $S_3$,
and its three roots are multiplicatively independent. For every $m\ge2$,
the ordinary trace recurrence is nondegenerate of exact order $3$, as is
its odd-index subsequence.
\end{proposition}
\begin{proof}
Irreducibility follows from Lemma~\ref{R16:lem:floor}. The discriminant is
\[
 \Delta(a)=16D(a),\qquad D(a)=a^3+a^2+18a-11.
\]
For even $m$, write $a=b^2$ with $b=7^{m/2}$ odd. If $m\ge4$, then
$b\ge49$, and the consecutive integers
\[
 L=b^3+(b-1)/2,\qquad U=L+1=b^3+(b+1)/2
\]
satisfy
\begin{align*}
 4(D(b^2)-L^2)&=4b^3+71b^2+2b-45>0,\\
 4(U^2-D(b^2))&=4b^3-71b^2+2b+45>0.
\end{align*}
For the second inequality, $4b-71>0$ already at $b\ge18$.
At $m=2$, direct arithmetic gives
\[
 347^2=120409<D(49)=120921<121104=348^2.
\]
Thus the discriminant is nonsquare and the irreducible cubic has group
$S_3$.

The root magnitudes $\beta,w,v$ are distinct, so no quotient of distinct
roots is a root of unity. The three exponential coefficients in $T_n$
are nonzero, giving a nondegenerate recurrence of exact order $3$.
The same reasoning applies to
\[
 T_{2r+1}=\beta(\beta^2)^r-w(w^2)^r+v(v^2)^r.
\]
For multiplicative independence in the $S_3$ case, let
$\prod\alpha_i^{e_i}=1$ with integer $e_i$. Apply a Galois transposition
and divide the two relations. Then
$(\alpha_i/\alpha_j)^{e_i-e_j}=1$, forcing $e_i=e_j$. All exponents
are equal, and the norm $-4$ gives $(-4)^{e_1}=1$, hence $e_1=0$.
\end{proof}

\subsubsection{An inert odd-trace divisor that misses every classical phase}

\begin{proposition}\label{R16:prop:13miss}
For every $m\ge2$ with $m\equiv4\pmod{12}$, the cubic $f_{7^m}$ is irreducible
modulo $13$, and $13\mid T_{61}(7^m)$, but
\[
 13\nmid T_{3^k}(7^m)\qquad\text{for every }k\ge1.
\]
\end{proposition}
\begin{proof}
The coefficient is $a\equiv9\pmod{13}$. In
$\F_{13}[x]/(x^3-9x^2-4x+4)$, testing all $13$ residues proves
irreducibility. Direct quotient-ring powering gives
\[
 x^{61}=(11,7,9),\quad\tau(x^{61})=0,\qquad x^{183}=9.
\]
Thus the ordinary divisible exponent is odd, not an even exponent omitted
by the floor bridge. The ten successive cubing states are given in
Table~\ref{R16:tab:13miss}. All traces are nonzero. The next state is
$u_{11}=(1,12,1)=3u_1$. Cubing preserves nonzero scalar proportionality,
so by induction every later phase is a nonzero scalar multiple of one
of the ten listed states. This proves the nondivisibility for all $k$.
\end{proof}

\begin{table}[ht]
\centering\small
\begin{tabular}{rcr@{\qquad}rcr}
\toprule
$k$&$x^{3^k}$&$\tau$&$k$&$x^{3^k}$&$\tau$\\
\midrule
1&$(9,4,9)$&6&6&$(5,1,0)$&11\\
2&$(11,4,1)$&2&7&$(4,1,11)$&12\\
3&$(1,9,1)$&4&8&$(1,8,1)$&8\\
4&$(7,0,12)$&10&9&$(11,6,3)$&3\\
5&$(6,8,11)$&3&10&$(7,7,7)$&5\\
\bottomrule
\end{tabular}
\caption{Complete projective classical phases for $p=13$, $a=9$.}
\label{R16:tab:13miss}
\end{table}

Equivalently, after removing the factor $3$ from $183=3\cdot61$, the
classical residual orbit is
\[
 \langle3\rangle\bmod61=\{1,3,9,27,20,60,58,52,34,41\}.
\]
The trace-zero classes among exponents divisible by $3$, divided by $3$,
are $18,51,53$. They miss this orbit. The inertness of a prime and an
ordinary odd trace divisor therefore do not suffice for the clock
intersection in \eqref{R16:eq:IW}.

\subsubsection{A primitive ordinary divisor need not be inert}

At $a=7^4=2401$, the integer
\[
 p=T_3(a)=13841316001
\]
is prime, as certified by exhaustive trial division in the portable
verifier. The preceding ordinary terms are
$T_0=3$, $T_1=2401$, and $T_2=5764809$, all positive and smaller than
$p$. Thus $p$ is primitive in the first-occurrence sense for this
ordinary trace sequence. Nevertheless, modulo $p$,
\begin{equation}\label{R16:eq:mixedexample}
 f_a(X)=(X-1977685529)
 (X^2+1977683128X-6128330130).
\end{equation}
The quadratic discriminant is a nonresidue; equivalently
$\Delta(a)^{(p-1)/2}\equiv-1\pmod p$. Hence the splitting type is
$(1,2)$, not $(3)$. This does not disprove a hypothetical eventual
inert-primitive-divisor theorem. It shows that primitivity alone is not
the required inertness condition, and a prime seed is of course not a
composite term at that seed.

\subsection{No rational iteration in the trace coordinate}

\begin{theorem}\label{R16:thm:rational}
For any fixed $a>8$, put $C_k=T_{3^k}(a)$ and
$B_k=T_{2\cdot3^k-1}(a)$. For every fixed rational function
$F\in\R(X)$, each equality
\[
 C_{k+1}=F(C_k),\qquad B_{k+1}=F(B_k)
\]
holds for at most finitely many indices. In particular neither trace clock
is eventually a scalar rational-function orbit in its own trace coordinate.
\end{theorem}
\begin{proof}
For the classical clock, write $n=3^k$, $t=\beta^n$, $u=w^n$,
$z=v^n$. Then $C_k=t-u+z$ and the exact identity
\begin{equation}\label{R16:eq:defect}
 D_k:=C_{k+1}-C_k^3=3(u-z)(t-u)(t+z)
\end{equation}
gives
\[
 D_k\sim3(\beta^2w)^n,\qquad
 D_k/C_k^2\longrightarrow0,\qquad D_k/C_k\longrightarrow+\infty.
\]
The final limit uses $\beta w=4/v>1$. If a fixed rational function
matched infinitely many pairs, then on that unbounded subsequence
$F(C_k)/C_k^3\to1$. Its Laurent expansion at infinity must be
\[
 F(X)=X^3+b_2X^2+b_1X+O(1).
\]
The first defect ratio forces $b_2=0$; the second contradicts the finite
limit $b_1$. This proves the classical assertion.

For the negative clock put $n=2\cdot3^k-1$, whose successor is $3n+2$.
The corresponding identity is
\begin{align*}
 B_{k+1}-\beta^2 B_k^3
 ={}&3\beta^2(u-z)(t-u)(t+z)\\
 &+(\beta^2-w^2)u^3+(v^2-\beta^2)z^3.
\end{align*}
The last two terms are $O(w^{3n})$, so the difference is asymptotic to
$3\beta^2(\beta^2w)^n$. Dividing by $B_k^2$ and $B_k$ again gives
$0$ and $+\infty$. A putative $F$ must have leading term
$\beta^2X^3$, and the same Laurent-expansion contradiction follows.
\end{proof}

The proof extends to any cubic algebraic integer with conjugates
$\beta,-\rho,\gamma$ satisfying
\[
 \beta>1,\qquad 0<|\gamma|<\rho<1.
\]
For odd $n$, set $u=\rho^n$ and $z=\gamma^n$, allowing $\gamma<0$.
Then $u-z\sim\rho^n$, while the nonzero integral norm gives
$\beta\rho=|\Norm(\beta)|/|\gamma|>1$. Both arguments are unchanged.

The conclusion is deliberately coordinate-specific. Classical conjugates
still evolve by coordinatewise cubing, or by a polynomial map in three
symmetric coordinates. Higher-dimensional dynamics and extra auxiliary
coordinates remain available. Nor does the theorem forbid interpolation
through finitely many pairs.

The scheduled sequences also cannot be fixed-order constant-coefficient
linear recurrences in $k$: such a recurrence has growth at most
$A R^k k^d$, whereas these traces grow as
$\beta^{3^k}$ or $\beta^{2\cdot3^k-1}$. Ordinary-index recurrence
results must therefore be applied with their actual indexing hypotheses.

\subsection{Applicable literature, reproducibility, and the remaining gap}

\subsubsection{A prime-factor theorem that does transfer}

For each fixed $a=7^m$, the ordinary recurrence is nondegenerate with
three nonzero roots by Proposition~\ref{R16:prop:S3}. Evertse's Corollary~3
\cite[p.~230]{R16-Evertse} implies that its greatest prime factor tends to
infinity; the precise formulation is also recorded in
\cite[equation (2), p.~342]{R16-Stewart}. Consequently
\[
 P^+(T_{3^k}(a))\longrightarrow\infty,\qquad
 P^+(T_{2\cdot3^k-1}(a))\longrightarrow\infty.
\]
This is an application of a published theorem, not a new general
prime-factor result. It says that eventually each term has a factor
outside any fixed finite prime set. It supplies no splitting-type control
and no primitive divisor at every index. In particular it does not by
itself establish \eqref{R16:eq:IW}.

The classical primitive-divisor theorem for Lucas and Lehmer sequences
\cite[Theorem 1.4]{R16-BHV} cannot be applied simply by naming our trace a
recurrence: even its odd-index subsequence has exact order $3$, and no
factor identity reducing it to the relevant two-root form has been
established here. Roskam's higher-order positive-density result
\cite{R16-Roskam} assumes a generalized Artin conjecture and concerns
unrestricted ordinary indices. For fixed subgroup index $h$, Roskam
also proves ordinary-index divisibility outside finitely many inert
reduction primes whose root-generated subgroup has that index; the
positive-density supply of such primes is conditional, and neither
statement locates an index on our clocks. J\"arviniemi's Theorems 1.1--1.2
\cite{R16-Jarviniemi} assume GRH; in the cubic case the structural theorem
provides ordinary-term divisors with mixed splitting type, and its
Section~8 distinguishes the irreducible-reduction difficulty. Algebraic
genericity alone does not remove these hypotheses or impose either clock.

Arithmetic-dynamics primitive-divisor results such as
\cite{R16-IngramSilverman} concern iterates of a fixed rational map with a
suitable periodic target. Theorem~\ref{R16:thm:rational} rules out the direct
scalar trace-coordinate interpretation. It does not rule out encoding
traces as reduced numerators of another rational orbit with varying
denominators, nor a higher-dimensional reduction. No such additional
representation is established here. These are bounded applicability checks,
not an exhaustive literature or historical-priority claim.

\subsubsection{Portable exact verification}

The original source archive contains six Python 3 programs requiring only the standard
library, each with its exact JSON output and text log:
\begin{enumerate}
\item \texttt{frobenius\_coset\_coverage\_verify.py}: primality by exhaustive
 division, irreducible reductions, complete $D$-sets and coverage sets,
 all displayed short phases and missing strata. A finite consistency scan
 of \eqref{R16:eq:cubiclaw} is not substituted for its reciprocity proof.
\item \texttt{frobenius\_coset\_matrix\_audit.py}: an independent
 companion-matrix recomputation of the coverage sets, using the preceding
 JSON only for comparison, not its quotient-ring implementation.
\item \texttt{mixed\_output\_prime\_norm\_verify.py}: $720$ exact norm
 identities, finite-field norm/projective orders, and the local
 roots-of-unity models. Numerical samples do not prove the height bounds.
\item \texttt{cyclotomic\_norm\_height\_verify.py}: homogeneous polynomial
 normalizations, $27$ actual finite-field exact-order norm certificates,
 and rational checks of the unit-bound constants and parity rounding.
\item \texttt{inert\_bridge\_verify.py}: the complete $p=13$ projective-phase
 obstruction, the primitive-but-mixed prime, and polynomial identities
 in the Galois and rational-dynamics arguments.
\item \texttt{inert\_bridge\_independent\_verify.py}: independent
 companion-matrix returns and an Eratosthenes-sieve trial-divisor
 certification of the primitive-but-mixed prime.
\end{enumerate}
Run each program from the \texttt{certificates} directory. The universal
statements about norms, asymptotics, and all even $m$ are proved in the
text, not extrapolated from finite scans. Checksums identify the delivered
sources and outputs. 

\subsubsection{Scope and the next arithmetic problem}

The positive coverage certificates supply usable auxiliary primes whenever
the fixed cubic is irreducible at one of the listed primes. They prove
neither that every coefficient meets one of these cells nor that a suitable
prime exists universally. The missing step for the inert-witness strategy
is an intersection theorem between irreducible reductions of a fixed
cubic and divisors occurring on its fixed sparse clock.

The conditional output-prime bounds are different: they use equality to a
growing prime output and exclude overly small norm-component orders.
They leave ample admissible group order and yield no contradiction.
No coefficient-density theorem is transferred to a prescribed powering
orbit, no separate-phase statement is promoted to a common-phase
statement, and no universal Mills or Pisot exclusion is claimed.




\endgroup

\clearpage\part{Actual cubic candidates and analytic prime interfaces}

\section{Both residue signs and all norm depth}\label{src:R17}
The positive residue branch is genuinely present and requires its own analysis. The generalized two-sign depth theorem and positive-branch resolvent statements below supplement, rather than erase, the negative-branch constraints in the preceding section.
\medskip
\begingroup




\subsection{The classical minimum and the two genuinely remaining signs}

Let $\xi$ denote the least real number greater than $1$ for which
$\floor{\xi^{3^j}}$ is prime for every $j\ge1$. We do not assume that
$\xi$ is algebraic. We study consequences \emph{if it is}.
Saito's reduction \cite[Theorem 1.7 and Section 10]{Saito} supplies a
fixed power $\beta=\xi^{3^{j_0}}$ that is a totally real cubic Pisot
number with conjugates
\begin{equation}\label{R17:eq:shape}
 \beta,-\rho,\gamma,\qquad
 \beta>1,\quad0<\sigma:=|\gamma|<\rho<1,
 \quad\rho\le\beta^{-17/40}.
\end{equation}
The theorem has additional conjugate inequalities; only those displayed
are needed for the basic statements here. Write the irreducible minimal
polynomial as
\[
 f(X)=X^3-aX^2+bX-c,\qquad
 a,b,c\in\Z,\quad c\ne0,
 \qquad\Delta=\operatorname{disc}(f)\ne0.
\]
At clock index $k\ge0$ put $M=3^k$ and
\begin{align*}
 A_k&=\Tr(\beta^M),\\
 B_k&=\sum_{i<j}(\beta_i\beta_j)^M,\qquad C_k=c^M,
\end{align*}
where $\beta_1=\beta$, $\beta_2=-\rho$, $\beta_3=\gamma$.
All three are integers. The exact coefficient-cubing recurrence is
\begin{equation}\label{R17:eq:Newton}
 \begin{split}
 A_{k+1}&=A_k^3-3A_kB_k+3C_k,\\
 B_{k+1}&=B_k^3-3A_kB_kC_k+3C_k^2,\\
 C_{k+1}&=C_k^3,
 \end{split}
 \qquad (A_0,B_0,C_0)=(a,b,c).
\end{equation}
For all sufficiently large $k$, the shadow
$\rho^M-\gamma^M$ is positive and below $1$. Hence
\begin{equation}\label{R17:eq:tail}
 p_k:=A_k=\floor{\beta^{3^k}}
\end{equation}
is prime, and these outputs strictly increase. All later arithmetic
arguments require only the fixed cubic, \eqref{R17:eq:shape}, and this
prime trace-floor tail; they make no additional use of minimality.
Indices are always relative to this chosen fixed power $\beta$.

\subsubsection{What the source arguments actually say}

Saito's theorem specifies the real conjugate shape and sizes, not the
residue $p_k\pmod3$. The supplied manuscript \cite{K031} also
retains both signs in Proposition \texttt{mil:prop:cubic-split-three}.
Its theorem \texttt{f13:thm:r2-splitbranch} explicitly treats a
\emph{positive-sign} norm-minus-one unit branch and, under an additional
exact-depth hypothesis, forces complete splitting rather than excluding
that branch. The subsequent discussion leaves those split outputs open.

By contrast, the formula $a\equiv2\pmod3$ at
\texttt{ft:eq:cubic-a-two} occurs in the nonreal-conjugate argument.
That argument uses infinitely recurring floor shifts of both values $0$
and $1$. In the surviving real shape, the floor equals the trace
eventually, so that premise is unavailable. A negative real conjugate is
not a negative trace residue. These are scope distinctions, not evidence
that an algebraic all-prime candidate exists in either branch.

We now prove the common arithmetic statements directly, so the subsequent
bounds do not depend on a citation to an unstated sign restriction.

\subsection{Original projective order and trace-limit rigidity}

Call a prime $p$ \emph{good} if $p\nmid6c\Delta$. Every sufficiently
late increasing output is good. The reduced cubic is then separable
with nonzero roots $\alpha_1,\alpha_2,\alpha_3$ in a finite field.
Let $Q_p$ be the order of the subgroup generated by all ratios
$\alpha_i/\alpha_j$. Equivalently, it is the order of the residue of
$X$ in the unit group of the original cubic quotient modulo nonzero
scalar units. Write
\[
 Q_p=3^{w_p}q_p,\qquad3\nmid q_p.
\]
This definition uses the original separable algebra, not a possibly
nonreduced polynomial obtained after powering.

\begin{lemma}[Scalar return]\label{R17:lem:return}
For every sufficiently late good output $p=p_k$ of
\eqref{R17:eq:tail},
\[
 w_p\ge k+1.
\]
Such an output cannot have irreducible cubic reduction. Its reduction is
split or mixed, and
\[
 v_3(p^2-1)\ge k+1.
\]
\end{lemma}
\begin{proof}
If $w_p\le k$, choose $L\ge1$ with $3^L\equiv1\pmod{q_p}$;
use $L=1$ if $q_p=1$. Then $Q_p$ divides
$3^k(3^{jL}-1)$ for every $j\ge0$. All residual roots at exponent
$3^{k+jL}$ are a common nonzero scalar multiple of their values at
$3^k$. Their sum is zero because $p=A_k$. Thus $p$ divides every
$A_{k+jL}$, contradicting the strictly increasing prime tail.

For an irreducible cubic reduction, the projective order divides
$p^2+p+1$, whose $3$-adic valuation is at most $1$. This is incompatible
with $w_p\ge k+1$ at $k\ge1$. In the split case $Q_p\mid p-1$;
in the mixed case $Q_p\mid p^2-1$. This proves the remaining assertion.
\end{proof}

\begin{proposition}[Both signs and the initial coefficient congruence]
\label{R17:prop:sign}
There is a fixed $\varepsilon\in\{1,-1\}$ such that
\[
 p_k\equiv\varepsilon\pmod{3^{k+1}}
\]
for every sufficiently late output. Moreover
\begin{equation}\label{R17:eq:residues}
 a\equiv\varepsilon\pmod3,\qquad c\equiv ab\pmod3,
 \qquad f(X)\equiv(X-\varepsilon)(X^2+b)\pmod3.
\end{equation}
The argument does not choose between the two values of $\varepsilon$.
\end{proposition}
\begin{proof}
The recurrence gives $A_k\equiv a\pmod3$ for every $k$.
An increasing prime tail excludes $a\equiv0\pmod3$. Let
$\varepsilon\in\{1,-1\}$ represent this fixed residue.
Lemma~\ref{R17:lem:return} implies that the factor $p_k-\varepsilon$
is divisible by $3^{k+1}$, since the other factor of $p_k^2-1$ is
prime to $3$.

For completeness, the standard trace congruence follows directly from
\eqref{R17:eq:Newton}. Denote its three-coordinate polynomial map by $F$.
For every integral vector $z$, $F(z)\equiv z\pmod3$.
If $u-v\in3^r\Z^3$ with $r\ge1$, expansion of each cube and of the
terms with coefficient $3$ gives
$F(u)-F(v)\in3^{r+1}\Z^3$. Induction therefore yields
\[
 A_{k+1}-A_k\in3^{k+1}\Z.
\]
The $3$-adic limit is $\varepsilon$ by the output congruence, and
$A_1\equiv\varepsilon\pmod9$. But
$A_1=a^3-3ab+3c$ and $a^3\equiv\varepsilon\pmod9$.
Hence $c\equiv ab\pmod3$. Expanding the residual factorization proves
\eqref{R17:eq:residues}.
\end{proof}

For $\varepsilon=-1$, completely split outputs are impossible because
$3\nmid p-1$, so only mixed type remains. The norm-component bounds
for that branch in the preceding research round have precisely this
hypothesis. In the next subsection we focus on
\begin{equation}\label{R17:eq:positive}
 \varepsilon=1,\qquad a\equiv1\pmod3,
 \qquad c\equiv b\pmod3.
\end{equation}
Neither the negative-branch mixed classification nor its prime-to-$3$
norm-component interpretation may be transferred without modification.

\subsection{The positive branch: exact splitting and a cyclotomic sum}

Fix a sufficiently late good output $p=p_k\equiv1\pmod3$. Write
\[
 t=v_3(p-1),\quad Q_p=3^wq,\quad3\nmid q,
 \quad M=3^k,\quad v=w-k\ge1.
\]

\begin{proposition}\label{R17:prop:split}
Both split and mixed reduction remain unexcluded by the current necessary
conditions, and in either case
\[
 k+1\le w\le t.
\]
In mixed type, $c$ is a cube modulo $p$ if and only if the normalized
quadratic root is a cube in $\F_{p^2}^*$. If $c$ is a cube, then
$w\le t-1$ and $t\ge k+2$; if it is not, then $w=t$ exactly.
\end{proposition}
\begin{proof}
In split type $Q_p\mid p-1$. In mixed type write the roots as
$\alpha\in\F_p^*$ and $\lambda,\lambda^p\in\F_{p^2}^*$, and put
$y=\lambda/\alpha$. Then $Q_p=\ord(y)\mid p^2-1$.
Since $p\equiv1\pmod3$, one has $v_3(p^2-1)=t$. The lower bound is
Lemma~\ref{R17:lem:return}.

The field norm of $y$ is $c/\alpha^3$. A generator of
$\F_{p^2}^*$ maps under the norm to a generator of $\F_p^*$.
Thus the exponent of $y$ is divisible by $3$ if and only if its norm
is a cube, which is equivalent to $c$ being a cube. The group
$\F_{p^2}^*$ has $3$-part $3^t$; its cubes have $3$-part at most
$3^{t-1}$, and a noncube has order with $3$-part exactly $3^t$.
\end{proof}

The cube improvement is specific to mixed type. It is automatic for
unit norms $c=\pm1$ or a fixed integral cube norm, but not for an
arbitrary nonunit output. Since inert type is already excluded, the
Legendre symbol $(\Delta/p)$ now distinguishes the surviving types:
$+1$ means complete splitting and $-1$ means mixed splitting.
This is a classification at actual outputs, not a fixed congruence
criterion for complete splitting of every prime.

\begin{lemma}[A small nonzero pair coefficient]\label{R17:lem:B}
For all odd $M=3^k$, the integer $B_k$ is negative. At every sufficiently
late output,
\begin{equation}\label{R17:eq:Bsmall}
 0<|B_k|<p_k.
\end{equation}
\end{lemma}
\begin{proof}
If $\gamma>0$, then
\[
 B_k=-(\beta\rho)^M+(\beta\sigma)^M-(\rho\sigma)^M<0.
\]
If $\gamma<0$, then
\[
 B_k=-(\beta\rho)^M-(\beta\sigma)^M+(\rho\sigma)^M<0.
\]
The inequalities follow from $\rho>\sigma$ and $\beta>\sigma$.
Also
\[
 |B_k|\le\beta^M(\rho^M+\sigma^M)+(\rho\sigma)^M,
 \qquad p_k>\beta^M-1,
\]
so $|B_k|/p_k\to0$.
\end{proof}

\begin{theorem}[Nonvanishing cyclotomic-sum bound]\label{R17:thm:cyclosum}
Set
\[
 N=Q_p/3^k=3^v q.
\]
For every sufficiently late positive-residue good output,
\begin{equation}\label{R17:eq:sum-bound}
 p^{f_p}\le3^{\varphi(N)}
       =3^{\varphi(3^v)\varphi(q)},
 \qquad
 f_p=\begin{cases}1,&\text{split type},\\2,&\text{mixed type}.
 \end{cases}
\end{equation}
The residue degrees stated here are exact.
\end{theorem}
\begin{proof}
Let $Y_i=\alpha_i^M$. Their sum is zero modulo $p$, and their ratio
subgroup has exact order $Q_p/\gcd(Q_p,M)=N$.
Choose a primitive $N$th root $\xi$ in the finite splitting field and
write the normalized roots as $1,\xi^A,\xi^B$.
For a complex primitive $N$th root $\zeta_N$, put
\[
 Z=1+\zeta_N^A+\zeta_N^B.
\]
We first prove that $Z\ne0$. If two unit complex numbers sum to $-1$,
then $|1+\zeta_N^A|=1$ forces
$\zeta_N^{2A}+\zeta_N^A+1=0$, and the other term is its square.
The normalized roots would be the three cube roots of unity. Reducing
modulo $p\ne3$ would then give $\sum_{i<j}Y_iY_j=0$ modulo $p$,
contradicting \eqref{R17:eq:Bsmall}. The nonzero second coefficient is
essential: trace zero alone would not exclude this equilateral case.

Evaluation $\zeta_N\mapsto\xi$ gives a prime ideal above $p$ in
$\Z[\zeta_N]$ at which $Z$ vanishes. Because $p\nmid N$, its residue
degree is $\ord_N(p)$. The norm of that prime ideal is
$p^{\ord_N(p)}$, so this integer divides the nonzero rational norm
of $Z$. Every complex conjugate of $Z$ has absolute value at most $3$.
Therefore
\[
 p^{\ord_N(p)}\le
 \bigl|\Norm_{\Q(\zeta_N)/\Q}(Z)\bigr|\le3^{\varphi(N)}.
\]

In split type $N\mid p-1$, giving degree $1$.
In mixed type normalize by the linear root, so $N=\ord(y^M)$ for
$y=\lambda/\alpha$. If $y^M\in\F_p$, then
$(\lambda/\lambda^p)^M=1$. That ratio has order dividing $p+1$,
which is prime to $3$. Raising to $M=3^k$ cannot kill it, so
$\lambda=\lambda^p$, contrary to the irreducible quadratic factor.
Thus $N\mid p^2-1$ but $N\nmid p-1$, giving degree exactly $2$.
\end{proof}

The powered roots are geometric elements of the original separable
algebra. Possible collisions in a separately defined powered polynomial
do not alter their exact ratio-subgroup order used above. No equality
between a nonreduced polynomial quotient and its image has been assumed.

% Standalone elementary closure, for integration with the two-sign theorem.
% No literature-priority claim: this is an extension relative to the supplied drafts.
\subsection{Multiplicative independence and strict candidate bounds}

\begin{lemma}[Multiplicative relations among totally real cubic roots]
\label{R17:lem:new-cubic-independence}
Let $z_1,z_2,z_3$ be the roots of a monic irreducible totally real cubic
in $\mathbb Z[X]$, and put $C=|z_1z_2z_3|\ge1$.
If $C>1$, its three roots are multiplicatively independent modulo roots
of unity. If $C=1$, every multiplicative relation modulo roots of unity
has exponent vector in $\mathbb Z(1,1,1)$. In particular every pair of
distinct roots is multiplicatively independent modulo roots of unity,
in either case.
\end{lemma}
\begin{proof}
For $\ell_i=\log|z_i|$, apply a Galois three-cycle to a relation
$\prod_i z_i^{e_i}$ equal to a root of unity. Taking absolute values and
real logarithms makes $(e_1,e_2,e_3)$ orthogonal to the three cyclic
permutations of $(\ell_1,\ell_2,\ell_3)$.
The corresponding circulant determinant is, up to sign,
\[
 \log C\left(\ell_1^2+\ell_2^2+\ell_3^2
       -\ell_1\ell_2-\ell_2\ell_3-\ell_3\ell_1\right).
\]
The second factor is positive, since three distinct real roots cannot
all have the same absolute value. If $C>1$, the matrix is invertible.
If $C=1$, its rank is two and its kernel is the line spanned by
$(1,1,1)$. This gives the assertions; a vector supported on at most two
coordinates cannot be a nonzero constant vector.
\end{proof}

\begin{corollary}[Independence in the entire explicit family]
\label{R17:cor:new-family-independence}
For every integer $a>8$, the three roots of
$X^3-aX^2-4X+4$ are multiplicatively independent modulo roots of unity.
In particular no parity restriction on $m$ is needed when $a=7^m$,
$m\ge2$.
\end{corollary}
\begin{proof}
The interval/root argument proves irreducibility and total reality,
and the absolute norm is $4$. Apply
Lemma~\ref{R17:lem:new-cubic-independence}.
\end{proof}

\begin{corollary}[Strict shadows for the conditional classical candidate]
\label{R17:cor:new-strict-shadow}
Suppose an irreducible totally real cubic Pisot number has conjugates
$\beta,-\rho,\gamma$, with $\beta>1$ and
$0<\sigma=|\gamma|<\rho<1$, and satisfies the Saito inequalities
\[
 \rho\le\beta^{-17/40},\qquad \rho\le\sigma^{17/23}.
\]
Then both inequalities are strict. Moreover
\[
 \frac12<\kappa:=\frac{2+\log\rho/\log\beta}{3}<\frac{21}{40}.
\]
Consequently the mixed-positive anti-norm constant
$3\log\beta/\log(\beta^2\rho)$ is strictly greater than $40/21$.
\end{corollary}
\begin{proof}
Equality in the first inequality would give
$\beta^{17}(-\rho)^{40}=1$, a forbidden two-root relation.
Equality in the second would imply
$(-\rho)^{46}\gamma^{-34}=1$, again forbidden; the even exponents
remove all sign issues. For the lower endpoint, the nonzero integral
norm gives
$1\le|\operatorname{Norm}(\beta)|=\beta\rho\sigma<\beta\rho^2$.
Thus $\rho>\beta^{-1/2}$. Substituting this and the strict upper shadow
bound into the definition of $\kappa$ proves its two strict endpoints.
Finally, writing $u=-\log\rho/\log\beta$ gives
$17/40<u<1/2$ and the anti-norm constant is $3/(2-u)>40/21$.
\end{proof}

% Integration fragment: audited and extended from Round 17, Sections 2--4.
% Requires amsmath/amsthm and theorem/corollary environments.
\subsection{A two-sign depth theorem for arbitrary totally real cubics}
\label{R17:sec:new-two-sign-depth}

The positivity of the dominant root and the Pisot inequalities are not
needed for the exact-depth argument. This observation both broadens that
argument and closes its sign restriction.

\begin{theorem}[Effective two-sign depth for totally real cubics]
\label{R17:thm:new-two-sign-depth}\label{R17:thm:depth}
Let
\[
 f(X)=X^3-aX^2+bX-c\in\mathbb Z[X]
\]
be irreducible and totally real, with $c\ne0$. Let
$\varepsilon\in\{1,-1\}$ and assume
\[
 a\equiv\varepsilon\pmod3,\qquad c\equiv\varepsilon b\pmod3.
\]
For the roots $\beta_1,\beta_2,\beta_3$ put
$A_k=\sum_i\beta_i^{3^k}$. The following finite positive integral
 offsets exist and can be computed by terminating modular searches.
\begin{enumerate}
\item If $b\equiv0\pmod3$, put $s=v_3(f(\varepsilon))$. Then
\[
 v_3(A_k-\varepsilon)=k+s
 \qquad(k\ge\max\{3,s-1\}).
\]
\item If $b\equiv2\pmod3$, there is one offset $s\ge1$ such that
\[
 v_3(A_k-\varepsilon)=k+s\qquad(k\ge s).
\]
\item If $b\equiv1\pmod3$, there are offsets $s_0,s_1\ge1$ such that
\[
 v_3(A_k-\varepsilon)=k+s_j
 \quad(k\equiv j\pmod2,\ k\ge s_j-1).
\]
\end{enumerate}
Only the third case uses the proved Baker--Brumer theorem.
\end{theorem}

\begin{proof}
Set
\[
 g(X)=\varepsilon f(\varepsilon X)
   =X^3-\varepsilon aX^2+bX-\varepsilon c.
\]
It is again a monic irreducible totally real cubic. Its roots are
$\varepsilon\beta_i$, its trace at the odd exponent $3^k$ is
$\varepsilon A_k$, and its coefficients satisfy
$a'=\varepsilon a\equiv1$ and $c'=\varepsilon c\equiv b$ modulo $3$.
Also
$v_3(\varepsilon A_k-1)=v_3(A_k-\varepsilon)$ and
$v_3(g(1))=v_3(f(\varepsilon))$.
We therefore prove the assertion for $g$ and suppress the primes.
All local valuations below are normalized by $v_3(3)=1$.

If $b\equiv c\equiv0$, Hensel factorization of
$g\equiv(X-1)X^2$ gives a root $\alpha\in1+3\mathbb Z_3$;
the other roots have valuations at least $1/2$. The latter assertion
follows by comparing valuations in their monic quadratic, whose other
coefficients are divisible by $3$. Put $L=\log\alpha$.
Irreducibility implies $\alpha\ne1$, so
\[
 v_3(L)=v_3(\alpha-1)=v_3(g(1))=s\ge1.
\]
For $k\ge3$, the two nonunit powers have valuations at least
$3^k/2\ge2k+2$, while the exponential remainder also has valuation
at least $2k+2$. Consequently
\[
 A_k-1=3^kL+R_k,\qquad v_3(R_k)\ge2k+2.
\]
For $k\ge s-1$ the leading term has strictly smaller valuation.

If $b\equiv c\equiv2$, then
$g\equiv(X-1)^2(X+1)$. Let
$\eta\in-1+3\mathbb Z_3$ be the simple-root lift, and let $u,v$
be the other roots. They satisfy
$v_3(u-1),v_3(v-1)\ge1/2$.
The logarithms $U=\log u,V=\log v$ have valuations at least $1/2$.
Cubing first gives $v_3(u^3-1),v_3(v^3-1)\ge3/2$, so
logarithm and exponential are inverse after cubing, even in the ramified
quadratic case. With $L=\log(-\eta)$ and $k\ge1$, this gives
\[
 A_k-1=3^k\lambda+R_k,\qquad
 \lambda=U+V-L=\log c-2\log(-\eta)\in3\mathbb Z_3,
 \qquad v_3(R_k)\ge2k+1.
\]
Here $v_3(n!)\le(n-1)/2$ bounds the exponential remainders.
If $\lambda=0$, then the rational local unit $c/\eta^2$, congruent
to $-1$ modulo $3$, lies in the logarithm kernel. It must equal $-1$,
so $\eta^2=-c$, contrary to the cubic irreducibility of $g$.
Thus $s=v_3(\lambda)\ge1$ is finite, and $k\ge s$ makes the
leading valuation strictly below the error valuation.

If $b\equiv c\equiv1$, the separable reduction
$(X-1)(X^2+1)$ splits in the unramified quadratic field
$K=\mathbb Q_3(i)$. Label the roots
\[
 \alpha=e^L,\qquad u=i e^U,\qquad v=-i e^V,
 \qquad L,U,V\in3\mathcal O_K.
\]
Local conjugation fixes $L$, swaps $U,V$, and sends $i$ to $-i$.
Thus
\[
 \lambda_j=L+(-1)^j i(U-V)\in3\mathbb Z_3,
 \qquad A_k-1=3^k\lambda_{k\bmod2}+R_k,
 \qquad v_3(R_k)\ge2k+2.
\]
It remains to prove both $\lambda_j$ nonzero, without a Pisot
hypothesis. By Lemma~\ref{R17:lem:new-cubic-independence}, the three global roots are multiplicatively independent modulo roots of unity when $C=|c|>1$, and every pair is independent when $C=1$.

All three local roots are units. The kernel of their local logarithm
consists of roots of unity, so the preceding argument gives rational
linear independence of $L,U,V$ when $C>1$. When $C=1$, $U,V$ are
rationally independent and $L+U+V=\log c=0$.
The Baker--Brumer theorem \cite[Theorem~2.1]{R17-Maksoud} promotes these rational independence
statements to independence over the algebraic numbers. It excludes
$\lambda_j=0$ in the first case; in the second it excludes
\[
 (-1+(-1)^ji)U+(-1-(-1)^ji)V=0.
\]
Hence $s_j=v_3(\lambda_j)\ge1$ is finite. The leading valuation
$k+s_j$ is strictly smaller than $2k+2$ for $k\ge s_j-1$.

Finally all offsets are effectively recognizable using the integral
coefficient-cubing recurrence modulo powers of $3$. In the second case,
search $k\ge1$ until $A_k-1$ is nonzero modulo $3^{2k+1}$; its
valuation $d$ gives $s=d-k$. In the third case search each parity
until the corresponding residue is nonzero modulo $3^{2k+2}$ and
set $s_j=d-k$. In the first case use $v_3(g(1))$ directly.
The proved nonvanishing establishes termination, without asserting a
uniform practical running time.
\end{proof}

\subsubsection{One cyclotomic order bound for both output signs}

Retain a fixed cubic Pisot candidate with real conjugates
$\beta,-\rho,\gamma$, where $0<|\gamma|<\rho<1$, and suppose that
$A_k=\lfloor\beta^{3^k}\rfloor=p_k$ is an increasing prime tail.
For a good output $p=p_k\nmid6c\operatorname{disc}(f)$ let
$Q_p=3^wq$ be the projective order, with $3\nmid q$.
The scalar-return argument gives $w\ge k+1$, excludes irreducible
reduction, and yields a fixed sign $\varepsilon$ with
$p_k\equiv\varepsilon\pmod{3^{k+1}}$ and
$a\equiv\varepsilon$, $c\equiv\varepsilon b\pmod3$.
Use Theorem~\ref{R17:thm:new-two-sign-depth} and denote its offsets by
$s_j$, repeating a common offset on both parities where appropriate.

\begin{corollary}[Unified bounded excess and prime-to-$3$ growth]
\label{R17:cor:new-two-sign-order}
At every sufficiently late good output of parity $j$, set
$N=Q_p/3^k=3^{w-k}q$ and $F_j=\varphi(3^{s_j})$.
Then
\[
 1\le w-k\le s_j,\qquad
 p^{f_p}\le3^{\varphi(N)},\qquad
 \varphi(q)\ge\frac{f_p\log p}{F_j\log3}
 >\frac{f_p(3^k\log\beta-\log2)}{F_j\log3},
\]
where $f_p=1$ for split outputs and $f_p=2$ for mixed outputs.
If $\varepsilon=-1$, every such output is mixed. In particular both
signs satisfy $\varphi(q)\gg_f3^k$ and $Q_p\gg_f9^k$.
\end{corollary}

\begin{proof}
The exact depth gives $v_3(p-\varepsilon)=k+s_j$. As $Q_p$ divides
$p^2-1$, this bounds $w$ above. Let $Y_i$ be the residual roots
raised to $3^k$. Their ratio group has order $N$ and their sum is
zero. Normalize to $1,\xi^A,\xi^B$, with $\xi$ a primitive $N$th
root, and lift to $Z=1+\zeta_N^A+\zeta_N^B$ in characteristic zero.
The second powered coefficient $B_k$ satisfies
$0<|B_k|<p$ for all sufficiently late outputs. If $Z=0$, its three
unit terms form an equilateral triple, whose second coefficient
vanishes after reduction; this contradicts $p\nmid B_k$.
Thus $Z\ne0$ and
\[
 p^{\operatorname{ord}_N(p)}
 \le|\operatorname{Norm}_{\mathbb Q(\zeta_N)/\mathbb Q}(Z)|
 \le3^{\varphi(N)}.
\]
The residue degree is one for split outputs. In mixed type put
$y=\lambda/\alpha$, using the linear root $\alpha$ and one quadratic
root $\lambda$. When $p\equiv1\pmod3$, the nontrivial anti-norm
ratio $\lambda/\lambda^p$ has order dividing $p+1$, prime to $3$;
its $3^k$th power remains nontrivial. Hence $y^{3^k}\notin\mathbb F_p$.
When $p\equiv-1\pmod3$, the order $N$ of $y^{3^k}$ has a nontrivial
$3$-part because $w>k$, whereas $\mathbb F_p^*$ has order prime to
$3$. Again $y^{3^k}\notin\mathbb F_p$. Thus the residue degree is
exactly two in either mixed branch. Finally
$\varphi(N)\le F_j\varphi(q)$,
$p>\beta^{3^k}/2$ eventually, and
$Q_p\ge3^{k+1}q\ge3^{k+1}\varphi(q)$.
\end{proof}

These results are necessary conditions for a hypothetical prime tail.
They do not force an inert output or a composite term, and the available
order $p^2-1$ remains far larger than the required $9^k$ scale.

\begin{corollary}[Prime-to-$3$ growth for every positive candidate]
\label{R17:cor:positivegrowth}
For a fixed candidate satisfying the positive prime tail, use the same
constant $s$ on both parities in the first two cases of
Theorem~\ref{R17:thm:depth}, and put
$F_j=\varphi(3^{s_j})=2\cdot3^{s_j-1}$.
At every sufficiently late good output of parity $j$,
\begin{equation}\label{R17:eq:positivegrowth}
 \varphi(q)\ge\frac{f_p\log p_k}{F_j\log3}
 >\frac{f_p(3^k\log\beta-\log2)}{F_j\log3},
\end{equation}
where $f_p=1$ in split type and $f_p=2$ in mixed type.
If $c$ is a cube modulo a mixed output, then necessarily $s_j\ge2$ and
\begin{equation}\label{R17:eq:cubegrowth}
 \varphi(q)\ge
 \frac{2\log p_k}{\varphi(3^{s_j-1})\log3}
 =\frac{6\log p_k}{F_j\log3}.
\end{equation}
In particular, for every fixed surviving positive-residue candidate,
\[
 \varphi(q)\gg_\beta3^k,\qquad Q_{p_k}\gg_\beta9^k.
\]
\end{corollary}
\begin{proof}
The depth theorem gives $t=k+s_j$. Proposition~\ref{R17:prop:split} implies
$1\le v=w-k\le s_j$, so $\varphi(3^v)\le F_j$.
Insert this in Theorem~\ref{R17:thm:cyclosum}; then use
$p_k>\beta^{3^k}/2$ eventually.
In the mixed cube case, $w\le t-1$ gives
$1\le v\le s_j-1$, hence $s_j\ge2$ and the improved denominator.
Finally $Q_{p_k}\ge3^{k+1}q\ge3^{k+1}\varphi(q)$.
\end{proof}

The improvement by a factor of six in \eqref{R17:eq:cubegrowth}, relative to
the split bound with the same $s_j$, consists of the exact residue-degree
factor $2$ and the smaller $3$-primary cyclotomic degree factor $3$.
For unit norms it is automatic in mixed type. If $s_j=1$, mixed type
with a cube norm is excluded on that parity, but complete splitting is
not. For arbitrary nonunits the cube condition remains an additional
condition at the particular output prime.


\subsection{A separate anti-norm bound for every mixed positive output}

The next condition does not need unit norm or a bounded excess
$3$-adic depth. Retain the mixed roots of Proposition~\ref{R17:prop:split}
and define
\[
 h=\ord(\lambda/\lambda^p).
\]
Then $h>1$, $h\mid p+1$, and $h\mid Q_p$. Since $p\equiv1\pmod3$,
\begin{equation}\label{R17:eq:anti-divides}
 3\nmid h,\qquad h\mid q.
\end{equation}
This is the anti-norm component. The norm-component order used in the
negative residue branch need not be prime to $3$ in the present branch.

For $j\ge1$, let $\Phi_j(X,Y)$ be the homogeneous cyclotomic polynomial,
and put $F(r)=r/(1-r)^2$ for $0<r<1$.
The elementary estimate
\begin{equation}\label{R17:eq:cyclo-est}
 \left|\log|\Phi_j(z)|\right|\le F(|z|)\quad(0<|z|<1)
\end{equation}
is uniform in $j$. Indeed, the M\"obius product gives a signed sum of
$\log|1-z^d|$, each bounded in absolute value by
$-\log(1-|z|^d)\le|z|^d/(1-|z|)$. Summing over all positive $d$
proves the claim. Reciprocity gives the corresponding estimate outside
the unit circle.

\begin{theorem}[Anti-norm height constraint]\label{R17:thm:anti}
Define
\[
 E_{\rm pair}=F(\rho/\beta)+F(\sigma/\beta)+F(\sigma/\rho).
\]
For every sufficiently late good mixed positive output,
\begin{equation}\label{R17:eq:anti-bound}
 \varphi(h)>
 \frac{3^k\log\beta-\log2-E_{\rm pair}}
      {\log(\beta^2\rho)}.
\end{equation}
The same lower bound holds for $\varphi(q)$. If there are infinitely
many such outputs, then along them
\begin{equation}\label{R17:eq:anti-liminf}
 \liminf\frac{Q_{p_k}}{9^k}
 \ge\frac{3\log\beta}{\log(\beta^2\rho)}.
\end{equation}
Under the strict size inequality $\rho<\beta^{-17/40}$, the last
candidate-specific constant is strictly greater than $40/21$.
The weak inequality in \eqref{R17:eq:shape} alone gives at least $40/21$.
\end{theorem}
\begin{proof}
For $h>1$, $\Phi_h(X,Y)$ is symmetric. Hence
\[
 K_h=\prod_{i<j}\Phi_h(\beta_i,\beta_j)
\]
is a rational algebraic integer, and therefore an integer.
It is nonzero because the three real conjugates have distinct absolute
values, so no pair ratio is a root of unity.
The exact order of $\lambda/\lambda^p$ and $p\nmid h$ give
$\Phi_h(\lambda,\lambda^p)=0$ modulo $p$; thus $p\mid K_h$.
Applying \eqref{R17:eq:cyclo-est} to the three ratios of smaller to larger
absolute value yields
\[
 0<|K_h|\le e^{E_{\rm pair}}(\beta^2\rho)^{\varphi(h)}.
\]
The base exceeds $1$, since the nonzero integral norm gives
$\beta\rho=|c|/\sigma>1$. The actual output satisfies
$p_k>\beta^{3^k}/2$ eventually, proving \eqref{R17:eq:anti-bound}.
If $h\mid q$, the prime factor formula for Euler's function gives
$\varphi(h)\le\varphi(q)$. Finally
$Q_{p_k}\ge3^{k+1}q\ge3^{k+1}\varphi(h)$, proving
\eqref{R17:eq:anti-liminf}. Write $u=-\log\rho/\log\beta$; the constant
is $3/(2-u)$, which is at least $40/21$ when $u\ge17/40$, and
strictly larger under the strict inequality.
\end{proof}

\subsection{Powered resolvents preserve the splitting character}

This section is a limitation of a proposed test, not an additional
exclusion. Restrict to a totally real cubic Pisot \emph{unit}, so
$c\in\{1,-1\}$ and $C_j=c$ for every odd powering. Let
$z_i=\beta_i^{3^j}$ at stage $j$. All these powers remain cubic:
a rational value would have equal conjugates, impossible when exactly
one conjugate has modulus greater than $1$.

\subsubsection{Three fields and an exact oriented covariance}

Choose an ordering of the roots and write
\[
 \delta_j=\prod_{i<\ell}(\beta_i^{3^j}-\beta_\ell^{3^j}),\quad
 \Delta_j=\delta_j^2,\quad I_j=\delta_j/\delta_0\in\Z.
\]
The ratio is a symmetric algebraic integer, because each difference
quotient is a polynomial with integral coefficients; hence it is an
integer. In particular $\Delta_j=\Delta_0I_j^2$.
Let $\omega^2+\omega+1=0$ and
$\eta_0=(\omega-\omega^2)\delta_0$, so $\eta_0^2=-3\Delta_0$.
Distinguish
\[
 K=\Q(\sqrt{\Delta_0}),\qquad
 K'=\Q(\sqrt{-3\Delta_0}),\qquad E=\Q(\omega,\delta_0).
\]
The quadratic resolvent is $K$, which may be $\Q$ in the cyclic case.
The Cardano radicands lie in $K'$, whereas $E$ contains $\mu_3$ and
is the appropriate Kummer base. A cubic reciprocity formula for
Eisenstein integers cannot simply be applied to an arbitrary element
of $K'$.

Define the Fourier resolvents and two rational coefficients by
\[
 \begin{aligned}
 U_j&=z_1+\omega z_2+\omega^2z_3,&
 V_j&=z_1+\omega^2z_2+\omega z_3,\\
 P_j&=A_j^2-3B_j,&
 H_j&=2A_j^3-9A_jB_j+27c.
 \end{aligned}
\]
Expansion gives
\[
 U_jV_j=P_j,\qquad U_j^3+V_j^3=H_j,
 \qquad U_j^3-V_j^3=3\eta_0I_j.
\]
Here $P_j=\frac12\sum_{i<\ell}(z_i-z_\ell)^2>0$, since the roots
are distinct and real. Thus neither resolvent vanishes. Put
\[
 R_j^+=(H_j+3\eta_0I_j)/2=U_j^3,
 \qquad R_j^-=(H_j-3\eta_0I_j)/2=V_j^3,
 \qquad R_j^+R_j^-=P_j^3.
\]
The weighted-square identity
$\sum_i\omega^{i-1}z_i^2=(2A_jU_j+V_j^2)/3$, together with
$z_i^3=A_jz_i^2-B_jz_i+c$, yields
\begin{equation}\label{R17:eq:covstep}
 U_{j+1}=D_jU_j,\qquad
 D_j=\frac{2A_j^2-3B_j}{3}
       +\frac{A_jR_j^-}{3P_j}\in K'^*.
\end{equation}
Consequently, with $\mathcal C_k=\prod_{0\le j<k}D_j$,
\begin{equation}\label{R17:eq:covariance}
 R_k^+=R_0^+\mathcal C_k^3,\qquad\mathcal C_k\in K'^*.
\end{equation}
This is equality of \emph{oriented} cube classes, not merely equality
up to inversion or equality of splitting fields. Reversing the initial
orientation consistently swaps the two classes. Standard Kummer
classification over a base containing $\mu_3$ is described in
\cite[Propositions 5.27--5.28]{R17-Milne}; the explicit multiplier here is
proved by \eqref{R17:eq:covstep}, not inferred from that classification.

\subsubsection{The powered test is the original splitting test}

Let $p\nmid6\Delta_0P_0$, assume $p\equiv1\pmod3$ and
$(\Delta_0/p)=1$, and choose a degree-one prime of $E$ over $p$.
Then $\omega$ and $\delta_0$ reduce to $\F_p$. If $r_0$ is the
nonzero reduction of $R_0^+$, one has the exact test
\begin{equation}\label{R17:eq:split-test}
 f\text{ splits completely modulo }p
 \quad\Longleftrightarrow\quad r_0^{(p-1)/3}=1.
\end{equation}
If $f$ splits, $U_0\in\F_p$ proves that $r_0$ is a cube.
Conversely, choose $u\in\F_p^*$ with $u^3=r_0$ and set $v=P_0/u$.
Then $v^3=R_0^-$, and the inverse Fourier expressions
\[
 \frac{A_0+\omega^{-\ell}u+\omega^\ell v}{3}
 \qquad(\ell=0,1,2)
\]
are all three roots of $f$ in $\F_p$. They are distinct since
$p\nmid\Delta_0$.

If also $p\nmid P_k$, both $R_0^\pm$ and $R_k^\pm$ are local units,
by their products $P_0^3$ and $P_k^3$. Equation~\eqref{R17:eq:covariance}
therefore makes $\mathcal C_k$ a local unit, even if an intermediate
formula for $D_j$ has a denominator divisible by $p$. It follows that
\[
 (R_k^+)^{(p-1)/3}=(R_0^+)^{(p-1)/3}\pmod{\mathfrak p}
\]
at the chosen prime $\mathfrak p$ of degree one. The powered character
is exactly the original character, not a second independent condition.
In the $S_3$ case this remains nonabelian splitting information over the
base containing the quadratic resolvent and $\mu_3$, not generally a
fixed rational congruence criterion. Cyclic cubic fields are different:
their abelian splitting may be described by conductor congruences.

At actual outputs the apparent degeneracy $P_k=0\pmod p$ disappears.
There is even an integral unit-cubic nonvanishing fact independent of the
real shape: $B_j\ne0$ for every $j\ge1$. For if the second recurrence
in \eqref{R17:eq:Newton} were zero, it would force $3\mid B_{j-1}$;
its right side would then be $3\pmod9$, a contradiction.
For our real candidate, Lemma~\ref{R17:lem:B} further gives
$0<|B_k|<p_k$ eventually. Thus at $p=A_k$,
\[
 P_k\equiv-3B_k\not\equiv0\pmod p,
 \qquad R_k^+\equiv(27c+3\eta_0I_k)/2\pmod{\mathfrak p}.
\]
The latter nonzero trace-zero radicand has the same cube character as
$R_0^+$. An additional restriction on that \emph{fixed} character along
the actual prime outputs would still be needed for an exclusion.

There is also a precise one-step index statement. Direct expansion gives
\[
 J_j:=I_{j+1}/I_j=-A_j^3c+A_j^2B_j^2-B_j^3
 =\prod_{i<\ell}(z_i^2+z_iz_\ell+z_\ell^2).
\]
Hence $\Delta_{j+1}=\Delta_jJ_j^2$. At a sufficiently late actual
output $p=A_k$, one has $J_k\equiv-B_k^3\ne0\pmod p$, so
\[
 v_p(\Delta_{k+1})=v_p(\Delta_k).
\]
This output cannot first enter the powering discriminant at the very next
step. Earlier or later collisions are not excluded; this is a statement
about polynomial discriminants, not an assertion of field ramification.

\subsubsection{One split actual prime seed, followed by a composite trace}

The following finite diagnostic prevents an overstrong individual-prime
inference. Take
\[
 f(X)=X^3-3529X^2-X+1,\qquad c=-1.
\]
Both $3529$ and
\[
 p=A_1=43949615473
\]
are prime, certified by exhaustive trial division, with the latter trial
bound $209641$. Its three distinct roots modulo $p$ are
\[
 36438575596,\qquad33092793291,\qquad18367865588.
\]
The complete projective order of this triple is
\[
 Q=21974807736=(p-1)/2,\qquad v_3(Q)=3\ge1+1,
 \qquad p-1=2^4\cdot3^3\cdot7^2\cdot2076229.
\]
The last factor is prime. Exact ratio exponent tests for every prime
factor of $Q$ certify that it is the full projective order, not only an
upper bound. The seed thus satisfies the required projective-depth
inequality at $k=1$.

The real roots lie in $(-1,0),(0,1),(3529,3530)$; none is integral,
so the cubic is irreducible. Their sum makes the negative small root
have larger modulus. Consequently every positive odd floor is its trace.
Furthermore
\[
 f(-1/50)=-48951/125000<0,\qquad50^{40}>3530^{17}.
\]
Thus its negative small modulus satisfies
$w<1/50<\beta^{-17/40}$. Since $\beta wv=1$, this also gives
$w<v^{17/23}$. These are strict Pisot shadow bounds, not merely a
freely prescribed finite-field root configuration.
Nevertheless the next classical trace is composite:
\[
 \begin{split}
 A_2=T_9&=84891701634551866699236663173929\\
       &=991\cdot85662665625178472955839216119.
 \end{split}
\]
This is one prime seed, not an all-prime tail, a minimum, or a
counterexample to any necessary-condition theorem. It shows only that
unit norm, strict shadow, complete splitting, and the individual seed
projective-depth condition are compatible.


\subsection{Reproducibility and the current arithmetic gap}

\subsubsection{Portable verification}

The original source archive supplies five Python 3 programs using only the standard
library, each with JSON output and a text log:
\begin{enumerate}
\item \texttt{positive\_branch\_verify.py}: seven mixed square-prime
cyclotomic norm certificates, four split certificates, $79$ exactly
certified totally real Pisot examples, and $1355$ modular depth
regressions across all three residue cases.
\item \texttt{positive\_branch\_independent\_verify.py}: an independent
implementation checking $540$ mixed cube-equivalence cases, $30$ mixed
norm certificates, $18$ split certificates, unit-depth examples, and an
exact anti-norm integer at an isolated actual prime trace.
\item \texttt{positive\_allnorm\_depth\_independent\_verify.py}:
$49$ certified Pisot inputs and $1225$ companion-matrix/Newton
comparisons, including modular offset extraction and unequal parity
offsets. No logarithm approximation is used to guess an offset.
\item \texttt{powered\_resolvent\_verify.py}: five formal polynomial
identities, all $162$ relevant modulo-$9$ inputs for the nonvanishing
lemma, and exact primality, splitting, order, shadow, covariance, and
next-composite certificates for the displayed seed.
\item \texttt{powered\_resolvent\_independent\_verify.py}:
an independent exact-grid verification of the formal identities, with
a degree bound justifying polynomial interpolation, and a separate
sieved trial-divisor certification of the prime seed.
\end{enumerate}
These are exact finite regressions or certificates, not samples from which
the universal depth theorem is inferred. The global logarithmic
nonvanishing and terminating-search assertions are proved in the depth theorem above.
The source programs contain no probable-prime proof oracle and require
neither a computer algebra package nor network access. Checksums bind the
delivered files. 

\subsubsection{What the partial results can be used for}

For a specified totally real cubic Pisot candidate with a prime tail of either residue sign, the exact depth algorithm computes the offsets that limit its
possible projective $3$-part. The cyclotomic-sum bound then forces a
prime-to-$3$ order component of exponential size in the clock index;
for a mixed output, the independent anti-norm certificate can give another
candidate-specific bound. A concrete proposed output or local model that
violates any one of these necessary conditions is ruled out.

The vanishing-sum proof also identifies a useful constraint that is absent
from arbitrary local root configurations: the actual second powered
coefficient is a nonzero integer strictly smaller than the prime output.
The resolvent calculation, in contrast, prevents counting the same
splitting character twice as if powering supplied an independent
obstruction. These are usable tests and method limitations, not only
formal changes of notation.


\subsubsection{What remains unresolved}

The available finite-field order can still be as large as $p^2-1$, while
$p$ itself grows exponentially in $3^k$. The lower bounds here therefore
leave ample room and produce no asymptotic contradiction. They do not
force an inert output, a particular discriminant character, a recurrent
divisor, or a composite future term for every candidate. The single
composite successor in the resolvent subsection is an exact isolated example, not a
universal mechanism.

Both residue signs remain logically distinct. The negative-branch bounds
from the preceding round do not cover the positive branch; the positive-branch refinements here complement the two-sign depth theorem, including nonunit norms. No density
statement for varying coefficients is transferred to this fixed prime
trace orbit, and Chebotarev is not invoked to select prime divisors of
fixed trace values. The all-norm depth theorem uses a proved logarithm
theorem, but neither it nor the cyclotomic bounds settle transcendence of
the classical Mills minimum. This paper also makes no closure claim for
the different schedule $2\cdot3^k-1$.




\endgroup

\section{Three analytic routes and their exact missing hypotheses}\label{src:R18}
The final analytic comparison uses the preceding arithmetic structure to test three specific interfaces. Each positive deduction is kept distinct from the missing assertion needed to reach every fixed candidate.
\medskip
\begingroup




\subsection{Conclusion and scope}

\textbf{The main targets remain open in this investigation.}
The three routes do produce rigorous information, but none supplies the
missing statement for every surviving fixed candidate. The distinction
between a finite configuration, a prescribed local orbit, and a single
global minimum is essential.

\begin{enumerate}
\item \textbf{Tao--Ziegler prime patterns.} Finite cubic trace shapes can
occur after a common additive translation. Such translated prime patterns
can coexist with recurring composite \emph{unshifted} Pisot floors.
There are exact finite obstructions to identifying the translation with
a new Pisot trace sequence.
\item \textbf{Finite-field sum--product.} The subgroup of exponents is a
valid input to BKT/BGK. Separate additive Fourier bounds do not control
its correlation with the particular Singer trace-zero set. The
phase-blind estimates studied here have a precise obstruction; this is
not a prohibition on every possible sum--product argument.
\item \textbf{Gafni--Tao short intervals.} A candidate gap forces a block
of exceptional starting points. The currently inspected envelope is too
weak to exclude that block. A stronger exponent or a theorem forcing
many disjoint bad intervals at a common scale would be sufficient, but
neither is established here.
\end{enumerate}

The first two routes use the explicit family
\begin{equation}\label{R18:eq:family}
 f_a(X)=X^3-aX^2-4X+4,\qquad a\in\Z,\quad a>8.
\end{equation}
The third route concerns a \emph{conditional reduction for the actual
classical Mills minimum}; its gap hypotheses are stated separately in
Section~\ref{R18:sec:interval}. We do not identify an arbitrary member of
\eqref{R18:eq:family} with that minimum.

The proofs below are deductions from the cited theorems and elementary
algebra, together with exact finite certificates. No claim of literature
priority is made. Analytic infinitude statements are not inferred from
the finite computations.

\subsection{Trace normalization and actual floors}

Write $T_n(a)$ for the sum of the $n$th powers of the three roots of
$f_a$. Newton's identities give
\begin{equation}\label{R18:eq:newton}
\begin{split}
 T_0=3,\quad T_1=a,\quad T_2=a^2+8,\\
 T_n=aT_{n-1}+4T_{n-2}-4T_{n-3}\qquad(n\ge3).
\end{split}
\end{equation}
For $n\ge1$, $T_n$ is a monic integer polynomial of degree $n$.

\begin{lemma}\label{R18:lem:floor}
For $a>8$, the roots of $f_a$ are $-w,v,\beta_a$, where
$0<v<w<1$ and $a<\beta_a<a+1$. The polynomial is irreducible over
$\mathbb Q$, and for every positive odd integer $n$,
\[
 \floor{\beta_a^n}=T_n(a),\qquad
 0<\beta_a^n-T_n(a)=w^n-v^n<1.
\]
Moreover, $T_n(1)=1$ for every positive odd $n$.
\end{lemma}
\begin{proof}
At $-1,0,1,a,a+1$, the values of $f_a$ are respectively
$7-a,4,1-a,4(1-a),(a-1)^2$. They locate the three roots in the
stated intervals. Their sum gives $w-v=\beta_a-a>0$. None of the
roots is an integer, so the rational-root test proves irreducibility.
The identity for odd powers proves the floor formula. Finally,
$f_1=(X-1)(X^2-4)$, so its odd power sums are $1+2^n+(-2)^n=1$.
\end{proof}

We use the two positive odd exponent clocks
\[
 n_+(k)=3^k,\qquad n_-(k)=2\cdot3^k-1\qquad(k\ge1).
\]
A recurring congruence $\ell\mid T_{n_\pm(k)}(a)$ produces infinitely
many composite actual floors once the trace exceeds $\ell$.

\subsection{Tao Ziegler patterns and finite rigidity}\label{R18:sec:tz}

\subsubsection{The common-zero theorem gives translated trace shapes}

The common-zero polynomial progression theorem of Tao--Ziegler
\cite{TZ2008} applies to any finite family of integer polynomials
vanishing at zero. It allows unequal degrees and supplies infinitely
many simultaneous prime translates, with the polynomial parameter at
most any prescribed positive power of the translating parameter.

\begin{proposition}\label{R18:prop:translated}
Let $I$ be a finite nonempty set of positive odd integers and let $B$
be any real lower bound. There are infinitely many pairs of positive
integers $(x,m)$ such that, for a fixed integer
$L\ge\max(8,B)$ and
\[
 a=1+Lm>B,\qquad c=x-1\ge0,
\]
all the numbers
\[
 T_n(a)+c=\floor{\beta_a^n}+c\qquad(n\in I)
\]
are prime. If $d=\max I$, these pairs can be chosen with
$m\le x^\varepsilon$ for any fixed $0<\varepsilon<1/d$.
Along an unbounded subsequence of such pairs,
\begin{equation}\label{R18:eq:dominating-shift}
 c\longrightarrow\infty,\qquad
 \max_{n\in I}|T_n(a)|=o(x),\qquad
 \frac{T_n(a)+c}{x}\longrightarrow1.
\end{equation}
\end{proposition}
\begin{proof}
Use $P_n(m)=T_n(1+Lm)-1$. Lemma~\ref{R18:lem:floor} gives $P_n(0)=0$,
and the cited theorem makes $x+P_n(m)$ prime for every $n\in I$.
Bounded $x$ permits only finitely many $m\le x^\varepsilon$, so
infinitely many pairs force an unbounded subsequence of $x$. Polynomial
growth gives $T_n(1+Lm)=O_{I,L}(x^{\varepsilon d})=o(x)$.
\end{proof}

This is a common translation of an entire finite vector, not unrelated
coordinatewise shifts. Nevertheless, in this application the translation
dominates every actual trace. It does not provide $c=0$, and choosing
$L$ large does not produce one fixed coefficient that survives all
increasing prefix lengths.

The unequal-constant theorem in Tao--Ziegler's later paper
\cite[Theorems 1.2 and 1.4, Corollary 1.5]{R18-TZ2016} has a different
hypothesis: every pairwise polynomial difference must have the maximal
degree. A raw tuple of three or more distinct trace degrees fails it,
since two members below the maximal degree have a lower-degree
difference. Subtracting a common polynomial does not change this defect.

\subsubsection{A valid prime pair inside an obstructed family}

\begin{proposition}\label{R18:prop:pair}
For infinitely many distinct integers $a\equiv43\pmod{186}$,
$a\ge229$, there are primes $q_1,q_2$ with
\[
 q_2-q_1=T_3(a)-a=a^3+11a-12.
\]
At each such coefficient, the actual floors are composite on the
recurring index classes
\begin{align*}
 \floor{\beta_a^{3^k}}&\quad\text{for }k\equiv1\pmod4,\\
 \floor{\beta_a^{2\cdot3^k-1}}&\quad\text{for }k\equiv3\pmod4.
\end{align*}
\end{proposition}
\begin{proof}
Put $a(m)=43+186m$ and apply the 2016 theorem to
$P_1(m)=a(m)$ and $P_2(m)=T_3(a(m))=a(m)^3+12a(m)-12$.
Their degree-three coefficients are $0$ and $186^3$, so the required
degree condition holds.

We check local admissibility. For a prime $\ell\nmid186$, choose $m$
modulo $\ell$ with $a(m)=1$; then $P_1=P_2=1$, and any
$x\not\equiv-1\pmod\ell$ works. At $\ell=2$, both values are odd;
at $\ell=3$, both are $1$, so $x=0$ works in either case. At
$\ell=31$, at most two residues of $x$ are forbidden. Hence the
singular series is positive and the theorem gives prime pairs
$x+P_1(m),x+P_2(m)$.

The coefficients $a(m)$ are unbounded, not merely the translations:
Theorem 1.4 provides a positive asymptotic weighted count of order $NM$
with $M$ near $N^{1/3}$. A fixed finite set of $m$ contributes only
$O(N\log^2N)=o(NM)$. Prime powers can be removed as in the cited
Corollary 1.5. Thus infinitely many distinct $m$ remain.

For compositeness, $a\equiv12\pmod{31}$ and
\[
 f_a(X)\equiv(X+6)^2(X+7)\pmod{31}.
\]
For odd $n$, this gives
\[
 T_n(a)\equiv2(-6)^n+(-7)^n=-6^n(2-4^n)\pmod{31}.
\]
Since $\ord_{31}(4)=5$ and $4^3\equiv2\pmod{31}$, the trace
vanishes exactly when $n\equiv3\pmod5$. The two exponent clocks
yield the displayed classes of $k$. Their traces exceed $31$ and
equal the actual floors by Lemma~\ref{R18:lem:floor}.
\end{proof}

For a small exact illustration, take $a=229,c=42$. Then
\[
 a+c=271,\qquad T_3(a)+c=12{,}011{,}767
\]
are prime, whereas
\[
 T_3(a)=12{,}011{,}725=5^2\cdot11\cdot31\cdot1409.
\]
The theorem, not this numerical example, proves infinitude. The example
shows concretely why translated prime geometry can occur in an already
obstructed unshifted family.

\subsubsection{Seven odd trace coordinates force trivial translation}

\begin{lemma}\label{R18:lem:zeros}
A nonzero real exponential sum
$F(t)=\sum_{j=1}^s b_j e^{\lambda_jt}$, with distinct real
$\lambda_j$ and nonzero real $b_j$, has at most $s-1$ distinct real zeros.
\end{lemma}
\begin{proof}
Induct on $s$. Divide by $e^{\lambda_1t}$ and differentiate. The
derivative has $s-1$ exponential terms. If the original sum had $s$
distinct zeros, Rolle's theorem would give at least $s-1$ derivative
zeros, contradicting the induction hypothesis.
\end{proof}

\begin{theorem}\label{R18:thm:rigidity}
Let $\beta,\gamma$ be totally real cubic Pisot numbers and $c\in\R$.
If
\[
 \Tr(\gamma^n)=\Tr(\beta^n)+c
\]
at seven distinct positive odd integers $n$, then $\beta=\gamma$ and
$c=0$. In particular, seven distinct exponents from either of the
actual clocks suffice.
\end{theorem}
\begin{proof}
Let $\beta_i,\gamma_i$ be the real conjugates. At odd integers,
$z^n=\operatorname{sgn}(z)|z|^n$. Thus the assumed equalities are
zeros of the real function
\[
 F(t)=\sum_{i=1}^3\operatorname{sgn}(\gamma_i)|\gamma_i|^t
      -\sum_{i=1}^3\operatorname{sgn}(\beta_i)|\beta_i|^t-c.
\]
Merge equal positive bases and remove zero coefficients. At most seven
distinct exponential terms remain, including the constant base $1$.
Lemma~\ref{R18:lem:zeros} forces $F$ to vanish identically. No Pisot
conjugate has magnitude $1$, so the coefficient at that base is $-c$,
whence $c=0$. If $\beta\ne\gamma$, the larger dominant root supplies
a unique largest base greater than $1$ with nonzero coefficient, a
contradiction.
\end{proof}

No common norm, common polynomial family, or consecutive-index
assumption is used. The theorem does not address cubics with nonreal
conjugates or algebraic numbers of unbounded degree. Nor does this
trace statement alone rule out every realization of an arbitrary
translated finite vector as floors of a non-Pisot base.

\subsection{BKT and BGK in the exponent field}\label{R18:sec:bkt}

\subsubsection{The subgroup is genuine after the correct change of field}

Let $p\ge5$ be prime and suppose
\[
 Q=p^2+p+1
\]
is also prime. The group $\F_{p^3}^*/\F_p^*$ is cyclic of order $Q$.
Choose $x$ generating it and define
\[
 D=\{d\in\Z/Q\Z:\Tr_{\F_{p^3}/\F_p}(x^d)=0\}.
\]
The trace-zero plane gives $|D|=p+1$, $0\notin D$, and $pD=D$.
For every $r\ne0$ there is exactly one ordered representation
$r=d-e$ with $d,e\in D$: the two distinct trace planes meet in one
projective point. This is the standard Singer difference-set
construction; see \cite{MRS}.

Write the unnormalized additive Fourier transform as
\[
 \widehat A(r)=\sum_{a\in A}e^{-2\pi ira/Q}.
\]
The difference-set identity gives
\begin{equation}\label{R18:eq:flat}
 |\widehat D(r)|=\sqrt p\qquad(r\ne0).
\end{equation}

Now set
\[
 H=\langle3\rangle\subset\F_Q^*,\qquad h=|H|,\qquad
 m=(Q-1)/h.
\]
For a nonscalar element $\gamma=cx^u$, the classical local question is
\begin{equation}\label{R18:eq:hit}
 N(u)=|D\cap uH|>0.
\end{equation}
Within its stated size range, BKT's prime-field sum--product theorem
\cite[Theorem 1.1]{R18-BKT} can be applied to $H$; here $HH=H$ supplies
multiplicative closure. BGK's subgroup exponential-sum
theorem \cite[Theorem 6]{R18-BGK} legitimately yields, for $h\ge Q^\delta$,
\[
 \max_{r\ne0}|\widehat H(r)|\le hQ^{-\eta(\delta)},\qquad
 \eta(\delta)=\exp(-C_1/\delta^{C_2}).
\]
The BGK paper is by Bourgain, Glibichuk and Konyagin; Tao is a coauthor
of the precursor BKT paper. Crucially, $H$ is a subgroup in the
\emph{exponent field} $\F_Q$. The set $\{x^{3^k}\}$ is not in
general a multiplicative subgroup of $\F_{p^3}^*$, and we do not apply
the subgroup theorem to that set.

\subsubsection{The exact limitation of phase-blind Fourier bounds}

Fourier inversion and Parseval give
\begin{align}
 N(u)&=\frac{(p+1)h}{Q}
 +\frac1Q\sum_{r\ne0}\widehat D(r)
                       \overline{\widehat H(ur)},\label{R18:eq:mixed}\\
 \sum_{r\ne0}|\widehat H(r)|^2&=h(Q-h).\label{R18:eq:parseval}
\end{align}

\begin{proposition}\label{R18:prop:cs}
The Cauchy--Schwarz estimate from \eqref{R18:eq:flat}--\eqref{R18:eq:parseval} is
\[
 \left|N(u)-\frac{(p+1)h}{Q}\right|
 \le\frac{\sqrt{ph(Q-1)(Q-h)}}{Q}.
\]
Its main term exceeds its displayed error if and only if $h>p^2$.
Consequently it does not certify positivity for any proper subgroup $H$.
\end{proposition}
\begin{proof}
The estimate is Cauchy--Schwarz. Positivity of the resulting lower bound
is equivalent to
$(p+1)^2h>p(Q-1)(Q-h)$.
Since $(p+1)^2+p(Q-1)=(p+1)Q$, this simplifies to $h>p^2$.
A proper subgroup satisfies $h\le(Q-1)/2<p^2$.
\end{proof}

The next statement is stronger, but still concerns only a specified
bounding method, not the actual error or every sum--product method.

\begin{theorem}\label{R18:thm:triangle}
For every proper subgroup $H\subset\F_Q^*$ and $p\ge5$, the exact
termwise triangle-inequality allowance
\[
 E_{\mathrm{tri}}=\frac{\sqrt p}{Q}
                       \sum_{r\ne0}|\widehat H(r)|
\]
is larger than the main term $\mu=(p+1)h/Q$ in \eqref{R18:eq:mixed}.
Thus even knowing every individual amplitude $|\widehat H(r)|$
exactly would not make that phase-blind bound certify a hit.
\end{theorem}
\begin{proof}
Expanding $1_H$ in the $m$ multiplicative characters trivial on $H$
and using the usual Gauss-sum magnitude gives
\[
 \max_{r\ne0}|\widehat H(r)|
 \le\frac{1+(m-1)\sqrt Q}{m}\le\sqrt Q.
\]
Together with Parseval, this implies
\[
 \sum_{r\ne0}|\widehat H(r)|\ge\frac{h(Q-h)}{\sqrt Q},\qquad
 \frac{E_{\mathrm{tri}}}{\mu}
 \ge\frac{\sqrt p(Q-h)}{(p+1)\sqrt Q}.
\]
For a proper subgroup $Q-h>Q/2$. Also $pQ>4(p+1)^2$ for $p\ge5$:
the difference is $p^3-3p^2-7p-4>0$. The claimed ratio is therefore
strictly greater than $1$.
\end{proof}

Directly inserting a BGK supremum gives an error of order
$\sqrt p\,hQ^{-\eta}$. Since the main term is of order
$hQ^{-1/2}$, this would require $\eta>3/4$ asymptotically.
For a proper subgroup of size $Q^\delta$, Parseval forces
$\max|\widehat H|\ge\sqrt{h/2}$, limiting any relative saving to
$\eta\le\delta/2+o(1)\le1/2+o(1)$.
The issue is therefore not a modest improvement to the BGK constant:
one needs information about the \emph{mixed phases} in
\eqref{R18:eq:mixed}, or a different structural argument.

\subsubsection{An exact local counterexample at 71}

\begin{proposition}\label{R18:prop:p71}
At $p=71$, $Q=5113$ is prime and $H=\langle3\rangle$ has order
$213>p$. For the Singer set formed from a root of
$f_0(X)=X^3-4X+4$ over $\F_{71}$, ten of the $24$ multiplicative
$H$-cosets miss $D$. In particular $D\cap13H=\varnothing$.

Inside the original family, $f_{17}$ is irreducible modulo $71$, but
$\Tr(x^{3^k})\ne0$ for every $k\ge0$. Thus prime projective order,
irreducibility, and $|H|>p$ are not sufficient for a classical local hit.
\end{proposition}
\begin{proof}[Exact certificate]
Trial division proves that $71$ and $5113$ are prime. Root testing
proves $f_0$ and $f_{17}$ irreducible over $\F_{71}$.
In the quotient by $f_0$, exact powering gives $x^{5113}=-4$;
the nonscalar projective class therefore has order $5113$.
The trace recurrence enumerates all $72$ elements of $D$, and
quotient-ring multiplication checks the same traces independently.
All $5112$ ordered nonzero differences occur exactly once.

The complete coset intersection multiset is
\[
 \underbrace{0,\ldots,0}_{10},\quad
 \underbrace{3,\ldots,3}_{6},\quad
 \underbrace{6,\ldots,6}_{6},\quad 9,9.
\]
The least-unused missing representatives are
$13,19,22,26,43,46,50,59,67,89$.
Order checks show $\ord_{5113}(3)=213$.

For $f_{17}$, all $213$ traces on the projective cubing clock are
nonzero, and the state after $213$ steps is $43x$. Subsequent cycles
differ by nonzero scalars, so no trace zero can appear. The independent
companion-matrix certificate reproduces these computations.
\end{proof}

This is a counterexample to a proposed \emph{local sufficient
condition}, not to the global composite-floor target. The same
$f_{17}$ has negative-clock zeros at $k=77,79,211\pmod{213}$ at
this auxiliary prime, and another auxiliary prime may work on its
classical clock. Also $D-D=\F_Q$ in this example: maximal additive
expansion itself does not imply the prescribed multiplicative-coset hit.

\subsubsection{A positive statement for almost all additive translates}

\begin{proposition}\label{R18:prop:shifts}
For any fixed $a\in\F_Q^*$, put $R_a(b)=|D\cap(b+aH)|$.
Then
\begin{align*}
 \sum_bR_a(b)&=(p+1)h,\\
 \sum_bR_a(b)(R_a(b)-1)&=h(h-1),\\
 \sum_bR_a(b)^2&=h^2+ph.
\end{align*}
In particular,
\begin{equation}\label{R18:eq:shift-bound}
 \#\{b:R_a(b)=0\}\le\frac{p(Q-h)}{h+p}.
\end{equation}
\end{proposition}
\begin{proof}
The first identity counts pairs in $D\times H$.
For two distinct elements of $H$, the perfect-difference property
provides exactly one corresponding ordered pair in $D$ with their
$a$-scaled difference. This gives the second identity, and adding
the first gives the third. Cauchy--Schwarz on the support of $R_a$
yields
\[
 \#\{b:R_a(b)>0\}\ge\frac{(p+1)^2h}{h+p}.
\]
Subtract from $Q$ and use $(p+1)^2=Q+p$.
\end{proof}

If $h/p\to\infty$, the fraction of missing additive translates tends
to zero. But the classical clock prescribes $b=0$, and the negative
clock prescribes the coupled values $a=2u,b=-u$; neither grants a
freely random translation. At $p=71,a=13$, exactly $262$ of the
$5113$ shifts miss, including $b=0$, whereas
\eqref{R18:eq:shift-bound} gives the valid upper bound $1225$.

For the classical clock, the missing input can also be formulated
using multiplicative characters. Let $H^\perp$ be the $m$ characters
trivial on $H$, and set $T_D(\chi)=\sum_{d\in D}\chi(d)$. Then
\[
 N(u)=\frac1m\left(p+1+
 \sum_{\substack{\chi\in H^\perp\\\chi\ne1}}
 \overline{\chi(u)}T_D(\chi)\right).
\]
Thus $(m-1)\max_{\chi\ne1}|T_D(\chi)|<p+1$ would suffice.
A hypothetical uniform bound $|T_D(\chi)|\le C\sqrt p$ would
give hits in a range roughly $h>Cp^{3/2}$. No such bound is
proved here. The identity $pD=D$ only removes characters with
$\chi(p)\ne1$; equivalently it allows $H$ to be enlarged to
$\langle3,p\rangle$ without changing whether a coset hits.

Even a successful local correlation theorem would still need a
prime-supply argument ensuring the required irreducibility, order,
and scheduled-orbit hypotheses for every fixed coefficient.

\subsection{Gafni Tao exceptional intervals}\label{R18:sec:interval}

\subsubsection{A precise conditional candidate interface}

Gafni--Tao \cite[Definition 1.1 and Theorem 1.3]{R18-GT2026} study
the exceptional starting points for the prime number theorem in short
intervals. For fixed $\delta>0$ and $\theta$, let
\[
 E_\delta(X,\theta)=
 \left\{x\in[X,2X]:
 \left|\sum_{x<n\le x+x^\theta}\Lambda(n)-x^\theta\right|
 \ge\delta x^\theta\right\}.
\]
Their estimates bound its Lebesgue measure by
$X^{\mu(\theta)+\varepsilon}$, with constants and onset depending
on fixed parameters. The explicit recipe uses zero-density and
zero-additive-energy inputs. These are asymptotic exponents, not
numerical constants at a prescribed finite height.

For clarity, the candidate-dependent hypotheses used here are stated
explicitly. Assume a fixed candidate supplies late prime traces
$p_r=\Tr(\beta^{3^r})$ and centers $x_r=p_r^3$, with
\begin{equation}\label{R18:eq:candidate}
\begin{gathered}
 p_{r+1}\text{ is the least prime above }x_r,\
 G_r=p_{r+1}-x_r\sim3x_r^\kappa,\qquad
 \tfrac12<\kappa<\tfrac{21}{40},\qquad
 x_{r+1}\sim x_r^3.
\end{gathered}
\end{equation}
These are the inputs at the present conditional classical-minimum
reduction. We do not assert \eqref{R18:eq:candidate} for an arbitrary
cubic Pisot number or for every member of \eqref{R18:eq:family}.

\begin{proposition}\label{R18:prop:exceptional-block}
Under \eqref{R18:eq:candidate}, fix $1/2<\theta<\kappa$.
At arbitrarily large scales $X\asymp x_r$,
\[
 |E_{1/2}(X,\theta)|\gg X^\kappa.
\]
Consequently an exceptional-set exponent strictly below $\kappa$
at such a $\theta$ would exclude this candidate.
\end{proposition}
\begin{proof}
The interval $(x_r,p_{r+1})$ is prime-free. Choose, for instance,
$X=2x_r/3$, so this interval lies inside $[X,2X]$ away from its
endpoints for large $r$. Every starting point in this gap at distance
at least $(2X)^\theta$ from its right endpoint has a prime-free
window $(x,x+x^\theta]$.

The total contribution of proper prime powers below $3X$ is
$O(X^{1/2}\log^2X)=o(X^\theta)$. Thus every such window is
$1/2$-exceptional for sufficiently large $r$. The set of its starts
has measure at least $G_r-O(X^\theta)\gg X^\kappa$.
If an upper exponent is $\mu(\theta)<\kappa$, choose
$\varepsilon<\kappa-\mu(\theta)$ to obtain a contradiction.
\end{proof}

\subsubsection{What the displayed numerical envelope does not prove}

Let $U(\theta)$ denote the upper-bound envelope produced by the
displayed majorants and optimization in \cite[Tables 1--2 and
Theorem 1.3]{R18-GT2026}. The following calculation is about this
\emph{recipe}, not a lower bound for the unknown true exceptional
exponent $\mu(\theta)$.

At $\sigma=7/10$, the displayed inputs are
\[
 \widetilde A=\frac{30}{13},\qquad
 \widetilde A^*=\frac{235}{39}.
\]
For $1/2<\theta<21/40$, this point is eligible since
$30/13>1/(1-\theta)$. Its two objective values are
\[
 u_2=\frac{9(1-\theta)}{13}+\frac25,\qquad
 u_4=\frac{47(1-\theta)}{26}-\frac15.
\]
Here $u_4<u_2$, so the supremum furnished by that recipe obeys
\begin{equation}\label{R18:eq:envelope}
 U(\theta)\ge u_4\ge\frac{137}{208}
 >\frac{21}{40}>\kappa,
 \qquad
 \frac{137}{208}-\frac{21}{40}=\frac{139}{1040}.
\end{equation}
Thus those specific inputs do not yield the sufficient exponent in
Proposition~\ref{R18:prop:exceptional-block}. This does not prove
$\mu(\theta)\ge137/208$, nor exclude improvements to the
zero-density input, moment bounds, sieve treatment, or
candidate-specific analysis. At the inspected $\sigma$, an improved
recipe would have to remove its eligibility or lower the relevant
objective; reindexing the same candidate does neither.

\subsubsection{The missing packing theorem}

If $N$ disjoint prime-free intervals in $[X,2X]$ each have length
at least $cX^\kappa$, their interior start-sets give, by the same
argument,
\[
 NX^\kappa\ll|E_{1/2}(X,\theta)|.
\]
An upper estimate $\ll X^{U(\theta)+\varepsilon}$ then implies
\begin{equation}\label{R18:eq:packing}
 N\ll X^{U(\theta)-\kappa+\varepsilon}.
\end{equation}
A theorem forcing one candidate to generate
\[
 N\ge X^{U(\theta)-\kappa+\eta}
\]
\emph{disjoint} bad intervals at arbitrarily large common scales,
for fixed $\eta>0$, would contradict \eqref{R18:eq:packing} after
choosing $\varepsilon<\eta$. The current hypotheses provide no
such theorem. The lacunarity $x_{r+1}\sim x_r^3$ instead gives at
most one late center in any $[X,2X]$. Replacing $\beta$ by
$\beta^{3^j}$ merely reindexes the same centers.

Even the favorable endpoint in \eqref{R18:eq:envelope} leaves a required
positive-power multiplicity gap of at least $139/1040$. Prime
patterns at unrelated coefficients or translations do not manufacture
disjoint \emph{prime-free} intervals for this one candidate.

\subsubsection{A sparse deletion diagnostic}

For $\theta>1/2$, replacing the von Mangoldt sum by the weighted
prime sum changes it by a uniformly negligible proper-prime-power
contribution. Suppose primes are deleted in intervals $J_r$ around
$x_r$, of length $O(x_r^\kappa)$, with $x_{r+1}\sim x_r^3$.
For a fixed range $[X,2X]$, only $O(1)$ late deletion intervals can
affect its windows. All affected starting points lie in neighborhoods
of total measure
\[
 O\bigl(X^{\max(\theta,\kappa)}\bigr).
\]
Outside these neighborhoods the original and retained weighted prime
sums are identical. Hence an exceptional-set estimate $O(X^U)$
survives these deletions whenever $U>\max(\theta,\kappa)$, with
the usual harmless adjustments of constants and tolerance.

Conditionally on an all-prime Pisot tower, deleting competing-child
primes while retaining the tower can make its base the minimum for
a \emph{modified subset of primes}. The inspected exceptional
envelope is stable under those sparse deletions. This is neither an
actual Mills counterexample nor a construction of an algebraic
all-prime tower. It identifies information absent from the
exceptional-measure bound alone.

\subsection{Verification and the remaining statements}

The portable scripts in the accompanying \texttt{certificates/}
directory use the Python standard library only. They check:
\begin{itemize}
\item the odd trace normalization, modular repeated-root identity,
both index classes, and exact prime-pair illustration;
\item the Singer difference sets, subgroup orders, complete coset
counts, and the $p=71$ counterexample, including an independent
companion-matrix reconstruction;
\item the exact translated-coset moments and finite shift counts;
\item the rational numerical comparison in \eqref{R18:eq:envelope}.
\end{itemize}
They do not numerically prove an analytic infinitude theorem, the
real-exponential zero-count lemma, or a universal prime-supply claim.
Independent proof reviews and exact regressions were used before packaging.

The remaining requirements are specific. The polynomial-pattern route
would need control of the unshifted, fixed-candidate slice rather than
free translations. The local finite-field route would need an adequate
mixed-correlation statement and a suitable auxiliary prime for every
candidate. The exceptional-interval route would need a sufficiently
small exponent or candidate-driven same-scale packing. None of those
requirements is met here. Conversely, the phase-blind obstruction and
the numerical-envelope obstruction are scoped results, not assertions
that all future uses of the cited methods are impossible.



\endgroup
\clearpage
\section{Final status and the remaining arithmetic problem}\label{sec:final}
The proved restrictions and infinite subfamilies do not constitute a proof
of the universal cubic Pisot recurring-composite statement. They also do
not prove the transcendence of the classical Mills minimum $\xi(3^k)$ or
of $\xi(2\cdot3^k-1)$. The surviving candidate space has not been reduced
to a finite list. Both residue signs must be respected in the classical
analysis, and reindexing a powering orbit does not reset an original seed
index or preserve a fixed coefficient line.

A completion by recurrent divisors would require a theorem that assigns
to every fixed surviving candidate an auxiliary prime with an admissible
trace-zero phase on the required sparse clock. Neither ordinary divisors
of a recurrence, nor density-one coefficient results, nor finitely many
successful tests supply this universal assertion. A completion by prime
intervals would require a genuinely stronger candidate-specific interval
input or a packing theorem at common scales. A completion by finite-field
analysis must control the relevant mixed correlation rather than only
separate Fourier magnitudes. These are precise sufficient directions,
not claims that all possible proof methods have been classified.

\appendix
\section{Source coverage and reproducibility}\label{app:coverage}
The twenty source packages are represented below. The first package
contains proof notes and certificates rather than a standalone PDF.
Copied predecessor sources in later packages are dependencies, not
additional papers. The initial positive divisibility exclusions
$28s$ and $403s$ are subsumed by the later complete positive-shift
elimination in Section~\ref{src:P01}; the initial large negative cutoff
is strengthened by the later negative-shift results. Such obsolete
restrictions are not counted as additional final theorems.

\begin{longtable}{@{}rp{8.5cm}l@{}}
\toprule
Source & Mathematical contribution & Location\\\midrule\endhead
1 & Scaling consequences local minima and normalized tails & \ref{src:INIT}\\
2 & Positive shifts and elimination of the exceptional unit candidates & \ref{src:P01}\\
3 & Negative shifts and integral fixed point rigidity & \ref{src:N01}\\
4 & The multiplier thirteen and a residual branch exclusion & \ref{src:R02}\\
5 & Prime cells height bounds and Pell families & \ref{src:R03}\\
6 & Odd multiplier extensions & \ref{src:R04}\\
7 & Finite prime forests and explicit continuation criteria & \ref{src:R05}\\
8 & Nonunit cubic families and fixed modulus limitations & \ref{src:R06}\\
9 & Full prefix positive extensions and cube rescaling & \ref{src:R07}\\
10 & Density on fixed norm lines & \ref{src:R08}\\
11 & Shifted prime patterns and compactness criteria & \ref{src:R09}\\
12 & Prime coefficients and simultaneous composite floors & \ref{src:R10}\\
13 & Signature controlled density and quantitative exceptional sets & \ref{src:R11}\\
14 & The thin prime power family & \ref{src:R12}\\
15 & A prime seed outlier and certified returns & \ref{src:R13}\\
16 & Adaptive returns and ancestral minimality & \ref{src:R14}\\
17 & Projective return criteria and finite horizon limitations & \ref{src:R15}\\
18 & Coset coverage and output prime order bounds & \ref{src:R16}\\
19 & Both residue signs and all norm depth & \ref{src:R17}\\
20 & Three analytic routes and their exact missing hypotheses & \ref{src:R18}\\
\bottomrule\end{longtable}
The standalone \LaTeX{} source contains the complete typesetting input;
no external bibliography, figure, or input file is needed to compile it.
Portable certificate procedures printed in the text reproduce the core
finite algebra. Commands and relative certificate paths quoted in the verification sections refer to the corresponding extracted original archive, not to additional files required by this standalone source. The original archives also retain their more extensive
machine-readable logs, prime certificates, and checker files. Claims
whose verification relies on a stated analytic theorem retain that
theorem as an explicit mathematical input rather than being described
as computationally certified.

\subsection{Independent checks of the strengthened statements}
The scalar-exponent identity was independently checked in $55\,222$
monic cubic quotient algebras with nonzero constant term at the primes
$5,7,11,13,17,19,23,31$. Each check used exact polynomial reduction
and repeated squaring, including reducible and repeated-root cases.
An additional implementation checked $14\,754$ cubic reductions and
$7\,264$ returned clock states. These are consistency checks for the
universal proof, not an exhaustive verification over all primes.
The two-sign depth formulas were also tested on $536$ irreducible
totally real cubics at indices $0$ through $7$ using integer coefficient
cubing and exact $3$-adic valuations. These finite tests do not replace
the nonvanishing arguments or establish uniform practical bounds for
the modular searches. The previously supplied source archives remain
the record for the older computational certificates; the consolidation
is not represented as a rerun of every original search.

\begin{thebibliography}{99}
\bibitem{BHP}
R. C. Baker, G. Harman and J. Pintz,
The difference between consecutive primes, II,
\emph{Proc. London Math. Soc. (3)} \textbf{83} (2001), 532--562.
\bibitem{BW}
A. Baker and G. W\"ustholz,
Logarithmic forms and group varieties,
\emph{J. reine angew. Math.} \textbf{442} (1993), 19--62.
\bibitem{CH}
M. Cully-Hugill,
\emph{Primes between consecutive powers},
arXiv:2107.14468v3, 14 February 2023.
\url{https://arxiv.org/html/2107.14468v3}.
Theorem 2, Table 2, and equations (2), (9), (15), (19)--(21).
\bibitem{Saito}
K. Saito,
\emph{Transcendency of variants of Mills' constant},
arXiv:2508.16068v3 (2025).
\url{https://arxiv.org/html/2508.16068v3}.
The numbering used here is that of this version, especially
Theorem 1.3 and Proposition 3.1.
\bibitem{Siegel}
C. L. Siegel,
Algebraic integers whose conjugates lie in the unit circle,
\emph{Duke Math. J.} \textbf{11} (1944), 597--602.
\bibitem{K031}
KEIO,
\emph{Prime-representing constants and related structures
  (encodings, arithmetic, geometry, and modular obstructions)},
unpublished research manuscript, 18 October 2026,
item [031] of the Verifiable Mathematical Reasoning Corpus.
\bibitem{K064}
KEIO,
\emph{Ternary trace dynamics and the transcendence of shifted Mills
constants}, unpublished research manuscript, 25 October 2026,
item [064] of the Verifiable Mathematical Reasoning Corpus.
The inherited inputs are its Mills-shadow corollary, approximation-magnitude
lemma, positive-fixed-point lemma, affine trace-return and twisted-limit
theory, and exceptional-table proposition.
\url{https://github.com/QEEDDD/Verifiable-Mathematical-Reasoning-Corpus}.
\bibitem{N01}
KEIO, Negative shifts and integral fixed point rigidity. Source research manuscript, 2026; integrated in Section~\ref{src:N01} of this article.
\bibitem{SaitoIntervals}
K.~Saito, \emph{Intervals without primes near an iterated linear
recurrence sequence}, arXiv:2504.14968v1, 21 April 2025.
Problem~1.1 and Theorem~1.3.
\url{https://arxiv.org/html/2504.14968v1}.
\bibitem{R02}
KEIO, The multiplier thirteen and a residual branch exclusion. Source research manuscript, 2026; integrated in Section~\ref{src:R02} of this article.
\bibitem{R05-AFGW}
K.~Aggarwal, R.~Frot, H.~Gou and H.~Wang,
\emph{Explicit bounds for the graphicality of the prime gap sequence},
arXiv:2512.24230v3, 5 August 2026.
\url{https://arxiv.org/abs/2512.24230v3}.
\bibitem{R05-CHL}
M.~Cully-Hugill and E.~S.~Lee,
\emph{Explicit interval estimates for prime numbers},
corrected version, arXiv:2103.05986v3, 14 September 2022.
\url{https://arxiv.org/pdf/2103.05986v3}.
All numerical interval parameters cited here are from this corrected
version, not the earlier tables.
\bibitem{Maynard}
J.~Maynard, \emph{On the difference between consecutive primes},
arXiv:1201.1787v3, 6 November 2012.
\url{https://arxiv.org/pdf/1201.1787v3}.
\bibitem{R05-OHP}
T.~Oliveira e Silva, S.~Herzog and S.~Pardi,
\emph{Empirical verification of the even Goldbach conjecture and
computation of prime gaps up to $4\cdot10^{18}$},
Mathematics of Computation \textbf{83} (2014), 2033--2060.
Author's publication record: \url{https://sweet.ua.pt/tos/bib/4.12.html}.
Computational results: \url{https://sweet.ua.pt/tos/gaps.html}.
\bibitem{R05-Stadlmann}
J.~Stadlmann, \emph{On the mean square gap between primes},
arXiv:2212.10867v2, 25 October 2026.
\url{https://arxiv.org/pdf/2212.10867v2}.
Page and lemma references in Section~\ref{R05:sec:analytic} refer to
this version.
\bibitem{R06-Lang}
S.~Lang, \emph{Fundamentals of Diophantine Geometry}, Springer, 1983,
Chapter~VIII, ``Siegel's theorem and integral points,'' pp.~188--224.
\url{https://doi.org/10.1007/978-1-4757-1810-2_8}.
\bibitem{R03}
KEIO, Prime cells height bounds and Pell families. Source research manuscript, 2026; integrated in Section~\ref{src:R03} of this article.
\bibitem{R07-K008}
KEIO, \emph{Additive clocks, non-reversible recurrences, and transcendental
prime-representing constants}, unpublished research manuscript,
5 September 2026, item [008] of the Verifiable Mathematical Reasoning
Corpus. Section~6, Theorem~6.1.
\bibitem{P01}
KEIO, Positive shifts and elimination of the exceptional unit candidates. Source research manuscript, 2026; integrated in Section~\ref{src:P01} of this article.
\bibitem{R08-Perucca}
A. Perucca, Prescribing valuations of the order of a point in the
reductions of abelian varieties and tori,
\emph{J. Number Theory} \textbf{129} (2009), 469--476.
\url{https://doi.org/10.1016/j.jnt.2008.07.004}.
Primary version inspected: arXiv:0712.2812v4, 11 October 2008,
Theorem~3.1 and Section~2.
\url{https://arxiv.org/pdf/0712.2812v4}.
\bibitem{Kedlaya}
K. S. Kedlaya, \emph{Notes on analytic number theory},
Chapter~4, Theorem~4.11 and its proof.
\url{https://kskedlaya.org/ant/chap-primes-in-ap.html}.
\bibitem{R08-DLMF}
NIST Digital Library of Mathematical Functions,
Section~27.11, Equation~27.11.9.
\url{https://dlmf.nist.gov/27.11.E9}.
\bibitem{R06}
KEIO, Nonunit cubic families and fixed modulus limitations. Source research manuscript, 2026; integrated in Section~\ref{src:R06} of this article.
\bibitem{GreenTao}
B.~Green and T.~Tao,
\emph{The primes contain arbitrarily long arithmetic progressions},
Annals of Mathematics \textbf{167} (2008), 481--547.
Theorems 1.1--1.2, printed p.~2 of
\href{https://arxiv.org/pdf/math/0404188v6}{arXiv:math/0404188v6}.
\bibitem{R09-LW}
T.~H. L\^e and J.~Wolf, \emph{Polynomial configurations in the primes},
\href{https://arxiv.org/abs/1210.4659}{arXiv:1210.4659}.
Theorem 4 and the following paragraph, printed p.~2;
\href{https://home.olemiss.edu/~leth/papers/appd.pdf}{author-hosted full text}.
\bibitem{TZ2008}
T.~Tao and T.~Ziegler,
\emph{The primes contain arbitrarily long polynomial progressions},
Acta Mathematica \textbf{201} (2008), 213--305;
\href{https://arxiv.org/abs/math/0610050}{arXiv:math/0610050}.
\bibitem{Lenstra}
H. Lenstra, \emph{The Chebotarev density theorem},
lecture notes, Fact~2.1 and Theorem~3.1.
\url{https://websites.math.leidenuniv.nl/algebra/Lenstra-Chebotarev.pdf}.
\bibitem{R08}
KEIO, Density on fixed norm lines. Source research manuscript, 2026; integrated in Section~\ref{src:R08} of this article.
\bibitem{R11-AKK}
S. Arango-Pi\~neros, D. Keliher, and C. Keyes,
\emph{Mertens' theorem for Chebotarev sets},
arXiv:2103.14747v2 (2021), Theorem~A.
\url{https://arxiv.org/abs/2103.14747v2}.
\bibitem{R11-BS}
J. Br\"udern and K. Soundararajan,
Common divisors of totients of polynomial sequences,
\emph{Mathematische Zeitschrift} \textbf{299} (2021), 527--542,
Section~2, Lemma~3.
\url{https://doi.org/10.1007/s00209-020-02674-7}.
In arXiv:1909.10808v1, the corresponding statement is Lemma~2:
\url{https://arxiv.org/abs/1909.10808v1}.
\bibitem{R11-HIKK}
S. Hahn, D. Ismailescu, G. Kim, and M. Kim,
\emph{Explicit algebraic numbers all whose integer parts of powers
are always composite},
arXiv:2608.05309v1 (2026), Section~4, Theorem~4.
\url{https://arxiv.org/abs/2608.05309v1}.
\bibitem{R10}
KEIO, Prime coefficients and simultaneous composite floors. Source research manuscript, 2026; integrated in Section~\ref{src:R10} of this article.
\bibitem{MRS}
Tam\'as M\'esz\'aros, Lajos R\'onyai, and Tibor Szab\'o,
\emph{Singer difference sets and the projective norm graph},
arXiv:1908.05591v1 (2019).
\url{https://arxiv.org/abs/1908.05591}.
See Section 1.2 and Corollary 1.2 for the classical trace-zero cosets
and their norm-one polynomial description.
\bibitem{R16-Sun}
Z.-H.~Sun, \emph{On the theory of cubic residues and nonresidues},
Acta Arithmetica \textbf{84} (1998), 291--335,
equations (1.1)--(1.2).
\url{https://matwbn.icm.edu.pl/ksiazki/aa/aa84/aa8441.pdf}.
\bibitem{R16-Evertse}
J.-H.~Evertse, \emph{On sums of S-units and linear recurrences},
Compositio Mathematica \textbf{53} (1984), 225--244, Corollary 3.
\url{https://www.numdam.org/item/CM_1984__53_2_225_0/}.
\bibitem{R16-Stewart}
C.~L.~Stewart, \emph{On prime factors of terms of linear recurrence
sequences}, in \emph{Number Theory and Related Fields},
Springer Proceedings in Mathematics \& Statistics \textbf{43} (2013),
equation (2), p.~342.
\url{https://doi.org/10.1007/978-1-4614-6642-0_18}.
\bibitem{R16-BHV}
Y.~Bilu, G.~Hanrot, and P.~M.~Voutier,
\emph{Existence of primitive divisors of Lucas and Lehmer numbers},
Journal f\"ur die reine und angewandte Mathematik \textbf{539} (2001),
75--122, Theorem 1.4.
\url{https://doi.org/10.1515/crll.2001.080}.
\bibitem{R16-Roskam}
H.~Roskam, \emph{Prime divisors of linear recurrences and Artin's
primitive root conjecture for number fields}, Journal de Th\'eorie des
Nombres de Bordeaux \textbf{13} (2001), 303--314.
\url{https://doi.org/10.5802/jtnb.323}.
\bibitem{R16-Jarviniemi}
O.~J\"arviniemi, \emph{Positive lower density for prime divisors of
generic linear recurrences}, Mathematical Proceedings of the Cambridge
Philosophical Society \textbf{175} (2023), 467--478,
Theorems 1.1--1.2 and \S8.
\url{https://doi.org/10.1017/S0305004123000257}.
\bibitem{R16-IngramSilverman}
P.~Ingram and J.~H.~Silverman, \emph{Primitive divisors in arithmetic
dynamics}, Mathematical Proceedings of the Cambridge Philosophical Society
\textbf{146} (2009), 289--302.
\url{https://arxiv.org/abs/0707.2505}.
\bibitem{R17-Maksoud}
A.~Maksoud, \emph{On the rank of Leopoldt's and Gross's regulator maps},
Documenta Mathematica \textbf{28} (2023), 1441--1471,
Theorem 2.1, p.~1445 (Baker--Brumer theorem).
\url{https://doi.org/10.4171/DM/935}.
\url{https://ems.press/content/serial-article-files/31139}.
\bibitem{R17-Milne}
J.~S.~Milne, \emph{Fields and Galois Theory}, version 5.10,
September 2022, Propositions 5.27--5.28, pp.~74--75.
\url{https://www.jmilne.org/math/CourseNotes/FT.pdf}.
\bibitem{R18-TZ2016}
T.~Tao and T.~Ziegler,
\emph{Polynomial patterns in the primes},
arXiv:1603.07817, Theorems 1.2 and 1.4 and Corollary 1.5.
\url{https://arxiv.org/abs/1603.07817}.
\bibitem{R18-BKT}
J.~Bourgain, N.~Katz and T.~Tao,
\emph{A sum-product estimate in finite fields, and applications},
Geometric and Functional Analysis \textbf{14} (2004), 27--57.
\url{https://arxiv.org/abs/math/0301343}.
\bibitem{R18-BGK}
J.~Bourgain, A.~A.~Glibichuk and S.~V.~Konyagin,
\emph{Estimates for the number of sums and products and for exponential
sums in fields of prime order},
Journal of the London Mathematical Society \textbf{73} (2006), 380--398.
\url{https://doi.org/10.1112/S0024610706022721}.
\bibitem{R18-GT2026}
A.~Gafni and T.~Tao,
\emph{On the number of exceptional intervals to the prime number
theorem in short intervals},
Essential Number Theory \textbf{5} (2026), 221--241.
The inspected theorem and numerical tables are in arXiv:2505.24017v1.
\url{https://arxiv.org/html/2505.24017v1}.
\bibitem{RamachandraSrinivasan1983}
K. Ramachandra and S. Srinivasan,
A note to a paper by Ramachandra on transcendental numbers,
\emph{Hardy--Ramanujan Journal} \textbf{6} (1983), 37--44.
\url{https://doi.org/10.46298/hrj.1983.98}.
\end{thebibliography}
\end{document}
